%%%%%%%%%%%%%%%%%%%%%%%%%%%%%%%%%%%%%%%%%%%%%%%%%%%%%%%%%%%%%%%%%%%%%%%%%%%%%
%%
%%  Pierre-Henri Chaudouard and Gerard Laumon
%%  "Le lemme fondamental pondere. II. Enonces cohomologiques"
%%  Annals of Mathematics 176 (2012), 1647-1781
%%  http://dx.doi.org/10.4007/annals.2012.176.3.6
%%
%%  ENGLISH TRANSLATION, following the JOURNAL version (not the arXiv preprint).
%%
%%  Numbering conventions reproduce the journal exactly:
%%   * paragraphs are numbered  <section>.<n>          (environment "paragr")
%%   * displayed equations      <section>.<n>.<m>      (counter "equation")
%%   * theorems, lemmas, propositions, definitions, corollaries, remarks,
%%     hypotheses and examples share ONE counter        (counter "theo"),
%%     numbered <section>.<n>.<m>, independent of the equation counter.
%%     (The journal has both "(1.4.1)" and "Theoreme 1.4.1" inside 1.4.)
%%
%%  Every numbered item carries a label equal to its journal number, so that
%%  \ref always reproduces the journal's number:
%%     \label{sec:9}      section 9
%%     \label{S:4.9}      paragraph 4.9
%%     \label{eq:4.7.1}   equation (4.7.1)
%%     \label{thm:9.1.3}  Theorem/Lemma/Proposition/... 9.1.3
%%  Descriptive aliases from the earlier draft are kept alongside them.
%%
%%%%%%%%%%%%%%%%%%%%%%%%%%%%%%%%%%%%%%%%%%%%%%%%%%%%%%%%%%%%%%%%%%%%%%%%%%%%%

\documentclass[12pt,a4paper,oneside,english,leqno]{smfart}
\usepackage[T1]{fontenc}
\usepackage[utf8]{inputenc}
\usepackage[english]{babel}
\usepackage[a4paper,top=2.5cm,bottom=2.5cm,left=2.5cm,right=2.5cm]{geometry}

% New TX supplies the AMS symbols. Loading amssymb as well produces a
% duplicate \Bbbk definition on installations where both are present.
\usepackage{newtxtext}
\makeatletter
\let\Bbbk\@undefined
\makeatother
\usepackage{newtxmath}

\usepackage{amscd}
\usepackage[all]{xy}
\usepackage{xcolor}

%------------------------------------------------------------------
%  Numbering
%------------------------------------------------------------------
%  Equations: one counter, reset at every paragraph (= subsection).
\numberwithin{equation}{subsection}
\renewcommand{\theequation}{\thesubsection.\arabic{equation}}

%  Theorem-like environments: a SEPARATE counter, also reset at every
%  paragraph.  This is what the journal does.
\newcounter{theo}
\numberwithin{theo}{subsection}
\renewcommand{\thetheo}{\thesubsection.\arabic{theo}}

%------------------------------------------------------------------
%  Paragraph environments (the numbered units "1.1", "1.2", ... )
%------------------------------------------------------------------
\makeatletter
\newenvironment{paragr}[1][]{\medskip\refstepcounter{subsection}%
  \par\noindent\hspace*{\parindent}\thesubsection.\ \emph{#1}%
  \@ifundefined{@empty}{}{}\ignorespaces}{\medskip}
\makeatother
\newenvironment{sparagr}[1][]{\refstepcounter{subsubsection}%
  \noindent\textbf{\thesubsubsection.\ \ #1}}{\medskip}

%------------------------------------------------------------------
%  Theorem environments -- all on the shared counter "theo"
%------------------------------------------------------------------
\theoremstyle{plain}
\newtheorem{theorem}[theo]{Theorem}
\newtheorem{theoreme}[theo]{Theorem}
\newtheorem{lemma}[theo]{Lemma}
\newtheorem{lemme}[theo]{Lemma}
\newtheorem{proposition}[theo]{Proposition}
\newtheorem{corollary}[theo]{Corollary}
\newtheorem{corollaire}[theo]{Corollary}
\newtheorem{conjecture}[theo]{Conjecture}

\theoremstyle{definition}
\newtheorem{definition}[theo]{Definition}
\newtheorem{appdefinition}[theo]{Definition}

\theoremstyle{remark}
\newtheorem{remark}[theo]{Remark}
\newtheorem{remarks}[theo]{Remarks}
\newtheorem{exemple}[theo]{Example}
\newtheorem{example}[theo]{Example}
\newtheorem{hypothese}[theo]{Hypothesis}
\newtheorem{hypothesis}[theo]{Hypothesis}
\newtheorem{numero}[theo]{}

%  Legacy aliases used in the first draft.
\newenvironment{remarque}{\begin{remark}}{\end{remark}}
\newenvironment{remarques}{\begin{remarks}}{\end{remarks}}
\newtheorem{apptheoreme}[theo]{Theorem}
\newtheorem{appcorollaire}[theo]{Corollary}
\newtheorem{appproposition}[theo]{Proposition}
\newtheorem{subpro}[theo]{Proposition}

%------------------------------------------------------------------
%  Proof environment
%------------------------------------------------------------------
\newenvironment{preuve}[1][]{\medskip\noindent\textit{Proof.} #1 ---\ }%
  {\hfill\ensuremath{\square}\medskip}
\DeclareMathOperator{\Lie}{Lie}
\DeclareMathOperator{\nd}{nd}
\DeclareMathOperator{\nr}{nr}
\DeclareMathOperator{\reg}{reg}
\DeclareMathOperator{\freg}{freg}
\DeclareMathOperator{\vol}{vol}
\DeclareMathOperator{\rang}{rank}
\DeclareMathOperator{\Ad}{Ad}
\DeclareMathOperator{\AD}{AD}
\DeclareMathOperator{\End}{End}
\DeclareMathOperator{\ad}{ad}
\DeclareMathOperator{\ab}{ab}
\DeclareMathOperator{\supp}{supp}
\DeclareMathOperator{\mes}{meas}
\DeclareMathOperator{\der}{der}
\DeclareMathOperator{\scnx}{sc}
\DeclareMathOperator{\SCNX}{SC}
\DeclareMathOperator{\Norm}{Norm}
\DeclareMathOperator{\Cent}{Cent}
\DeclareMathOperator{\el}{ell}
\DeclareMathOperator{\Out}{Out}
\DeclareMathOperator{\Aut}{Aut}
\DeclareMathOperator{\Gal}{Gal}
\DeclareMathOperator{\Hom}{Hom}
\DeclareMathOperator{\homc}{\mathcal{H}\mathit{om}}
\DeclareMathOperator{\inv}{inv}
\DeclareMathOperator{\Int}{Int}
\DeclareMathOperator{\Ind}{Ind}
\DeclareMathOperator{\Indp}{Ind-p}
\DeclareMathOperator{\Id}{Id}
\DeclareMathOperator{\stab}{stab}
\DeclareMathOperator{\aff}{aff}
\DeclareMathOperator{\loc}{loc}
\DeclareMathOperator{\geom}{geom}
\DeclareMathOperator{\Tran}{Tran}
\DeclareMathOperator{\Res}{Res}
\DeclareMathOperator{\cusp}{cusp}
\DeclareMathOperator{\disc}{disc}
\DeclareMathOperator{\Ker}{Ker}
\DeclareMathOperator{\Cok}{Cok}
\DeclareMathOperator{\Coker}{Coker}
\DeclareMathOperator{\Kil}{Kil}
\DeclareMathOperator{\rss}{rank_{ss}}
\DeclareMathOperator{\vect}{vect}
\DeclareMathOperator{\val}{val}
\DeclareMathOperator{\SL}{SL}
\DeclareMathOperator{\Spec}{Spec}
\DeclareMathOperator{\Spf}{Spf}
\DeclareMathOperator{\Div}{Div}
\DeclareMathOperator{\Tor}{Tor}
\DeclareMathOperator{\Ext}{Ext}
\DeclareMathOperator{\trace}{trace}
\DeclareMathOperator{\Sym}{Sym}
\DeclareMathOperator{\car}{\mathfrak{car}}
\DeclareMathOperator{\lng}{long}
\DeclareMathOperator{\Def}{Def}
\DeclareMathOperator{\card}{card}
\DeclareMathOperator{\tors}{tors}
\DeclareMathOperator{\parab}{par}
\DeclareMathOperator{\sep}{sep}
%----------------------------------
%
%
%
%    MACROS MATHS
%
%
%
%
%
%----------------------------------
%-----------------------------Ensembles ---------------------
\newcommand{\ZZ}{\mathbb{Z}}
\newcommand{\Gm}{\mathbb{G}_m}
\newcommand{\Gmk}{\mathbb{G}_{m,k}}
\newcommand{\SG}{\mathfrak{S}}
\newcommand{\NN}{\mathbb{N}}
\newcommand{\RR}{\mathbb{R}}
\newcommand{\AAA}{\mathbb{A}}
\newcommand{\CC}{\mathbb{C}}
\newcommand{\BBB}{\mathbb{B}}
\newcommand{\QQ}{\mathbb{Q}}
\newcommand{\FF}{\mathbb{F}}
\newcommand{\Fq}{\mathbb{F}_q}
\newcommand{\Qp}{\mathbb{Q}_p}
\newcommand{\Ql}{\mathbb{Q}_{\ell}}
\newcommand{\Qlb}{\overline{\mathbb{Q}}_{\ell}}
\newcommand{\Qb}{\overline{\mathbb{Q}}}
\newcommand{\GA}{G(\mathbb{A})}
\newcommand{\CGA}{C^{\infty}_{c}(G(\mathbb{A}))}
\newcommand{\GQ}{G(\mathbb{Q})}
\newcommand{\ga}{\gamma}

%---------------Lettres calligraphi\'{e}es ------------------



\newcommand{\oc}{\mathcal{O}}
\newcommand{\rc}{\mathcal{R}}
\newcommand{\Sc}{\mathcal{S}}
\newcommand{\tc}{\mathcal{T}}
\newcommand{\uc}{\mathcal{U}}
\newcommand{\vc}{\mathcal{V}}
\newcommand{\wc}{\mathcal{W}}
\newcommand{\ec}{\mathcal{E}}
\newcommand{\gc}{\mathcal{G}}
\newcommand{\mc}{\mathcal{M}}
\newcommand{\yc}{\mathcal{Y}}
\newcommand{\lc}{\mathcal{L}}
\newcommand{\dc}{\mathcal{D}}
\newcommand{\fc}{\mathcal{F}}
\newcommand{\pc}{\mathcal{P}}
\newcommand{\hc}{\mathcal{H}}
\newcommand{\nc}{\mathcal{N}}
\newcommand{\Ac}{\mathcal{A}}
\newcommand{\bc}{\mathcal{B}}
\newcommand{\kc}{\mathcal{K}}
\newcommand{\Ic}{\mathcal{I}}
\newcommand{\Jc}{\mathcal{J}}
\newcommand{\Qc}{\mathcal{Q}}
\newcommand{\xc}{\mathcal{X}}

%--------------------Gros chapeaux
\newcommand{\Gc}{\widehat{G}}
\newcommand{\Tc}{\widehat{T}}
\newcommand{\Pc}{\widehat{P}}
\newcommand{\Bc}{\widehat{B}}
\newcommand{\Mc}{\widehat{M}}
\newcommand{\Lc}{\widehat{L}}
\newcommand{\Rc}{\widehat{R}}
\newcommand{\Sch}{\widehat{S}}
\newcommand{\Uc}{\widehat{U}}
\newcommand{\Hc}{\widehat{H}}
\newcommand{\Zc}{\widehat{Z}}

%--------------------Lettres grasses math\'{e}matiques ---------
\newcommand{\GG}{\mathbf{G}}
\newcommand{\HH}{\mathbf{H}}
\newcommand{\PP}{\mathbf{P}}
\newcommand{\QQQ}{\mathbf{Q}}
\newcommand{\MM}{\mathbf{M}}
\newcommand{\XX}{\mathbf{X}}
\newcommand{\LL}{\mathbf{L}}
\newcommand{\CCC}{\mathbf{C}}
\newcommand{\RRR}{\mathbf{R}}
\newcommand{\TT}{\mathbf{T}}
\newcommand{\BB}{\mathbf{B}}

%----------------Lettres gothiques ------------------------
\newcommand{\ggo}{\mathfrak{g}}
\newcommand{\kgo}{\mathfrak{k}}
\newcommand{\gA}{\mathfrak{g}(\mathbb{A})}
\newcommand{\gQ}{\mathfrak{g}(\mathbb{Q})}
\newcommand{\SgA}{\mathcal{S}(\mathfrak{g}(\mathbb{A}))}
\newcommand{\of}{\mathfrak{o}}
\newcommand{\Ogo}{\mathfrak{O}}
\newcommand{\mgo}{\mathfrak{m}}
\newcommand{\ngo}{\mathfrak{n}}
\newcommand{\ago}{\mathfrak{a}}
\newcommand{\cgo}{\mathfrak{c}}
\newcommand{\pgo}{\mathfrak{p}}
\newcommand{\qgo}{\mathfrak{q}}
\newcommand{\sgo}{\mathfrak{s}}
\newcommand{\rgo}{\mathfrak{r}}
\newcommand{\hgo}{\mathfrak{h}}
\newcommand{\tgo}{\mathfrak{t}}
\newcommand{\lgo}{\mathfrak{l}}
\newcommand{\ugo}{\mathfrak{u}}
\newcommand{\zgo}{\mathfrak{z}}
\newcommand{\Zgo}{\mathfrak{Z}}
\newcommand{\bgo}{\mathfrak{b}}
\newcommand{\Xgo}{\mathfrak{X}}
\newcommand{\Fgo}{\mathfrak{F}}

%--------------lettres avec une barre ---------------

\newcommand{\gb}{\bar{g}}
\newcommand{\qb}{\bar{q}}
\newcommand{\mb}{\bar{m}}
\newcommand{\nb}{\bar{n}}
\newcommand{\hb}{\bar{h}}
\newcommand{\xb}{\bar{x}}

%----------------------------------------------

\newcommand{\Gg}{G_{\gamma}}
\newcommand{\GAun}{G(\mathbb{A})^1}

\newcommand{\npa}{\ngo_P(\AAA)}
\newcommand{\sig}{\sigma_1^2}
\newcommand{\quo}{\GQ \backslash \GAun}
\newcommand{\quop}{P_1(\QQ) \backslash \GAun}

\newcommand{\rgoe}{\mathfrak{r}^\ast}
\newcommand{\mti}{\tilde{\mathfrak{\mgo}}}
\newcommand{\nbar}{\overline{\mathfrak{n}}}
\newcommand{\Oreg}{\mathcal{O}_{\reg}}

\newcommand{\hgoe}{\mathfrak{h}^\ast}
\newcommand{\Yb}{\overline{Y}}
\newcommand{\CT}{\mathcal{C}_0^2 (T_1)}


%-------------------Lettres grecques ---------------------------
\newcommand{\al}{\alpha}
\newcommand{\be}{\beta}
\newcommand{\Si}{\Sigma}
\newcommand{\zreg}{\zgo_{M,\reg}}
\newcommand{\am}{a_M}
\newcommand{\ams}{a_M^\ast}
\newcommand{\wm}{\wc(\am)}
\newcommand{\om}{\omega}
\newcommand{\Om}{\Omega}
\newcommand{\Lo}{\Lambda_\om}

\newcommand{\la}{\lambda}

\newcommand{\Ga}{\Gamma}
\newcommand{\back}{\backslash}
\newcommand{\ie}{i.e.\ }
\newcommand{\Cc}{C_c^\infty}

%-----------------------lettres à petit chapeau -------------
\newcommand{\ic}{\hat{i}}
\newcommand{\jc}{\hat{j}}
\newcommand{\ac}{\hat{a}}
\newcommand{\rf}{\hat{r}}





\newcommand{\geg}{\Ga^{\ec}(\ggo)}
\newcommand{\bg}{\langle}
\newcommand{\bd}{\rangle}

%--------------------lettres etoilees -----------------------
\newcommand{\Le}{L^\ast}
\newcommand{\Ge}{G^\ast}
\newcommand{\Me}{M^\ast}
\newcommand{\Te}{T^\ast}
\newcommand{\Se}{S^\ast}
\newcommand{\gem}{\Ga^{\ec}_{\Ge}(\mgo)}
\newcommand{\mgoe}{\mgo^\ast}
\newcommand{\ggoe}{\ggo^\ast}
\newcommand{\lgoe}{\lgo^\ast}
\newcommand{\sgoe}{\sgo^\ast}


\newcommand{\lce}{\lc^\ast}
\newcommand{\geme}{\Ga^{\ec}_{\Ge}(\mgoe)}
\newcommand{\gege}{\Ga^{\ec}(\ggoe)}
\newcommand{\gegell}{\Ga^{\ec}_{\el}(\ggo)}
\newcommand{\gag}{\Ga(\ggo)}
\newcommand{\gge}{\Ga(\ggoe)}
\newcommand{\fe}{f^\ast}
\newcommand{\dge}{\Delta_{\Ge}}
\newcommand{\dme}{\Delta_{\Me}}
\newcommand{\Fcl}{\overline{F}}
\newcommand{\kcl}{\overline{l}}


%------------------- L-groupes ---------------
\newcommand{\LG}{{}^LG}
\newcommand{\LM}{{}^LM}
\newcommand{\LT}{{}^LT}
\newcommand{\LLL}{{}^LL}
%---------------------Lettres pointees ---------------
\newcommand{\Gp}{\dot{G}}
\newcommand{\Mp}{\dot{M}}
\newcommand{\eps}{\varepsilon}
\newcommand{\epse}{\epsilon^\ast}
\newcommand{\Qe}{Q^\ast}
\newcommand{\Pe}{P^\ast}
\newcommand{\ome}{\omega^\ast}
\newcommand{\Ome}{\Omega^\ast}
\newcommand{\Gps}{G_{\eps}}
\newcommand{\Geps}{\Ge_{\epse}}

%--------------------------------
%
%
%   nouveaux <= et >=
%
%
%
%----------------------------
\renewcommand{\leq}{\leqslant}
\renewcommand{\geq}{\geqslant}


\usepackage[
  colorlinks=true,
  linkcolor={blue!45!black},
  citecolor={green!35!black},
  urlcolor={blue!45!black},
  linktocpage=true,
  pdfborder={0 0 0},
  bookmarksnumbered=true,
  pdftitle={The Weighted Fundamental Lemma. II. Cohomological Statements},
  pdfauthor={Pierre-Henri Chaudouard and Gerard Laumon}
]{hyperref}
\IfFileExists{bookmark.sty}{\usepackage{bookmark}}{}

% The journal resets both counters at each numbered paragraph.  Give hyperref
% the full journal number as well, so distinct equations and statements never
% receive duplicate PDF destinations.
\renewcommand{\theHequation}{\thesection.\arabic{subsection}.\arabic{equation}}
\newcommand{\theHtheo}{\thesection.\arabic{subsection}.\arabic{theo}}

\pagestyle{plain}
\setlength{\parindent}{1.5em}
\setlength{\parskip}{0pt}
\setlength{\emergencystretch}{3em}
\tolerance=2000
\hbadness=10000
\vbadness=10000

\title[The Weighted Fundamental Lemma. II]{The Weighted Fundamental Lemma. II. Cohomological Statements}
\alttitle{Le lemme fondamental pond\'er\'e. II. \`Enonc\'es cohomologiques}
\author{Pierre-Henri Chaudouard \and G\'erard Laumon}
\thanks{English translation of the journal version}
\date{\emph{Annals of Mathematics} \textbf{176} (2012), 1647--1781\\
\href{https://doi.org/10.4007/annals.2012.176.3.6}{doi:10.4007/annals.2012.176.3.6}}

\begin{document}

\begin{abstract}
In this paper, we study the cohomology of the truncated Hitchin fibration, which was introduced in a previous paper. We extend Ng\^o's main theorems on the cohomology of the elliptic part of the Hitchin fibration. As a consequence, we get a proof of Arthur's weighted fundamental lemma.
\end{abstract}

\maketitle

\begin{center}\small
Section, paragraph, equation and statement numbers are those of the journal version.
\end{center}
\bigskip

\begin{center}
\begin{minipage}{0.86\textwidth}\small
\textbf{Note on the text.} This is a translation of the journal version, not of the arXiv preprint; the two differ in several places. Section, paragraph, equation and statement numbers are those of the journal, so a reference such as "Theorem 10.7.3" points at the same result in both. A handful of obvious misprints in the printed French (unbalanced parentheses, a missing accent on a subscript, and similar) have been silently corrected where leaving them would mislead; all other idiosyncrasies of the original notation have been reproduced as printed.
\end{minipage}
\end{center}
\bigskip

\renewcommand{\contentsname}{Contents}
\tableofcontents
\bigskip
\section{Introduction}\label{sec:1}

\begin{paragr}\label{S:1.1} The aim of this work, which follows \cite{LFPI}, is to prove the \emph{weighted fundamental lemma}. This statement was conjectured by Arthur in his work on the stabilization of the Arthur--Selberg traces (cf. conjecture 5.1 of \cite{Arthur-stable-I}). It generalizes the Langlands–Shelstad fundamental lemma and consists of a family of combinatorial identities between weighted $p$-adic orbital integrals. After Ng\^o's work on the fundamental lemma, it is the only missing ingredient in the stabilization of the trace formula. According to a strategy due to Langlands, one should obtain, by means of the stable trace formula, all endoscopic functorialities. For example, Arthur’s ongoing work should provide a classification of the automorphic spectrum of classical groups in terms of that of general linear groups. Apart from work of Kottwitz (cf. \cite{Kottwitz-basechange}) and Whitehouse (cf. \cite{Whitehouse}), the weighted fundamental lemma had until now been a completely open problem.
\end{paragr}

\begin{paragr}\label{S:1.2} More precisely, we prove a variant in positive characteristic and for the Lie algebras of Arthur’s statement. Moreover, in order to simplify the exposition as much as possible, we restrict ourselves in this article to the case of split groups. The general case, which also includes the “non-standard” forms of the weighted fundamental lemma due to Waldspurger, can be obtained by similar methods and will be treated later. Work of Waldspurger (cf. \cite{Waldspurger-tordu} and \cite{Waldspurger-endoscopie}) shows that all these variants of the weighted fundamental lemma imply Arthur’s original statement for groups over $p$-adic fields.
\end{paragr}

\begin{paragr}\label{S:1.3} From now on, we reserve the term fundamental lemma or weighted fundamental lemma for the variants for Lie algebras in positive characteristic of the original statements. In \cite{Ngo-lemme}, Ng\^o proved the ordinary fundamental lemma. In his approach, the fundamental lemma results from a cohomological statement on the elliptic part of the Hitchin fibration. Similarly here, the weighted fundamental lemma is a consequence of a cohomological theorem (cf. theorem \ref{thm:10.7.3}) on the truncated Hitchin fibration introduced in \cite{LFPI}. The key result in Ng\^o's approach is his theorem describing the supports of the different constituents of the cohomology of the elliptic part of the Hitchin fibration. In an informal sense, this theorem shows that the global orbital integrals are the “limits” of the orbital integrals associated to the most regular elements possible. In particular, to prove the global variant of the fundamental lemma, it suffices to do so for these very regular elements for which a “by hand” verification is possible. The local statement then follows from the global one. Our main cohomological result is that, for the cohomology of the truncated Hitchin fibration, no new support appears. The most striking consequence is that the global weighted orbital integrals are also “limits” of ordinary orbital integrals.
\end{paragr}

\begin{paragr}[The weighted fundamental lemma.]\label{S:1.4} We now state the local theorem that we prove. Let $\Fq$ be a finite field with $q$ elements and $G$ a connected reductive group, defined and split over $\Fq$. Let $T$ be a maximal torus of $G$, defined and split over $\Fq$. We assume that the order of the Weyl group $W^G$ of $(G,T)$ is invertible in $\Fq$. Let $M$ be a Levi subgroup of $G$ defined over $\Fq$ containing the torus $T$. We denote by Gothic letters the corresponding Lie algebras. Consider the local field $K=\Fq((\eps))$ and $\oc=\Fq[[\eps]]$, its ring of integers. Let $X\in \mgo(K)$ be a semisimple and $G$-regular element. Its centralizer in $G\times_{\Fq}K$ is a maximal $K$-torus denoted $J_X$, which is contained in $M\times_{\Fq}K$. In this context, an Arthur weighted orbital integral is given by
$$J_M^G(X)=|D^G(X)|^{1/2} \int_{J_X(K)\back G(K)} \mathbf{1}_{\ggo(\oc)}(\Ad(g^{-1})X) v_M^G(g) dg,$$
where
\begin{itemize}
\item $|D^G(X)|$ is the usual absolute value of the Weyl discriminant; \item $\mathbf{1}_{\ggo(\oc)}$ is the characteristic function of $\ggo(\oc)$; \item $\Ad$ denotes the adjoint action of $G$ on $\ggo$; \item $v_M^G(g)$ is Arthur’s “weight” function which is left $M(F)$-invariant; \item $dg$ is the quotient measure deduced from Haar measures on $G(K)$ and $J_X(K)$.
\end{itemize}
The weight and the measure $dg$ depend on choices that can be normalized. With appropriate measure conventions, the integral $J_M^G(X)$ depends on $X$ only up to $M(K)$-conjugation.

Let $\Gc$ and $\Tc$ be the complex dual groups in the sense of Langlands of $G$ and $T$. The group $\Tc$ identifies with a maximal torus of $\Gc$. The dual $\Mc$ of $M$ identifies with a Levi subgroup of $\Gc$ which contains $\Tc$. Let $s\in \Tc$ and let $\Mc'$ be the neutral component of the centralizer of $s$ in $\Mc$. This is a connected reductive subgroup of $\Mc$ which contains $\Tc$. We assume that the neutral components of the centers of $\Mc$ and $\Mc'$ are equal. Let $M'$ be the group over $\Fq$ dual to $\Mc'$. The torus $T$ identifies with a maximal torus of $M'$ and the Weyl group $W^{M'}$ identifies with a subgroup of the Weyl group $W^M$. Thus we have a canonical morphism
\begin{equation}\label{eq:1.4.1}
  \mgo'//M' \simeq \tgo //W^{M'} \to \tgo //W^M \simeq \mgo//M,
\end{equation}
where the isomorphisms are those of Chevalley and the quotients are the categorical quotients. The open subset of semisimple $G$-regular elements defines, via the canonical map $\mgo\to \mgo//M$, an open subset in the quotient $\mgo//M$, whose preimage under the composite map $\mgo'\to \mgo'//M'\to\mgo//M$ is the open subset of semisimple regular elements of $G$. Following Waldspurger (cf. \cite{Waldspurger-transfert}), one can adapt the definition of Langlands–Shelstad transfer factors (cf. \cite{Langlands-Shelstad}) to Lie algebras. By considering $M'$ as an “endoscopic group” of $M$, one obtains a function
$$\Delta_{M',M} : \mgo'(K) \times \mgo(K) \to \CC$$
defined on the open set of pairs $(Y,X)$ where each component is a semisimple and $G$-regular element. The transfer factor is in fact invariant under the adjoint action of $M'(K)\times M(K)$. For a pair $(Y,X)$ with each component semisimple and $G$-regular, the factor $\Delta_{M',M}(Y,X)$ is nonzero if and only if $Y$ and $X$ have the same image in $(\mgo//M)(K)$ via the canonical maps $\mgo'(K)\to (\mgo//M)(K)$ (induced by (\ref{eq:1.4.1})) and $\mgo(K)\to (\mgo//M)(K)$. In this case, when moreover $M'=M$, the transfer factor equals $1$.

For every element $Y\in \mgo'(K)$ that is semisimple and regular, one can form the linear combination
$$J^G_{M',M}(Y)=\sum_{X} \Delta_{M',M}(Y,X)J_M^G(X),$$
where the sum is taken over the set of $M(K)$-conjugacy classes of semisimple and $G$-regular elements $X \in \mgo(K)$. The sum has finite support.

We now assume that $s$ \emph{is not central} in $\Mc$, that is, $M'\not=M$. Let $\ec(s)$ be the finite set of connected reductive subgroups $\Hc$ of $\Gc$ which satisfy
\begin{itemize}
\item there exists an integer $i\geq 2$ and elements $s_1\in s Z_{\Mc}$ (where $Z_{\Mc}$ is the center of $\Mc$) and $s_2, \dots, s_n$ in the center of $\Mc'$ such that $\Hc$ is the neutral component of the intersection of the centralizers of the $s_i$ in $\Gc$; \item the neutral components of the centers of $\Gc$ and $\Hc$ are equal.
\end{itemize}
Every $\Hc\in\ec(s)$ contains $\Mc'$ as a Levi subgroup. Note that the set $\ec(s)$ does not contain $\Gc$. For every $s'\in \Tc$ and every subgroup $\Hc$ of $\Gc$, denote by $\Hc_{s'}$ the neutral component of the centralizer of $s'$ in $\Hc$. We say that $s'$ is elliptic in $\Hc$ if the neutral components of the centers of $\Hc_{s'}$ and $\Hc$ are equal. For all $\Hc\in \ec(s)$ and $s'\in Z_{\Mc'}$ elliptic in $\Hc$, the group $\Hc_{s}$ belongs to $\ec(s)$. By Langlands duality, the elements of $\ec(s)$ are identified with a family of connected reductive groups defined over $\Fq$ that contain $M'$ as a Levi subgroup. For every $Y\in \mgo'(K)$ semisimple and regular, one defines for every $H\in \ec(s)$ the “stable” variant $S_{M'}^H(Y)$ of the weighted orbital integrals by the following recursive formula
$$S_{M'}^H(Y)=J_{M',M'}^H(Y)-\sum_{s'\in Z_{\Mc'}/Z_{\Hc}, \ s'\not=1 } |Z_{\Hc_{s'}}/Z_{\Hc}|^{-1} S_{M'}^{H_{s'}}(Y),$$
where by convention $|Z_{\Hc_{s'}}/Z_{\Hc}|^{-1}$ is taken to be $0$ when $s'$ is not elliptic in $\Hc$. One can show that the sum has finite support.

For every $s'\in sZ_{\Mc}$ that is elliptic in $\Gc$, the group $\Gc_{s'}$ belongs to $\ec(s)$. In particular, the expression $S_{M'}^{G_{s'}}(Y)$ is well defined. One may now state the weighted fundamental lemma that we prove.

\begin{theoreme}\label{thm:1.4.1}
For every $Y\in \mgo'(K)$ semisimple and $G$-regular, we have the following identity
$$  J^G_{M',M}(Y)=\sum_{s'\in Z_{\Mc}/Z_{\Gc}} |Z_{\Gc_{s'}}/Z_{\Gc}|^{-1} S_{M'}^{G_{s'}}(Y),$$
where, as before, the sum has finite support on the elements $s'$ elliptic in $\Gc$.
\end{theoreme}

\end{paragr}

\begin{paragr}[The support theorem.]\label{S:1.5} As mentioned above, the previous theorem relies on a cohomological study of the suitably truncated Hitchin fibration. The cohomological counterpart of the previous theorem is theorem \ref{thm:10.7.3} below, whose statement we will not reproduce in this introduction. Instead, we state our support theorem, which is the key to theorem \ref{thm:10.7.3} in that it reduces the cohomological study of the truncated Hitchin fibration to that of its elliptic part. On the latter, theorem \ref{thm:10.7.3} is known by a fundamental theorem of Ng\^o (cf. \cite{Ngo-lemme}).

The situation is as follows. Let $k$ be an algebraic closure of a finite field. Let $G$ be a semisimple group defined over $k$ and $T\subset G$ a maximal torus. We assume that the characteristic of $k$ does not divide the order of the Weyl group $W$ of $(G,T)$. Let $C$ be a smooth, connected, projective curve endowed with an effective, even divisor $D$ of degree greater than the genus of $C$. Let $\infty$ be a closed point of $C$ outside the support of $D$. The Hitchin stack $\mc$ over $k$ classifies Hitchin triples $(\ec,\theta,t)$ consisting of a $G$-torsor $\ec$, a section $\theta$ of the vector bundle $\Ad(\ec)(D)$ obtained by pushing forward the torsor $\ec$ via the adjoint representation twisted by $D$, and a $G$-regular element $t\in \tgo$ such that $t$ and $\theta(\infty)$ have the same image under the canonical map
  \begin{equation}\label{eq:1.5.1}
    \ggo \to \ggo//G\simeq \tgo//W.
  \end{equation}
Let $\Ac$ be the $k$-scheme that classifies pairs $(h_a,t)$ consisting of a section $h_a$ of the bundle $\tgo//W \otimes_k \oc(D)$ and a $G$-regular element $t\in \tgo$ such that the image of $t$ under $\tgo \to \tgo//W$ equals $h_a(\infty)$. The Hitchin fibration $f : \mc \to \Ac$ is the morphism which associates to a triple $(\ec,\theta,t)$ the pair $(\chi(\theta),t)$, where $\chi_D$ is a twisted version of the morphism (\ref{eq:1.5.1}). In \cite{LFPI} we introduced a truncated open substack $\mc^\xi$ of $\mc$, consisting of “$\xi$-semistable” triples (for reminders, see section \ref{sec:4} below), which depends on a parameter $\xi$ taking values in a real vector space. When $\xi$ is “in general position” (in the sense of remark \ref{thm:4.9.2}), the restriction of the morphism $f^\xi$ to $\mc^\xi$ is proper (cf. theorem 6.2.2 of \cite{LFPI} recalled in \ref{thm:4.9.1}). We define in section \ref{sec:9} an open subset $\Ac^{\mathrm{bon}}$ of $\Ac$. Let $\Ac^{\el}$ denote the elliptic open subset of $\Ac$ and
$$j:\mathcal{A}^{\mathrm{ell}}\cap \mathcal{A}^{\mathrm{bon}} \hookrightarrow \mathcal{A}^{\mathrm{bon}}$$
be the canonical open immersion. Let $\ell$ be a prime number invertible in $k$ and $\Qlb$ an algebraic closure of $\Ql$. For every integer $i$, let ${}^{\mathrm{p}}\mathcal{H}^{i}(Rf_{\ast}^\xi \overline{\mathbb{Q}}_{\ell})$ denote the $i$th perverse cohomology sheaf of $Rf_{\ast}^\xi \overline{\mathbb{Q}}_{\ell}$. Here is our main cohomological theorem (cf. theorem \ref{thm:10.5.1} below):

\begin{theoreme}\label{thm:1.5.1}
For every integer $i$, there is a canonical isomorphism
$${}^{\mathrm{p}}\mathcal{H}^{i}(Rf_{\ast}^\xi \overline{\mathbb{Q}}_{\ell})\cong j_{!\ast}j^{\ast}{}^{\mathrm{p}}\mathcal{H}^{i}(Rf_{\ast}^\xi \overline{\mathbb{Q}}_{\ell})$$
over $\mathcal{A}^{\mathrm{bon}}$.
\end{theoreme}

It is known that the stack $\mc^\xi$ is smooth over $k$. By Deligne’s theorem, the properness of the morphism $f^\xi$ implies that the complex $Rf_{\ast}^\xi \overline{\mathbb{Q}}_{\ell}$ is pure. It follows that the perverse sheaves ${}^{\mathrm{p}}\mathcal{H}^{i}(Rf_{\ast}^\xi \overline{\mathbb{Q}}_{\ell})$ are semisimple. The preceding theorem asserts that on $\Ac^{\mathrm{bon}}$, the supports of the irreducible constituents all meet the elliptic locus.
\end{paragr}

\begin{paragr}[Description of the article.]\label{S:1.6} Let us now give a brief description of the content of the article. The first sections (sections \ref{sec:2} and \ref{sec:3}) are devoted to some notations and supplementary material. In section \ref{sec:4} we introduce the truncated Hitchin fibration. In section \ref{sec:5} we recall the action, introduced by Ng\^o, of a Picard stack on the fibers of Hitchin. As this action does not preserve the truncated fibers, in section \ref{sec:6} we introduce a substack that does respect the truncation. Section \ref{sec:7} provides some reminders on cameral curves and the sheaf of connected components of the Picard stack introduced by Ng\^o. In section \ref{sec:8} we deduce a decomposition into $s$-parts of the cohomology of the truncated Hitchin fibration which generalizes Ng\^o's decomposition on the elliptic part (cf. theorem \ref{thm:8.5.1}). In section \ref{sec:9} we introduce a certain “good” open subset of the base of the Hitchin fibration and give a sufficient condition for a point of the base to lie in this “good” open subset (theorem \ref{thm:9.1.3}). Section \ref{sec:10} contains the main cohomological results: theorem \ref{thm:10.5.1} on “support” and theorem \ref{thm:10.7.3}, which is the cohomological analogue of the weighted fundamental lemma. In section \ref{sec:12} we express the global weighted orbital $s$-integrals introduced in section \ref{sec:11} in terms of the Frobenius trace on the $s$-part of the cohomology (cf. theorem \ref{thm:12.1.1}). In section \ref{sec:15} we deduce, from the preceding considerations, a global variant of the weighted fundamental lemma (cf. theorem \ref{thm:15.1.1}). The remainder of the article is devoted to passing from this global form to the local statement.
\end{paragr}

\begin{paragr}[Acknowledgements.]\label{S:1.7} We thank Luc Illusie, Ng\^o Bao Ch\^au, Michel Raynaud and Jean-Loup Waldspurger for useful discussions.

Part of this article was written during a stay of the first-named author at the \emph{Institute for Advanced Study} in Princeton in the autumn of 2008. He wishes to thank that institute for its hospitality, as well as the \emph{National Science Foundation} for the support (agreement No.\ DMS-0635607) that made this stay possible.

Finally, we wish to thank the anonymous referees for their careful reading.
\end{paragr}

\section{Notations}\label{sec:2}
\begin{paragr}\label{S:2.1}
Let $k$ be an algebraic closure of a finite field $\FF_q$ with $q$ elements. Let $G$ be a connected reductive algebraic group over $k$. Let $G_{\der}$ denote the derived group of $G$. Let $T$ be a maximal torus of $G$. Define
$$W=W^G=W^G_T,$$
$$\Phi=\Phi^G=\Phi^G_T \quad \text{and} \quad \Phi^\vee=\Phi^{G,\vee}_T$$
to be, respectively, the Weyl group, the set of roots, and the set of coroots of $T$ in $G$. The map that sends a root to its coroot induces a bijection between $\Phi$ and $\Phi^\vee$. For any set $X$, we denote by $|X|$ its cardinality (finite or infinite). In the sequel, we assume that the order of the Weyl group $|W^G_T|$ is invertible in $k$.

The Lie algebra of a group denoted by an uppercase letter is written using the corresponding Gothic letter. Thus $\ggo$ and $\tgo$ denote the Lie algebras of $G$ and $T$, respectively. Let
$$X^*(G)=\Hom_k(G,\Gmk)$$
be the group of characters of $G$. Let
$$X_*(G)=\Hom_{\ZZ}(X^*(G),\ZZ)$$
be the dual group. In the case of a torus, this group canonically identifies with the group of cocharacters. We will also use the real vector spaces
$$\ago_G=X_*(G)\otimes_{\ZZ} \RR$$
and
$$\ago_G^*=X^*(G)\otimes_{\ZZ} \RR.$$
\end{paragr}

\begin{paragr}\label{S:2.2}
Let $\Int$ denote the action of $G$ on itself by inner automorphisms and $\Ad$ the adjoint representation. Let
$$\car=\car_G=\ggo//G$$
be the categorical quotient of $\ggo$ by the adjoint action of $G$. Let
$$ \chi=\chi_G \ :\ \ggo \to \car$$
be the canonical morphism, which is called the \emph{characteristic morphism}. Let
$$\chi^T=\chi_G^T \ :\ \tgo \to \car_G$$
be the restriction of $\chi_G$ to $\tgo$. This is a finite Galois covering of $\car_G$ with Galois group $W$. The morphism $\chi_G^T$ is $W$-invariant and induces, by Chevalley's isomorphism, an isomorphism
$$\tgo//W\simeq \car.$$
It follows that the scheme $\car$ is an affine space of dimension $n$, where $n$ is the rank of $G$, with coordinates given by $n$ homogeneous $G$-invariant polynomials on $\ggo$ whose degrees, denoted $e_{1}\leq \ldots \leq e_{n}$, are uniquely determined.
\end{paragr}

\begin{paragr}[Kostant Section.]\label{S:2.3}  An element of $\ggo$ is said to be $G$-\emph{regular} if the dimension of its centralizer in $G$ equals the rank of $G$. In the sequel, when the context is clear, we simply speak of regular elements. Let $\ggo^{\mathrm{reg}}\subset\ggo$ denote the open subset of regular elements in $\ggo$.

Let $B$ be a Borel subgroup of $G$ containing $T$, and let $\Delta\subset \Phi$ be the set of simple roots of $T$ in $B$. Complete these data to a pinning $(B,T,\{X_\al\}_{\al\in \Delta})$ of $G$. Set
$$X_{\pm}=\sum_{\al \in \pm \Delta} X_\al,$$
where we have defined $X_{-\al}$, for $\al\in \Delta$, as the unique root vector associated with the root $-\al$ satisfying $[X_{\al},X_{-\al}]=\al^\vee$. Let $\ggo_{X_+}$ denote the centralizer of $X_+$ in $\ggo$. Kostant showed that the affine space $X_- +\ggo_{X_+}$ is contained in the regular open subset $\ggo^{\mathrm{reg}}$ and that the restriction of the morphism $\chi : \ggo \to \car$ induces an isomorphism
$$X_- +\ggo_{X_+} \to \car.$$
Let
$$\eps=\eps_G  \ : \ \car \to  X_- +\ggo_{X_+}$$
be the inverse morphism.
\end{paragr}

\begin{paragr}[Parabolic Subgroups and Levi Subgroups.]\label{S:2.4}
For any subgroups $M$ and $Q$ of $G$, denote by $\fc^Q(M)$ the set of parabolic subgroups $P$ of $G$ satisfying $M\subset P\subset Q$. The elements of $\fc=\fc^G(T)$ are called the \emph{semi-standard} parabolic subgroups of $G$.

A \emph{Levi subgroup} of $G$ is defined as a Levi factor of a parabolic subgroup of $G$. Let $\lc=\lc^G(T)$ be the set of \emph{semi-standard} Levi subgroups of $G$ (in the sense that they contain $T$).

For every $P\in\fc$, let $N_P$ denote the unipotent radical of $P$ and let $M_P\in \lc$ be the unique Levi subgroup of $P$ that contains $T$. For every Levi subgroup $M\in \lc$, denote by $\pc(M)$ (respectively, $\lc(M)$) the subset of those $P\in  \fc$ for which $M_P=M$, (respectively, those $L\in \lc$ such that $M\subset L$). When it is necessary to indicate the ambient group $G$, we add an exponent $G$ to these notations.

Let $M\in \lc$. Denote by $\ago_T^M$ the subspace of $\ago_T$ spanned by the coroots of $T$ in $M$. By considering the roots, one similarly defines the subspace $(\ago_T^M)^*\subset \ago_T^*$. The canonical morphism $\ago_T \to \ago_M$ has kernel $\ago_T^M$. This kernel has, as a complement in $\ago_T$, the orthogonal of $(\ago_T^M)^*$, which identifies $\ago_M$ with a subspace of $\ago_T$. In this way, we obtain a decomposition
$$\ago_T=\ago_T^M \oplus \ago_M.$$
Define $\ago_M^G=\ago_M \cap \ago_T^G$. Then one has $\ago_M=\ago_M^G \oplus \ago_G$.

\medskip

We equip the real vector space $\ago_T$ with an inner product invariant under the natural action of the Weyl group $W^G$. The direct sum decompositions above are then orthogonal. All subspaces of $\ago_T$ are endowed with the Euclidean measure (that is, the Lebesgue measure which assigns covolume $1$ to any orthonormal basis lattice).
\end{paragr}

\begin{paragr}[Langlands Dual.]\label{S:2.5}
Let $\Tc$ be the torus over $\CC$ defined by
$$\Tc(\CC)=\Hom_{\ZZ}(X_*(T),\CC^\times).$$
Let $\Gc$ be a connected reductive group over $\CC$ together with an embedding of $\Tc$ as a maximal torus so that we have the following identification between the root data
$$(X^*(T),\Phi^{G}_T,X_*(T),\Phi^{G,\vee}_T)=(X_*(\Tc),\Phi^{\Gc,\vee}_{\Tc},X^*(\Tc),\Phi^{\Gc}_{\Tc}).$$
We say that $\Gc$ is “the” Langlands dual of $G$. Conversely, every connected reductive subgroup $\Mc$ of $\Gc$ that contains $\Tc$ determines, up to isomorphism, a connected reductive group $M$ over $k$ together with an embedding of $T$ in $M$ for which one has an identification as above between the root data of $(M,T)$ and those of $(\Mc,\Tc)$. The groups $\Mc$ and $M$ are said to be dual. Note that the Weyl group $W^M$ is then naturally a subgroup of $W^G$. Moreover, when $\Mc$ is a Levi subgroup of $\Gc$, the group $M$ equipped with the embedding of $T$ identifies with the Levi subgroup $M\in \lc^G(T)$ defined by the condition
$$\Phi^M_T=\Phi^{\Mc,\vee}_{\Tc}.$$

In what follows, by an abuse of notation, we shall not distinguish between a complex algebraic group and its group of complex points.
\end{paragr}

\begin{paragr}\label{S:2.6}
Let $M\in \lc$ and $P\in \pc(M)$. Let $\pgo$ be the Lie algebra of $P$. Let
   $$\chi_P \ :\ \pgo \to \car_P=\pgo//P$$
be the canonical morphism from $\pgo$ to its adjoint quotient. The restriction of this morphism to $\mgo$ is $M$-invariant and induces an isomorphism from $\car_M=\mgo //M$ onto $\car_P$ (see \cite{LFPI} Lemma 2.7.1). Moreover, there is a commutative diagram in which the second horizontal arrow is the projection given by the canonical decomposition $\pgo=\mgo\oplus \ngo_P$.

 \begin{equation}\label{eq:2.6.1}
    \xymatrix{ \mgo \ar[r]  \ar[d]^{\chi_M} & \pgo \ar[r]  \ar[d]^{\chi_P} & \mgo \ar[d]^{\chi_M}  \\
 \car_M \ar[r]^{\sim} &\car_P\ar[r]^{\sim} & \car_M }
\end{equation}

Let $\pgo^{\reg}=\pgo\cap \ggo^{\reg}$ be the open subset of $\pgo$ consisting of $G$-regular elements.

\begin{lemme}\label{thm:2.6.1} Let $X_1$ and $X_2$ be elements of $\pgo$ such that $\chi_P(X_1)=\chi_P(X_2)$. If $X_1$ and $X_2$ are either both semisimple or both $G$-regular, then there exists $p\in P$ that conjugates $X_1$ to $X_2$.
\end{lemme}

\begin{preuve}
Every semisimple element of $\pgo$ is clearly conjugate under $P$ to an element of $\tgo$. Hence the result is immediate when $X_1$ and $X_2$ are both semisimple.

Let $X\in \pgo$ and write its Jordan decomposition as $X=X_s+X_n$. We show that $\chi_P(X)=\chi_P(X_s)$. By conjugating $X$ by an element of $P$, we may assume that $X_s\in \mgo$. Let $X'_n$ be the projection of $X_n$ onto $\mgo$ (with respect to the decomposition $\pgo=\mgo\oplus\ngo_P$). Then the projection of $X$ onto $\mgo$ has Jordan decomposition $X_s+X'_n$. We know that $\chi_M(X_s+X'_n)=\chi_M(X_s)$. It follows that
$$\chi_P(X_s)=\chi_M(X_s)=\chi_M(X_s+X'_n)=\chi_P(X)$$
by the commutative diagram (\ref{eq:2.6.1}).

Now, let $X_1$ and $X_2$ be two $G$-regular elements in $\pgo$ such that $\chi_P(X_1)=\chi_P(X_2)$. We must show that they are conjugate under $P$. From the above, the semisimple parts in the Jordan decompositions of $X_1$ and $X_2$ are conjugate under $P$. By conjugating $X_2$ by an element of $P$, we may assume that their semisimple parts are equal to some $Y\in \pgo$.

Replacing $G$ by $G_Y$, the centralizer of $Y$ in $G$ (which is known to be a connected reductive group), and $P$ by the parabolic subgroup $G_Y\cap P$, we reduce to the case where $X_1$ and $X_2$ are nilpotent and $G$-regular. It is known that such elements are conjugate by an element $g\in G$. We now show that necessarily $g\in P$, which will complete the proof. Let $B_1$ and $B_2$ be Borel subgroups such that, for $i=1,2$, one has $X_i\in \bgo_i$ and $B_i\subset P$. These subgroups are uniquely determined. In particular, since $\Ad(g)X_1=X_2$, we must have $gB_1g^{-1}=B_2$. Thus, the parabolic subgroups $gPg^{-1}$ and $P$, which are conjugate and contain the same Borel subgroup $B_2$, are equal. It follows that $g\in P$, as required.

\end{preuve}
\end{paragr}

\begin{paragr}[Factorization of the Discriminant.]\label{S:2.7}
Let $\Mc$ be a connected reductive subgroup of $\Gc$ that contains $\Tc$. Let $M$ be a connected reductive group over $k$ equipped with an embedding of $T$ and dual to $\Mc$. We then have inclusions $\Phi^M\subset \Phi^G$ and $W^M\subset W^G$. Hence, there is a morphism
  $$\chi_G^M \ : \ \car_M \to \car_G$$
which satisfies $\chi_G^M\circ \chi_M^T=\chi_G^T$. Moreover, when $M\in \lc^G(T)$, the morphism $\chi_G^M\circ \chi_M$ is the restriction of $\chi_G$ to $\mgo$.

Let $B\in \pc(T)$ be a Borel subgroup and let $\Phi^B$ denote the roots of $T$ in $B$. For any root $\al\in \Phi$, let $d\al\in k[\tgo]$ be its derivative. The elements of $k[\tgo]$ defined by
$$D^M=\prod_{\al\in \Phi^M} d\al$$
and
$$R_M^B=\prod_{\al\in \Phi^B-\Phi^M} d\al$$
are $W^M$-invariant: hence they descend to regular functions on $\car_M$. This is clear for the first one, which is nothing other than the Weyl discriminant. For the second, this follows from the formula
$$w(R_T^B)=(-1)^{\ell(w)} R_T^B$$
for all $w\in W$, where $\ell(w)$ denotes the length of $w$ (see \cite{Bourbaki}, Chapter V, §5.4 Proposition 5) and from the formula
$$R_T^B = R_M^B R_{T}^{M\cap B}.$$
Note that $R_M^B$ depends on the choice of $B$ only up to a sign. Let $\car^{M\textrm{-}\reg}_M$ be the open subset where $D^M$ does not vanish. This is precisely the locus where the morphism $\chi^T_M$ is étale. We shall not confuse it with the open subset $\car^{G\textrm{-}\reg}_M$, where the discriminant for $G$
$$D^G=D^M (R_M^B)^2$$
does not vanish. Only the latter will appear in what follows, and we set
$$\car_M^{\reg}=\car_M^{G\textrm{-}\reg}.$$
Note that the fiber of $\chi_M$ above $\car_M^{\reg}$ is the open subset of $\mgo$ consisting of semisimple $G$-regular elements. The open subset $\tgo^{\reg}$ is the set of $G$-regular elements.
\end{paragr}

\section{Complements on Centralizers}\label{sec:3}

\begin{paragr}[Centralizers.]\label{S:3.1}
For every parabolic subgroup $P\in \pc^G(T)$, let $I^P$ be the group scheme of centralizers on $\pgo$ defined by
$$I^P=\{(p,X)\in P \times_k \pgo \mid \Ad(p)X=X\}.$$
We have the following lemma.
\end{paragr}

\begin{lemma}\label{thm:3.1.1}
The canonical morphism $I^P \to I^G$ induces an isomorphism
\begin{equation}\label{eq:3.1.1}
  I^P_{|\pgo^{\reg}} \buildrel\sim\over\longrightarrow  I^G_{|\pgo^{\reg}}
\end{equation}
over the open set
$$\pgo^{\reg}=\pgo\cap \ggo^{\reg}.$$
Moreover, the group schemes $I^P_{|\pgo^{\reg}}$ and $I^G_{|\pgo^{\reg}}$ are commutative and smooth over $\pgo^{\reg}$.
\end{lemma}

\begin{proof}
The morphism (\ref{eq:3.1.1}) is an isomorphism if and only if, for every $X\in \pgo^{\reg}$, the centralizer $G_X$ of $X$ in $G$ is contained in $P$. By using the Jordan decomposition $X=X_s+X_n$, we have $G_X=G_{X_s}\cap G_{X_n}$ and $G_{X_s}$ is a connected reductive group since $X_s$ is semisimple. Since $P\cap G_{X_s}$ is a parabolic subgroup of $G_{X_s}$ and $X_n$ is $G_{X_s}$-regular, we reduce to the case where $X$ is nilpotent, which we now assume. Let $B\subset P$ be a Borel subgroup whose Lie algebra contains $X$. By the $G$-regularity of $X$, such a Borel subgroup is unique, so $G_{X}$ is contained in the normalizer of $B$ in $G$, which is $B$, hence contained in $P$.

It is well known that $I^G_{|\ggo^{\reg}}$ is a commutative group scheme, smooth over $\ggo^{\reg}$. By base change, the same holds for $I^G_{|\pgo^{\reg}}$ and therefore for $I^P_{|\pgo^{\reg}}$.
\end{proof}

\begin{paragr}[Regular Centralizers.]\label{S:3.2}
Let $M\in \lc^G(T)$ be a semi-standard Levi subgroup. Let $\eps_M$ be a Kostant section of the characteristic morphism $\chi_M  : \mgo \to\car_M$ (see § \ref{S:2.3}). This section takes values in the open set of $M$-regular elements. Let
$$J^M=\eps_M^*I^M,$$
which we call the group scheme of regular centralizers over $\car_M$. It is a commutative group scheme, smooth over $\car_M$ (cf. Lemma \ref{thm:3.1.1}).

Following Ng\^o, we introduce the isomorphism
\begin{equation}\label{eq:3.2.1}
   (\chi_G^*J^G)_{|\ggo^{\reg}} \to I^G_{\ggo^{\reg}}
\end{equation}
which, to a pair $(g,X)$ in $G\times_k \ggo^{\reg}$ such that $\Ad(g)\eps_G(\chi_G(X))=\eps_G(\chi_G(X))$, associates $(hgh^{-1},X)$, where $h$ is an element of $G$ such that $\Ad(h)\eps_G(\chi_G(X))=X$. Since both $X$ and $\eps_G(\chi_G(X))$ are $G$-regular and have the same image in $\car_G$, such an $h$ exists and is well defined up to left translation by an element of the centralizer $G_X$ of $X$. By the $G$-regularity of $X$, we know that $G_X$ is commutative. Therefore, the morphism (\ref{eq:3.2.1}) is well defined and is clearly an isomorphism. Here is a proposition that slightly generalizes Proposition 3.2 of \cite{Ngo-endoscopie}.
\end{paragr}

\begin{proposition}\label{thm:3.2.1}
For every $P\in \fc^G(T)$, there exists a unique morphism of group schemes
$$(\chi_G^* J^G)_{|\pgo} \to I^P$$
which extends the isomorphism
$$(\chi_G^*J^G)_{|\pgo^{\reg}} \to I^G_{|\pgo^{\reg}}=I^P_{|\pgo^{\reg}}$$
obtained by composing the restriction of the morphism (\ref{eq:3.2.1}) to $\pgo^{\reg}$ with the isomorphism of Lemma \ref{thm:3.1.1}.

Moreover, the diagram below is commutative
\begin{equation}\label{eq:3.2.2}
 \xymatrix{    (\chi_G^* J^G)_{|\pgo}   \ar[r]     \ar[rd] &   I^P  \ar[d] \\ & I^G_{|\pgo}}
\end{equation}
where the vertical arrow is the canonical injection and the diagonal arrow is the restriction to $\pgo$ of the morphism $\chi_G^* J^G\to I^G$.
\end{proposition}

\begin{proof}
Since the scheme $(\chi_G^* J^G)_{|\pgo}$ is normal, the complement $\pgo-\pgo^{\reg}$ has codimension at least $2$ in $\pgo$, and $I^P$ is affine, the morphism
$$(\chi_G^*J^G)_{|\pgo^{\reg}} \to I^G_{|\pgo^{\reg}}=I^P_{|\pgo^{\reg}}$$
extends uniquely to a morphism of group schemes over $\pgo$. The same argument shows that, since the triangle (\ref{eq:3.2.2}) is commutative on $\pgo^{\reg}$, it is also commutative on $\pgo$.
\end{proof}

\begin{paragr}\label{S:3.3}
In this paragraph, we prove the following proposition.

\begin{proposition}\label{thm:3.3.1}
Let $M\in \lc^{G}(T)$ be a Levi subgroup. There exists a unique morphism of group schemes over $\car_M$
\begin{equation}\label{eq:3.3.1}
  (\chi_G^M)^* J^G \to X_*(M)\otimes_{\ZZ} \mathbb{G}_{m,\car_M}
\end{equation}
such that for every $P\in \pc^G(M)$ the following diagram commutes
\begin{equation}\label{eq:3.3.2}
 \xymatrix{   \chi_P^*(\chi_G^M)^* J^G=(\chi_G^*J^G)_{|\pgo}   \ar[r]     \ar[rd] &   I^P  \ar[d] \\  & X_*(M)\otimes_{\ZZ} \mathbb{G}_{m,\pgo}}
\end{equation}
where
\begin{itemize}
\item the diagonal arrow is deduced from the morphism (\ref{eq:3.3.1}) by base change
$$\chi_P : \pgo\to \car_M \, ;$$
\item the horizontal arrow is that defined in Proposition \ref{thm:3.2.1}; \item the vertical arrow is induced by the canonical morphism
$$P \to X_*(P)\otimes_{\ZZ}\Gmk =X_*(M)\otimes_{\ZZ}\Gmk.$$
\end{itemize}
\end{proposition}

\begin{remark}\label{thm:3.3.2}
The $\ZZ$-module $X_*(M)$ is free of finite rank equal to the dimension of the center of $M$. Consequently, $X_*(M)\otimes_{\ZZ}\Gmk$ is a torus over $k$ of the same rank.
\end{remark}

\begin{proof}
Consider the two projections $p_1$ and $p_2$
$$ \pgo\times_{\car_M} \pgo \longrightarrow \pgo.$$
Let $p=\chi_P\circ p_1=\chi_P\circ p_2$. By faithfully flat descent via the morphism $\chi_P$, it suffices to show that the morphism
\begin{equation}\label{eq:3.3.3}
  \chi_P^*(\chi_G^M)^* J^G\to X_*(M)\otimes_{\ZZ} \mathbb{G}_{m,\pgo},
\end{equation}
which makes diagram (\ref{eq:3.3.2}) commute, belongs to the kernel of
$$\Hom_{\pgo}\Bigl( (\chi_G^* J^G)_{|\pgo},  X_*(M)\otimes_{\ZZ} \mathbb{G}_{m,\pgo}\Bigr)  \longrightarrow  \Hom_{\pgo\times_{\car_M} \pgo} \Bigl( p^* (\chi_G^M)^* J^G,  X_*(M)\otimes_{\ZZ} \mathbb{G}_{m,\pgo\times_{\car_M} \pgo}\Bigr).$$
It suffices to verify this over the open set $\pgo^{\reg}\times_{\car_M}\pgo^{\reg}$. Now, for $i=1,2$, the image of the morphism (\ref{eq:3.3.3}) under $p_i^*$ is given on $\pgo^{\reg}\times_{\car_M}\pgo^{\reg}$ by the map which, to $(g,X_1,X_2) \in G\times_k \pgo^{\reg}\times_{\car_M}\pgo^{\reg}$ satisfying $\Ad(g)\eps_G(\chi_G(X_i))=\eps_G(\chi_G(X_i))$, associates
$$\lambda\in X^*(P) \mapsto  \lambda\Bigl(h_i g h_i^{-1}\Bigr),$$
where $h_i$ conjugates $\eps_G(\chi_G(X_i))$ to $X_i$. Note that $h_i g h_i^{-1}$ belongs to the centralizer $G_{X_i}$ of $X_i$ in $G$, hence lies in $P$ (recall the inclusion $G_{X_i}\subset P$, cf. Lemma \ref{thm:3.1.1}). But we know that $X_1$ and $X_2$ are conjugate in $P$ (cf. Lemma \ref{thm:2.6.1}). It follows that $h_2\in P\,G_{X_1}\,h_1 \subset P\,h_1$. Therefore, for every $\lambda\in X^*(P)$,
$$ \lambda\Bigl(h_1 g h_1^{-1}\Bigr)=\lambda\Bigl(h_2 g h_2^{-1}\Bigr),$$
which is the desired equality.

Thus, we have shown the existence of the morphism (\ref{eq:3.3.1}) which makes diagram (\ref{eq:3.3.2}) commute for a parabolic subgroup $P\in\pc(M)$. We now verify that the morphism (\ref{eq:3.3.1}) we have just defined does not depend on the choice of $P$. Over the dense open set $\car_M^{\reg}$ where $D^G$ does not vanish, the morphism (\ref{eq:3.3.1}) is the map which, to a pair $(g,a)\in G\times_k \car_M$ satisfying $\Ad(g)\eps_G(\chi_G^M(a))=\eps_G(\chi_G^M(a))$, associates
$$\lambda\in X^*(P) \mapsto \lambda\Bigl(hgh^{-1}\Bigr),$$
where $h\in G$ is such that $\Ad(h)\eps_G(\chi_G^M(a))$ lies in the fiber of $\chi_P$ over $a$. Such an element $h$ exists: let $Y\in \mgo$ be such that $\chi_P(Y)=a$. Then $Y$ is semisimple and $G$-regular, hence $G$-conjugate to $\eps_G(\chi_G^M(a))$. Such an $h$ is well defined up to right multiplication by an element of $P$, since the fiber of $\chi_P$ over $a$ is a $P$-conjugacy class (cf. Lemma \ref{thm:2.6.1}). In particular, $\lambda\Bigl(hgh^{-1}\Bigr)$ is well defined for every $\lambda\in X^*(P)$ (by an argument already mentioned, we have $hgh^{-1}\in P$). Clearly, one could equally define $h$ as an element of $G$ such that $\Ad(h)\eps_G(\chi_G^M(a))$ lies in the fiber of $\chi_M$ over $a$, hence the independence of $P$.
\end{proof}
\end{paragr}

\begin{paragr}[Galois Description of $J^G$.]\label{S:3.4} We leave it to the reader to verify the following lemma.

\begin{lemma}\label{thm:3.4.1}
The morphism
$$(\chi_G^T)^*J^G \to X_*(T)\otimes_\ZZ \mathbb{G}_{m,\tgo}$$
from Proposition \ref{thm:3.3.1} is $W$-equivariant, where the action on the target is the diagonal action of $W$.
\end{lemma}

For every $\tgo$-scheme $X$, let $(\chi_G^T)_*X$ be the $\car_G$-scheme obtained from $X$ by Weil restriction of scalars (recall that the $k[\car_G]$-module $k[\tgo]$ is free of finite rank, equal to the order of $W$). The adjunction morphism $J^G\to (\chi_G^T)_*(\chi_G^T)^*J^G$ factors through the morphism
$$J^G \longrightarrow \bigl((\chi_G^T)_*(\chi_G^T)^*J^G\bigr)^W,$$
where the exponent $W$ denotes the subscheme of fixed points under $W$. Using the $W$-equivariance of the morphism $(\chi_G^T)^*J^G \to X_*(T)\otimes_\ZZ \mathbb{G}_{m,\tgo}$ (cf. Lemma \ref{thm:3.4.1}), one obtains a morphism
\begin{equation}\label{eq:3.4.1}
  J^G \longrightarrow \bigl(X_*(T)\otimes_\ZZ (\chi_G^T)_*\mathbb{G}_{m,\tgo}\bigr)^W,
\end{equation}
whose target is the closed subscheme of fixed points under $W$ in $X_*(T)\otimes_\ZZ \chi^T_*\mathbb{G}_{m,\tgo}$; it is moreover a smooth commutative group scheme over $\car_G$ (cf.\ \cite[Lemma~2.4.1]{Ngo-lemme}) whose Lie algebra, viewed as an $\oc_{\car_G}$-module, is
$$\Lie\bigl((X_*(T)\otimes_\ZZ (\chi_G^T)_*\mathbb{G}_{m,\tgo})^W\bigr)=\bigl(X_*(T)\otimes_\ZZ (\chi_G^T)_*\oc_{\tgo}\bigr)^W.$$
Let us recall the following proposition, due to Ng\^o, which provides a Galois description of the Lie algebra of $J^G$.

\begin{proposition}\label{thm:3.4.2}
The morphism
$$J^G \longrightarrow \bigl(X_*(T)\otimes_\ZZ (\chi_G^T)_*\mathbb{G}_{m,\tgo}\bigr)^W$$
is an isomorphism over the open subset $\car_G^{\reg}$. It also induces an isomorphism between the open subschemes of the neutral components of the fibers. In particular, this morphism induces an isomorphism between the Lie algebras viewed as $\oc_{\car_G}$-modules
$$\Lie(J^G)=\bigl(X_*(T)\otimes_\ZZ (\chi_G^T)_*\oc_{\tgo}\bigr)^W.$$
\end{proposition}

\begin{preuve}
We refer the reader to \cite[Prop.~2.4.2 and Cor.~2.4.8]{Ngo-lemme}.
\end{preuve}
\end{paragr}

\section{The Truncated Hitchin Fibration}\label{sec:4}

\begin{paragr}\label{S:4.1}
Let $C$ be a projective, smooth, and connected curve over $k$ of genus $g$. Let $D=2D'$ be an effective divisor of degree $>2g$ and let $\infty$ be a point in $C(k)$ that is not in the support of $D$. Let $\lc_D$ be the $\Gmk$-torsor on $C$ associated to $D$, equipped with its canonical trivialization on the complement of the support of $D$. For any $k$-scheme $V$ with a $\Gmk$-action, define
$$V_D=\lc_D\times_k^{\Gmk} V$$
to be the contracted product; this is a scheme over $C$ which is a locally trivial fibration with fiber type $V$. Let $M\in \lc(T)$. The homothety action of the multiplicative group $\Gmk$ on $\tgo$ induces an action of $\Gmk$ on $\car_M$, for which the characteristic morphism $\chi_M^T$ is $\Gmk$-equivariant. By forming the contracted product with $\lc_D$, we obtain $C$-morphisms
$$\chi_{M,D}^T : \tgo_D \to \car_{M,D}$$
and, for every $P\in \pc(M)$,
$$\chi_{P,D} : \pgo_D  \to \car_{M,D}.$$
\end{paragr}

\begin{paragr}[Hitchin Stack.]\label{S:4.2} This is the algebraic stack denoted $\mc=\mc_G$ which classifies Hitchin triples $m=(\mathcal{E},\theta,t)$ where:
\begin{itemize}
\item $\mathcal{E}$ is a $G$-torsor on $C$; \item $\theta \in H^{0}(C,\Ad_D(\ec))$; \item $t\in \mathfrak{t}^{G\textrm{-reg}}$ with $\chi(t)=\chi(\theta_{\infty})$ in $\mathfrak{car}_G$.
\end{itemize}
We denote
$$\Ad_D(\ec)=\ec\times_k^{G,\Ad_D}\ggo_D$$
the vector bundle on $C$ obtained by pushing forward the $G$-torsor $\ec$ via the morphism
$$\Ad_D \ : \ G\to \Aut(\ggo_D)$$
induced by the adjoint action.

Let $[\ggo_D/G]$ be the stack over $C$ defined as the quotient of $\ggo_D$ by $G\times_k C$. By definition, for every $C$-test scheme $S$, the groupoid $[\ggo_D/G](S)$ is the groupoid of pairs $(\ec,\theta)$ where $\ec$ is a $G$-torsor over $S$ and $\theta$ is a section over $S$ of $\Ad_D(\ec)=\ec\times_k^{G,\Ad_D}\ggo_D$.

Since the morphism $\chi_G$ is both $\Gmk$-equivariant and $G$-invariant, it induces a morphism of $C$-stacks
\begin{equation}\label{eq:4.2.1}
   [\chi_{G,D}] \ :\ [\ggo_D/G] \to \car_{G,D}.
\end{equation}

One may then give an equivalent definition of $\mc$. Let $S$ be a $k$-test scheme. The data of a triple $m=(\ec,\theta,t)\in \mc(S)$ is equivalent to the data of a pair $m=(h_m,t)$ consisting of a $C$-morphism
$$h_m  \ : \ C\times_k S \longrightarrow [\ggo_D/G]$$
and a point $t\in \tgo^{\reg}(S)$ such that
$$ ([\chi_{G,D}]\circ h_m)(\infty_S)=\chi^T_G(t),$$
where $\infty_S$ is the base change of $\infty$ via $S\to k$.
\end{paragr}

\begin{paragr}[Smoothness of $\mc$ over $k$.]\label{S:4.3} We have the following proposition due to Biswas and Ramamanan.

\begin{proposition}\label{thm:4.3.1} (cf. \cite{Biswas-Ramanan} and \cite{Ngo-lemme} Theorem 4.14.1)
The algebraic stack $\mathcal{M}$ is smooth over $k$.
\end{proposition}
\end{paragr}

\begin{paragr}[Characteristic Space.]\label{S:4.4}
In the (second) definition of $\mc$, if one replaces the quotient stack $[\ggo_D/G]$ by the categorical quotient of $\ggo_D$ by $G$, which is none other than $\car_{G,D}$, one obtains a quasi-projective scheme, smooth over $k$, denoted $\Ac=\Ac_G$ and called the \emph{characteristic space}, which classifies pairs $a=(h_a,t)$ where
\begin{itemize}
\item $h_a\in H^0(C,\car_{G,D})$; \item $t\in \mathfrak{t}^{G\textrm{-reg}}$ with $\chi(t)=h_a(\infty)$.
\end{itemize}

For each $M\in \mathcal{L}(T)$, there is the base $\mathcal{A}_{M}$ of the Hitchin fibration for $M$ and an open subset
$$
\mathcal{A}_{M}^{G\textrm{-reg}}\subset \mathcal{A}_{M},
$$
where we require that $t\in \mathfrak{t}^{G\textrm{-reg}}\subset \mathfrak{t}^{M\textrm{-reg}}$. The morphism
\begin{equation}\label{eq:4.4.1}
  \chi^M_{G} : \mathcal{A}_{M}^{G\textrm{-reg}}\hookrightarrow \mathcal{A}
\end{equation}
which sends $(h,t)$ to $(\chi_{G,D}^M \circ h,t)$ is a closed immersion (cf. \cite{LFPI} Proposition 3.5.1).

\textit{Henceforth, only the open subset $\mathcal{A}_{M}^{G\textrm{-reg}}$ will play a role, which we shall denote simply by $\mathcal{A}_{M}$.}

\begin{definition}\label{thm:4.4.1}
The elliptic part $\mathcal{A}^{\mathrm{ell}}=\mathcal{A}_{G}^{\mathrm{ell}}$ of $\mathcal{A}$ is defined as the complement of the union of the $\mathcal{A}_{M}$ for $M\in \mathcal{L}(T)$ with $M\neq G$.
\end{definition}

Note that the elliptic open subset is nonempty for dimensional reasons. Similarly, one defines $\mathcal{A}_{M}^{\mathrm{ell}}$. We have the following proposition (cf. \cite{LFPI} Proposition 3.6.2).

\begin{proposition}\label{thm:4.4.2}
The scheme $\Ac_G$ is the disjoint union of the locally closed subschemes $\Ac_{M}^{\el}$ for $M\in \lc(T)$:
$$\Ac_G=\bigcup_{M\in \lc^G(T)} \Ac_{M}^{\el}.$$
\end{proposition}
\end{paragr}

\begin{paragr}[Hitchin Fibration.]\label{S:4.5} We now define the Hitchin fibration using the morphism $[\chi_{G,D}]$ in (\ref{eq:4.2.1}).

\begin{definition}\label{thm:4.5.1}
The Hitchin fibration is the morphism
$$
f=f_{G}:\mathcal{M}\rightarrow \mathcal{A}
$$
which sends $m=(h_m,t)$ to $a=([\chi_{G,D}]\circ h_m,t)$.
\end{definition}

Let $a=(h_a,t)\in \Ac$. The Hitchin fiber $\mc_a=f^{-1}(a)$ classifies the sections $h_m$ of the stack $[\ggo_D/G]$ making the diagram below commutative
\begin{equation}\label{eq:4.5.1}
  \xymatrix{ C \ar[r]^{h_m}  \ar[rd]_{h_a} & [\ggo_D/G] \ar[d]^{[\chi_{G,D}]}  \\ &\car_{G,D} }
\end{equation}
\end{paragr}

\begin{paragr}[Reduction of a Hitchin Triple.]\label{S:4.6} Let $m=(h_m,t)\in \mc_G$ and let $P\in \pc(M)$ be a parabolic subgroup of $G$ with Levi factor $M\in \lc(T)$. A reduction of $m$ to $P$ is, by definition, the data of a section
$$h_{P} :   C \to [\pgo_D/P]$$
of the $C$-stack $[\pgo_D/P]$ (the quotient of $\pgo_D$ by $P\times_k C$) satisfying the following two conditions:
\begin{enumerate}
\item We have $h_m=i_P\circ h_P$, where
$$i_P : [\pgo_D/P] \to [\ggo_D/G]$$
is the canonical morphism; \item The pair $(h_{a_P},t)$ defined by
$$h_{a_P}=[\chi_{P,D}]\circ h_P,$$
belongs to $\Ac_M$, that is,
$$h_{a_P}(\infty)=\chi_P^T(t).$$
\end{enumerate}

Thus, if $h_P$ is a reduction of $m$ to $P$, the following diagram commutes:
\begin{equation}\label{eq:4.6.1}
\xymatrix{      C \ar[r]^{h_P}   \ar@/^3pc/[rr]_{h_m}  \ar@/_4pc/[rrd]_{h_a} \ar[rd]_{h_{a_P}} & [\pgo_D/P] \ar[d]^{[\chi_{P,D}]} \ar[r]^{i_P} &[\ggo_D/G] \ar[d]^{[\chi_{G,D}]}  \\  &\car_{P,D} \ar[r]_{\chi^M_{G,D}}     &\car_{G,D} }.
\end{equation}

It is clear from the above diagram that for $m$ to admit a reduction to $P$, it is necessary that $f(m)$ belongs to $\Ac_M$. In fact, this condition is sufficient. More precisely, we have the following proposition (cf. Proposition 4.10.1 of \cite{LFPI}).

\begin{proposition}\label{thm:4.6.1}
Let $m\in \mc_G$ and $M\in \lc(T)$ such that $f(m)\in \Ac_{M}^{\el}$. Then for every $P\in \fc(M)$, there exists a unique reduction $h_P$ of $m$ to $P$.
\end{proposition}
\end{paragr}

\begin{paragr}[Degree $\deg_P$.]\label{S:4.7}
For any group scheme $H$ over $C$, let $\mathrm{B}(H)$ be the $C$-classifying stack of $H$, so that for every $C$-test scheme $S$, the groupoid $\mathrm{B}(H)(S)$ is the groupoid of $H\times_C S$-torsors over $S$. This stack is also the quotient of $C$ by the trivial action of $H$.

The structural morphism $\pgo_D \to C$ induces a morphism of $C$-stacks
\begin{equation}\label{eq:4.7.1}
  [\pgo_D/P] \longrightarrow \mathrm{B}(P\times_k C).
\end{equation}
Moreover, the natural morphism $P\to X_*(P)\otimes_\ZZ \Gmk$ induces a morphism of $C$-stacks
\begin{equation}\label{eq:4.7.2}
\mathrm{B}(P\times_k C) \longrightarrow \mathrm{B}\Bigl(X_*(P)\otimes_\ZZ \mathbb{G}_{m,C}\Bigr).
\end{equation}
There is a morphism
$$\mathrm{B}(\mathbb{G}_{m,C})(C)  \longrightarrow \ZZ$$
given by the degree of a $\mathbb{G}_{m,C}$-torsor. This immediately yields a morphism
\begin{equation}\label{eq:4.7.3}
   \mathrm{B}\Bigl(X_*(P)\otimes_\ZZ \mathbb{G}_{m,C}\Bigr)(C)  \longrightarrow X_*(P).
\end{equation}
By composing (\ref{eq:4.7.3}) with (\ref{eq:4.7.1}) and (\ref{eq:4.7.2}), we obtain a degree morphism
\begin{equation}\label{eq:4.7.4}
  \deg_P : [\pgo_D/P](C)  \longrightarrow X_*(P).
\end{equation}
\end{paragr}

\begin{paragr}[Convex Set Associated to $m\in\mc$.]\label{S:4.8}
Let $M\in \lc(T)$ and $m\in \mc$ such that $f(m)\in \Ac_M^{\el}$. For every $P\in \pc(M)$, let $h_P$ be the unique reduction of $m$ to $P$ (cf. Proposition \ref{thm:4.6.1}). In (\ref{eq:4.7.4}) we defined the degree $\deg_P(h_P)$ of the section $h_P$ of $[\pgo_D/P]$.

Define $\mathcal{C}_m\subset \ago_T$ to be the convex set obtained by translating the convex hull of the points
$$-\deg_P(h_P)\in X_*(P)=X_*(M) \subset \ago_M,$$
for $P\in \pc(M)$, by the vectors in the subspace $\ago_T^M\oplus \ago_G$.
\end{paragr}

\begin{paragr}[$\xi$-Stability.]\label{S:4.9}
Let $\xi\in \ago_T$ and $m\in \mc$. We say that $m$ is $\xi$-semi-stable if the convex set $\mathcal{C}_m$ from paragraph \ref{S:4.8} contains $\xi$. Let $\mc^\xi$ denote the substack of $\mc$ consisting of $\xi$-stable points and define
$$\mc^{\el}=\mc\times_{\Ac_G} \Ac_G^{\el},$$
the elliptic open subset of $\mc$. Denote by
$$f^\xi : \mc^\xi \to \Ac_G$$
and
$$f^{\el} : \mc^{\el}\to\Ac_G^{\el}$$
the restrictions of the Hitchin morphism $f$ to $\mc^\xi$ and $\mc^{\el}$ respectively.

Here is the main result of \cite{LFPI} (see Proposition 6.1.4 and Theorem 6.2.2).

\begin{theorem}\label{thm:4.9.1}
The stack $\mc^\xi$ is an open substack of $\mc$, hence smooth over $k$, and it contains the elliptic open subset $\mc^{\el}$.

If $G$ is semisimple and if $\xi$ is in \emph{general position}, the stack $\mc^\xi$ is a Deligne--Mumford stack and the morphism $f^\xi$ is proper.
\end{theorem}

\begin{remark}\label{thm:4.9.2}
The phrase $\xi$ \emph{in general position} means that $\xi$ does not belong to certain rational hyperplanes in $\mathfrak{a}_T$ (cf. \cite{LFPI} Definition 6.1.3). In particular, it implies that for a Hitchin triple, the notions of $\xi$-semi-stability and $\xi$-stability are equivalent.
\end{remark}
\end{paragr}

\section{Action of the Picard Stack \texorpdfstring{$\Jc_a$}{J-a}}\label{sec:5}

\begin{paragr}[Action of $\Gmk$ on the Centralizers.]\label{S:5.1}
In §\ref{S:3.2} we defined the group scheme $J^G$ of regular centralizers over $\ggo$. Let $(g,a)\in J^G$ and $x\in \Gmk$. The elements of $\ggo^{\reg}$ defined by $\eps_G(x\cdot a)$ and $x\cdot \eps_G(a)$ (where $\cdot$ denotes, respectively, the action of $\Gmk$ on $\car_G$ defined in §\ref{S:4.1} and the scaling action of $\Gmk$ on $\ggo$) are $G$-conjugate since they have the same image via $\chi_G$. Therefore, there exists an element $h_x\in G$ such that
$$\eps_G(x\cdot a)=\Ad(h_x)(x\cdot \eps_G(a)).$$
One may check that $h_x g h_x^{-1}$ centralizes $\eps_G(x\cdot a)$ and that this element is independent of the choice of $h_x$ (since the regular centralizers are commutative). We then define
$$x\cdot (g,a)= (h_x g h_x^{-1},\, x\cdot a),$$
which defines an action of $\Gmk$ on $J^G$ for which the canonical morphism $J^G \to \car_G$ is $\Gmk$–equivariant. By the construction of §\ref{S:4.1}, one deduces a group scheme $J^G_D$ over $\car_{G,D}$.

For every parabolic subgroup $P\in \fc(T)$, the group scheme $I^P$ over $\pgo$ of centralizers is endowed with an action of $\Gmk$ defined for all $x\in \Gmk$ and $(p,X)\in I^P$ by
$$x\cdot(p,X)=(p,\, x\cdot X),$$
where $\cdot$ denotes the scaling action of $\Gmk$ on $\pgo$. Again, the construction of §\ref{S:4.1} yields a group scheme $I^P_D$ over $\pgo_{D}$.

It follows from Proposition \ref{thm:3.2.1} that the morphism
$$(\chi_G^* J^G)_{|\pgo} \to I^P$$
defined there is $\Gmk$–equivariant. Hence one obtains a morphism of group schemes
\begin{equation}\label{eq:5.1.1}
  (\chi_{G,D}^* J^G_D)_{|\pgo_D} \to I^P_D
\end{equation}
over $\pgo_D$.

Let $M\in \lc(T)$ and let $\chi^M_{G,D}: \car_{M,D}\to \car_{G,D}$ be the morphism deduced from the $\Gmk$–equivariant morphism $\chi^M_G : \car_M \to \car_G$. Proposition \ref{thm:3.3.1} shows that there exists a unique morphism of group schemes
\begin{equation}\label{eq:5.1.2}
(\chi^M_{G,D})^*J_D^G\to X_*(M)\otimes_\ZZ \mathbb{G}_{m,\car_{M,D}}
\end{equation}
which, for every $P\in \pc(M)$, is the composition of (\ref{eq:5.1.1}) with the obvious morphism
$$I^P_D \to X_*(M)\otimes_\ZZ \mathbb{G}_{m,\car_{M,D}}.$$
\end{paragr}

\begin{paragr}[The Picard Stack $\Jc_a$.]\label{S:5.2}
Let $a=(h_a,t)\in\Ac$. Define
$$J_a=h_a^* J^G_D.$$
This is a smooth, commutative group scheme over $C$. Let $\Jc_a$ denote the $k$–stack of $J_a$–torsors over $C$. More precisely, for any test $k$–scheme $S$, the groupoid $\Jc_a(S)$ is the groupoid of $C$–morphisms
$$C\times_k S \to \mathrm{B}(J_a).$$

Since the group scheme $J_a$ is commutative, the contracted product of two $J_a$–torsors is again a $J_a$–torsor. This product endows $\Jc_a$ with the structure of a Picard stack over $\Spec(k)$.
\end{paragr}

\begin{paragr}[Degree Morphism.]\label{S:5.3}
Let $M\in \lc(T)$ and $a_M=(h_{a_M},t)\in \Ac_M$. Let $a$ be the image of $a_M$ in $\Ac_G$ via the morphism $\chi^M_G$, cf. line (\ref{eq:4.4.1}). Thus we have $h_a=\chi^M_{G,D}\circ  h_{a_M}$ and
$$J_a=h_{a_M}^* (\chi^M_{G,D})^* J^G_D.$$
By the base change $h_{a_M} : C \to \car_{M,D}$, one deduces from the morphism (\ref{eq:5.1.2}) a morphism of group schemes
$$J_a \to X_*(M)\otimes_\ZZ \mathbb{G}_{m,C},$$
and a morphism of $C$–Picard stacks between the classifying stacks
\begin{equation}\label{eq:5.3.1}
  \mathrm{B}(J_a) \to \mathrm{B}( X_*(M)\otimes_\ZZ \mathbb{G}_{m,C}).
\end{equation}
Since we have a degree morphism, cf. (\ref{eq:4.7.3}),
$$\mathrm{B}( X_*(M)\otimes_\ZZ \mathbb{G}_{m,C})(C) \to X_*(M),$$
one obtains, by composition, a degree morphism
\begin{equation}\label{eq:5.3.2}
  \deg_M : \Jc_a \to X_*(M)
\end{equation}
for every $a\in \Ac_M$, which is a homomorphism of Picard stacks (if one views $X_*(M)$ as the constant $k$–Picard stack).

\begin{proposition}\label{thm:5.3.1}
Let $M\subset L$ be two Levi subgroups in $\lc^G(T)$. Let
$$p^M_L : X_*(M)\to X_*(L)$$
be the dual projection of the restriction morphism
$$X^*(L)\to X^*(M).$$
Let $a\in \Ac_M$. Then one has the following equality of morphisms from $\Jc_a$ to $X_*(L)$:
$$p_L^M\circ \deg_M=\deg_L.$$
\end{proposition}

\begin{preuve}
Let $Q\in \pc(L)$ and $P\in \pc(M)$ with $P\subset Q$. Consider the commutative diagram
$$\xymatrix{
\pgo \ar[r] \ar[d]_{\chi_P}  & \qgo \ar[d]^{\chi_Q }\\
\car_M \ar[rd]_{\chi^M_G} \ar[r]^{\chi^M_L} & \car_L \ar[d]^{\chi_G^L} \\
& \car_G
}.$$
Recall that in §\ref{S:3.2} we defined a group scheme $J^G$ over $\car_G$. From the diagram above, one has the equality
$$\chi_P^* (\chi_G^M)^* J^G=(\chi_Q^* (\chi_G^L)^* J^G)_{|\pgo}.$$
Thus one obtains a diagram
\begin{equation}\label{eq:5.3.3}
  \xymatrix{
    \chi_P^* (\chi_G^M)^* J^G  \ar[r]  \ar[rd] &  X_*(M)\otimes_{\ZZ} \mathbb{G}_{m,\pgo} \ar[d] \\
    & X_*(L)\otimes_{\ZZ} \mathbb{G}_{m,\pgo}
  }
\end{equation}
where the horizontal arrow is the one defined in (\ref{eq:3.3.1}) (Proposition \ref{thm:3.3.1}), the diagonal arrow is the one defined in (\ref{eq:3.3.1}) relative to $L$ and restricted to $\pgo$, and the vertical arrow is induced by $p_M^L$. We now show that diagram (\ref{eq:5.3.3}) is commutative. It suffices to check it over the dense open subset $\pgo^{\reg}$. The horizontal (resp. diagonal) morphism in (\ref{eq:5.3.3}) factors through $I^P$, (resp. through $I^Q_{\pgo}$) (cf. Proposition \ref{thm:3.3.1}). Thus we have a diagram
$$
\xymatrix{
(\chi_P^* (\chi_G^M)^* J^G)_{|\pgo^{\reg}}  \ar[r]  \ar[rd] & I^P_{|\pgo^{\reg}} \ar[r] \ar[d] &  X_*(M)\otimes_{\ZZ} \mathbb{G}_{m,\pgo^{\reg}} \ar[d] \\
&   I^Q_{|\pgo^{\reg}} \ar[r]   &  X_*(L)\otimes_{\ZZ} \mathbb{G}_{m,\pgo^{\reg}} }
$$
whose commutativity is what has to be checked. Now the first triangle consists only of isomorphisms and is clearly commutative (cf.\ Lemma \ref{thm:3.1.1} and (\ref{eq:3.2.1})). The commutativity of the right-hand square is easily verified.

For every $a\in \Ac_M$, the commutativity of (\ref{eq:5.3.3}) yields a commutative diagram
$$
\xymatrix{
\mathrm{B}(J_a) \ar[r] \ar[rd] &  \mathrm{B}(X_*(M)\otimes_{\ZZ} \mathbb{G}_{m,C}) \ar[d]\\
&  \mathrm{B}(X_*(L)\otimes_{\ZZ} \mathbb{G}_{m,C}) }
$$
where the vertical arrow is induced by $p^M_L$ and the other arrows are given by (\ref{eq:5.3.1}). To conclude, one invokes the commutativity of the diagram
$$
\xymatrix{
\mathrm{B}(X_*(M)\otimes_{\ZZ} \mathbb{G}_{m,C})(C) \ar[r] \ar[d] & X_*(M) \ar[d] \\
\mathrm{B}(X_*(L)\otimes_{\ZZ} \mathbb{G}_{m,C})(C)  \ar[r] & X_*(L) }
$$
where the horizontal morphisms are the degree maps and the vertical morphisms are induced by $p_L^M$.
\end{preuve}
\end{paragr}

\begin{paragr}[Stack Quotient of the Centralizers.]\label{S:5.4}
Let $P\in \fc^G(T)$ be a semi-standard parabolic subgroup. Define an action of the group $P$ on the group scheme $I^P$ as follows: for any $h\in P$ and any $(p,X)\in I^P$, set
$$h\cdot(p,X)= (hph^{-1},\, \Ad(h)X).$$
This action commutes with the scaling action of $\Gmk$ on the first factor. One thus obtains an action of $P$ on $I^P_D$ for which the morphism $I^P_D \to \pgo_D$ is $P$–equivariant.

The stack quotient $[I^P_D/P]$ defines a group scheme over the stack $[\pgo_D/P]$. Let $h_P$ be a section over $C$ of the stack $[\pgo_D/P]$. By pulling back the stack $[I^P_D/P]$ along the section $h_P$, one obtains a group scheme over $C$, denoted $h_P^*[I^P_D/P]$. For a concrete description of this group: the section $h_P$ corresponds to the data of a $P$–torsor on $C$ together with a section $\theta$ of $\Ad{P,D}(\ec)=\ec\times^{P,\Ad}\pgo_D$. The group scheme $h_P^*[I^P_D/P]$ is the group scheme $\Aut_P(\ec,\theta)$ of automorphisms of $\ec$ that centralize $\theta$.

We now define an action of the group $G$ on the group scheme $\chi_G^* J^G$ as follows: for every $h\in G$ and every $(g,X)\in \chi_G^* J^G$, set
$$h\cdot(g,X)= (g,\, \Ad(h)X).$$
One checks that this action commutes with that of $\Gmk$ (induced from the one defined on $J^G$ in Paragraph \ref{S:5.1}). Hence one deduces an action of $G$ on $\chi_{G,D}^* J^G_D$. One verifies that the morphism
$$(\chi_{G,D}^* J^G_D)_{|\pgo_D} \to I^P_D$$
defined in Paragraph \ref{S:5.1} is $P$–equivariant.

The stack quotient $[(\chi_{G,D}^* J^G_D)/G]$ defines a smooth, commutative group scheme over $[\ggo_D/G]$. Since the action of $G$ on the first factor of $\chi_G^* J^G$ is trivial, one has the natural identification
\begin{equation}\label{eq:5.4.1}
  [(\chi_{G,D}^* J^G_D)/G]=[\chi_{G,D}]^*J_D^G,
\end{equation}
where we distinguish the morphisms $\chi_{G,D} : \ggo_D \to \car_{G,D}$ (cf. §\ref{S:4.1}) and $[\chi_{G,D}] : [\ggo_D/G] \to \car_{G,D}$, cf. (\ref{eq:4.2.1}).
\end{paragr}

\begin{paragr}[Action of $\Jc_a$ on $\mc_a$.]\label{S:5.5}
Let $a=(h_a,t)\in \Ac$ and $m=(h_m,t)\in \mc_a$. In §\ref{S:5.2} we defined a group scheme $J_a$ over $C$. According to diagram (\ref{eq:4.5.1}) and line (\ref{eq:5.4.1}) in Paragraph \ref{S:5.4}, one also has
\begin{equation}\label{eq:5.5.1}
J_a=h_m^* [\chi_{G,D}]^*J_D^G = h_m^*\bigl[(\chi_{G,D}^* J^G_D)/G\bigr].
\end{equation}
The morphism
$$\chi_{G,D}^* J^G_D \to I^G_D$$
defined in (\ref{eq:5.1.1}) in Paragraph \ref{S:5.1} is $G$–equivariant (cf. §\ref{S:5.4}). Hence one obtains a morphism of group schemes
\begin{equation}\label{eq:5.5.2}
  J_a = h_m^*\bigl[(\chi_{G,D}^* J^G_D)/G\bigr] \to h_m^*[I^G_D/G] = \Aut(m).
\end{equation}
The last equality identifies $h_m^*[I^G_D/G]$ with the group scheme $\Aut(m)$ of automorphisms of $m$ (cf. §\ref{S:5.4}).

We then let the Picard stack $\Jc_a$ act on the Hitchin fiber $\mc_a$ as follows: to any $m\in \mc_a$ and any $J_a$–torsor $j$, we associate the contracted product
$$j\cdot m = j\times^{J_a} m,$$
where $J_a$ acts on $m$ via the morphism $J_a \to \Aut(m)$ (cf. (\ref{eq:5.5.2})). One verifies that $j\times^{J_a} m \in \mc_a$.

\begin{proposition}\label{thm:5.5.1}
Let $a\in \Ac_M^{\el}$ and $m\in \mc_a$. For every $j\in \Jc_a$, the convex sets associated to $m$ and to $j\cdot m$ (cf. §\ref{S:4.8}) are translations of one another. More precisely, one has
$$\mathcal{C}_m = \mathcal{C}_{\, j\cdot m} + \deg_M(j),$$
where $\deg_M$ is the degree morphism defined in (\ref{eq:5.3.2}) in §\ref{S:5.3}.
\end{proposition}

\begin{preuve}
Let $P\in \pc(M)$ and let $h_P$ be the corresponding $P$–reduction of $m$ (cf. Proposition \ref{thm:4.6.1}). Examining Diagram \ref{eq:4.6.1}, one sees that
$$ J_a = h_P^* i_P^* [\chi_{G,D}]^* J_D^G.$$
One verifies that there is the following identification of group schemes over $[\pgo_D/P]$:
$$ i_P^* [\chi_{G,D}]^* J_D^G = \bigl[(J^G_D)_{|\pgo_D}/P\bigr].$$

The group scheme morphism over $[\ggo_D/G]$ deduced from the $G$–equivariant morphism (\ref{eq:5.1.1}) (for $P=G$) gives a morphism (see also (\ref{eq:5.4.1}))
$$ [\chi_{G,D}]^* J_D^G \to [I^G_D/G].$$
It follows from Proposition \ref{thm:3.2.1} that the induced morphism
$$ i_P^* [\chi_{G,D}]^* J_D^G = \bigl[(J^G_D)_{|\pgo_D}/P\bigr] \to i_P^*[I^G_D/G] $$
factors through
$$ \bigl[(J^G_D)_{|\pgo_D}/P\bigr] \to [I^P_D/P],$$
where the latter arrow is induced by the $P$–equivariant morphism (\ref{eq:5.1.1}) (the equivariance is discussed in §\ref{S:5.4}). Consequently, the morphism
$$J_a \to h_m^*[I^G_D/G] = \Aut(m)$$
factors through
\begin{equation}\label{eq:5.5.3}
  J_a \to h_P^*[I^P_D/P] = \Aut(h_P).
\end{equation}

In particular, if $j$ is a $J_a$–torsor, the action of $j$ on $m$ sends the $P$–reduction $h_P$ to $j\times^{J_a} h_P$ (where $J_a$ acts on $h_P$ via (\ref{eq:5.5.3})), which is necessarily the unique $P$–reduction of $j\cdot m$. The proposition then follows from the obvious equality on definitions,
$$\deg_P(j\times^{J_a} h_P) = \deg_M(j) + \deg_P(h_P).$$
\end{preuve}
\end{paragr}

\bigskip
\section{Degree Morphism and Picard Stack \texorpdfstring{$\Jc^1$}{J1}}\label{sec:6}

\begin{paragr}[Topological Space $\bg\lc^G\bd$.]\label{S:6.1} Let $\bg\lc^G\bd$ be the sober topological space whose underlying set is $\lc^G=\lc^G(T)$ and whose topology is the coarsest one for which the subsets $\lc^M$, for $M\in \lc^G$, are closed.

Let $M\subset L$ be two elements of $\lc^G(T)$. Let $X_*(L)_{\lc^M}$ be the constant sheaf on $\bg\lc^M\bd$ with value $X_*(L)$. Let
$$i_M : \bg\lc^M\bd \hookrightarrow \bg\lc^G\bd$$
be the canonical inclusion. Let $i_{M,*}(X_*(L)_{\lc^M})$ be the sheaf on $\bg\lc^G\bd$ supported on $\lc^M$ whose fiber at every point of $\lc^M$ is $X_*(L)$. There is a morphism
\begin{equation}\label{eq:6.1.1}
  i_{M,*}(X_*(M)_{\lc^M})\to i_{M,*}(X_*(L)_{\lc^M})
\end{equation}
induced by the morphism $p^M_L :  X_*(M)\to X_*(L)$ (cf. \ref{thm:5.3.1}). On the other hand, there is a morphism
\begin{equation}\label{eq:6.1.2}
  i_{L,*}(X_*(L)_{\lc^L})\to i_{M,*}(X_*(L)_{\lc^M})
\end{equation}
which, fiberwise, is the identity morphism of $X_*(L)$ at a point of $\lc^M$ and is zero outside $\lc^M$.

We then define a double morphism
\begin{equation}\label{eq:6.1.3}
  \bigoplus_{M\in \lc^G(T)}  i_{M,*}(X_*(M)_{\lc^M}) \stackrel{\displaystyle \longrightarrow}{\longrightarrow}  \bigoplus_{\stackrel{ M,L\in \lc^G(T)}{ M\subset L}}  i_{M,*}(X_*(L)_{\lc^M})
\end{equation}
as follows: on each component $i_{M,*}(X_*(M)_{\lc^M})$, for $M\in \lc^G$, the first arrow is obtained by composing the sum of the morphisms (\ref{eq:6.1.1}) for $L\in \lc^G(M)$ with the canonical inclusion
$$ \bigoplus_{L\in \lc^G(M)} i_{M,*}(X_*(L)_{\lc^M})\to  \bigoplus_{\stackrel{ M,L\in \lc^G(T)}{ M\subset L}}  i_{M,*}(X_*(L)_{\lc^M}); $$
the second arrow is obtained, on the component $i_{L,*}(X_*(L)_{\lc^L})$, for $L\in \lc^G$, by composing the sum of the morphisms (\ref{eq:6.1.2}) for $M\in \lc^L(T)$ with the canonical inclusion
$$ \bigoplus_{M \in \lc^L(T)} i_{M,*}(X_*(L)_{\lc^M})\to  \bigoplus_{\stackrel{ M,L\in \lc^G(T)}{ M\subset L}}  i_{M,*}(X_*(L)_{\lc^M}).$$
Note that in the double arrow (\ref{eq:6.1.3}), each component $i_{M,*}(X_*(L)_{\lc^M})$ of the target (with $M\subset L$) receives from each arrow only one component of the source, namely $  i_{M,*}(X_*(M)_{\lc^M})$ by the first arrow and $i_{L,*}(X_*(L)_{\lc^L})$ by the second.

Let $\mathcal{K}$ be the kernel of the double morphism (\ref{eq:6.1.3}). We leave to the reader the verification of the following lemma.

\begin{lemme}\label{thm:6.1.1}  The sheaf $\mathcal{K}$ is the constructible sheaf on $\bg \lc^G \bd$ whose fiber $\mathcal{K}_M$ at $M\in \lc^G$ is the $\ZZ$–module $X_*(M)$ and whose specialization map $\mathcal{K}_M \to \mathcal{K}_L$, for $M$ and $L$ in $\lc^G$ such that $M\subset L$, is the morphism $p^M_L : X_*(M) \to X_*(L)$ (cf. Proposition \ref{thm:5.3.1}).
\end{lemme}

\end{paragr}

\begin{paragr}[The Sheaf $\xc$ on $\Ac_G$.]\label{S:6.2} Let
  \begin{equation}\label{eq:6.2.1}
    \Ac_G \to  \bg \lc^G \bd
  \end{equation}
be the map which assigns the value $M$ on each locally closed subset $\Ac^{\el}_M$ for $M\in \lc^G(T)$ (cf. Proposition \ref{thm:4.4.2}).

  \begin{lemme}\label{thm:6.2.1}
The map (\ref{eq:6.2.1}) is continuous.
  \end{lemme}

  \begin{preuve}
Let $M\in \lc^G(T)$. The inverse image of the closed subset $\lc^M(T)$ is the union of the $\Ac_L^{\el}$ for $L\in \lc^M(T)$, which identifies with the closed set $\Ac_M$ (cf. (\ref{eq:4.4.1}) and Proposition \ref{thm:4.4.2}).
  \end{preuve}

Let $\xc$ be the constructible sheaf on $\Ac_G$ defined as the pullback by the continuous map (\ref{eq:6.2.1}) (cf. Lemma \ref{thm:6.2.1}) of the sheaf $\mathcal{K}$ on $\bg \lc^G \bd$. Thus, this sheaf is constant with value $X_*(M)$ on $\Ac^{\el}_M$, for $M\in \lc^G(T)$.
\end{paragr}

\begin{paragr}[The Picard Stack $\Jc$.]\label{S:6.3} Let
  \begin{equation}\label{eq:6.3.1}
    h_\Ac :    C\times_k \Ac \longrightarrow \car_{G,D}
  \end{equation}
be the canonical map which to a triple $(c,h_a,t)\in C\times_k \Ac$ associates $h_a(c)\in \car_{G,D}$. Let
$$J_{\Ac}=h^*_\Ac J^G_D$$
be the commutative group scheme, smooth over $C\times_k\Ac$, obtained as the pullback of $J^G_D$ by the morphism $h_\Ac$ in (\ref{eq:6.3.1}). Every $a=(h_a,t)\in \Ac(k)$ defines by base change a morphism $a_C : C\to C\times_k \Ac$ which, composed with $h_\Ac$, yields $h_a$. The group scheme over $C$ defined by $a_C^* J_{\Ac}$ is none other than the group scheme denoted $J_a$ defined in §\ref{S:5.2}.

Let $\mathrm{B}(J_{\Ac})$ be the $C\times_k \Ac$–classifying stack of $J_{\Ac}$. Let $\Jc$ be the stack over $\Ac$ whose associated groupoid $\Jc(S)$ for any test $\Ac$–scheme $S$ is the groupoid of morphisms
$$C\times_k S \to \mathrm{B}(J_{\Ac}).$$
Since $J_\Ac$ is commutative, the contracted product endows $\Jc$ with the structure of a Picard stack over $\Ac$.
\end{paragr}

\begin{paragr}[The Degree of $\Jc$.]\label{S:6.4} It is defined by the following proposition.

  \begin{proposition}\label{thm:6.4.1}
There exists a unique morphism
$$\deg : \Jc \to \xc$$
which, for every $M\in \lc^G(T)$ and every $a\in \Ac_M^{\el}$, induces the morphism
$$\deg_M : \Jc_a \to \xc_a=X_*(M)$$
defined in (\ref{eq:5.3.2}) in §\ref{S:5.3}.
  \end{proposition}

  \begin{preuve}
Uniqueness being evident, we concentrate on existence. Let $M\in \lc^G(T)$. The homomorphism $\deg_M$ from (\ref{eq:5.3.2}) induces a homomorphism
$$(\chi_G^M)^* \Jc  \longrightarrow X_*(M)_{\Ac_M},$$
where $\chi_G^M$ is the closed immersion of $\Ac_M$ into $\Ac_G$ defined in (\ref{eq:4.4.1}) and $X_*(M)_{\Ac_M}$ is the constant sheaf on $\Ac_M$ with value $X_*(M)$. By adjunction, one thus obtains a homomorphism
$$\Jc \longrightarrow (\chi_G^M)_* (X_*(M)_{\Ac_M}).$$

Taking the sum over $M\in \lc^G(T)$, one obtains a homomorphism
$$\Jc \longrightarrow \bigoplus_{M\in \lc^G(T)} (\chi_G^M)_* (X_*(M)_{\Ac_M}).$$
We will show that this homomorphism takes values in the subsheaf $\xc$ and satisfies the stated properties. To see that it takes values in $\xc$, it suffices to check that it lands in the kernel of the double arrow
$$\bigoplus_{M\in \lc^G(T)} (\chi_G^M)_* (X_*(M)_{\Ac_M})  \stackrel{\displaystyle \longrightarrow}{\longrightarrow}  \bigoplus_{\stackrel{ M,L\in \lc^G(T)}{ M\subset L}}  (\chi_G^M)_* (X_*(L)_{\Ac_M}),$$
which is the pullback of the double arrow (\ref{eq:6.1.3}) by the continuous map $\Ac_G\to \bg \lc^G \bd$ defined in (\ref{eq:6.2.1}). This follows from Proposition \ref{thm:5.3.1}. The other properties follow from Lemma \ref{thm:6.1.1}.
  \end{preuve}

\begin{paragr}[Action of $\Jc$ on $\mc$.]\label{S:6.5} This is simply to assemble the construction of §\ref{S:5.5}. Let $S$ be a test $k$–scheme. Let $m=(h_m,t)\in \mc(S)$ and $f(m)=(h_a,t)\in \Ac(S)$. Let $J_{\Ac,S}=J_{\Ac}\times_{\Ac}S$ and let $j\in \Jc(S)$ be a $J_{\Ac,S}$–torsor. We have, cf. line (\ref{eq:5.5.1}) of §\ref{S:5.5},
$$ J_{\Ac,S} = h_a^* J^G_D= h_m^* [\chi_{G,D}]^*J_D^G = h_m^*\bigl[(\chi_{G,D}^* J^G_D)/G\bigr].$$
As in (\ref{eq:5.5.2}), one deduces a morphism
$$J_{\Ac,S}\longrightarrow   h_m^*[I^G_D/G]=\Aut(m).$$
The action of $\Jc$ is then given by the contracted product
$$j\cdot m = j \times^{J_{\Ac,S}} m,$$
where $J_{\Ac,S}$ acts on $m$ via the morphism to $\Aut(m)$ just introduced.
\end{paragr}

\begin{paragr}[Action of $\Jc^1$ on $\mc^\xi$.]\label{S:6.6}  Let $\Jc^1$ be the kernel of the degree homomorphism defined in Proposition \ref{thm:6.4.1}. It is an open Picard substack of $\Jc$ which, fiberwise, contains the neutral component of $\Jc$. Since $\Jc$ acts on $\mc$, so does $\Jc^1$.

    \begin{proposition}\label{thm:6.6.1}
For every $\xi\in \ago_T$, the action of $\Jc^1$ preserves the open subset $\mc^\xi$.
    \end{proposition}

    \begin{preuve}
Let $m\in \mc$ and $a=f(m)\in \Ac$. Let $j\in \Jc_a$. By Proposition \ref{thm:5.5.1}, we have $\mathcal{C}_m = \mathcal{C}_{\, j\cdot m} + \deg_M(j)$. Moreover, if $j\in \Jc_a^1$ then $\deg_M(j)=0$ (cf. Proposition \ref{thm:6.4.1}). It follows that, for every $j\in \Jc_a^1$, we have
 $$\mathcal{C}_m = \mathcal{C}_{\, j\cdot m},$$
and, in particular, $j$ preserves the $\xi$–stability, which is defined by the condition $\xi\in \mathcal{C}_m$.
    \end{preuve}
\end{paragr}
\end{paragr}

\section{Cameral Curve and the Sheaf \texorpdfstring{$\pi_0(\Jc)$}{pi0(J)}}\label{sec:7}

\begin{paragr}[Cameral Curve.]\label{S:7.1} Let
$$\infty_\Ac : \Ac \to C\times_k \Ac$$
be the morphism deduced from $\infty$ by base change $\Ac\to \Spec(k)$. Following Donagi and Gaitsgory (cf. \cite{Donagi-Gaitsgory}), we introduce the universal cameral curve, denoted $Y$, which is the cover of $C\times_k \Ac$ defined by the Cartesian diagram
$$\xymatrix{
Y\ar[r]\ar[d]_{\pi_{Y}}&\tgo_{D}\ar[d]^{\chi^T_{G,D}}\\
C\times_k \Ac\ar[r]_{h_{\Ac}}&\car_{G,D}
}$$
where $h_\Ac$ is the canonical morphism (cf. (\ref{eq:6.3.1})) and $\chi^T_{G,D}$ is the morphism $\chi^T_G$ from §\ref{S:2.7} “twisted by $D$”. Since the morphism $\chi^T_{G,D}$ is $W^G$–invariant, there is a natural action of $W^G$ on $Y$. The covering $\pi_Y$ is finite, Galois with group $W^G$, and étale over a neighborhood of the image of $\infty_\Ac$. For every $a=(h_a,t)\in \Ac$, the point $t$ satisfies $\chi_G^T(t)=h_a(\infty)$. Hence there is a canonical morphism
$$\infty_Y : \Ac \to Y$$
over $\infty_\Ac$. Let $U\subset C\times_k \Ac$ be the inverse image by $h_\Ac$ of the open subset $\car_{G,D}^{\reg}$ where the discriminant $D^G$ does not vanish (cf. §\ref{S:2.7}). Let $V$ be the open subset of $Y$ defined by $V=\pi_Y^{-1}(U)$. Note that $\infty_\Ac$ takes values in $U$ and $\infty_Y$ takes values in $V$. The covering $\pi_Y$ is precisely étale over $U$.

Let $a=(h_a,t)\in \Ac$ and let $k(a)$ be the residue field of $a$. By pulling back the previous construction via the morphism $C\times_k k(a)\to C\times_k\Ac$ deduced from $a$ by base change $C\to \Spec(k)$, one obtains the cameral curve $Y_a$ over $k(a)$ defined by the Cartesian diagram
$$\xymatrix{
Y_{a}\ar@{^{(}->}[r]\ar[d]_{\pi_{Y_a}}&\mathfrak{t}_{D}\times_k k(a)
\ar[d]^{\chi^T_{G,D}}\\
C\times_k k(a)\ar@{^{(}->}[r]_{h_{a}}&\mathfrak{car}_{G,D}\times_k k(a)
}.$$
The group $W^G$ preserves $Y_a$ and the covering $\pi_{Y_a}$ is finite, Galois with group $W^G$, and étale over the point $\infty$. In its fiber over this point, we have the marked point $\infty_{Y_a}=\infty_Y\circ a$. Let $U_a$ and $V_a$ be the open subsets of $C\times_k k(a)$ and $Y_a$, respectively, deduced from $U$ and $V$ by base change.
\end{paragr}

\begin{paragr}[The Function $W_a$.]\label{S:7.2}  In this entire paragraph, the letter $W$ denotes an \emph{arbitrary} subgroup of the Weyl group $W^G$. For every $a\in \Ac$, let $W_a$ be the subgroup of $W^G$ defined as the stabilizer (normalizer) of the irreducible component of $Y_a$ passing through $\infty_{Y_a}$. Equivalently, it is the stabilizer of the connected component of $V_a$ which contains $\infty_{Y_a}$.

Let $\bg W^G \bd$ be the sober topological space whose underlying set consists of the subgroups $W$ of $W^G$ and whose topology is the coarsest one for which the subsets formed by those subgroups $W'$ contained in a given subgroup $W$ are closed.

\begin{proposition}\label{thm:7.2.1}
The map $\Ac \to \bg W^G \bd$ defined by $a\mapsto W_a$ is continuous.
\end{proposition}

\begin{preuve}
We first show the constructibility of this map. The composite morphism $Y\to C\times_k \Ac \to \Ac$ induces a surjective morphism $p : V \to \Ac$ for which $\infty_Y$ is a section. For every $a\in \Ac$, let $V_a^0$ denote the connected component of $V_a$ that contains $\infty_Y(a)$. Let $V^0$ be the union of the $V_a^0$ as $a$ runs through $\Ac$. According to Grothendieck (cf. \cite{EGAIV3} Proposition 15.6.4), $V^0$ is an open subset. For every subgroup $W$ of $W^G$, consider the constructible set
$$\Ac_W= p\Bigl( \bigcap_{w\in W} wV^0\Bigr) - \bigcup_{w\in W^G-W} p\Bigl( V^0\cap  wV^0\Bigr).$$
We will show the equality
\begin{equation}\label{eq:7.2.1}
  \Ac_W=\{a\in \Ac \mid W_a=W \},
\end{equation}
which implies the constructibility of the map $W_a$. Let $a\in \Ac_W$. There exists $y\in \bigcap_{w\in W} wV^0$ such that $p(y)=a$. Consequently, the union $\bigcup_{w\in W} wV^0_a$, each of whose components contains $y$, is connected. It follows that $\bigcup_{w\in W} wV^0_a=V^0_a$ and hence $W\subset W_a$. Now, let $w\in W_a$. Then $wV_a^0=V_a^0$ and $y\in V^0_a\cap  wV^0_a$. Necessarily $w\in W$, so that $W_a\subset W$, and finally $W_a=W$. Conversely, let $a\in \Ac$ be such that $W_a=W$. For any $y\in V_a^0$, since $V_a^0$ is stable under $W_a=W$, we have $y\in \bigcap_{w\in W} wV^0$. Moreover, if $w\notin W=W_a$, then $V^0_a\cap  wV^0_a=\emptyset$. Hence $a\in \Ac_W$. This proves (\ref{eq:7.2.1}).

Next, we show continuity. It suffices to show that for every subgroup $W$ of $W^G$, the union $\bigcup \Ac_{W'}$ taken over all subgroups $W'$ of $W$ is closed. As this union is constructible, it is enough to show that it is stable under specialization. Let $a$ be a point of $\Ac$ and $a'\in \Ac$ a generalization of $a$ in $\Ac$. We will show that $W_{a}\subset W_{a'}$, which implies the desired stability. By definition of $W_{a'}$, we have
$$V_{a'}^0 \subset V- \bigcup_{w\in W^G-W_{a'}} wV^0.$$
By the aforementioned Grothendieck result (cf. \cite{EGAIV3} Proposition 15.6.1), the set $V- \bigcup_{w\in W^G-W_{a'}} wV^0$ is closed in $V$. Therefore,
$$\overline{V_{a'}^0} \subset V- \bigcup_{w\in W^G-W_{a'}} wV^0,$$
where $\overline{V_{a'}^0}$ denotes the closure of $V_{a'}^0$ in $V$. Again by Grothendieck, we have $V^0_{a}\subset \overline{V_{a'}^0}$ and hence
$$V_{a}^0 \subset V- \bigcup_{w\in W^G-W_{a'}} wV^0,$$
which immediately implies that $W_a\subset W_{a'}$, as required.
\end{preuve}
\end{paragr}

\begin{paragr}[The Morphism from $X_*(T)_{\Ac}$ to $\Jc$.]\label{S:7.3} Let $S$ be an affine scheme over $k$ and $a=(h_a,t)\in \Ac(S)$. Write $C_S=C\times_k S$ and let $\mathcal{O}_{C_S,\infty}$ be the completion of $\mathcal{O}_{C_S}$ along $\{\infty\}\times_k S \subset C_S$, which can be identified with $\mathcal{O}_{S}[[\varpi_{\infty}]]$ by choosing a uniformizer $\varpi_{\infty}$ on $C$ at the point $\infty$. Let $F_{C_S,\infty}=\mathcal{O}_{C_S,\infty}[1/\varpi_{\infty}]$ (which does not depend on the choice of $\varpi_{\infty}$). Let $C^\infty= C-\{\infty\}$ and $C^\infty_S =C^\infty\times_k S$.

\begin{proposition}\label{thm:7.3.1}
There exists a canonical isomorphism
\begin{equation}\label{eq:7.3.1}
   J_a \times_C  \Spec(\mathcal{O}_{C_S,\infty}) \simeq  T\times_k  \Spec(\mathcal{O}_{C_S,\infty}).
\end{equation}
\end{proposition}

\begin{preuve}
In a neighborhood of $\infty$, the bundle $\car_{G,D}$ is canonically isomorphic to the trivial bundle with fiber $\car_G$. It follows that the section $h_a : C_S \to \car_{G,D}\times_k S$ restricts to a section
$$h_{a,\infty} : \Spec(\mathcal{O}_{C_S,\infty}) \to \car^{\reg}_G\times_k \Spec(\mathcal{O}_{C_S,\infty}).$$
Since the morphism $\chi_G^T : \tgo^{\reg} \to \car^{\reg}_G$ is étale, there exists a unique morphism
$$h_{a_T,\infty} : \Spec(\mathcal{O}_{C_S,\infty}) \to \tgo^{\reg}\times_k \Spec(\mathcal{O}_{C_S,\infty})$$
lifting $t$ and making the diagram below commute
$$\xymatrix{
\Spec(\mathcal{O}_{C_S,\infty}) \ar[rr]^{h_{a_T,\infty}} \ar[drr]_{h_{a,\infty}} & &  \tgo^{\reg}\times_k \Spec(\mathcal{O}_{C_S,\infty}) \ar[d]^{\chi_G^T} \\
& & \car^{\reg}_G\times_k \Spec(\mathcal{O}_{C_S,\infty}).
}$$
The isomorphism (\ref{eq:3.2.1}) from §\ref{S:3.2} induces the following isomorphism over $\tgo^{\reg}$:
$$\bigl((\chi_G^T)^*J^G\bigr)_{|\tgo^{\reg}}\simeq I^G_{\tgo^{\reg}}=T\times_k \tgo^{\reg}.$$
From this isomorphism and the equality
$$ J_a \times_C  \Spec(\mathcal{O}_{C_S,\infty}) = h_{a,\infty}^*J^G = h_{a_T,\infty}^*\bigl((\chi_G^T)^*J^G\bigr),$$
we deduce the desired isomorphism.
\end{preuve}

By Beauville--Laszlo formal descent (cf.\ \cite{Beauville-Laszlo}; for a recent reference better adapted to our situation, cf.\ \cite{Heinloth}), for every $\lambda \in X_*(T)$, there exists a triple $(j,\al,\beta)$, unique up to unique isomorphism, consisting of a $J_a$–torsor $j$ over $C_S$ and of $J_a$–equivariant isomorphisms
$$\al : j_{|\Spec(\mathcal{O}_{C_S,\infty}) } \to J_a \times_C  \Spec(\mathcal{O}_{C_S,\infty})$$
and
$$\beta : j_{|C^\infty_S} \to J_a \times_C   C_S^\infty$$
such that the induced isomorphism over $\Spec(F_{C_S,\infty})$ given by $\beta\circ \al^{-1}$ on
$$J_a \times_C  \Spec(F_{C_S,\infty}) \simeq T\times_k  \Spec(F_{C_S,\infty})$$
(isomorphism deduced from (\ref{eq:7.3.1})) is given by left translation by $\varpi_\infty^\lambda$.

From this construction one obtains a homomorphism
$$X_*(T)_{\Ac} \to \Jc,$$
where $X_*(T)_{\Ac}$ is the constant sheaf on $\Ac$ with value $X_*(T)$.
\end{paragr}

\begin{paragr}[Return to Some Results of Ng\^o.]\label{S:7.4}  Let $K$ be an extension of $k$ that is separably closed. Let $a=(h_a,t)\in \Ac(K)$ be a “geometric” point. In this entire paragraph, $\tc$ denotes a torus scheme over the open subset $U_a$ of $C_K=C\times_k K$ (cf. §\ref{S:7.1}). Such a torus scheme extends to a smooth group scheme over $C_K$ with \emph{connected fibers} and every such extension maps to the unique maximal extension denoted $\tilde{\tc}$ and called the connected Néron model of $\tc$. We introduce the following notations: let $\tc_{\infty_K}$ be the fiber of $\tc$ at the point $\infty_K=\infty\times_k K$. Let $X_*(\tc)=X_*(\tc_{\infty_K})=\Hom_{K\textrm{-gp}}(\mathbb{G}_{m,K},\tc_{\infty_K})$. The fundamental group of $U_a$, pointed by $\infty_K$, is denoted $\pi_1(U_a,\infty_K)$ and acts on $X_*(\tc_{\infty_K})$. Define
$$A(\tc)=X_*(\tc_{\infty_K})_{\pi_1(U_a,\infty_K)}$$
where the subscript $\pi_1(U_a,\infty_K)$ denotes the group of coinvariants under $\pi_1(U_a,\infty_K)$. Let $B(\tc)$ be the group of connected components of the Picard stack of $\tilde{\tc}$–torsors. Let $A$ and $B$ be the functors from the category of torus schemes on $U_a$ to the category of abelian groups defined by $\tc \mapsto A(\tc)$ and  $\tc \mapsto B(\tc)$. Using a variant of a result of Kottwitz (cf. \cite{Kottwitz-isocrystals} Lemma 2.2), Ng\^o showed that the functors $A$ and $B$ are isomorphic. It is important for us to specify such an isomorphism. For this, we use the following “Ng\^o-style” variant of the aforementioned Kottwitz lemma, where we denote by $X_*$ the functor $\tc \mapsto X_*(\tc_{\infty_K})$.

\begin{lemme}\label{thm:7.4.1}
The canonical morphisms
$$\Hom(A,B) \to \Hom(X_*,B)$$
and
$$ \Hom(X_*,B) \to \Hom(X_*(\mathbb{G}_{m,U_a}),B(\mathbb{G}_{m,U_a}))= B(\mathbb{G}_{m,U_a})$$
are isomorphisms. The group $B(\mathbb{G}_{m,U_a})$ is free of rank $1$, and any element of $\Hom(A,B)$ whose image under the composition of the two preceding isomorphisms is a generator of $B(\mathbb{G}_{m,U_a})$ is an isomorphism.
\end{lemme}

\begin{preuve}
The fundamental properties of the functor $B$ are proved by Ng\^o (cf. \cite{Ngo-endoscopie}). The rest is merely a transposition of results of Kottwitz (cf. \cite{Kottwitz-isocrystals} pp.206-207).
\end{preuve}

We now define a morphism of functors from $X_*$ to $B$. For every torus scheme $\tc$ on $U_a$, we thus define a morphism from $X_*(\tc_{\infty_K})$ into the stack of $\tilde{\tc}$–torsors. With the notations of §\ref{S:7.3}, for every $\lambda \in X_*(\tc_{\infty_K})$, one glues the trivial $\tilde{\tc}$–torsor on $C^\infty_K$ with the trivial $\tc$–torsor on $\Spec(\oc_{C_K,\infty})$ via Beauville–Laszlo using the gluing data $\varpi_\infty^\lambda$ (the torus $\tc$ is in fact constant on $\oc_{C_K,\infty}$, isomorphic to $X_*(\tc_{\infty_K})\otimes_\ZZ \mathbb{G}_{m,\oc_{C_K,\infty}}$). This yields a morphism
$$X_* \to B$$
by composing with the canonical morphism from the stack of $\tilde{\tc}$–torsors to its group of connected components.

\begin{lemme}\label{thm:7.4.2}
The functor morphism $X_* \to B$ factors canonically as in the diagram
$$\xymatrix{
X_* \ar[rd] \ar[r] & A \ar[d] \\
&  B
}$$
where the horizontal morphism is the canonical one and the vertical morphism is an isomorphism.
\end{lemme}

\begin{preuve}
By Lemma \ref{thm:7.4.1}, it suffices to see that the functor $X_* \to B$ sends the canonical generator of $X_*(\mathbb{G}_{m,\infty_K})$ (given by the identity homomorphism) to a generator of $B(\mathbb{G}_{m,U_a})$. But when $\tc= \mathbb{G}_{m,U_a}$, we have $\tilde{\tc}=\mathbb{G}_{m,C_K}$, and the degree morphism of the Picard stack of $\tilde{\tc}$–torsors (which is nothing other than the Picard stack of locally free $\oc_{C_K}$–modules of rank $1$) identifies its group of connected components with $\ZZ$. The canonical generator of $X_*(\tc_{\infty_K})$, defined by the gluing explained above, yields a $\mathbb{G}_{m,C_K}$–torsor of degree $1$, hence the lemma.
\end{preuve}
\end{paragr}

\begin{paragr}[Ng\^o's Factorization.]\label{S:7.5} Let $\mathcal{Y}$ be the sheaf on $\Ac$ defined as the pullback via the continuous map $\Ac\rightarrow \bg W^G \bd$ of the sheaf whose fiber over $W\subset W^G$ is $X_{*}(T)_{W}$, with the natural transition morphisms. By definition, this sheaf is a quotient of the constant sheaf $X_{*}(T)_{\Ac}$. In §\ref{S:7.3} we defined a morphism $X_{*}(T)_{\Ac} \to \Jc$. According to Ng\^o, there exists a constructible sheaf on $\Ac$, denoted $\pi_0(\Jc)$, whose fiber at a geometric point $a$ of $\Ac$ is the group $\pi_0(\Jc_a)$ of connected components of $\Jc_a$ (cf. Proposition 6.2 of \cite{Ngo-endoscopie}). By composing with the canonical morphism $\Jc\to \pi_0(\Jc)$, one obtains a morphism $X_{*}(T)_{\Ac}\to \pi_0(\Jc)$. Ng\^o showed the following factorization (cf. Proposition 5.5.1 of \cite{Ngo-lemme}).

\begin{proposition}\label{thm:7.5.1}
We have the following factorization
$$\xymatrix{
   X_{*}(T)_{\Ac}  \ar[r] \ar[rd] & \yc \ar[d]\\
& \pi_0(\Jc)
}$$
where the horizontal arrow is the canonical morphism and the other two arrows are epimorphisms.
\end{proposition}

\begin{preuve}
For the reader’s convenience, we recall some points of Ng\^o's proof. It suffices to prove the factorization fiberwise. Let $a=(h_a,t)\in \Ac(K)$ be a geometric point, where $K$ is an extension of $k$ that is separably closed. Let $\Jc^0_a$ be the stack of $J^0$–torsors on $C\times_k K$, where $J^0_a$ is the connected group scheme which, fiberwise, is the connected component of the identity section of $J_a$. Since the fiber of $J_a$ at $\infty_K$ is connected, the morphism $X_*(T)\to \Jc_a$ factors through the evident morphism $X_*(T)\to \Jc^0_a$. According to Ng\^o (cf. \cite{Ngo-endoscopie} Proposition 6.3), the canonical morphism $\Jc^0_a \to \Jc_a$ induces a surjective morphism $\pi_0(\Jc^0_a) \to \pi_0(\Jc_a)$. Moreover, the canonical morphism $J^0_a \to \tilde{J}_a$ from $J^0_a$ to the connected Néron model $\tilde{J}_a$ of the torus scheme $J_{|U_a}$ induces a morphism from $\Jc^0_a$ to the stack $\tilde{\Jc}_a$ of $\tilde{J}_a$–torsors on $C\times_k K$ and, by Ng\^o (cf. \cite{Ngo-endoscopie} Proposition 6.4), an isomorphism $\pi_0(\tilde{\Jc}_a)=\pi_0(\Jc^0_a)$.

It thus suffices to prove the desired factorization for the morphism $X_*(T)\to \pi_0(\tilde{\Jc}_a)$, that is, the factorization through the epimorphism $X_*(T)_{W_a}\to \pi_0(\tilde{\Jc}_a)$ (in fact, this will be an isomorphism).

With the notations of §\ref{S:7.1}, we have a commutative diagram
$$\xymatrix{
\Spec(K) \ar[rr]^{\infty_{Y_a}} \ar[rrd]_{\infty} &  & V_{a}\ar@{^{(}->}[rr]^{i_a} \ar[d]^{\pi_{Y_a}}&& \mathfrak{t}^{\reg}_{D,K} \ar[d]^{\chi^T_{G,D}}\\
& & U_a\ar@{^{(}->}[rr]_{h_{a}}& &\car^{\reg}_{G,D,K}
}$$
where the square is Cartesian and where we denote by a subscript $K$ the base change from $k$ to $K$.

Let $(Z,z)$ be a connected, étale, Galois cover of $U_a$ pointed by a point $z : \Spec(K)\to Z$ lying over $\infty$. We may assume, without loss of generality, that for each connected component of $V_a$, there exists a $U_a$–morphism from $Z$ onto that component. The composition on the right with $z$ induces a bijection from $\Hom_{U_a}(Z,V_a)$ onto the fiber of $\pi_{Y_a}$ over $\infty$. Denote by $\alpha : Z\to V_a$ the morphism corresponding to $\infty_{Y_a}$. With the notations of §\ref{S:5.2}, we have
$$J_a=J_D^G\times_{\car^{\reg}_{G,D,K}} U_a.$$
But according to the isomorphism (\ref{eq:3.2.1}) in §\ref{S:3.2}, we have the following isomorphism
$$J_D^G \times_{\car^{\reg}_{G,D,K}}\mathfrak{t}^{\reg}_{D,K}\simeq  T \times_k \mathfrak{t}^{\reg}_{D,K},$$
whence
$$J_a\times_{U_a}V_a=J_D^G\times_{\car^{\reg}_{G,D,K}} V_a \simeq  T \times_k V_a.$$
Using the morphism $\alpha$, one obtains
\begin{equation}\label{eq:7.5.1}
  J_a\times_{U_a} Z = (J_a\times_{U_a} V_a)\times_{V_a,\alpha} Z \simeq T \times_k Z.
\end{equation}
For every $\Phi\in \Aut_{U_a}(Z)$, there exists a unique element $w_{\Phi}\in W^G$ such that $\alpha\circ \Phi = w_{\Phi}\cdot\alpha$ (where $\cdot$ denotes the action of $W^G$ on $V_a$). Let $\pi_1(U_a,\infty)$ be the fundamental group of $U_a$ pointed by $\infty$. The assignment $\Phi \mapsto w_{\Phi}$ defines a morphism from $\pi_1(U_a,\infty)$ into $W^G$, whose image is precisely the subgroup $W_a$. The obvious action of $\pi_1(U_a,\infty)$ on $J_a\times_{U_a} Z$ corresponds, via the isomorphism (\ref{eq:7.5.1}), to the diagonal action of $\pi_1(U_a,\infty)$ on $T \times_k Z$, where the action on the first factor is given by the morphism into $W^G$. In particular, we have the identification
$$\Hom_{Z\textrm{-gpe}}( \mathbb{G}_{m,Z},  J_a\times_{U_a} Z)=X_*(T)$$
which is equivariant under the action of $\pi_1(U_a,\infty)$ when this group acts on $X_*(T)$ via its morphism into $W_a$. One deduces the identifications
$$X_*(J_a)_{\infty}=X_*(T)$$
and
\begin{equation}\label{eq:7.5.2}
\bigl(X_*(J_a)_{\infty}\bigr)_{\pi_1(U_a,\infty) }=X_*(T)_{W_a}.
\end{equation}
Note that this identification also holds for $\tilde{J}_a$ since $J_a\simeq \tilde{J}_a$ over $U_a$.

Thus, the morphism $X_*(T)\to \pi_0(\tilde{\Jc}_{a})$ identifies with the morphism $X_*(\tilde{J}_a)_{\infty} \to \pi_0(\tilde{\Jc}_{a})$ considered in §\ref{S:7.4}. Lemma \ref{thm:7.4.2} and line (\ref{eq:7.5.2}) imply that this morphism factors through $X_*(T)_{W_a}$ and that the morphism $X_*(T)_{W_a}\to  \pi_0(\tilde{\Jc}_{a})$ is an isomorphism. This concludes the proof.
\end{preuve}
\end{paragr}
\section{The \texorpdfstring{$s$}{s}-decomposition}\label{sec:8}

\begin{paragr}\label{S:8.1}
In \cite[\S8]{Ngo-endoscopie}, Ngô introduced a decomposition of the perverse cohomology sheaves of the Hitchin fibration over the elliptic part. The aim of this section is to generalize Ngô's decomposition to open subsets that are not necessarily elliptic.
\end{paragr}

\begin{paragr}\label{S:8.2}
The following proposition will be useful in what follows.

\begin{proposition}\label{thm:8.2.1}
Let $a\in\Ac_G$ and let $M\in\lc^G$ be a Levi subgroup. The following two assertions are equivalent\,:
\begin{enumerate}
\item $a\in\Ac_M$\,; \item $W_a\subset W^M$.
\end{enumerate}
\end{proposition}

\begin{preuve}
Let $a=(h_a,t)\in\Ac_G$ and let $k(a)$ be the residue field of $a$. Let us show that assertion 2 implies the first one. With the notation of Section~\ref{S:7.1}, we have a commutative diagram
\begin{equation}\label{eq:8.2.1}
\xymatrix{
\Spec(k(a))\ar[r]^-{\infty_{Y_a}}\ar[dr]_-{\infty} & V_a\ar@{^{(}->}[r]^-{i_a}\ar[d]^{\pi_{Y_a}} & \tgo^{\reg}_{D,k(a)}\ar[d]^{\chi^T_{G,D}}\\
& U_a\ar@{^{(}->}[r]_-{h_a} & \car^{\reg}_{G,D,k(a)}\,,
}
\end{equation}
where the square is Cartesian. Let $V_a^0$ be the connected component of $V_a$ that contains $\infty_{Y_a}$. By definition of $W_a$, the restriction of $\pi_{Y_a}$ to $V_a^0$ makes $V_a^0$ into a connected, étale and Galois covering of group $W_a$. Since the inclusion $i_a$ is $W$-equivariant, we obtain by passing to the quotient by $W_a$ a morphism
$$
U_a\to\tgo^{\reg}_{D,k(a)}//W_a
$$
over $\car^{\reg}_{G,D,k(a)}$. Since $W_a\subset W^M$, we may compose this morphism with the canonical morphism $\tgo^{\reg}_{D,k(a)}//W_a\to\car^{\reg}_{M,D,k(a)}$, so that we obtain a morphism $h_{a_M}$ which fits into the commutative diagram
$$
\xymatrix{
U_a\ar@{^{(}->}[r]^-{h_{a_M}}\ar[d] & \car_{M,D,k(a)}\ar[d]^{\chi^M_{G,D}}\\
C\times_k k(a)\ar@{^{(}->}[r]_-{h_a} & \car_{G,D,k(a)}\,. }
$$
By properness of the morphism $\chi^M_{G,D}$, the morphism $h_{a_M}$ extends to $C\times_k k(a)$. One checks that $(h_{a_M},t)$ belongs to $\Ac_M$ and has image $a$ in $\Ac_G$.

The converse follows from the next lemma.
\end{preuve}

\begin{lemma}\label{thm:8.2.2}
Let $a_M\in\Ac_M$ and let $a\in\Ac_G$ be its image under the closed immersion $\Ac_M\to\Ac_G$ of~(\ref{eq:4.4.1}). Let $W_{a_M}$ be the subgroup of $W^M$ defined in Section~\ref{S:7.2} relative to the group $M$. We have
$$
W_{a_M}=W_a.
$$
\end{lemma}

\begin{preuve}
Let $a_M=(h_{a_M},t)\in\Ac_M$ and let $a=(h_a,t)$ be its image in $\Ac_G$. We adopt the notation of the proof of Proposition~\ref{thm:8.2.1}, but relative to the group $M$. We thus have a diagram analogous to~(\ref{eq:8.2.1})
$$
\xymatrix{
\Spec(k(a_M))\ar[r]^-{\infty_{Y_{a_M}}}\ar[dr]_-{\infty} & V_{a_M}\ar@{^{(}->}[r]\ar[d] & \tgo^{\reg}_{D,k(a_M)}\ar[d]^{\chi^T_{M,D}}\\
& U_{a_M}\ar@{^{(}->}[r]_-{h_{a_M}} & \car^{\reg}_{G,D,k(a_M)}\,, }
$$
where the exponent reg means $G$-regular (in accordance with the conventions of~\S\ref{S:2.7}). Moreover we have $U_a=U_{a_M}$ and an immersion $V_{a_M}\hookrightarrow V_a$ which is both open and closed, which is $W^M$-equivariant and which sends $\infty_{Y_{a_M}}$ to $\infty_{Y_a}$. In particular, the connected component of $\infty_{Y_a}$ in $V_a$ is equal to the image of that of $\infty_{Y_{a_M}}$ in $V_{a_M}$. If $w\in W^G$ satisfies $wV_{a_M}\cap V_{a_M}\not=\emptyset$, then necessarily $w\in W^M$. It follows that $W_a\subset W^M$. The result is then clear.
\end{preuve}
\end{paragr}

\begin{paragr}\label{S:8.3}
In Propositions~\ref{thm:6.4.1} and~\ref{thm:7.5.1} we defined respectively a morphism $\deg:\Jc\to\xc$, which obviously factors through a morphism
$$
\deg:\pi_0(\Jc)\to\xc,
$$
and a morphism $\yc\to\pi_0(\Jc)$. Composing these two morphisms, we obtain $\yc\to\xc$ and, composing further with the morphism $X_*(T)_{\Ac}\to\yc$, we obtain a morphism $X_*(T)_{\Ac}\to\xc$. Let $\yc^1$ and $X_*(T)^1_{\Ac}$ be the constructible sheaves on $\Ac$ defined respectively by
$$
\yc^1=\Ker(\yc\to\xc)
$$
and
$$
X_*(T)^1_{\Ac}=\Ker(X_*(T)_{\Ac}\to\xc).
$$
Let $\pi_0(\Jc^1)$ be the sheaf of connected components of $\Jc^1$: it is a subsheaf of $\pi_0(\Jc)$ whose fiber at a geometric point $a$ is the subgroup of $\pi_0(\Jc_a)$ formed by the components of degree $0$ (where the degree is the morphism above). The morphism $X_*(T)_{\Ac}\to\Jc$ then restricts to a morphism $X_*(T)^1_{\Ac}\to\Jc^1$ which, by Proposition~\ref{thm:7.5.1}, fits into the commutative diagram
\begin{equation}\label{eq:8.3.1}
\xymatrix{
X_*(T)^1_{\Ac}\ar[r]\ar[dr] & \yc^1\ar[d]\\
& \pi_0(\Jc^1),
}
\end{equation}
where the two arrows with target $\pi_0(\Jc^1)$ are epimorphisms.

\begin{proposition}\label{thm:8.3.1}
Let $a\in\Ac_M^{\el}$ with $M\in\lc^G$ and let $a'\in\Ac_G$ be a specialization of $a$. Let $L\in\lc^G$ be defined by the condition $a'\in\Ac_L^{\el}$.
\begin{enumerate}
\item The sheaves $X_*(T)^1_{\Ac}$ and $\yc^1$ have as fiber at $a$ respectively the group $X_*(T\cap M_{\der})$ and the image group of the canonical morphism
$$
X_*(T\cap M_{\der})_{W_a}\to X_*(T)_{W_a}\,;
$$
\item we have $L\subset M$ and $W_{a'}\subset W_a$\,; \item the specialization arrows $(X_*(T)^1_{\Ac})_{a'}\to(X_*(T)^1_{\Ac})_a$ and $\yc^1_{a'}\to\yc^1_a$ are respectively given by the canonical morphisms
$$
X_*(T\cap L_{\der})\to X_*(T\cap M_{\der})
$$
and
$$
\mathrm{Im}(X_*(T\cap L_{\der})_{W_{a'}}\to X_*(T)_{W_{a'}})\to\mathrm{Im}(X_*(T\cap M_{\der})_{W_a}\to X_*(T)_{W_a})\,;
$$
\item the fibers of the sheaf $\yc^1$ are finite\,; \item the fibers of the sheaf $\pi_0(\Jc^1)$ are finite.
\end{enumerate}
\end{proposition}

\begin{preuve}
Let us prove assertion 1. By Lemma~\ref{thm:8.3.2} below, the fiber at $a\in\Ac_M^{\el}$ of $X_*(T)^1_{\Ac}$ is the kernel of the canonical morphism $X_*(T)\to X_*(M)$: it is therefore indeed $X_*(T\cap M_{\der})$. We know by Proposition~\ref{thm:7.5.1} that this arrow factors through $X_*(T)_{W_a}$, which is the fiber at $a$ of $\yc$. The fiber at $a$ is therefore the kernel of the arrow $X_*(T)_{W_a}\to X_*(M)$ (which is well defined since $W_a\subset W^M$ by Proposition~\ref{thm:8.2.1}). But this kernel is also the image of the morphism $X_*(T\cap M_{\der})_{W_a}\to X_*(T)_{W_a}$.

Assertion 2 follows from the continuity of the maps $\Ac_G\to\bg\lc^G\bd$ and $\Ac_G\to\bg W^G\bd$ (cf.\ Lemma~\ref{thm:6.2.1} and Proposition~\ref{thm:7.2.1}).

We leave assertion 3 to the reader. Assertion 4 follows from Lemma~\ref{thm:8.3.3} below. Assertion 5 is a consequence of assertion 4, since we have an epimorphism $\yc^1\to\pi_0(\Jc^1)$.
\end{preuve}

\begin{lemma}\label{thm:8.3.2}
Let $M\in\lc^G$. The morphism $X_*(T)_{\Ac}\to\xc$ has as fiber at $a\in\Ac_M^{\el}$ the canonical morphism
$$
X_*(T)\to X_*(M).
$$
\end{lemma}

\begin{preuve}
Let $a\in\Ac_M^{\el}$ and let $K$ be its residue field. Let $h_{a_M}$ be the morphism from $C_K$ to $\car_{M,D}$ deduced from $a$. We take up again the notation of Section~\ref{S:7.3}. Let $\lambda\in X_*(T)$ and let $j_\lambda$ be the $J_a$-torsor on $C_K=C\times_k K$ obtained by gluing by means of $\varpi^\lambda_\infty\in T(\oc_{C_K,\infty})\simeq J_a(\oc_{C_K,\infty})$ (cf.\ the isomorphism of Proposition~\ref{thm:7.3.1}) which defines the morphism $X_*(T)\to\Jc_a$. The degree of the $J_a$-torsor $j_\lambda$ is the map which, to $\mu\in X^*(M)$, associates the degree of the $\mathbb{G}_{m,C_K}$-torsor $\mu(j_\lambda)$ obtained when one pushes $j_\lambda$ forward by the composite morphism
\begin{equation}\label{eq:8.3.2}
J_a\longrightarrow X_*(M)\otimes_{\ZZ}\mathbb{G}_{m,C_K}\stackrel{\mu}{\longrightarrow}\mathbb{G}_{m,C_K},
\end{equation}
where the first arrow is obtained by pulling back by $h_{a_M}$ the morphism
$$
(\chi^M_{G,D})^*J^G_D\to X_*(M)\otimes_{\ZZ}\mathbb{G}_{m,\car_{M,D}}
$$
defined in~(\ref{eq:5.1.2}) and deduced from Proposition~\ref{thm:3.3.1}. The degree of $j_\lambda$ is therefore the valuation of the image of $\varpi^\lambda_\infty$ under the composite morphism~(\ref{eq:8.3.2}).

In fact, on $\Spec(\oc_{C_K,\infty})$, the section $h_{a_M}$ takes regular values and factors through a canonical section, denoted $h_{a_T,\infty}$, of $\tgo_D^{\reg}$ (cf.\ the proof of Proposition~\ref{thm:7.3.1}). The isomorphism $J_a\times_C\Spec(\oc_{C_S,\infty})\simeq T\times_k\Spec(\oc_{C_S,\infty})$ is then obtained by pullback under $h_{a_T,\infty}$ of the isomorphism
$$
\bigl((\chi^T_G)^*J^G\bigr)_{|\tgo^{\reg}}\simeq X_*(T)\otimes_{\ZZ}\mathbb{G}_{m,\tgo^{\reg}}.
$$
As can be seen from the definitions, the morphism
$$
X_*(T)\otimes_{\ZZ}\mathbb{G}_{m,\tgo^{\reg}}=\bigl((\chi^T_G)^*J^G\bigr)_{\tgo^{\reg}}\to X_*(M)\otimes_{\ZZ}\mathbb{G}_{m,\tgo^{\reg}}
$$
which one obtains by pulling back by $\chi^T_M$ the morphism of Proposition~\ref{thm:3.3.1} is deduced from the canonical morphism $X_*(T)\to X_*(M)$. Consequently, we have
$$
\deg(\mu(j_\lambda))=\mu(\lambda)
$$
(the pairing is the one between characters and cocharacters of $T$) and the degree of $\mu(j_\lambda)$ is indeed the image of $\lambda$ under the projection $X_*(T)\to X_*(M)$.
\end{preuve}

\begin{lemma}\label{thm:8.3.3}
For every $M\in\lc^G$ and every $a\in\Ac_M^{\el}$ the group $X_*(T\cap M_{\der})_{W_a}$ is finite. In particular, the image of the morphism
$$
X_*(T\cap M_{\der})_{W_a}\to X_*(T)_{W_a}
$$
is finite.
\end{lemma}

\begin{preuve}
We shall first prove that the group $X_*(T\cap M_{\der})^{W_a}$ of invariants under $W_a$ is trivial. Let $\lambda$ be an element of this group, which we view as a morphism from $\Gmk$ to $M$. Let $L$ be the Levi subgroup in $\lc^M$ which is the centralizer of the image torus. We thus have $W_a\subset W^L$ and, by Proposition~\ref{thm:8.2.1}, we have $a\in\Ac_L$. Since $a$ is elliptic in $M$, it follows that $M=L$. Hence $\lambda$ is central in $M_{\der}$, which is possible only if $\lambda=0$, which is what had to be proved.

Since $X_*(T\cap M_{\der})$ is of finite type, the finiteness of $X_*(T\cap M_{\der})_{W_a}$ is equivalent to the vanishing of $X_*(T\cap M_{\der})_{W_a}\otimes_{\ZZ}\QQ$. Now averaging over $W_a$ induces an isomorphism of this $\QQ$-vector space onto $X_*(T\cap M_{\der})^{W_a}\otimes_{\ZZ}\QQ$, whose vanishing we have proved. This concludes the proof.
\end{preuve}
\end{paragr}

\begin{paragr}[The open subset $\uc_M$.]\label{S:8.4}
Let $M\in\lc^G$ and let $\uc_M\subset\Ac_G$ be the open subset defined by
$$
\uc_M=\bigcup_{L\in\lc^G(M)}\Ac_L^{\el}.
$$
It is also the complement of the closed subset
$$
\bigcup_{L\in\lc^G-\lc^G(M)}\Ac_L.
$$
The open subset $\uc_G$ is nothing but the elliptic open subset $\Ac_G^{\el}$. We define the $M$-parabolic part of $\mc^\xi$, denoted $\mc^{\xi,M-\parab}$, by the base change
$$
\mc^{\xi,M-\parab}=\mc^\xi\times_{\Ac}\uc_M.
$$
Let
$$
f^{\xi,M-\parab}:\mc^{\xi,M-\parab}\to\uc_M
$$
be the restriction of the Hitchin morphism $f^\xi$ to $\mc^{\xi,M-\parab}$. When $M=G$, we replace ``$G$-par'' by ``ell''.

Let $\Gamma(\uc_M,\cdot)$ be the functor of sections over $\uc_M$.

\begin{lemma}\label{thm:8.4.1}
The restriction of $\yc^1$ to $\uc_M$ is a quotient of the constant sheaf with value $X_*(T\cap M_{\der})$. In particular, we have group homomorphisms
$$
X_*(T\cap M_{\der})\longrightarrow\Gamma(\uc_M,\yc^1)
$$
and
\begin{equation}\label{eq:8.4.1}
X_*(T\cap M_{\der})\longrightarrow\Gamma(\uc_M,\pi_0(\Jc^1)).
\end{equation}
\end{lemma}

\begin{preuve}
The first assertion follows from the description of $\yc^1$ given in Proposition~\ref{thm:8.3.1}. The rest follows from it.
\end{preuve}
\end{paragr}

\begin{paragr}[The $s$-decomposition.]\label{S:8.5}
With the notation of Section~\ref{S:2.5}, we denote by $\Tc$ the complex torus with character lattice $X^*(\Tc)=X_*(T)$. In all that follows, we always identify in the notation the algebraic groups over $\CC$ with their groups of $\CC$-valued points. Let $Z^0_{\Mc}$ be the connected center of the Langlands dual $\Mc$ of $M$. It is also the neutral component of the subgroup of $\Tc$ defined as the intersection of the kernels of the characters $\alpha^\vee$, where $\alpha^\vee$ runs over the coroots of $T$ in $M$. Let us note that we have
$$
\Tc/Z^0_{\Mc}=\Hom_{\ZZ}(X_*(T\cap M_{\der}),\CC^\times).
$$
Let $(\Tc/Z^0_{\Mc})_{\tors}$ be the subgroup of torsion points of $\Tc/Z^0_{\Mc}$. Let $\Qb$ be the algebraic closure of $\QQ$ in $\CC$. We fix an embedding of $\Qb$ into $\Qlb$. We then have the natural inclusions
$$
(\Tc/Z^0_{\Mc})_{\tors}\subset\Hom_{\ZZ}(X_*(T\cap M_{\der}),\Qb^\times)\subset\Hom_{\ZZ}(X_*(T\cap M_{\der}),\Qlb^\times).
$$

The action of the Picard stack $\Jc^1$ on $\mc^\xi$ over $\Ac$ induces an action of $\Jc^1$ on $\mc^{\xi,M-\parab}$ over $\uc_M$, hence an action on each perverse cohomology sheaf ${}^{\mathrm{p}}\hc^i(Rf_*^{\xi,M-\parab}\Qlb)$ of $Rf_*^{\xi,M-\parab}\Qlb$. By the homotopy lemma (cf.\ \cite[Lemma~3.2.3]{Laumon-Ngo}), this action factors through an action of the sheaf $\pi_0(\Jc^1)$ on ${}^{\mathrm{p}}\hc^i(Rf_*^{\xi,M-\parab}\Qlb)$.

\begin{theorem}\label{thm:8.5.1}
For all $s\in(\Tc/Z^0_{\Mc})_{\tors}$ and $i\in\ZZ$, let
$$
{}^{\mathrm{p}}\hc^i(Rf_*^{\xi,M-\parab}\Qlb)_s
$$
be the sum of the $s'$-isotypic components of the finite commutative group
$$
\Gamma(\uc_M,\pi_0(\Jc^1))
$$
acting on the perverse cohomology sheaf ${}^{\mathrm{p}}\hc^i(Rf_*^{\xi,M-\parab}\Qlb)$, for the characters $s'\in\Hom_{\ZZ}(\Gamma(\uc_M,\pi_0(\Jc^1)),\Qlb^\times)$ which induce the character $s$ on $X_*(T\cap M_{\der})$ via the morphism~(\ref{eq:8.4.1}).

We then have the following assertions\,:
\begin{enumerate}
\item if $L$ is an element of $\lc^G(M)$ and if $a$ is a geometric point of $\Ac_L^{\el}$ such that the fiber ${}^{\mathrm{p}}\hc^i(Rf_*^{\xi,M-\parab}\Qlb)_{a,s}$ is nonzero, then there exists a lift $\tilde s$ of $s$ to $\Tc^{W_a}$. \item there are only finitely many $s$ for which
$$
{}^{\mathrm{p}}\hc^i(Rf_*^{\xi,M-\parab}\Qlb)_s\not=0\,;
$$
\item we have the following decomposition\,:
\begin{equation}\label{eq:8.5.1}
{}^{\mathrm{p}}\hc^i(Rf_*^{\xi,M-\parab}\Qlb)=\bigoplus_{s\in(\Tc/Z^0_{\Mc})_{\tors}}{}^{\mathrm{p}}\hc^i(Rf_*^{\xi,M-\parab}\Qlb)_s.
\end{equation}
\end{enumerate}
\end{theorem}

\begin{remark}\label{thm:8.5.2}
When $M=T$, we have $Z_{\Mc}=\Tc$ and the decomposition~(\ref{eq:8.5.1}) is tautological. When $M=G$, the above theorem is due to Ngô.
\end{remark}

\begin{preuve}
Since $\Gamma(\uc_M,\pi_0(\Jc^1))$ is a finite commutative group (cf.\ Proposition~\ref{thm:8.3.1}, assertion 5) which acts on ${}^{\mathrm{p}}\hc^i(Rf_*^{\xi,M-\parab}\Qlb)$, we have a decomposition of the cohomology into a direct sum of isotypic components indexed by the characters of $\Gamma(\uc_M,\pi_0(\Jc^1))$ with values in $\Qb$. Such a character induces a character of finite order of $X_*(T\cap M_{\der})$ via the morphism~(\ref{eq:8.4.1}) of Lemma~\ref{thm:8.4.1}, hence an element of $(\Tc/Z^0_{\Mc})_{\tors}$. The decomposition~(\ref{eq:8.5.1}) follows from this, as does the finiteness in assertion 2.

Let $L\in\lc^G(M)$ and let $a$ be a geometric point of $\Ac_L^{\el}$. We thus have $W_a\subset W^L$ by virtue of Proposition~\ref{thm:8.2.1}. The group $\Gamma(\uc_M,\pi_0(\Jc^1))$ acts on the fiber ${}^{\mathrm{p}}\hc^i(Rf_*^{\xi,M-\parab}\Qlb)_{a,s}$ via the restriction morphism $\Gamma(\uc_M,\pi_0(\Jc^1))\to\pi_0(\Jc^1_a)$. Let $s'$ be a character of the finite group $\pi_0(\Jc^1_a)$ such that ${}^{\mathrm{p}}\hc^i(Rf_*^{\xi,M-\parab}\Qlb)_{a,s}$ has a nonzero $s'$-isotypic component. Then, by definition of the $s$-part, the character $s'$ composed with the morphism $X_*(T\cap M_{\der})\to\pi_0(\Jc^1_a)$ gives back $s$. Now this last morphism factors through $\yc^1_a$ (cf.\ the commutative diagram~(\ref{eq:8.3.1})). But we have (cf.\ Proposition~\ref{thm:8.3.1} and its proof)
$$
\yc^1_a=\mathrm{Im}(X_*(T\cap L_{\der})_{W_a}\to X_*(T)_{W_a})=\Ker(X_*(T)_{W_a}\to X_*(L)).
$$
The morphism $X_*(T\cap M_{\der})\to\yc^1_a$ is dual to the morphism
$$
\Tc^{W_a}/Z^0_L\to\Tc/Z^0_M.
$$
The character $s'$ of $\pi_0(\Jc^1_a)$ induces a character of $\yc^1_a$, denoted $\tilde s$, which is an element of $\Tc^{W_a}/Z^0_L$ whose image in $\Tc/Z^0_M$ is $s$. Assertion 1 follows from this.
\end{preuve}

\begin{corollary}\label{thm:8.5.3}
Let $s\in(\Tc/Z^0_{\Mc})_{\tors}$ and $j:\Ac_G^{\el}\hookrightarrow\Ac_G$. We then have the following decomposition\,:
$$
j^*({}^{\mathrm{p}}\hc^i(Rf_*^{\xi,M-\parab}\Qlb)_s)=\bigoplus_{s'}{}^{\mathrm{p}}\hc^i(Rf_*^{\el}\Qlb)_{s'},
$$
where the sum is taken over the $s'\in(\Tc/Z^0_{\Gc})_{\tors}$ with image $s$ under the canonical map
$$
\Tc/Z^0_{\Gc}\to\Tc/Z^0_{\Mc}.
$$
\end{corollary}

\begin{preuve}
Over the open subset $\Ac^{\el}$, the group $\Gamma(\uc_M,\pi_0(\Jc^1))$ acts on
$$
j^*({}^{\mathrm{p}}\hc^i(Rf_*^{\xi,M-\parab}\Qlb))
$$
via the restriction morphism $\Gamma(\uc_M,\pi_0(\Jc^1))\to\Gamma(\Ac^{\el},\pi_0(\Jc^1))$. The corollary is then an obvious consequence of the commutative diagram
$$
\xymatrix{
X_*(T\cap M_{\der})\ar[r]\ar[d] & X_*(T\cap G_{\der})\ar[d]\\
\Gamma(\uc_M,\pi_0(\Jc^1))\ar[r] & \Gamma(\Ac^{\el},\pi_0(\Jc^1)) }
$$
which follows from Lemma~\ref{thm:8.4.1}.
\end{preuve}
\end{paragr}
\section{The open subset \texorpdfstring{$\Ac^{\mathrm{bon}}_G$}{A-bon-G}}\label{sec:9}

\begin{paragr}[The main statement.]\label{S:9.1}
In this section, we keep the notation of the preceding sections. In particular, $C$ is a projective, smooth and connected curve over $k$ of genus $g$ equipped with an effective divisor of degree $>2g$ and with a closed point $\infty$ outside the support of $D$. The pair $(G,T)$ consists of a connected reductive group $G$ over $k$ and a maximal subtorus $T$. The Weyl group $W^G$ is denoted simply by $W$ in this section. We also recall that the characteristic of $k$ does not divide the order of $W^G$. Let $\Ac_G$ be the smooth $k$-scheme defined in Section~\ref{S:4.4}. Following Ng\^o, we are going to define an upper semi-continuous function on $\Ac_G$, which will be denoted by $\delta$ (cf.\ \S\ref{S:9.15}). At $a\in\Ac_G(k)$, this function is given by
\[
\delta_a=\dim_k(H^0(C,(\tgo\otimes_k\pi_{Y_a,*}(\rho_{a,*}\oc_{X_a}/\oc_{Y_a}))^W)),
\]
where $\pi_Y:Y\to C$ is the cameral covering of $C$ associated with $a$ (cf.\ \S\ref{S:7.1}) and $\rho_a:X_a\to Y_a$ is the normalization morphism. We then lay down the following definition.

\begin{definition}\label{thm:9.1.1}
Let $\Ac_G^{\mathrm{bon}}\subset\Ac_G$ be the largest open subset of $\Ac_G$ such that, for every Zariski point $a\in\Ac_G^{\mathrm{bon}}$, the following inequality holds:
\[
\operatorname{codim}_{\Ac_G}(a)\geq\delta_a.
\]
\end{definition}

\begin{remark}\label{thm:9.1.2}
Since $\delta$ is semi-continuous, the condition $\delta_a=\delta$ defines a locally closed subset $\Ac_\delta$ of $\Ac_G$. By the noetherianity of $\Ac_G$, there are only finitely many $\delta$ for which $\Ac_\delta$ is nonempty. The open subset $\Ac_G^{\mathrm{bon}}$ is the complement of the irreducible components $\Ac'_\delta$ of the strata $\Ac_\delta$ which satisfy
\[
\operatorname{codim}_{\Ac_G}(\Ac'_\delta)<\delta.
\]
\end{remark}

\emph{A priori}, $\Ac_G^{\mathrm{bon}}$ may be empty. If the base field is of characteristic zero, Ng\^o has shown that one has in fact $\Ac_G^{\mathrm{bon}}=\Ac_G$. For our base field $k$, Theorem~\ref{thm:9.1.3} (and the remark that follows) shows that the open subset $\Ac_G^{\mathrm{bon}}$ is large enough. Before stating this theorem we are going to introduce certain divisors. Let $M$ be a reductive group equipped with an embedding of $T$, which is dual to a subgroup of $\Gc$ containing $\Tc$. In this context, we defined in Section~\ref{S:2.7} the discriminant $D^M$ on $\car_M$. It determines a divisor of $\car_M$ which is $\Gmk$-invariant. We deduce from it a divisor $\mathfrak{D}^M$ of $\car_{M,D}$. Let $B\in\pc(T)$ be a Borel subgroup of $G$. Again in~\S\ref{S:2.7}, we defined a homogeneous function $R_M^B$ on $\car_M$. Let $\mathfrak{R}_M^B$ be the divisor on $\car_{M,D}$ deduced from it. Since $R_M^B$ depends on the choice of $B$ only up to a sign, the divisor $\mathfrak{R}_M^B$ does not depend on the choice of $B$. For this reason, we denote it $\mathfrak{R}_M^G$.
\end{paragr}

\begin{theorem}\label{thm:9.1.3}
Let $M$ be a reductive group over $k$ dual to a subgroup of $\Gc$ containing $\Tc$ and let $a=(h_a,t)\in\Ac_M(k)$ be such that there exists a finite set $S$ of closed points of $C$ such that
\begin{enumerate}
\item the curve $h_a(C)$ traced in $\car_{M,D}$ satisfies, for every $c\in C-S$, one of the three following hypotheses:
  \begin{enumerate}
\item $h_a(c)$ does not belong to the divisor $\mathfrak{R}_M^G\cup\mathfrak{D}^M$\,; \item at the point $h_a(c)$, the curve $h_a(C)$ does not meet $\mathfrak{R}_M^G$ and meets transversally the divisor $\mathfrak{D}^M$, which is smooth at the point $h_a(c)$\,; \item at the point $h_a(c)$, the curve $h_a(C)$ does not meet $\mathfrak{D}^M$ and meets transversally the divisor $\mathfrak{R}_M^G$, which is smooth at the point $h_a(c)$\,;
  \end{enumerate}
\item for every $c$ in $S$, the intersection multiplicity of $h_a(C)$ with the divisor $\mathfrak{D}^G$ is strictly bounded above by
  \[
2|S|^{-1}|W|^{-1}(\deg(D)-2g+2).
  \]
\end{enumerate}
Then the image of $a$ in $\Ac_G$ belongs to the open subset $\Ac_G^{\mathrm{bon}}$.
\end{theorem}

\begin{remark}\label{thm:9.1.4}
Let $a=(h_a,t)\in\Ac_G(k)$ be such that $h_a(C)$ meets the divisor $\mathfrak{D}^G$ transversally. Then $a=(h_a,t)$ belongs to the open subset $\Ac_G^{\mathrm{bon}}$. In particular, $\Ac_G^{\mathrm{bon}}$ is nonempty.
\end{remark}

The rest of the section is devoted to the proof of this theorem.

\begin{paragr}[The moduli problem $\BBB_G$.]\label{S:9.2}
Inspired by the cameral curve of Donagi--Gaitsgory (cf.\ \cite{Donagi-Gaitsgory}), we are going to consider the following moduli problem $\BBB_G$. To a $k$-scheme $S$ we associate the set of isomorphism classes of commutative diagrams
\[
\xymatrix{
X \ar[r]^-{f} \ar[d]_{\pi} & \tgo_D\times_k S \ar[dl] \\
C\times_k S, }
\]
where $X$ is a projective and smooth curve over $S$, with not necessarily connected geometric fibers, equipped with an action of $W$, $\pi$ is a finite, $W$-invariant and flat morphism, and $f$ is a finite and $W$-equivariant morphism. Moreover, these data satisfy the following condition: the morphism $\pi$ is étale and Galois with Galois group $W$ above $\infty\times_k S$ and the morphism $f$ restricted to $\pi^{-1}(\infty\times_k S)$ is injective.

We leave it to the reader to verify the following representability statement.
\end{paragr}

\begin{lemma}\label{thm:9.2.1}
The moduli problem $\BBB_G$ is representable by a $k$-algebraic space of finite type.
\end{lemma}

\begin{paragr}[The scheme $\AAA_G$.]\label{S:9.3}
Let $\AAA_G$ be the smooth scheme such that for every $k$-scheme $S$ the set $\AAA_G(S)$ is the set of sections $a$ of $\car_D\times_k S$ over $C\times_k S$ which send the $S$-point $\infty\times_k S$ into the $G$-regular open subset $\car_D^{\reg}$. We have the following formula for the dimension of $\AAA_G$ (cf.\ \cite[Lemma 4.13.1]{Ngo-lemme}):
\begin{equation}\label{eq:9.3.1}
\dim_k(\AAA_G)=\deg(D)(\dim(G)+\rang(G))/2+\rang(G)(1-g).
\end{equation}
\end{paragr}

\begin{lemma}\label{thm:9.3.1}
Let $S$ be a $k$-scheme and $(X,\pi,f)\in\BBB_G(S)$.
\begin{enumerate}
\item The formation of the quotient scheme $X//W$ commutes with every base change $S'\to S$. \item The morphism $X//W\to C\times_k S$ induced by $\pi$ is an isomorphism. \item The morphism $f$ defines, by passing to the quotient by $W$, an element $a_f\in\AAA_G(S)$\,; the map $f\mapsto a_f$ defines a morphism, denoted $\psi$, from $\BBB_G$ to $\AAA_G$.
\end{enumerate}
\end{lemma}

\begin{preuve}
Let $S'\to S$. Since the order of $W$ is invertible in $k$, for every $S'$-algebra equipped with an action of $W$ over $S$, the functor of $W$-invariants is given by the average over $W$, which is $S'$-linear. Assertion 1 follows.

Let us next deal with assertion 2. The $\oc_{C\times_k S}$-module $\pi_*(\oc_X)$ is flat and of finite type, hence it is locally free. Since the order of $W$ is invertible in $k$, the $\oc_{C\times_k S}$-module $\lc=\pi_*(\oc_X)^W$ is a direct factor of $\pi_*(\oc_X)$. Consequently, $\lc$ is also a locally free $\oc_{C\times_k S}$-module. Since the morphism $\pi$ is étale Galois, with Galois group $W$, above $\infty\times_k S$, the unit section
\[
\eps:\oc_{C\times_k S}\to(\pi_*\oc_X)^W
\]
is an isomorphism in a neighbourhood of $\infty\times_k S$. In particular, $\lc$ is an invertible $\oc_{C\times_k S}$-module. Let us decompose the identity of $\lc$ as
\[
\lc\to\lc\otimes_{\oc_{C\times_k S}}\lc\to\lc,
\]
where the first arrow is $\eps\otimes 1$ and the second is the product of the algebra $\lc$. Since $\lc$ is invertible, so is $\lc\otimes_{\oc_{C\times_k S}}\lc$, and the two arrows above are necessarily isomorphisms. It follows that $\eps$ is an isomorphism and assertion 2 is proved.

Assertion 2 implies that $f$ defines, by passing to the quotient, an element $a_f\in\AAA_G(S)$. We indeed obtain a morphism by assertion 1, whence assertion 3.
\end{preuve}

\begin{paragr}[The invariant $\delta$ and the strata $\AAA_{G,\delta}$.]\label{S:9.4}
Let $S$ be a $k$-scheme and $(X,\pi,f)\in\BBB_G(S)$. By assertion 3 of Lemma~\ref{thm:9.3.1}, we deduce from it a section $a$ of $\car_D\times_k S$. Most of the time we shall leave $a$ implicit and shall write simply $\pi:Y\to C\times_k S$ for the cameral covering $\pi_a:Y_a\to C\times_k S$ of Section~\ref{S:7.1} (it is not to be confused with the universal cameral covering denoted $Y$ in \S\ref{S:7.1}). We thus have a commutative diagram
\begin{equation}\label{eq:9.4.1}
\xymatrix{
X \ar@/^1.4pc/[rr]^{f} \ar[r]_{\rho} \ar[dr]_{\pi}
  & Y \ar[r]_{f_Y} \ar[d]^{\pi_Y}
  & \tgo_D\times_k S \ar[d]^{\chi} \\
& C\times_k S \ar[r]^{a} \ar@{=}[dr]
  & \car_D\times_k S \ar[d]^{p} \\
& & C\times_k S
}
\end{equation}
with cartesian square, where
\begin{itemize}
\item[--] $\rho$ is the normalization morphism of the cameral curve $Y$\,; \item[--] $p$ is the canonical projection\,; \item[--] $\car_D=\car_{G,D}$\,; \item[--] $\chi$ is the morphism $\chi^T_{G,D}$ of Section~\ref{S:7.1}.
\end{itemize}
This diagram is obtained by the base change $S\to\BBB_G$ associated with $(X,\pi,f)$ from the universal diagram
\begin{equation}\label{eq:9.4.2}
\xymatrix{
\xc \ar@/^1.4pc/[rr]^{f_{\xc}} \ar[r]_{\rho_{\xc}} \ar[dr]_{\pi_{\xc}}
  & \yc \ar[r]_{f_{\yc}} \ar[d]^{\pi_{\yc}}
  & \tgo_D\times_k\BBB_G \ar[d]^{\chi} \\
& C\times_k\BBB_G \ar[r] \ar@{=}[dr]
  & \car_D\times_k\BBB_G \ar[d]^{p} \\
& & C\times_k\BBB_G,
}
\end{equation}
where $\xc$ the universal curve over $\BBB_G$ and $\yc$ is the universal cameral curve over $c\times_k\BBB_G$.
\end{paragr}

\begin{lemma}\label{thm:9.4.1}
The morphism $\psi:\BBB_G\to\AAA_G$ (defined in Lemma~\ref{thm:9.3.1}, assertion 3) induces a bijective map on the underlying topological spaces.
\end{lemma}

\begin{preuve}
Let us first show that the induced map on the underlying topological spaces is surjective on closed points. Let $a$ be such a closed point. Its residue field is the algebraically closed field $k$. In particular, $a$ is identified with a section of $\car_D$ over $C$. Let $Y$ be the associated cameral curve. It follows from the fact that $a$ is generically regular that $Y$ is geometrically reduced. Let $X$ be the normalization of $Y$. We then know that $X$ is smooth (cf.\ \cite[Prop.\ 17.15.14]{EGAIV4}). According to diagram (\ref{eq:9.4.1}), we obtain a triple $(X,\pi,f)$ which is a $k$-point of $\BBB_G$ with image $a$. This proves that all the closed points of $a$ are reached by the morphism $\BBB_G\to\AAA_G$.

Let $a$ be a non-closed point of the topological space underlying $\AAA_G$. Let $\bar a$ be its closure. We have to see that $a$ belongs to $\psi(\psi^{-1}(\bar a))$. By Chevalley's theorem, this is a constructible set. By \cite[Chap.\ 0, Prop.\ 9.2.3]{EGAIII}, we have the following alternative on $\bar a\cap\psi(\psi^{-1}(\bar a))$:
\begin{itemize}
\item[--] either this subset contains a nonempty open subset of $\bar a$ and hence $a$\,; \item[--] or this subset is rare in $\bar a$.
\end{itemize}
We have to exclude the second possibility. The latter means that the closure of $\psi(\psi^{-1}(\bar a))$ in $\bar a$ has empty interior. In particular, the closure of $\psi(\psi^{-1}(\bar a))$ in $\bar a$ is strictly smaller than $\bar a$. There is therefore a closed point of $\bar a$ which is not in the closure of $\psi(\psi^{-1}(\bar a))$. But this is not possible since, as we have seen, such a point always lies in $\psi(\psi^{-1}(\bar a))$. We have thus ruled out this second possibility and thereby finished the proof of surjectivity.

Injectivity is obvious since every triple $(X,\pi,f)$ of $\BBB_G$ with image $a\in\AAA_G$ is obtained, after a suitable extension of scalars, by normalization of the cameral curve associated with $a$.
\end{preuve}

Let $p_2:C\times_k\BBB_G\to\BBB_G$ be the projection onto the second factor. By the definition of $\BBB_G$, the $\oc_{\BBB_G}$-module
\[
p_{2,*}\pi_{\yc,*}(\rho_{\xc,*}\oc_{\xc}/\oc_{\yc})
\]
is locally free. Since the characteristic of $k$ does not divide the order of $W$, the same holds for the $\oc_{\BBB_G}$-module
\begin{equation}\label{eq:9.4.3}
(\tgo\otimes_k p_{2,*}\pi_{\yc,*}(\rho_{\xc,*}\oc_{\xc}/\oc_{\yc}))^W.
\end{equation}

\begin{definition}\label{thm:9.4.2}
Let $\delta:\BBB_G\to\NN$ be the locally constant function on $\BBB_G$ given by the rank of the locally free $\oc_{\BBB_G}$-module (\ref{eq:9.4.3}).
\end{definition}

Let $\delta:\AAA_G\to\NN$ be the constructible function on $\AAA_G$ defined for every $a\in\AAA_G$ by
\[
\delta(a)=\delta(b),
\]
where $b\in\BBB_G$ is the unique point such that $\psi(b)=a$.

\begin{remarks}\label{thm:9.4.3}
The existence and the uniqueness of $b$ were proved in Lemma~\ref{thm:9.4.1}. The constructibility of $\delta$ on $\AAA_G$ follows from Chevalley's theorem.

Let $K$ be an extension of $k$ and $(X,\pi,f)\in\BBB_G(K)$. Then, with the notation of diagram (\ref{eq:9.4.1}) for $S=\Spec(K)$, we have
\[
\delta(X,\pi,f)=\sum_c\delta_c(X,\pi,f),
\]
where the sum is taken over all the closed points $c$ of $C\times_k K$ and the local invariant $\delta_c$ is the following length:
\[
\delta_c(X,\pi,f)=\lng((\tgo\otimes_k\pi_{Y,*}(\rho_*\oc_X/\oc_Y))^W_c).
\]
This invariant vanishes at a point $c$ such that $Y$ is smooth at every point $y\in Y$ above $c$. We also have the following formula:
\[
\delta(X,\pi,f)=\dim_K(H^0(C\times_k K,(\tgo\otimes_k\pi_{Y,*}(\rho_*\oc_X/\oc_Y))^W)).
\]
\end{remarks}

\begin{proposition}\label{thm:9.4.4}
The function $\delta$ on $\AAA_G$ is upper semi-continuous.
\end{proposition}

\begin{preuve}
Since the function $\delta$ is constructible, it suffices to show that for every morphism from a trait $S$ to $\AAA_G$ one has $\delta(\eta)\leq\delta(s)$, where $s$ and $\eta$ are the closed and the generic points of $S$. Such a morphism from $S$ to $\AAA_G$ is interpreted as a section of $\car_{D,S}$ over $C\times_k S$. Let $Y$ be the relative cameral curve over $S$ deduced from it. Let $X$ be the normalization of $Y$. We are thus in the situation of diagram (\ref{eq:9.4.1}), whose notation we use. The generic fiber of $Y$ is generically reduced. Up to replacing $\eta$ by a radicial extension and $S$ by its normalization in this extension, we may assume that the generic fiber $X_\eta$ of $X$ is smooth (cf.\ \cite[Prop.\ 17.15.14]{EGAIV4}). The triple $(X_\eta,\pi_\eta,f_\eta)$ deduced from $(X,\pi,f)$ by base change to $\eta$ defines a point of $\BBB_G$, with values in the residue field $k(\eta)$ of $\eta$, above $a(\eta)$. We therefore have
\[
\delta(\eta)=\dim_{k(\eta)}((\tgo\otimes_k H^0(C_\eta,\pi_{Y_\eta,*}(\rho_{\eta,*}\oc_{X_\eta}/\oc_{Y_\eta})))^W).
\]
Since the $\oc_S$-module $\tgo\otimes_k H^0(C,\pi_{Y,*}(\rho_*\oc_X/\oc_Y))$ is flat, the $\oc_S$-module
\[
(\tgo\otimes_k H^0(C,\pi_{Y,*}(\rho_*\oc_X/\oc_Y)))^W,
\]
which is a direct factor of it, is flat too. It follows that we also have
\[
\delta(\eta)=\dim_{k(s)}((\tgo\otimes_k H^0(C_s,\pi_{Y_s,*}(\rho_{s,*}\oc_{X_s}/\oc_{Y_s})))^W),
\]
where we denote by an index $s$ the base change to $s$, and by $k(s)$ the residue field of $s$. Again, up to replacing $k(s)$ by a radicial extension, we may assume that the normalization, denoted $\tilde{X}_s$, of $Y_s$ is smooth. Let $\tilde{\rho}_s$ be the normalization morphism $\tilde{X}_s\to Y_s$. We then have
\[
\delta(s)=\dim_{k(s)}((\tgo\otimes_k H^0(C_s,\pi_{Y_s,*}(\tilde{\rho}_{s,*}\oc_{\tilde{X}_s}/\oc_{Y_s})))^W).
\]
Since we have $\rho_{s,*}\oc_{X,s}\subset\tilde{\rho}_{s,*}\oc_{\tilde{X}_s}$, we deduce the desired inequality $\delta(\eta)\leq\delta(s)$.
\end{preuve}

To conclude this paragraph, let us introduce the following definition.

\begin{definition}\label{thm:9.4.5}
For every natural number $\delta$, let $\AAA_{G,\delta}$ be the locally closed subset consisting of the $a\in\AAA_G$ such that $\delta(a)=\delta$.
\end{definition}

\begin{remarks}\label{thm:9.4.6}
By Proposition~\ref{thm:9.4.4}, the union $\bigcup_{\delta\geq\delta'}\AAA_{G,\delta}$ is closed. The open subset $\AAA_{G,0}$ contains in particular the open subset of the $a$ for which the cameral curve $Y_a$ is smooth.
\end{remarks}

In what follows, under certain conditions, we shall bound above the dimension of the strata $\AAA_{G,\delta}$.

\begin{paragr}[Cotangent complex for a point of $\BBB_G(k)$.]\label{S:9.5}
Let $(X,\pi,f)$ be a point of $\BBB_G(k)$. The deformation theory at this point for $\BBB_G$ is controlled by the complex
\[
R\!\Hom_{\oc_X}(L_{X/\tgo_D},\oc_X)^W,
\]
where the cotangent complex $L_{X/\tgo_D}$ is here a perfect complex of length $1$ concentrated in degrees $-1$ and $0$,
\[
L_{X/\tgo_D}=[f^*\Omega_{\tgo_D/k}\to\Omega^1_{X/k}]\cong[f^*\Omega_{\tgo_D/C}\to\Omega^1_{X/C}]
\]
whose arrow is $df$. This complex has two cohomology sheaves, namely
\[
\hc^{-1}(L_{X/\tgo_D})=\Ker(f^*\Omega^1_{\tgo_D/C}\to\Omega^1_{X/C})
\]
which is a vector bundle on $X$ of rank the rank of the group $G$, and
\[
\hc^{0}(L_{X/\tgo_D})=\Coker(f^*\Omega^1_{\tgo_D/C}\to\Omega^1_{X/C})=\Omega^1_{X/Y}
\]
which is a torsion $\oc_X$-module on $X$.
\end{paragr}

\begin{paragr}[Cotangent complex for a point of $\AAA_G(k)$.]\label{S:9.6}
In an analogous way, the deformation theory of $\AAA_G$ at a point $a\in\AAA_G(k)$ is controlled by the complex
\[
R\!\Hom_{\oc_Y}(L_{Y/\tgo_D},\oc_Y)^W,
\]
where $Y=Y_a$ is the cameral curve and
\[
L_{Y/\tgo_D}=[f_Y^*\Omega^1_{\tgo_D/k}\to\Omega^1_{Y/k}]\cong[f_Y^*\Omega^1_{\tgo_D/C}\to\Omega^1_{Y/C}]
\]
is a perfect complex concentrated in degrees $-1$ and $0$, the arrow being given by $df_Y$. This complex is quasi-isomorphic to $\hc^{-1}(L_{Y/\tgo_D})[1]$.
\end{paragr}

\begin{lemma}\label{thm:9.6.1}
We have
\[
\hc^{-1}(L_{Y/\tgo_D})=f_Y^*\chi^*\Omega^1_{\car_D/C}=\pi_Y^*a^*\Omega^1_{\car_D/C}.
\]
\end{lemma}

\begin{preuve}
The morphism $\chi$ induces the exact sequence
\begin{equation}\label{eq:9.6.1}
0\longrightarrow\chi^*\Omega^1_{\car_D/C}\longrightarrow\Omega^1_{\tgo_D/C}\longrightarrow\Omega^1_{\tgo_D/\car_D}\longrightarrow 0.
\end{equation}
Let us show the left exactness, the only one that is not obvious. Since $\chi$ is generically étale, the sequence is generically left exact. Since $\car_D$ is smooth over $C$, the $\oc_{\tgo_D}$-module $\chi^*\Omega^1_{\car_D/C}$ is locally free and cannot contain a nontrivial torsion module.

Since $\pi_Y$ is obtained from $\chi$ by the base change given by $a:C\to\car_D$, we have $a\pi_Y=\chi f_Y$ and
\begin{equation}\label{eq:9.6.2}
\Omega^1_{Y/C}=f_Y^*\Omega^1_{\tgo_D/\car_D}.
\end{equation}
Consequently, applying the functor $f_Y^*$ to the sequence (\ref{eq:9.6.1}), we obtain the exact sequence
\begin{equation}\label{eq:9.6.3}
0\longrightarrow f_Y^*\chi^*\Omega^1_{\car_D/C}\longrightarrow f_Y^*\Omega^1_{\tgo_D/C}\longrightarrow\Omega^1_{Y/C}\longrightarrow 0,
\end{equation}
where the left exactness is seen as before, whence the lemma.
\end{preuve}

\begin{lemma}\label{thm:9.6.2}
Let $a\in\AAA_G(k)$ and $Y=Y_a$.
\begin{enumerate}
\item There is no obstruction to deforming $Y\hookrightarrow\tgo_D$ and $\AAA_G$ is smooth at the point $a$; in other words we have
  \[
\Ext^2_{\oc_Y}(L_{Y/\tgo_D},\oc_Y)^W=0\,;
  \]
\item the tangent space to $\AAA_G$ at the point $a$ is the space
  \[
\Ext^1(L_{Y/\tgo_D},\oc_Y)^W
  \]
which is of dimension $\dim(\AAA_G)$.
\end{enumerate}
\end{lemma}

\begin{preuve}
Let $i\in\{0,1\}$. Since $\hc^{-1}(L_{Y/\tgo_D})$ is locally free, we have
\begin{align*}
\Ext^{i+1}(L_{Y/\tgo_D},\oc_Y)&=\Ext^i(\hc^{-1}(L_{Y/\tgo_D}),\oc_Y)\\
&=H^i(Y,\homc_{\oc_Y}(\hc^{-1}(L_{Y/\tgo_D}),O_Y)).
\end{align*}
By Lemma~\ref{thm:9.6.1}, we have
\[
\homc_{\oc_Y}(\hc^{-1}(L_{Y/\tgo_D}),\oc_Y)=\pi_Y^*\,\homc_{\oc_C}(a^*\Omega^1_{\car_D/C},\oc_C).
\]
Let $n$ be the rank of $G$. There exist $n$ homogeneous elements $(t_1,\dots,t_n)$ of $k[\tgo]$ of degrees $e_1,\dots,e_n$ such that $k[\tgo]^W=k[t_1,\dots,t_n]$ and $2\sum_{i=1}^n e_i=\dim(G)+n$ (cf.\ property 33 of the summary of \cite{Bourbaki}). Using these elements, we obtain an identification
\[
a^*\Omega^1_{\car_D/C}\cong\bigoplus_{i=1}^n\oc_C(-e_iD)
\]
hence an isomorphism
\[
\homc(L_{Y/\tgo_D},O_Y)=\pi_Y^*\bigoplus_{i=1}^n\oc_C(e_iD).
\]
Thus we have
\[
\Ext^{i+1}(L_{Y/\tgo_D},\oc_Y)=H^i(Y,\pi_Y^*\bigoplus_{i=1}^n\oc_C(e_iD))=H^i(C,\pi_{Y,*}\pi_Y^*\bigoplus_{i=1}^n\oc_C(e_iD)).
\]
By the projection formula, we have
\[
\pi_{Y,*}\pi_Y^*\bigoplus_{i=1}^n\oc_C(e_iD)=\pi_{Y,*}\oc_Y\otimes_{\oc_C}\bigoplus_{i=1}^n\oc_C(e_iD).
\]
Passing to the invariants under $W$, we find
\begin{align*}
\Ext^{i+1}(L_{Y/\tgo_D},\oc_Y)^W&=H^i(C,(\pi_{Y,*}\oc_Y)^W\otimes_{\oc_C}\bigoplus_{i=1}^n\oc_C(e_iD))\\
&=H^i(C,\bigoplus_{i=1}^n\oc_C(e_iD))
\end{align*}
since $(\pi_{Y,*}\oc_Y)^W=\oc_C$. The result follows from the Riemann--Roch theorem.
\end{preuve}

\begin{paragr}[The invariant $\kappa$.]\label{S:9.7}
Let us consider the universal situation of the commutative diagram (\ref{eq:9.4.2}). Let $L_{\xc/\tgo_D\times\BBB_G}$ be the cotangent complex of $\xc$ relative to $\tgo_D\times\BBB_G$. Let
\[
\kc=p_{2,*}(\pi_{\xc,*}(\hc^{-1}(L_{\xc/\tgo_D\times\BBB_G})\otimes_{\oc_{\xc}}\Omega^1_{\xc/C\times\BBB_G}))^W
\]
which we view as a coherent $\oc_{\BBB_G}$-module on $\BBB_G$ (cf.\ \cite[Cor.\ 2.3.11]{Illusie}).
\end{paragr}

\begin{definition}\label{thm:9.7.1}
Let $\kappa:\BBB_G\to\NN$ be the upper semi-continuous function defined for every $b\in\BBB_G$ by
\[
\kappa(b)=\dim_{k(b)}(\kc\otimes_{\oc_{\BBB_G}}k(b)),
\]
where $k(b)$ is the residue field of $b$.
\end{definition}

Let $\kappa:\AAA_G\to\NN$ be the constructible function defined by
\[
\kappa(a)=\kappa(b),
\]
where $b$ is the unique point of $\BBB_B$ with image $a$ under the bijection $\psi$ of Lemma~\ref{thm:9.4.1}.

\begin{remarks}\label{thm:9.7.2}
The semi-continuity of $\kappa$ on $\BBB_G$ is an obvious consequence of Nakayama's lemma. The constructibility of $\kappa$ on $\AAA_G$ follows from Chevalley's theorem. In general, the function $\kappa$ is not semi-continuous on $\AAA_G$. One can show however that it is semi-continuous on the strata $\AAA_{G,\delta}$ of Definition~\ref{thm:9.4.5}.
\end{remarks}

Let $(X,\pi,f)\in\BBB_G(k)$. For every closed point $c$ of $C$, let
\begin{equation}\label{eq:9.7.1}
\kappa_c(X,\pi,f)=\lng(\pi_*(\hc^{-1}(L_{X/\tgo_D})\otimes_{\oc_X}\Omega^1_{X/C})^W_c).
\end{equation}

\begin{remark}\label{thm:9.7.3}
For every $\oc_X$-module $\lc$, we denote, by abuse, by $\pi_*(\lc)^W$ what should more correctly be denoted $(\pi_*\lc)^W$.
\end{remark}

We then have
\begin{equation}\label{eq:9.7.2}
\kappa(X,\pi,f)=\sum_{c\in C}\kappa_c(X,\pi,f).
\end{equation}
This sum is of course with finite support. More precisely, we have the following proposition.

\begin{proposition}\label{thm:9.7.4}
For every closed point $c$ of $C$ at which one of the two following conditions is realized:
\begin{itemize}
\item[--] $X$ is étale above $c$\,; \item[--] $Y$ is smooth above $c$,
\end{itemize}
one has $\kappa_c(X,\pi,f)=0$.
\end{proposition}

\begin{preuve}
Let $x$ be a closed point of $X$ such that $\pi$ is étale at $x$. We thus have
\[
\Omega^1_{X/C,x}=0
\]
and $\kappa_c(X,\pi,f)=0$ for every closed point $c$ of $C$ such that $\pi$ is étale above $c$.

Let us next place ourselves at a point $c$ above which $Y$ is smooth. In this case, locally above $c$ we have $X=Y$, $\pi=\pi_Y$ and $\hc^{-1}(L_Y)=\hc^{-1}(L_X)$. Using Lemma~\ref{thm:9.6.1}, we thus have, locally above $c$,
\[
\hc^{-1}(L_{X/\tgo_D})=\pi^*a^*\Omega^1_{\car_D/C}.
\]
The projection formula then gives
\[
(\pi_*(\hc^{-1}(L_{X/\tgo_D})\otimes_{\oc_X}\Omega^1_{X/C}))^W_c=(a^*\Omega^1_{\car_D/C}\otimes_{\oc_C}(\pi_*(\Omega^1_{X/C}))^W)_c
\]
which vanishes since
\begin{equation}\label{eq:9.7.3}
\pi_*(\Omega^1_{X/C}))^W=0.
\end{equation}
Indeed, this last module is the quotient of $\pi_*(\Omega^1_{X/k}))^W=\Omega^1_{C/k}$ (cf.\ assertion 2 of Lemma~\ref{thm:9.11.1}) by
\[
\pi_*(\pi^*\Omega^1_{C/k}))^W=\Omega^1_{C/k}\otimes_{\oc_C}\pi_*(\oc_X)^W=\Omega^1_{C/k}
\]
since $\pi_*(\oc_X)^W=\oc_C$.
\end{preuve}

\begin{paragr}[Smoothness of $\BBB_G$.]\label{S:9.8}
Here is a condition, in terms of the function $\kappa$, sufficient for a point of $\BBB_G(k)$ to be smooth.
\end{paragr}

\begin{proposition}\label{thm:9.8.1}
Let $(X,\pi,f)\in\BBB_G(k)$ be such that
\[
\deg(D)>2g-2+\kappa(X,\pi,f).
\]
Then one has
\[
\Ext^2_{\oc_X}(L_{X/\tgo_D},\oc_X)^W=(0)
\]
and $\BBB_G$ is therefore smooth at such a point.
\end{proposition}

\begin{preuve}
We have
\[
\Ext^2_{\oc_X}(L_{X/\tgo_D},\oc_X)=\Ext^1_{\oc_X}(\hc^{-1}(L_{X/\tgo_D}),\oc_X).
\]
By Serre duality, this $k$-vector space is dual to $H^0(X,\hc^{-1}(L_{X/\tgo_D})\otimes_{\oc_X}\Omega^1_{X/k})$. Passing to the module of invariants under $W$ (which is isomorphic to the module of co-invariants since the order of $W$ is invertible in $k$), one sees that the desired vanishing is equivalent to
\[
H^0(X,\hc^{-1}(L_{X/\tgo_D})\otimes_{\oc_X}\Omega^1_{X/k})^W=0
\]
or again
\begin{equation}\label{eq:9.8.1}
H^0(C,F_L(\Omega^1_{X/k}))=0,
\end{equation}
where $F_L$ is the exact functor from the category of $\oc_X$-modules to that of $\oc_C$-modules defined by
\[
F_L(\lc)=\pi_*(\hc^{-1}(L_{X/\tgo_D})\otimes_{\oc_X}\lc)^W.
\]
Let us apply the functor $F_L$ to the exact sequence of $\oc_X$-modules
\begin{equation}\label{eq:9.8.2}
0\longrightarrow\pi^*\Omega^1_{C/k}\longrightarrow\Omega^1_{X/k}\longrightarrow\Omega^1_{X/C}\longrightarrow 0.
\end{equation}
We obtain the exact sequence of $\oc_C$-modules
\[
0\longrightarrow F_L(\pi^*\Omega^1_{C/k})\longrightarrow F_L(\Omega^1_{X/k})\longrightarrow F_L(\Omega^1_{X/C})\longrightarrow 0.
\]
The module $F_L(\Omega^1_{X/C})$ is torsion and its length is $\kappa(X,\pi,f)$, by the very definition of this invariant. There is therefore an effective divisor $K$ on $C$, of degree bounded above by $\kappa(X,\pi,f)$, such that
\[
F_L(\pi^*\Omega^1_{C/k})\subset F_L(\Omega^1_{X/k})\subset F_L(\pi^*\Omega^1_{C/k})(K).
\]
To prove (\ref{eq:9.8.1}), it therefore suffices to prove that one has
\begin{equation}\label{eq:9.8.3}
H^0(C,F_L(\pi^*\Omega^1_{C/k})(K))=0.
\end{equation}
Now, by the projection formula, we have
\begin{equation}\label{eq:9.8.4}
F_L(\pi^*\Omega^1_{C/k})=\pi_*(\hc^{-1}(L_{X/\tgo_D}))^W\otimes_{\oc_C}\Omega^1_{C/k}.
\end{equation}
From the obvious injection $\hc^{-1}(L_{X/\tgo_D})\hookrightarrow f^*\Omega^1_{\tgo_D/C}$, we deduce the injections
\begin{equation}\label{eq:9.8.5}
\pi_*(\hc^{-1}(L_{X/\tgo_D}))^W\hookrightarrow\pi_*(\hc^{-1}(L_{X/\tgo_D}))\hookrightarrow\pi_*(f^*\Omega^1_{\tgo_D/C})
\end{equation}
whence, by tensoring with $\Omega^1_{C/k}(K)$ and by the isomorphism (\ref{eq:9.8.4}), an injection
\begin{equation}\label{eq:9.8.6}
F_L(\pi^*\Omega^1_{C/k})(K)\hookrightarrow\pi_*(f^*\Omega^1_{\tgo_D/C})\otimes_{\oc_C}\Omega^1_{C/k}(K).
\end{equation}
We are going to prove that one has
\begin{equation}\label{eq:9.8.7}
H^0(C,\pi_*(f^*\Omega^1_{\tgo_D/C})\otimes_{\oc_C}\Omega^1_{C/k}(K))=0,
\end{equation}
which will prove (\ref{eq:9.8.3}) and hence (\ref{eq:9.8.1}). The above vanishing is equivalent to
\begin{equation}\label{eq:9.8.8}
H^0(X,f^*\Omega^1_{\tgo_D/C}\otimes_{\oc_X}\pi^*\Omega^1_{C/k}(K))=0.
\end{equation}

Let $\tgo^\vee$ be the dual of the $k$-vector space $\tgo$. We then have an isomorphism
\[
\Omega^1_{\tgo_D/C}=\tgo^\vee\otimes_k\chi^*p^*\oc_C(-D)
\]
whence
\[
f^*\Omega^1_{\tgo_D/C}=\tgo^\vee\otimes_k\pi^*\oc_C(-D).
\]
Hence the vanishing (\ref{eq:9.8.8}) is equivalent to the vanishing
\begin{equation}\label{eq:9.8.9}
H^0(X,\pi^*(\Omega^1_{C/k}(K-D)))=0.
\end{equation}
But, for this last one, it suffices to have
\[
\deg(\Omega^1_{C/k}(K-D))<0
\]
that is
\[
\deg(K)<\deg(D)-2g+2.
\]
Since $\deg(K)\leq\kappa(X,\pi,f)$, the condition
\[
\kappa(X,\pi,f)<\deg(D)-2g+2
\]
is sufficient in order to have the vanishing (\ref{eq:9.8.9}), hence the vanishing (\ref{eq:9.8.1}) and hence $\Ext^2_{\oc_X}(L_{X/\tgo_D},\oc_X)^W=(0)$.
\end{preuve}
\begin{paragr}[Estimating the dimension of $\BBB_G$ at smooth points.]\label{S:9.9}
The aim of this section is to prove the following proposition.

\begin{proposition}\label{thm:9.9.1}
Let $(X,\pi,f)\in\BBB_G(k)$ be such that
$$\Ext^2(L_{X/\tgo_D},\oc_X)^W=0.$$
We have the inequality
$$\dim_k(\AAA_G)-\dim_k(\Ext^1(L_{X/\tgo_D},\oc_X)^W)\geq\delta(X,\pi,f).$$
\end{proposition}

In the proof of Proposition~\ref{thm:9.9.1}, it will be convenient to have the following definition at our disposal.

\begin{definition}\label{thm:9.9.2}
Let $F_\Omega$ be the functor from the category of $\oc_X$-modules to that of $\oc_C$-modules defined, for every $\oc_X$-module $\lc$, by
$$F_\Omega(\lc)=\pi_*(\lc\otimes_{\oc_X}\Omega^1_{X/k})^W.$$
\end{definition}

\begin{remark}\label{thm:9.9.3}
The functor $F_\Omega$ is exact.
\end{remark}

\begin{preuve}
Under the condition
$$\Ext^2_{\oc_X}(L_{X/\tgo_D},\oc_X)^W=0,$$
the desired inequality is equivalent to
$$\deg(F_\Omega(\rho^*\Omega^1_{Y/C}))-\deg(F_\Omega(\Omega^1_{X/C}))\geq\delta(X,\pi,f)$$
by Proposition~\ref{thm:9.10.1} of \S\ref{S:9.10} below. Combining Lemmas~\ref{thm:9.10.2}, \ref{thm:9.10.3} and~\ref{thm:9.10.4} of \S\ref{S:9.10} below, we obtain
$$\deg(F_\Omega(\rho^*\Omega^1_{Y/C}))=\delta(X,\pi,f)+\deg((\tgo^\vee\otimes_k\pi_*\Omega^1_{X/C})^W).$$
The desired inequality is therefore equivalent to
$$\deg((\tgo^\vee\otimes_k\pi_*\Omega^1_{X/C})^W)-\deg(F_\Omega(\Omega^1_{X/C}))\geq 0.$$
For this it suffices to check that we have
$$\dim_k((\tgo^\vee\otimes_k\pi_*\Omega^1_{X/C})^W_c)\geq\dim_k(F_\Omega(\Omega^1_{X/C})_c)$$
for every $c\in C(k)$. Now this last inequality follows from assertions 3 and 4 of Lemma~\ref{thm:9.11.1} of \S\ref{S:9.11}.
\end{preuve}
\end{paragr}

\begin{paragr}\label{S:9.10}
In this paragraph we state and prove some results which occur in the proof of Proposition~\ref{thm:9.9.1}.

\begin{proposition}\label{thm:9.10.1}
Under the condition
$$\Ext^2_{\oc_X}(L_{X/\tgo_D},\oc_X)^W=0,$$
we have
$$\dim_k(\Ext^1_{\oc_X}(L_{X/\tgo_D},\oc_X)^W)=\dim(\AAA_G)+\deg(F_\Omega(\Omega^1_{X/C}))-\deg(F_\Omega(\rho^*\Omega^1_{Y/C})).$$
\end{proposition}

\begin{preuve}
For every $k$-vector space $V$, we denote by $V^\vee$ its dual. By Serre duality, we have
$$\Ext^i_{\oc_X}(L_{X/\tgo_D},\oc_X)^\vee=H^{1-i}(X,L_{X/\tgo_D}\otimes_{\oc_X}\Omega^1_{X/k})$$
and therefore
$$(\Ext^i_{\oc_X}(L_{X/\tgo_D},\oc_X)^\vee)^W=H^{1-i}(C,F_\Omega(L_{X/\tgo_D})).$$
The hypothesis of the statement is equivalent to
$$H^{-1}(C,F_\Omega(L_{X/\tgo_D}))=0.$$
We also have, for $i\geq 1$,
$$H^i(C,F_\Omega(L_{X/\tgo_D}))=0.$$
The dimension we are looking for, namely $\dim_k(\Ext^1_{\oc_X}(L_{X/\tgo_D},\oc_X)^W)$, is equal to the Euler-Poincar\'e characteristic, denoted $\chi(C,F_\Omega(L_{X/\tgo_D}))$, of $R\Gamma(C,F_\Omega(L_{X/\tgo_D}))$. We have a distinguished triangle
$$\rho^*L_{Y/\tgo_D}\to L_{X/\tgo_D}\to[\rho^*\Omega^1_{Y/C}\to\Omega^1_{X/C}]\to$$
and therefore a distinguished triangle
$$F_\Omega(\rho^*L_{Y/\tgo_D})\to F_\Omega(L_{X/\tgo_D})\to[F_\Omega(\rho^*\Omega^1_{Y/C})\to F_\Omega(\Omega^1_{X/C})]\to,$$
where the last complex is concentrated in degrees $-1$ and $0$. Consequently we have the equality between Euler-Poincar\'e characteristics
$$\chi(C,F_\Omega(L_{X/\tgo_D}))=\chi(C,F_\Omega(\rho^*L_{Y/\tgo_D}))+\chi(C,F_\Omega(\Omega^1_{X/C}))-\chi(C,F_\Omega(\rho^*\Omega^1_{Y/C})).$$
By the Riemann-Roch formula, the difference
$$\chi(C,F_\Omega(\Omega^1_{X/C}))-\chi(C,F_\Omega(\rho^*\Omega^1_{Y/C}))$$
is equal to $\deg(F_\Omega(\Omega^1_{X/C}))-\deg(F_\Omega(\rho^*\Omega^1_{Y/C}))$. It remains to compute
$$\chi(C,F_\Omega(\rho^*L_{Y/\tgo_D}))=-\chi(C,F_\Omega(\rho^*\hc^{-1}(L_{Y/\tgo_D}))).$$
By Lemma~\ref{thm:9.6.1}, we have
$$\rho^*\hc^{-1}(L_{Y/\tgo_D})=\rho^*\pi_Y^*a^*\Omega^1_{\car_D/C}=\pi^*a^*\Omega^1_{\car_D/C}.$$
By the projection formula, we therefore have
$$F_\Omega(\rho^*\hc^{-1}(L_{Y/\tgo_D}))=a^*\Omega^1_{\car_D/C}\otimes_{\oc_C}\pi_*(\Omega^1_{X/k})^W.$$
Now $(\pi_*\Omega^1_{X/k})^W=\Omega^1_{C/k}$ (cf.\ Lemma~\ref{thm:9.11.1}, assertion 2) and one may identify $a^*\Omega^1_{\car_D/C}$ with $\bigoplus_{i=1}^n\oc_C(-e_iD)$ (cf.\ the proof of~\ref{thm:9.6.2}). Hence, using the Riemann-Roch formula, we have
\begin{align*}
\chi(C,F_\Omega(\rho^*L_{Y/\tgo_D}))&=-\sum_{i=1}^n\chi(C,\Omega^1_{C/k}(-e_iD))\\
&=(e_1+\cdots+e_n)\deg(D)-n(g-1)\\
&=\dim(\AAA_G).
\end{align*}
This completes the proof.
\end{preuve}

\begin{lemma}\label{thm:9.10.2}
We have the following equality
$$\deg(F_\Omega(\rho^*\Omega^1_{Y/C}))=\deg(F_\Omega(f^*\Omega^1_{\tgo_D/C}))-\deg(F_\Omega(f^*\chi^*\Omega^1_{\car_D/C})).$$
\end{lemma}

\begin{preuve}
Let us apply the functor $f^*$ to the exact sequence~(\ref{eq:9.6.1}). We obtain the sequence
$$0\longrightarrow f^*\chi^*\Omega^1_{\car_D/C}\longrightarrow f^*\Omega^1_{\tgo_D/C}\longrightarrow f^*\Omega^1_{\tgo_D/\car_D}\longrightarrow 0$$
which is exact (left exactness is seen as that of sequence~(\ref{eq:9.6.3})). It follows from~(\ref{eq:9.6.2}) that $f^*\Omega^1_{\tgo_D/\car_D}=\rho^*\Omega^1_{Y/C}$. Applying the exact functor $F_\Omega$ to the above sequence, we obtain the desired relation.
\end{preuve}

\begin{lemma}\label{thm:9.10.3}
We have
\begin{multline*}
\deg(F_\Omega(f^*\Omega^1_{\tgo_D/C}))-\deg((\tgo^\vee\otimes_k\pi_*(\oc_X))^W\otimes_{\oc_C}\Omega^1_{C/k}(-D))\\
=\deg((\tgo^\vee\otimes_k\pi_*\Omega^1_{X/C})^W).
\end{multline*}
\end{lemma}

\begin{preuve}
Let us identify $\Omega^1_{\tgo_D/C}$ with $\chi^*p^*(\tgo^\vee\otimes_k\oc_C(-D))$ and $f^*\Omega^1_{\tgo_D/C}$ with $\tgo^\vee\otimes_k\pi^*(\oc_C(-D))$. By the projection formula, we have
\begin{equation}\label{eq:9.10.1}
F_\Omega(f^*\Omega^1_{\tgo_D/C})=(\tgo^\vee\otimes_k\pi_*(\Omega^1_{X/k}))^W\otimes_{\oc_C}\oc_C(-D).
\end{equation}
and
\begin{equation}\label{eq:9.10.2}
(\tgo^\vee\otimes_k\pi_*\oc_X)^W\otimes_{\oc_C}\Omega^1_{C/k}=\pi_*(\tgo^\vee\otimes_k\pi^*\Omega^1_{C/k})^W.
\end{equation}

Starting from the exact sequence~(\ref{eq:9.8.2}), we see that we have an exact sequence
$$
\begin{aligned}
0\longrightarrow\pi_*(\tgo^\vee\otimes_k\pi^*\Omega^1_{C/k})^W(-D)&\longrightarrow\pi_*(\tgo^\vee\otimes_k\Omega^1_{X/k})^W(-D)\\
&\longrightarrow\pi_*(\tgo^\vee\otimes_k\Omega^1_{X/C})^W(-D)\longrightarrow 0
\end{aligned}
$$
which we rewrite, using~(\ref{eq:9.10.1}) and~(\ref{eq:9.10.2}), as
$$
\begin{aligned}
0\longrightarrow F_\Omega(f^*\Omega^1_{\tgo_D/C})&\longrightarrow(\tgo^\vee\otimes_k\pi_*\oc_X)^W\otimes_{\oc_C}\Omega^1_{C/k}(-D)\\
&\longrightarrow\pi_*(\tgo^\vee\otimes_k\Omega^1_{X/C})^W(-D)\longrightarrow 0.
\end{aligned}
$$
The announced equality follows from this since we have
$$\deg(\pi_*(\tgo^\vee\otimes_k\Omega^1_{X/C})^W(-D))=\deg(\pi_*(\tgo^\vee\otimes_k\Omega^1_{X/C})^W)$$
in view of the fact that $\pi_*(\tgo^\vee\otimes_k\Omega^1_{X/C})^W$ is a torsion $\oc_C$-module.
\end{preuve}

\begin{lemma}\label{thm:9.10.4}
We have
$$\deg((\tgo^\vee\otimes_k\pi_*(\oc_X))^W\otimes_{\oc_C}\Omega^1_{C/k}(-D))-\deg(F_\Omega(f^*\chi^*\Omega^1_{\car_D/C}))=\delta(X,\pi,f).$$
\end{lemma}

\begin{preuve}
We shall show the following isomorphism
\begin{equation}\label{eq:9.10.3}
F_\Omega(f^*\chi^*\Omega^1_{\car_D/C})=(\tgo^\vee\otimes_k\pi_{Y,*}\oc_Y)^W\otimes_{\oc_C}\Omega^1_{C/k}(-D)
\end{equation}
which immediately gives the result.

Let us start from the relation (cf.\ Lemma~\ref{thm:9.11.1}, assertion 1)
$$\Omega^1_{\car_D/C}=\chi_*(\Omega^1_{\tgo_D/C})^W.$$
Now $\Omega^1_{\tgo_D/C}$ is identified with $\tgo^\vee\otimes_k\chi^*p^*\oc_C(-D)$. We therefore have
$$\Omega^1_{\car_D/C}=(\tgo^\vee\otimes_k\chi_*\oc_{\tgo_D})^W\otimes_{\oc_{\car_D}}p^*\oc_C(-D).$$
Consequently we have
\begin{align}
a^*\Omega^1_{\car_D/C}&=(\tgo^\vee\otimes_k a^*\chi_*\oc_{\tgo_D})^W\otimes_{\oc_C}\oc_C(-D)\label{eq:9.10.4}\\
&=(\tgo^\vee\otimes_k\pi_{Y,*}\oc_Y)^W\otimes_{\oc_C}\oc_C(-D)\nonumber
\end{align}
since $a^*\chi_*\oc_{\tgo_D}=\pi_{Y,*}\oc_Y$. Now we have
$$f^*\chi^*\Omega^1_{\car_D/C}=\pi^*a^*\Omega^1_{\car_D/C}.$$
Consequently, by the projection formula, we have
\begin{align}
F_\Omega(f^*\chi^*\Omega^1_{\car_D/C})&=a^*\Omega^1_{\car_D/C}\otimes_{\oc_C}\pi_*(\Omega^1_{X/k})^W\label{eq:9.10.5}\\
&=a^*\Omega^1_{\car_D/C}\otimes_{\oc_C}\Omega^1_{C/k},\nonumber
\end{align}
the last equality resulting from $\pi_*(\Omega^1_{X/k})^W=\Omega^1_{C/k}$ (cf.\ Lemma~\ref{thm:9.11.1}, assertion 2). Combining lines~(\ref{eq:9.10.4}) and~(\ref{eq:9.10.5}), we do obtain~(\ref{eq:9.10.3}).
\end{preuve}
\end{paragr}

\begin{paragr}[Some auxiliary lemmas.]\label{S:9.11}
We collect in the following lemma some results which have been used in what precedes.

\begin{lemma}\label{thm:9.11.1}
Let $c\in C(k)$. We have the following relations:
\begin{enumerate}
\item $\chi_*(\Omega^1_{\tgo_D/C})^W=\Omega^1_{\car_D/C}$; \item $\pi_*(\Omega^1_{X/k})^W=\Omega^1_{C/k}$; \item $\dim_k((\tgo^\vee\otimes_k\pi_*\Omega^1_{X/C})^W_c)=\dim(\tgo)-\dim(\tgo^{W_x})$ where $W_x\subset W$ is the stabilizer of a point $x\in X(k)$ above $c$; \item $\dim_k(F_\Omega(\Omega^1_{X/C})_c)=\left\{\begin{array}{ll}
0&\text{if $X$ is \'etale above $c$};\\
1&\text{otherwise.}
\end{array}\right.$
\end{enumerate}
\end{lemma}

\begin{preuve}
Let us prove assertion 1. We deduce from~(\ref{eq:9.6.1}) the exact sequence
$$0\longrightarrow(\chi_*\chi^*\Omega^1_{\car_D/C})^W\longrightarrow(\chi_*\Omega^1_{\tgo_D/C})^W.$$
But, by the projection formula, we see that the adjunction map $\Omega^1_{\car_D/C}\to\chi_*\chi^*\Omega^1_{\car_D/C}$ identifies $\Omega^1_{\car_D/C}$ with $(\chi_*\chi^*\Omega^1_{\car_D/C})^W$. Choosing a trivialization of the bundle $\lc_D$ in a neighborhood of a point of $C(k)$, we are reduced to proving
$$\chi_*(\Omega^1_{\tgo/k})^W=\Omega^1_{\car/k},$$
where we denote by $\chi$ the canonical morphism $\tgo\to\car$. Let $k(\tgo)$ be the field of fractions of the ring of regular functions on $\tgo$. Since $\chi$ is generically \'etale, we have
$$k(\tgo)\otimes_{k[\tgo]}\Omega^1_{\tgo/k}=k(\tgo)\otimes_{k[\tgo]^W}\Omega^1_{\car/k}$$
whence
\begin{equation}\label{eq:9.11.1}
(k(\tgo)\otimes_{k[\tgo]}\Omega^1_{\tgo/k})^W=(k(\tgo)\otimes_{k[\tgo]^W}\Omega^1_{\car/k})^W=k(\tgo)^W\otimes_{k[\tgo]^W}\Omega^1_{\car/k}.
\end{equation}
Let $t_1,\ldots,t_n$ and $\sigma_1,\ldots,\sigma_n$ be homogeneous elements of $k[\tgo]$ such that $k[\tgo]\simeq k[t_1,\ldots,t_n]$ and $k[\tgo]^W\simeq k[\sigma_1,\ldots,\sigma_n]$. Let $d\Sigma\in\chi_*(\Omega^1_{\tgo/k})^W$. By~(\ref{eq:9.11.1}), there exist $b_1,\ldots,b_n$ in $k(\tgo)^W$ such that $d\Sigma=\sum_{i=1}^nb_id\sigma_i$. Let $i$ be an integer $1\leq i\leq n$. The form
$$d\sigma_1\wedge\cdots\wedge d\Sigma\wedge\cdots\wedge d\sigma_n\in\bigwedge^n\Omega^1_{\tgo/k},$$
where $d\Sigma$ replaces $d\sigma_i$, is invariant under $W$. On the other hand, this form is equal to
$$b_id\sigma_1\wedge\cdots\wedge d\sigma_n=b_iJdt_1\wedge\cdots\wedge dt_n,$$
where $J$ is the Jacobian determinant $\det(\frac{\partial\sigma_i}{\partial t_j})$. We therefore have $b_iJ\in k[\tgo]$. The invariance under $W$ of the above form implies that $b_jJ$ is anti-invariant, that is to say we have, for every $w\in W$,
$$w(b_iJ)=\det(w)b_iJ,$$
where $\det(w)$ is the determinant of $w$ acting on $\tgo$. Now the set of anti-invariant elements of $k[\tgo]$ is $k[\tgo]^WJ$ (cf.\ \cite[Prop.~6, Chap.~V]{Bourbaki}). This proves that $b_i\in k[\tgo]^W$, whence assertion 1.

Let us pass to the other assertions. Throughout the proof we shall use the following notation. We denote by $c$ a point of $C(k)$. Let $\oc_{C,c}$ be the completion of the local ring of $C$ at $c$ and $F_c$ the field of fractions of $\oc_{C,c}$. The $F_v$-algebra $\oc_X\times_{\oc_C}F_c$ is a finite product of fields which are finite tamely ramified extensions of $F_v$. They correspond bijectively to the points of $X(k)$ above $c$. Let $x$ be such a point and $F_x$ the corresponding field. The integral closure of $\oc_{C,c}$ in $F_x$ is the completion $\oc_{X,x}$ of the local ring of $X$ at $x$. The Weyl group $W$ acts on the fiber of $X$ above $c$. Let $W_x\subset W$ be the stabilizer of $x$. Since the characteristic of $k$ does not divide the order of $W$, this subgroup is cyclic; it depends on the choice of $x$ only up to conjugation. The choice of a uniformizer $\varpi$ of $\oc_{X,x}$ makes it possible to identify this ring with the formal power series ring $k[[\varpi]]$, and the action of $W_x$ on $\oc_{X,x}$ is identified with the continuous action of $W_x$ on $k[[\varpi]]$ which acts by a primitive character $\zeta$ on the uniformizer $\varpi$. Let $m$ be the order of $W_x$. Then the subring $\oc_{C,c}$ is identified with the subring $k[[\varpi^m]]$.

Let us prove assertion 2. It suffices to prove it in the formal neighborhood $\Spec(\oc_{C,c})$ of the point $c$, in which case it reduces to the isomorphism
$$(\pi_*\Omega^1_{k[[\varpi]]/k})^{W_x}=\Omega^1_{k[[\varpi^m]]/k}.$$
Now $\pi_*\Omega^1_{k[[\varpi]]/k}$ is the $k[[\varpi^m]]$-module $k[[\varpi]]d\varpi$. Hence its module of invariants under $W_x$ is
$$k[[\varpi^m]]\varpi^{m-1}d\varpi=k[[\varpi^m]]d\varpi^m$$
which is indeed the module $\Omega^1_{k[[\varpi^m]]/k}$.

Let us prove assertion 3. We have
$$(\tgo^\vee\otimes_k\pi_*\Omega^1_{X/C})^W_c\cong(\tgo^\vee\otimes_k\pi_*\Omega^1_{\oc_{X,x}/\oc_{C,c}})^{W_x}$$
and $\pi_*\Omega^1_{\oc_{X,x}/\oc_{C,c}}$ is identified with the $k[[\varpi^m]]$-module $k[[\varpi]]d\varpi/\varpi^{m-1}k[[\varpi]]d\varpi$. Hence, as a representation of $W_x$, it is the sum of the characters $\zeta^i$ for $1\leq i<m$ with multiplicity 1. Let $\tgo_i\subset\tgo$ be the isotypic component of $\tgo$ on which $W_x$ acts by the character $\zeta^i$ for $0\leq i<m$. We therefore have
\begin{align*}
\dim_k((\tgo^\vee\otimes_k\pi_*\Omega^1_{\oc_{X,x}/\oc_{C,c}})^{W_x})&=\sum_{i=1}^{m-1}\dim(\tgo_{m-i})\\
&=\dim(\tgo)-\dim(\tgo^{W_x})
\end{align*}
whence assertion 3.

Let us prove assertion 4. If $X$ is \'etale above $c$, we have $\Omega^1_{X/C,x}=0$, whence the desired vanishing. If $X$ is not \'etale above $c$, we have $m\geq 2$. By definition, we have $F_\Omega(\Omega^1_{X/C})=\pi_*(\Omega^1_{X/C}\otimes_{\oc_X}\Omega^1_{X/k})^W$. We therefore have
\begin{align*}
F_\Omega(\Omega^1_{X/C})_c&=\pi_*(\Omega^1_{X/C}\otimes_{\oc_X}\Omega^1_{X/k})^W_c\\
&\simeq(k[[\varpi]]d\varpi\otimes_kd\varpi/\varpi^{m-1}k[[\varpi]]d\varpi\otimes_kd\varpi)^{W_x}\\
&\simeq k[[\varpi^m]]/\varpi^mk[[\varpi^m]](\varpi^{m-2}d\varpi\otimes_kd\varpi)
\end{align*}
whence the assertion on the dimension.
\end{preuve}
\end{paragr}

\begin{paragr}[The Ng\^o-Bezrukavnikov formula.]\label{S:9.12}
Ng\^o has related the local invariant $\delta_c$ to the dimension of an affine Springer fiber (cf.\ \cite[Prop.~3.7.5 and Cor.~3.8.3]{Ngo-lemme}). On the other hand, Bezrukavnikov has proved (cf.\ \cite{Bezrukavnikov}) a formula stated by Kazhdan-Lusztig (cf.\ \cite{Kazhdan-Lusztig}) for this dimension. Combining the two, Ng\^o obtained the following lemma, of which we give here a direct proof.

\begin{lemma}\label{thm:9.12.1}
Let $(X,\pi,f)\in\BBB_G(k)$ and let $a\in\AAA_G(k)$ be its image under the morphism defined in Lemma~\ref{thm:9.3.1}. For every $c\in C(k)$, we have the following formula:
$$\delta_c(X,\pi,f)=1/2(\val_c(D^G(a))-\dim_k(\tgo)+\dim_k(\tgo^{W_x})),$$
where $W_x\subset W$ is the stabilizer of a point $x\in X(k)$ above $c$.
\end{lemma}

\begin{preuve}
Let $(X,\pi,f)\in\BBB_G(k)$ and $a\in\AAA_G(k)$ be as above. Let $c\in C(k)$ and $x\in X(k)$ above $c$. We take up again the notation of the proof of Lemma~\ref{thm:9.11.1}. By the definition of $\delta_c(X,\pi,f)$, it suffices to prove the following two relations (where $m$ is the order of $W_x$):
\begin{equation}\label{eq:9.12.1}
\dim_k\Big(\frac{\tgo\otimes_k\oc_{X,x}}{\pi^*(\tgo\otimes_k\pi_{Y,*}\oc_Y)^W_x}\Big)=m\val_c(D^G(a))/2
\end{equation}
and
\begin{equation}\label{eq:9.12.2}
\dim_k\Big(\frac{\tgo\otimes_k\oc_{X,x}}{\pi^*(\tgo\otimes_k\pi_*\oc_X)^W_x}\Big)=m(\dim_k(\tgo)-\dim_k(\tgo^{W_x}))/2.
\end{equation}
Since one can find a non-degenerate bilinear form invariant under $W$, the above formulas are equivalent to those obtained when one replaces $\tgo$ by its dual $\tgo^\vee$. By formulas~(\ref{eq:9.10.3}) and~(\ref{eq:9.10.5}), we have
$$\pi^*(\tgo^\vee\otimes_k\pi_{Y,*}\oc_Y)^W=\pi^*a^*\Omega^1_{\car/C}=f^*\chi^*\Omega^1_{\car/C}$$
and the inclusion $\pi^*(\tgo^\vee\otimes_k\pi_{Y,*}\oc_Y)^W_x\hookrightarrow\tgo^\vee\otimes_k\oc_{X,x}$ is identified with the inclusion of $(f^*\chi^*\Omega^1_{\car/C})_x$ in $(f^*\Omega^1_{\tgo/C})_x$. The dimension we are looking for is therefore equal to $m$ times the valuation of the Jacobian determinant $J$ considered in the proof of Lemma~\ref{thm:9.11.1}. Now we know that, up to a factor in $k^\times$, the square of this Jacobian is equal to the discriminant $D^G(a)$ (cf.\ Proposition~5 of~\cite{Bourbaki}). Formula~(\ref{eq:9.12.1}) follows.

Let us now prove formula~(\ref{eq:9.12.2}). We have
$$\pi^*(\tgo\otimes_k\pi_*\oc_X)^W_x=k[[\varpi]]\tgo_0\oplus\bigoplus_{i=1}^{m-1}\varpi^{m-i}k[[\varpi]]\tgo_i.$$
We therefore have the following formula:
$$\dim_k\Big(\frac{\tgo\otimes_k\oc_{X,x}}{\pi^*(\tgo\otimes_k\pi_*\oc_X)^W_x}\Big)=\sum_{i=1}^{m-1}(m-i)\dim_k(\tgo_i)$$
which gives the desired result in view of the fact that $\dim_k(\tgo_i)=\dim_k(\tgo_{m-i})$ by the existence of a non-degenerate bilinear form invariant under $W$.
\end{preuve}
\end{paragr}

\begin{paragr}[A bound for $\kappa$ by the discriminant.]\label{S:9.13}
The aim of this paragraph is to obtain a bound for $\kappa$ by a multiple of the discriminant. Beforehand, we shall give another expression for $\kappa$.

\begin{lemma}\label{thm:9.13.1}
Let $(X,\pi,f)\in\BBB_G(k)$ and let $c$ be a closed point of $C$. Then
$$\kappa_c(X,\pi,f)=\lng(\pi_*(\hc^{-1}(L_{X/Y})\otimes_{\oc_X}\Omega^1_{X/C})^W_c).$$
\end{lemma}

\begin{preuve}
We have a distinguished triangle
$$\rho^*L_{Y/\tgo_D}\to L_{X/\tgo_D}\to L_{X/Y}\to,$$
where we have introduced the complex
$$L_{X/Y}=[\rho^*\Omega^1_{Y/C}\to\Omega^1_{X/C}]$$
placed in degrees $-1$ and $0$. We therefore have an exact sequence
$$0\to\hc^{-1}(\rho^*L_{Y/\tgo_D})\to\hc^{-1}(L_{X/\tgo_D})\to\hc^{-1}(L_{X/Y})\to 0$$
in view of the fact that $\hc^0(L_{Y/\tgo_D})=0$. It follows that we have an exact sequence
$$
\begin{aligned}
\pi_*(\hc^{-1}(\rho^*L_{Y/\tgo_D})\otimes_{\oc_X}\Omega^1_{X/C})^W&\to\pi_*(\hc^{-1}(L_{X/\tgo_D})\otimes_{\oc_X}\Omega^1_{X/C})^W\\
&\to\pi_*(\hc^{-1}(L_{X/Y})\otimes_{\oc_X}\Omega^1_{X/C})^W\to 0.
\end{aligned}
$$
By Lemma~\ref{thm:9.6.1}, we have
$$\hc^{-1}(\rho^*L_{Y/\tgo_D})=\pi^*a^*\Omega^1_{\car_D/C}$$
whence
$$\pi_*(\hc^{-1}(\rho^*L_{Y/\tgo_D})\otimes_{\oc_X}\Omega^1_{X/C})^W=a^*\Omega^1_{\car_D/C}\otimes_{\oc_X}\pi_*(\Omega^1_{X/C})^W=0$$
since $\pi_*(\Omega^1_{X/C})^W=0$ (cf.\ line~(\ref{eq:9.7.3})). It follows that we have
$$\pi_*(\hc^{-1}(L_{X/\tgo_D})\otimes_{\oc_X}\Omega^1_{X/C})^W\simeq\pi_*(\hc^{-1}(L_{X/Y})\otimes_{\oc_X}\Omega^1_{X/C})^W.$$
The lemma then follows from the definition of $\kappa_c$; cf.~(\ref{eq:9.7.1}).
\end{preuve}

The following bound, although rather crude, will suffice for what follows.

\begin{proposition}\label{thm:9.13.2}
Let $(X,\pi,f)\in\BBB_G(k)$ and let $a\in\AAA_G(k)$ be the corresponding point under the morphism $\psi$ defined in assertion 3 of Lemma~\ref{thm:9.3.1}. For every closed point $c$ of $C$, we have the inequality
$$\kappa_c(X,\pi,f)\leq|W|/2\val_c(D^G(a)).$$
\end{proposition}

\begin{preuve}
Let $c$ be a closed point of $C$. By Lemma~\ref{thm:9.13.1}, we have
\begin{align*}
\kappa_c(X,\pi,f)&=\lng(\pi_*(\hc^{-1}(L_{X/Y})\otimes_{\oc_X}\Omega^1_{X/C})^W_c)\\
&\leq\lng(\pi_*(\hc^{-1}(L_{X/Y})\otimes_{\oc_X}\Omega^1_{X/k})^W_c),
\end{align*}
where the inequality results from the obvious surjection
$$\hc^{-1}(L_{X/Y})\otimes_{\oc_X}\Omega^1_{X/k}\to\hc^{-1}(L_{X/Y})\otimes_{\oc_X}\Omega^1_{X/C}\to 0.$$
Since we also have an obvious injection
$$1\to\hc^{-1}(L_{X/Y})\otimes_{\oc_X}\Omega^1_{X/k}\to\rho^*\Omega^1_{Y/C}\otimes_{\oc_X}\Omega^1_{X/k}$$
and since we have (cf.~(\ref{eq:9.6.2}))
$$\rho^*\Omega^1_{Y/C}=f^*\Omega^1_{\tgo_D/\car_D},$$
it follows that
$$\kappa_c(X,\pi,f)\leq\lng(\pi_*(f^*\Omega^1_{\tgo_D/\car_D}\otimes_{\oc_X}\Omega^1_{X/k})^W_c).$$
Let $x$ be a point of $X$ above $c$. Let $W_x\subset W$ be the stabilizer of $x$ in $W$ and $m$ the order of $W_x$. One may therefore bound the above length by the length of the $\oc_{X,x}$-module $(f^*\Omega^1_{\tgo_D/\car_D}\otimes_{\oc_X}\Omega^1_{X/k})_x$ (where one has removed the invariants under $W_x$). Now the latter is the length of $(f^*\Omega^1_{\tgo_D/\car_D})_x$, hence it is equal to the valuation (at $x$) of the Jacobian determinant $J$ considered in the proof of Lemma~\ref{thm:9.11.1}, whose square is equal to the discriminant $D^G(a)$. This length is therefore
$$m/2\val_c(D^G(a)).$$
We then obtain the desired bound.
\end{preuve}
\end{paragr}

\begin{paragr}\label{S:9.14}
In this paragraph we shall prove the following lemma.

\begin{lemma}\label{thm:9.14.1}
Let $a\in\AAA_G(k)$ and $c\in C(k)$ be such that $\val_c(D^G(a))=1$. Then the cameral curve $Y_a$ associated with $a$ is smooth above $c$.
\end{lemma}

\begin{preuve}
By the Ng\^o-Bezrukavnikov formula (cf.\ Lemma~\ref{thm:9.12.1}), there exists an integer $\delta$ such that
$$1=\val_c(D^G(a))=2\delta+\dim_k(\tgo)-\dim_k(\tgo^{W_x}),$$
where $W_x\subset W$ is the stabilizer of a point $x\in X(k)$ above $c$. We therefore have $\delta=0$ and $\dim_k(\tgo)-\dim_k(\tgo^{W_x})=1$. Let $M$ be the Levi subgroup of $G$ defined as the centralizer of $\tgo^{W_x}$. Since $\tgo^{W_x}$ is of codimension 1 in $\tgo$ and is central in $M$, the rank of the derived group of $M$ is 1 (rank zero corresponds to $M=T$, which is excluded: the group $W_x$ would then be trivial). In particular, the group $W_x$ is equal to the Weyl group of $M$, of order 2. The morphism $X\to C$ in a formal completion of $x$ is given by the inclusion $k[[\varpi^2]]\hookrightarrow k[[\varpi]]$. With the notation of~(\ref{eq:9.4.1}), we have, by the choice of a trivialization of the line bundle associated with $D$, a morphism $f:\Spec(k[[\varpi]])\to\tgo\times_k\Spec(k[[\varpi^2]])$. Taking the quotient by $W^M$, we obtain a section $\Spec(k[[\varpi^2]])\to\car_M\times_k\Spec(k[[\varpi^2]])$. We are thus immediately reduced to the case where $G=M$. Using the decompositions $\tgo=\zgo_M\oplus\tgo_{M_{\der}}$ and $\car_M=\zgo_M\oplus\car_{M_{\der}}$, where we denote by $\zgo_M$ and $\tgo_{M_{\der}}$ the Lie algebras of the center of $M$ and of the torus $T\cap M_{\der}$, we are reduced to the case where $M$ is semi-simple of rank 1. A direct computation shows that the local germ of the curve $Y$ is of the form $\Spec(k[[u,y]]/(y^2=a(u)))$ with $\val_u(a(u))=1$, which is indeed smooth.
\end{preuve}
\end{paragr}

\begin{paragr}[Proof of Theorem~\protect\ref{thm:9.1.3}.]\label{S:9.15}
The morphism forgetting the datum $t$
$$\Ac_G\to\AAA_G$$
makes $\Ac_G$ a finite, \'etale and Galois covering of $\AAA_G$ with Galois group $W$. Composing with this morphism the functions $\delta$ and $\kappa$ of Definitions~\ref{thm:9.7.1}, we obtain functions still denoted by $\delta$ and $\kappa$. The first is upper semi-continuous (cf.\ Proposition~\ref{thm:9.4.4}) and the second is only constructible.

We are at last in a position to give the proof of Theorem~\ref{thm:9.1.3}. We take up again its notation: in particular, $a=(h_a,t)\in\Ac_G(k)$ satisfies assertions 1 and 2 of Theorem~\ref{thm:9.1.3}. The condition
$$\kappa<\deg(D)-2g+2$$
therefore defines a constructible subset $\Ac^{\kappa-\mathrm{bon}}$ of $\Ac_G$. We shall show that this subset contains a neighborhood of the point $a$. By \cite[Chap.~0, Prop.~9.2.5]{EGAIII}, it is necessary and sufficient that, for every irreducible closed subset $\yc$ of $\Ac_G$ containing $a$, the subset $\Ac^{\kappa-\mathrm{bon}}\cap\yc$ be dense in $\yc$. For this, it suffices to prove that $\Ac^{\kappa-\mathrm{bon}}\cap\yc$ contains the generic point of $\yc$. It suffices further to prove the following lemma.

\begin{lemma}\label{thm:9.15.1}
Let $K$ be an algebraically closed extension of $k$ and $b$ a morphism from $\Spec(K[[u]])$ to $\Ac_G$ with special fiber $b_s=a\times_kK$. Let $b_\eta$ be the morphism deduced from $b$ on the generic fiber. Then $b_\eta$ defines a point of $\Ac^{\kappa-\mathrm{bon}}$.
\end{lemma}

\begin{preuve}
Let $h_b$ be the section of the bundle $\car_D\times_kK[[u]]$ above the relative curve $C\times_kK[[u]]$ underlying the datum of $b$. Let $C_\eta$ and $C_s$ be the generic and special fibers of this curve. Let $S\subset C$ be the finite set of closed points of $C$ which satisfies assertions 1 and 2 of Theorem~\ref{thm:9.1.3}. In what follows, we identify $S$ with a set of closed points of $C_s=C\times_kK$.

\begin{lemma}\label{thm:9.15.2}
For every closed point $c_\eta$ of $C_\eta$ which specializes to a point of $C_s-S$ we have
$$\kappa_{c_\eta}(b_\eta)=0,$$
where $\kappa_{c_\eta}$ is the local invariant defined in~(\ref{eq:9.7.1}).
\end{lemma}

\begin{preuve}
If $c_\eta$ specializes to a point at which $h_a(C_s)$ does not meet $\mathfrak{R}^G_M$ and meets transversally or does not meet $\mathfrak{D}^M$, then $h(C_\eta)$ at $c_\eta$ does not meet or meets transversally $\mathfrak{D}^G$. In all cases, Lemma~\ref{thm:9.14.1} implies that the cameral curve associated with $b_\eta$ is smooth above $c_\eta$, whence the vanishing of $\kappa$ at this point (cf.\ Proposition~\ref{thm:9.7.4}).

If $c_\eta$ specializes to a point $c_s$ at which $h(C_s)$ does not meet $\mathfrak{D}^M$ and meets transversally or does not meet $\mathfrak{R}^G_M$, then $h(C_\eta)$ at this point meets the divisor $\mathfrak{D}^G$ with multiplicity less than or equal to 2. If this multiplicity is zero or equal to 1, we conclude as before that $\kappa$ at $c_\eta$ is zero. So let us suppose that this multiplicity is equal to 2. Let $Y$ be the cameral curve associated with $h_b$. Let $Y_s$ and $Y_\eta$ be its special and generic fibers. Let $y_s$ be a point above $c_s$. By hypothesis, there exists a unique root $\pm\alpha$ of $T$ in $G$ such that $Y_s$ meets at $y_s$ the divisor defined by $\alpha$. In particular, the stabilizer in $W$ of $y_s$ is of order 2. Let $Y_\eta$ be the cameral curve associated with $h_{b_\eta}$ and $y_\eta$ a point above $c_\eta$. This point $y_\eta$ specializes to a point $y_s\in Y_s$ above $c_s$. The stabilizer of $y_\eta$ in $W$ is contained in the stabilizer of $y_s$, hence its order is at most 2. Replacing if necessary $K((u))$ by a purely inseparable extension and $K[[u]]$ by its normalization in that extension, we may suppose that the normalization $X_\eta$ of $Y_\eta$ is smooth (cf.\ \cite[Prop.~17.15.14]{EGAIV4}). Let $x_\eta$ be a point of $X_\eta$ above $y_\eta$. The stabilizer of $x_\eta$ in $W$ is contained in the stabilizer of $y_\eta$, hence is a group of order at most 2. By the Ng\^o-Bezrukavnikov formula (Lemma~\ref{thm:9.12.1}), we know that the valuation of the discriminant and the difference $\dim(\tgo)-\dim(\tgo^{W_{x_\eta}})$ have the same parity. Now our hypothesis is that the valuation of the discriminant is 2. Hence the difference $\dim(\tgo)-\dim(\tgo^{W_{x_\eta}})$ is even. If $|W_x|=2$, this quantity is equal to 1. Hence $W_x=\{1\}$ and $X_\eta$ is \'etale above $c_\eta$. Consequently, $\kappa$ vanishes at $c_\eta$ (cf.\ Proposition~\ref{thm:9.7.4}).
\end{preuve}

By~(\ref{eq:9.7.2}) and the preceding lemma, we have
$$\kappa(b_\eta)=\sum_{c_\eta\in S_\eta}\kappa_{c_\eta}(b_\eta),$$
where $S_\eta$ is the set of points $c_\eta$ of $C_\eta$ which specialize to points of $S$. Proposition~\ref{thm:9.13.2} then gives the bound
$$\kappa(b_\eta)\leq|W|/2\sum_{c_\eta\in S_\eta}\val_{c_\eta}(D^G(b_\eta)).$$
The sum in the above bound is equal to $\sum_{c\in S}\val_c(D^G(b_s))$. Now $b_s=a\times_kK$. Hypothesis 2 of Theorem~\ref{thm:9.1.3} implies the new bound
$$\kappa(b_\eta)<\deg(D)-2g+2$$
which indeed proves that $b_\eta$ defines a point of $\Ac^{\kappa-\mathrm{bon}}$, and this completes the proof of Lemma~\ref{thm:9.15.1}.
\end{preuve}

We have thus proved that $\Ac^{\kappa-\mathrm{bon}}$ contains an open neighborhood of $a$, which we denote by $U$. Let us next show that $U\subset\Ac^{\mathrm{bon}}_G$. Let $\delta$ be an integer and $\Ac'_\delta$ an irreducible component of a stratum $\Ac_\delta$. We must show that if $U\cap\Ac'_\delta$ is non-empty then we have
$$\dim_k(\Ac'_\delta)\leq\dim(\Ac_G)-\delta.$$
Up to a finite \'etale covering, the intersection $U\cap\Ac'_\delta$ is the image under the morphism $\psi$ considered in Lemma~\ref{thm:9.4.1} of a locally closed subset of $\BBB_G$. At every point $b\in\BBB_G(k)$ of this subset, we have $\delta(b)=\delta$ and $\kappa(b)<\deg(D)-2g+2$. Propositions~\ref{thm:9.8.1} and~\ref{thm:9.9.1} show that $\BBB_G$ is smooth at such a point, of dimension $\leq\dim_k(\Ac_G)-\delta$. We therefore have
$$\dim_k(\Ac'_\delta)=\dim_k(U\cap\Ac'_\delta)\leq\dim(\Ac_G)-\delta$$
which is what we had to prove. This completes the proof of Theorem~\ref{thm:9.1.3}.
\end{paragr}
\section{The cohomological statement extending Ngô's}\label{sec:10}

\begin{paragr}[The decomposition theorem]\label{S:10.1}
In what follows, we denote with an index $0$ an object over $\Fq$, and the omission of the index $0$ indicates the extension of scalars to $k$. Let $A_0$ be an integral $\Fq$-scheme of finite type. Following our conventions, we have $A = A_0\otimes_{\Fq}k$. Let $d_A$ be the dimension of $A$. For every point $a$ of $A$, let $d_a$ be the dimension of the irreducible closed subset $\overline{\{a\}}$ of $A$. Let $i_a = \overline{\{a\}}\hookrightarrow X$ and $j_a = \{a\}\hookrightarrow\overline{\{a\}}$ be the inclusions. We will say that a local system of $\Qlb$-vector spaces on the scheme $\{a\}$ is \emph{admissible} if it extends to a local system $\fc_V$ on a dense open subset $V$ of $\overline{\{a\}}$. For such a local system $\fc_a$, we will denote by $j_{a,!*}\fc_a$ the intermediate extension $j_{V,!*}\fc_V$ of the local system $\fc_V$ to the whole of $\overline{\{a\}}$, for any small enough dense open subset $j_V : V\hookrightarrow\overline{\{a\}}$.

\begin{proposition}\label{thm:10.1.1}
(\cite[Theorem~5.3.8]{BBD}) Let $K_0$ be a pure $\Qlb$-perverse sheaf of weight $w$ on $A_0$. The perverse sheaf $K$ on $A$, obtained by extension of scalars from $\Fq$ to $k$, admits a canonical decomposition into a direct sum
$$K = \bigoplus_{a\in A} i_{a,*}j_{a,!*}\fc_a[d_a],$$
where for all $a$ but finitely many, $\fc_a = (0)$.
\end{proposition}

\begin{preuve}
Uniqueness follows from the fact that
$$\Hom(i_{a,*}j_{a,!*}\fc_a[d_a], i_{a',*}j_{a',!*}\fc_{a'}[d_{a'}]) = (0)$$
for all $a\neq a'$ in $A$. Indeed the restriction to $\overline{\{a\}}\cap\overline{\{a'\}}$ of $j_{a,!*}\fc_a[d_a]$ lies in ${}^{\mathrm{p}}D_c^{\leq-1}(\overline{\{a\}}\cap\overline{\{a'\}},\Qlb)$, and therefore $i_a^*i_{a',*}j_{a,!*}\fc_a[d_a]$ lies in ${}^{\mathrm{p}}D_c^{\leq-1}(\overline{\{a\}},\Qlb)$.

For existence, we proceed by induction on the dimension of $A$. Let $j : U\hookrightarrow A$ be an open subset whose reduction is smooth over $\Fq$, which contains all the generic points of the irreducible components of $A$ of dimension $d_A$ and on which $j^{*\mathrm{p}}\hc^{-d_A}(K)$ is a local system and $j^*K = j^{*\mathrm{p}}\hc^{-d_A}(K)[d_A]$. There then exists a filtration
$$(0) = K_0\subset K_1\subset K_2\subset K_3 = K$$
with $K_1/K_0$ and $K_3/K_2$ supported in $A-U$ and $K_2/K_1 = j_{!,*}j^{*\mathrm{p}}\hc^{-d_A}(K)[d_A]$. Now for every pure local system $\fc$ on $U$ and every pure $\Qlb$-perverse sheaf $L$ of the same weight on $A$ supported by $A-U$, we have
$$\Ext^1(L, j_{!,*}\fc) = \Ext^1(j_{!,*}\fc, L) = (0)$$
(on $A$ the $\Ext^0$ vanish, as we have already recalled, and the $\Ext^1$ cannot have $1$ as an eigenvalue of Frobenius, by Deligne's theorem). It follows that $j_{!,*}j^{*\mathrm{p}}\hc^{-d_A}(K)[d_A]$ is a direct factor of $K$ with a complement supported by $A-U$.
\end{preuve}
\end{paragr}

\begin{paragr}\label{S:10.2}
Let $f_0 : M_0\to A_0$ be a morphism of separated $\Fq$-schemes of finite type. Let $f : M\to A$ be the morphism obtained by base change from $\Fq$ to $k$. We assume that $f_0$ is proper, purely of relative dimension $d_f$, and that $M_0$ is smooth, purely of dimension $d_M = d_f+d_A$ over $\Fq$. We consider the graded perverse sheaf on $A_0$, pure in each degree,
$$K_0^\bullet = {}^{\mathrm{p}}\hc^\bullet(Rf_{0,*}\overline{\QQ}_{\ell,M_0}[d_M])
= \bigoplus_n {}^{\mathrm{p}}\hc^n(Rf_{0,*}\overline{\QQ}_{\ell,M_0}[d_M]).$$
Let $K^\bullet$ be the graded perverse sheaf on $A$ deduced from $K_0^\bullet$ by base change from $\Fq$ to $k$. By Proposition~\ref{thm:10.1.1}, we have
$$K^\bullet = \bigoplus_{a\in A} i_{a,*}j_{a,!*}\fc_a^\bullet[d_a].$$
where each $\fc_a^\bullet$ is an admissible graded local system on $\{a\}$. Poincaré duality ensures that $\fc_a^{-n} = (\fc_a^n)^\vee(d_a)$.

We define \emph{the socle} of $Rf_*\overline{\QQ}_{\ell,M}[d_M]$ as the finite set
$$\mathrm{Socle}(Rf_*\overline{\QQ}_{\ell,M}[d_M])\subset A$$
of the $a$ such that $\fc_a^\bullet\neq(0)$. For every $a$ in this socle, let $n_a$, resp. $n_a^-$, be the largest integer $n$, the smallest one, such that $\fc_a^n\neq(0)$. We define \emph{the amplitude} of $\fc_a^\bullet$ by
$$\mathrm{Ampl}(\fc_a^\bullet) = n_a - n_a^-\geq 0.$$

\begin{lemma}\label{thm:10.2.1}
Let $a\in\mathrm{Socle}(Rf_*\overline{\QQ}_{\ell,M}[d_M])$. We have the inequalities
$$\mathrm{Ampl}(\fc_a^\bullet)\leq 2(d_f-d_A+d_a)$$
and
$$\mathrm{codim}_A(a) := d_A-d_a\leq d_f.$$
\end{lemma}

\begin{preuve}
Let us write as above $\mathrm{Ampl}(\fc_a^\bullet) = n_a-n_a^-$. By Poincaré duality, we have $n_a = -n_a^-$, so that
$$\mathrm{Ampl}(\fc_a^\bullet) = 2n_a.$$
The restriction to $\{a\}$ of $R^{d_M+n_a-d_a}f_*\overline{\QQ}_{\ell,M}$, which contains $\fc_a^{n_a}$ as a direct factor, is nonzero. Consequently, we have $d_M+n_a-d_a\leq 2d_f$, or again
$$\mathrm{Ampl}(\fc_a^\bullet)\leq 2(d_f-d_A+d_a).$$
In particular, we necessarily have $\mathrm{codim}_A(a) := d_A-d_a\leq d_f$.
\end{preuve}
\end{paragr}

\begin{paragr}[Back to a key result of Ngô.]\label{S:10.3}
Let $g_0 : P_0\to A_0$ be a smooth commutative group scheme, purely of relative dimension $d_f$ over $A_0$, with connected fibers. We assume that $P_0$ acts on $M_0$ over $A_0$. Let
$$V_\ell P_0 = R^{2d_f-1}g_{0,!}\Qlb(d_f).$$
This is a sheaf on $A_0$ whose fiber at a geometric point $\bar{a}$ of $A_0$ is the $\Qlb$-Tate module $V_\ell P_{0,\bar{a}}$ of the fiber of $P_0$ at $\bar{a}$. We have an action of the commutative Lie algebra $V_\ell P_0$ on $Rf_{0,*}\Qlb$
$$V_\ell P_0\otimes Rf_{0,*}\Qlb\to Rf_{0,*}\Qlb[-1].$$
For every geometric point $\bar{a}$ of $A_0$, this action induces on the one hand a structure of
$$\wedge^\bullet V_\ell P_{0,\bar{a}} = \bigoplus_{i\in\NN}\wedge^i V_\ell P_{0,\bar{a}}$$
graded module on the cohomology $H^\bullet(M_{0,\overline{a}},\Qlb)$. On the other hand it induces a structure of
$$\wedge^\bullet V_\ell P_0 = \bigoplus_{i\in\NN}\wedge^i V_\ell P_0$$
module on
$${}^{\mathrm{p}}\hc^\bullet(Rf_{0,*}\Qlb) = \bigoplus_{n\in\ZZ}{}^{\mathrm{p}}\hc^n(Rf_{0,*}\Qlb).$$
In particular, for every $a$ in the socle of $Rf_*\overline{\QQ}_{\ell,M}[d_M]$, it induces a structure of $\wedge^\bullet V_\ell P_{0,a}$ on $\fc_a^\bullet$.

For each geometric point $\overline{a}$ of $A_0$ with residue field denoted $k(\bar{a})$, the fiber $P_{0,\overline{a}}$ of $P_0$ admits a canonical dévissage
\begin{equation}\label{eq:10.3.1}
1\to P_{0,\overline{a}}^{\mathrm{aff}}\to P_{0,\overline{a}}\to P_{0,\overline{a}}^{\mathrm{ab}}\to 1,
\end{equation}
where $P_{0,\overline{a}}^{\mathrm{ab}}$ is an abelian $k(\overline{a})$-scheme and $P_{0,\overline{a}}^{\mathrm{aff}}$ is an affine $k(\overline{a})$-group scheme. The dimension of $P_{0,\overline{a}}^{\mathrm{aff}}$ depends only on the image $a$ of $\overline{a}$ and will be denoted $\delta_a^{\mathrm{aff}}$. The function $a\mapsto\delta_a^{\mathrm{aff}}$ is upper semi-continuous. In particular, if $a$ is a point of $A$, the function $b\mapsto\delta_b^{\mathrm{aff}}$ is constant of value $\delta_a^{\mathrm{aff}}$ on a dense open subset of $\overline{\{a\}}$. Let $\overline{a}$ be a geometric point of $A_0$. We have an analogous dévissage at the level of the Tate modules
\begin{equation}\label{eq:10.3.2}
0\to V_\ell P_{0,\overline{a}}^{\mathrm{aff}}\to V_\ell P_{0,\overline{a}}\to V_\ell P_{0,\overline{a}}^{\mathrm{ab}}\to 0,
\end{equation}
where $V_\ell P_{0,\overline{a}}^{\mathrm{ab}}$ is of dimension $2(d_f-\delta_a^{\mathrm{aff}})$ and $V_\ell P_{0,\overline{a}}^{\mathrm{aff}}$ is of dimension equal to the dimension of the toric part of $P_{0,\overline{a}}^{\mathrm{aff}}$.

We have just seen above that the cohomology $H^\bullet(M_{0,\overline{a}},\Qlb)$ is endowed with a structure of $\wedge^\bullet V_\ell P_{0,\overline{a}}$-graded module. Every splitting of the exact sequence~(\ref{eq:10.3.2}) induces an action of
$$\wedge^\bullet V_\ell P_{\overline{a}}^{\mathrm{ab}}$$
on $H^\bullet(M_{0,\overline{a}},\Qlb)$. This action depends \emph{a priori} on the splitting. If the geometric point $\overline{a}$ is localized at a closed point $a$ of $A_0$, the residue field $k(a)$ is finite and the extension of $P_{0,\overline{a}}^{\mathrm{ab}}$ by $P_{0,\overline{a}}^{\mathrm{aff}}$ given by~(\ref{eq:10.3.1}) is of finite order. It is therefore split up to an isogeny, and the splitting obtained is necessarily unique. In particular the sequence~(\ref{eq:10.3.2}) is canonically split (cf.\ \cite[prop.\ 7.5.3]{Ngo-lemme}).

Let $a$ be a point of the socle of $Rf_*\overline{\QQ}_{\ell,M}[d_M]$. In fact, the dévissage~(\ref{eq:10.3.1}) exists for $a$ after a radicial base change. In particular, one can define the dévissage~(\ref{eq:10.3.2}) for $a$. By a weight argument (cf.\ \cite[\S7.4.8]{Ngo-lemme}), the action of $\wedge^\bullet V_\ell P_{0,a}$ on $\fc_a^\bullet$ factors through an action of $\wedge^\bullet V_\ell P_{0,a}^{\mathrm{ab}}$ on $\fc_a^\bullet$. Thus $\fc_a^\bullet$ is a graded module over $\wedge^\bullet V_\ell P_{0,a}^{\mathrm{ab}}$.

Following Ngô (cf.\ \cite[\S7.1.4]{Ngo-lemme}), one says that $V_\ell P_0$ is \emph{polarisable} if, locally for the étale topology of $A_0$, there exists an alternating bilinear form
$$V_\ell P_0\times V_\ell P_0\to\Qlb(1)$$
such that, for each geometric point $\overline{a}$ of $A_0$, the kernel of the bilinear form induced on $V_\ell P_{\overline{a}}$ is $V_\ell P_{\overline{a}}^{\mathrm{aff}}$ and the form induced on $V_\ell P_{\overline{a}}^{\mathrm{ab}}$ is nondegenerate. The following proposition is a reformulation of a result of Ngô.

\begin{proposition}\label{thm:10.3.1}
Let $f_0 : M_0\to A_0$ be a morphism of separated $\Fq$-schemes of finite type. Let $g_0 : P_0\to A_0$ be a smooth commutative group scheme, purely of relative dimension $d_f$ over $A_0$, with connected fibers. Let $P_0\times M_0\to M_0$ be an action of $P_0$ on $M_0$ over $A_0$.

The following hypotheses:
\begin{enumerate}
\item \begin{enumerate} \item[(a)] $f_0$ is proper, purely of relative dimension $d_f$; \item[(b)] $M_0$ is smooth, purely of dimension $d_M = d_f+d_A$ over $\Fq$;
  \end{enumerate}
\item the Tate module $V_\ell P_0$ is polarisable; \item for every geometric point $\bar{a}$ localized at a closed point $a$ of $A_0$, the action of $\wedge^\bullet V_\ell P_{\overline{a}}^{\mathrm{ab}}$ on $H^\bullet(M_{\overline{a}},\Qlb)$, induced by the canonical splitting of~(\ref{eq:10.3.2}) described above, makes $H^\bullet(M_{\overline{a}},\Qlb)$ a free $\wedge^\bullet V_\ell P_{\overline{a}}^{\mathrm{ab}}$-module;
\end{enumerate}
imply that for every point $a\in\mathrm{Socle}(Rf_*\overline{\QQ}_{\ell,M}[d_M])$, the fiber at $a$ of $\fc_a^\bullet$ is a free $\wedge^\bullet V_\ell P_{0,a}^{\mathrm{ab}}$-module.
\end{proposition}

\begin{preuve}
This is essentially the content of Proposition 7.4.10 of \cite{Ngo-lemme}. The hypotheses under which Ngô places himself in that proposition are not exactly the ones we adopt. More precisely, we do not assume that the morphism $f_0$ is projective. In fact, in Ngô's proof, the projectivity of $f_0$ only intervenes through Proposition 7.5.1 of \cite{Ngo-lemme}, whose conclusion is our hypothesis 3.
\end{preuve}

Following Ngô, one deduces the following corollary.

\begin{corollary}\label{thm:10.3.2}
Under the hypotheses of Proposition~\ref{thm:10.3.1}, for every point $a\in A$ which belongs to the socle of $Rf_*\overline{\QQ}_{\ell,M}[d_M]$, we have the inequalities
$$\mathrm{Ampl}(\fc_a^\bullet)\geq 2(d_f-\delta_a^{\mathrm{aff}})$$
and
$$\delta_a^{\mathrm{aff}}\geq\mathrm{codim}_A(a).$$
\end{corollary}

\begin{preuve}
The first inequality is an obvious consequence of Proposition~\ref{thm:10.3.1} and the second follows from Lemma~\ref{thm:10.2.1}.
\end{preuve}

The following lemma will be useful to us in checking the hypotheses of Proposition~\ref{thm:10.3.1}.

\begin{lemma}\label{thm:10.3.3}
Let $\al : B'\to B$ be an isogeny between abelian varieties defined over an algebraically closed field $K$. Let $X$ be a $K$-scheme (or more generally a Deligne-Mumford $K$-stack) endowed with an action of $B'$ and with a proper morphism $f : X\to B$ which is $B'$-equivariant when $B'$ acts on $B$ by translation via the isogeny $\al$. Then the action of
$$H_\bullet(B',\Qlb) = \wedge^\bullet V_\ell B'$$
on $H^\bullet(X,\Qlb)$ induced by the action of $B'$ on $X$ makes $H^\bullet(X,\Qlb)$ a free $\wedge^\bullet V_\ell B'$-module.
\end{lemma}

\begin{preuve}
We consider the cartesian diagram
$$\xymatrix{
X_0\times_K B' \ar[r]^-{\al'} \ar[d]_{\mathrm{pr}_{B'}} & X \ar[d]^{f}\\
B' \ar[r]^-{\al} & B,
}$$
where $\mathrm{pr}_{B'}$ is the projection onto the second factor, $X_0$ is the fiber of $f$ at the origin of $B$ and $\al'$ is the restriction to $X_0\times_K B'$ of the action $X\times_K B'\to X$. The cohomology of $X_0\times_K B'$, which is given by the tensor product $H^\bullet(X_0)\otimes H^\bullet(B')$, is obviously a free module over $\wedge^\bullet V_\ell B'$. The cohomology of $X$ is identified with the invariants under $\ker(\al)$ in the cohomology of $X\times_B B'\simeq X_0\times_K B'$. Now $\ker(\al)$ acts trivially on $H^\bullet(B')$ by the usual homotopy argument. Thus the cohomology of $X$ is identified with $H^\bullet(X)^{\ker(\al)}\otimes H^\bullet(B')$ and it is thus free as a module over $\wedge^\bullet V_\ell B'$.
\end{preuve}
\end{paragr}

\begin{paragr}[Néron model.]\label{S:10.4}
We return to the notation used since the beginning of the article. Let $a\in\Ac_G$ be a closed point. Let $\nu : J_a\hookrightarrow\tilde{J}_a$ be the Néron model of the group scheme $J_a$ over $C$. The arrow $\nu$ is an isomorphism over the open subset $U_a$, so that the quotient $\tilde{J}_a/J_a$ has punctual support. We have the following proposition, due to Ngô (cf.\ \cite[cor.\ 4.8.1]{Ngo-lemme}).

\begin{proposition}\label{thm:10.4.1}
Let $a = (h_a,t)\in\Ac_G(k)$. Let $\rho_a : X_a\to Y_a$ be the normalization of the curve $Y_a$. There exists a unique isomorphism
$$\tilde{J}_a\stackrel{\sim}{\longrightarrow}(X_*(T)\otimes_\ZZ\pi_{a,*}\rho_{a,*}\mathbb{G}_{m,X_a})^W,$$
where the exponent $W$ denotes the fixed points under $W$, which fits into the cartesian square
$$\xymatrix{
J_a \ar[r] \ar@{^{(}->}[d]_{\nu} & (X_*(T)\otimes_\ZZ\pi_{a,*}\mathbb{G}_{m,Y_a})^W \ar@{^{(}->}[d]\\
\tilde{J}_a \ar[r]^-{\sim} & (X_*(T)\otimes_\ZZ\pi_{a,*}\rho_{a,*}\mathbb{G}_{m,X_a})^W,
}$$
where
\begin{itemize}
\item[--] the top horizontal arrow is obtained by base change by $h_a$ from the morphism
  $$J_D^G\longrightarrow(X_*(T)\otimes_\ZZ(\chi_G^T)_*\mathbb{G}_{m,\tgo_D})^W$$
deduced from~(\ref{eq:3.4.1}); \item[--] the right vertical arrow is induced by the adjunction arrow
  $$\mathbb{G}_{m,Y_a}\hookrightarrow\rho_{a,*}\mathbb{G}_{m,X_a}.$$
\end{itemize}
In particular, the Lie algebra of $\tilde{J}_a$ is the commutative Lie algebra with underlying vector bundle
$$(X_*(T)\otimes_\ZZ\pi_{a,*}\rho_{a,*}\oc_{X_a})^W.$$
\end{proposition}
\end{paragr}

\begin{paragr}[Application to the Hitchin fibration.]\label{S:10.5}
From now on we assume that the group $G$ is \emph{semi-simple}. Let $(G_0,T_0)$ be a pair formed of a connected reductive group $G_0$ over $\Fq$ and of $T_0\subset G_0$ a maximal $\Fq$-subtorus, split over $\Fq$. Let $\tau$ be the Frobenius automorphism of $k$ given by $x\mapsto x^q$. In what follows we limit ourselves to the case where the pair $(G,T)$ is obtained from $(G_0,T_0)$ after extension of scalars from $\Fq$ to $k$. The pair $(G,T)$ is then endowed with the natural action of $\tau$.

We fixed in Section~\ref{S:4.1} a curve $C$, projective, smooth and connected over $k$, as well as a divisor $D$ and a point $\infty\in C(k)$. From now on we assume that $C$ comes by base change from a curve $C_0$ over $\Fq$. We again denote by $\tau$ the automorphism $1\times\tau$ of $C = C_0\times_{\Fq}k$. We assume that $D$ and $\infty$ are fixed under $\tau$.

The scheme $\Ac_G$ is then endowed with the natural action of the Frobenius $\tau$. The latter leaves stable the closed subschemes $\Ac_M$ for $M\in\lc^G(T)$. Likewise, the stack $\mc_G$ is endowed with the natural action of $\tau$. Let $\xi\in\ago_T$ \emph{in general position} (cf.\ Remark~\ref{thm:4.9.2}). Let $f^\xi : \mc_G^\xi\to\Ac_G$ be the truncated Hitchin morphism, which is proper with smooth source over $k$ (cf.\ \S\ref{S:4.9}). Then the open substack $\mc_G^\xi$ is stable under $\tau$ and the morphism $f^\xi$ is $\tau$-equivariant. Equivalently, one could have directly defined objects $f_0^\xi : \mc_{G_0}^\xi\to\Ac_{G_0}$ over $\Fq$ which give back, after extension of scalars from $\Fq$ to $k$, the analogous objects denoted without the index $0$.

Let $d_{\mc^\xi}$ be the relative dimension of $f^\xi$. Let us recall that we have defined an open subset $\Ac_G^{\mathrm{bon}}$ of $\Ac_G$ (cf.\ Definition~\ref{thm:9.1.1}).

\begin{theorem}\label{thm:10.5.1}
The socle of the restriction of $Rf^\xi_*\Qlb[d_{\mc^\xi}]$ to $\Ac^{\mathrm{bon}}$ is entirely contained in the elliptic open subset $\Ac^{\el}\cap\Ac^{\mathrm{bon}}$, so that if
$$j : \Ac^{\el}\cap\Ac^{\mathrm{bon}}\hookrightarrow\Ac^{\mathrm{bon}}$$
is the inclusion of this open subset, we have a canonical isomorphism
$${}^{\mathrm{p}}\hc^i(Rf^\xi_*\Qlb[d_{\mc^\xi}])\cong j_{!*}j^{*\mathrm{p}}\hc^i(Rf^\xi_*\Qlb[d_{\mc^\xi}])$$
on $\Ac^{\mathrm{bon}}$ for every integer $i$.
\end{theorem}

\begin{preuve}
We are going to apply Corollary~\ref{thm:10.3.2} to the Hitchin fibration $\mc^\xi$ endowed with the action of the Picard stack $\Jc^1$. We shall deduce from it that every $a$ which is in the socle of $Rf^\xi_*\Qlb[d_{\mc^\xi}]$ satisfies the inequality $\delta_a^{\mathrm{aff}}\geq\mathrm{codim}_\Ac(a)$. We shall then show that one has $\delta_a^{\mathrm{aff}}\leq\delta_a$ with equality if and only if $a$ is elliptic. Hence, if $a$ is not elliptic, $a$ cannot be in $\Ac^{\mathrm{bon}}$.

We shall first of all show that for every $a\in\Ac_G(k)$, the fiber $\mc_a^\xi$ satisfies the hypotheses required by Proposition~\ref{thm:10.3.1} and Corollary~\ref{thm:10.3.2}. For this, we apply Lemma~\ref{thm:10.3.3} not to $\mc_a$ but to a homeomorphic stack, which suffices to prove the freeness of the cohomology of $\mc_a$.

Let $M\in\lc^G$ be such that $a$ is the image of a point $(h_{a_M},t)\in\Ac_M^{\el}(k)^\tau$. Let $X_a = \eps_M(h_{a_M})\in\mgo(F)$ where $\eps_M$ is the Kostant section and $F$ is the function field of $C$. Let $\AAA$ be the ring of adèles of $F$ and $\oc$ the product of the completions of the local rings of the closed points of $C$. Let $S$ be the finite set of closed points of $C$ complementary to the open subset $h_{a_M}^{-1}(\car_{M,D}^{\reg})$. For every $v\in S$, let $\Sc_v$ be the affine Springer fiber associated with $v$ and $X_a$. Let $F_v$ be the fraction field of the completion of the local ring $\oc_v$ at $v$. Let $J$ be the centralizer of $X_a$ in $M\times_k F$. The ind-scheme $J(F_v)/J(\oc_v)$ acts on $\Sc_v$. Let $\Xgo$ be the contracted product of $J(\AAA)/J(\oc)\times\prod_{v\in S}\Sc_v$ by the diagonal action of $\prod_v J(F_v)/J(\oc_v)$. The stack quotient of $\Xgo$ by the action (on the first factor) of $J(F)$ is homeomorphic to the Hitchin fiber $\mc_a$. The homeomorphism between $[J(F)\back\Xgo]$ and $\mc_a$ is equivariant under the action of the Picard stack $\Jc_a = [J(F)\back J(\AAA)/J(\oc)]$.

The latter admits as a quotient $[J(F)\back J(\AAA)/J(\oc)\prod_v J(F_v)]$, which is isogenous to the stack $\Jc_a^{\mathrm{ab}}$ of $\tilde{J}_a$-torsors, where $\tilde{J}_a$ is the Néron model of $J_a$. The stack $\Jc_a$ acts on $[J(F)\back\Xgo]$ and the obvious morphism from $[J(F)\back\Xgo]$ to $[J(F)\back J(\AAA)/J(\oc)\prod_v J(F_v)]$ is $\Jc_a$-equivariant.

More generally, one can define a ``truncated'' $\Xgo^\xi$ whose quotient by $J(F)$ is endowed with an action by $\Jc_a^1$ and is homeomorphic, in a $\Jc_a^1$-equivariant way, to $\mc_a^\xi$. We have a $\Jc_a^1$-equivariant morphism from $[J(F)\back\Xgo^\xi]$ to $[J(F)\back J(\AAA)/J(\oc)\prod_v J(F_v)]$ whose image is isogenous to the stack $\Jc_a^{1,\mathrm{ab}}$, the image of $\Jc_a^1$ in $\Jc_a^{\mathrm{ab}}$. One obtains an action of $\Jc_a^{1,\mathrm{ab}}$ on $[J(F)\back\Xgo^\xi]$ by choosing, as is legitimate, a section up to isogeny of the canonical morphism $\Jc_a^{1,\mathrm{ab}}\to\Jc_a^1$.

The hypotheses of Lemma~\ref{thm:10.3.3} are therefore satisfied for the quotient
$$[J(F)\back\Xgo^\xi]$$
and one deduces from it the desired freeness of the cohomology of $[J(F)\back\Xgo^\xi]$, hence of $\mc_a^\xi$.

Let us now compare, for a point $a$ of $\Ac_G$, the invariants $\delta_a$ and $\delta_a^{\mathrm{aff}}$. The stack $\Jc_a^1$ is a Deligne-Mumford stack whose automorphism groups are all the finite group $H^0(C,J_a)$. Let $\Jc_a^0$ be the neutral component of $\Jc_a^1$. We recall that $\tilde{J}_a$ is the Néron model of $J_a$. Let $\Jc_a^{0,\mathrm{ab}}$ be the abelian scheme which represents the functor of isomorphism classes of $\tilde{J}_a$-torsors. The morphism $\Jc^0\longrightarrow\Jc^{0,\mathrm{ab}}$, which sends a $J_a$-torsor to the class of the $\tilde{J}_a$-torsor obtained when one pushes it out by the canonical morphism $J_a\to\tilde{J}_a$, fits into the exact sequence
$$1\longrightarrow\Jc^{0,\mathrm{aff}}\longrightarrow\Jc^0\longrightarrow\Jc^{0,\mathrm{ab}}\longrightarrow 0,$$
where the kernel $\Jc^{0,\mathrm{aff}}$ is the quotient of an affine algebraic group by the finite group $H^0(C,J_a)$ acting trivially. We are going to compute the dimension of these objects by a Lie algebra computation. The exact sequence
$$1\longrightarrow J_a\longrightarrow\tilde{J}_a\longrightarrow\tilde{J}_a/J_a\longrightarrow 1$$
gives the long exact sequence
$$1\longrightarrow H^0(C,J_a)\longrightarrow H^0(C,\tilde{J}_a)\longrightarrow H^0(C,\tilde{J}_a/J_a)$$
$$\longrightarrow H^1(C,J_a)\longrightarrow H^1(C,\tilde{J}_a)\longrightarrow 1$$
since $\tilde{J}_a/J_a$, having punctual support, has trivial $H^1$. Let $\delta_a^{\mathrm{aff}} = \dim(\Jc^{0,\mathrm{aff}})$. We therefore have
\begin{align*}
\delta_a^{\mathrm{aff}} &= \dim(\Jc^0)-\dim(\Jc^{0,\mathrm{ab}})\\
&= \dim(H^1(C,J_a))-\dim(H^0(C,J_a))-\dim(H^1(C,\tilde{J}_a))\\
&= \dim(H^0(C,\tilde{J}_a/J_a))-\dim(H^0(C,\tilde{J}_a)),
\end{align*}
the last equality coming from the long exact sequence above.

By~\ref{thm:10.4.1}, we have
\begin{align*}
\dim(H^0(C,\tilde{J}_a/J_a)) &= \dim_k(H^0(C,(\tgo\otimes_k\pi_{Y_a,*}(\rho_{a,*}\oc_{X_a}/\oc_{Y_a}))^W))\\
&= \delta_a,
\end{align*}
where $\delta_a$ is the value at $a$ of the function $\delta$ considered in Section~\ref{sec:9}. It also follows from~\ref{thm:10.4.1} that one has
$$H^0(C,\tilde{J}_a) = T^{W_a},$$
where $W_a$ is the subgroup of $W^G$ defined in Section~\ref{S:7.2}. Combining the previous results, we have
$$\delta_a^{\mathrm{aff}} = \delta_a-\dim(T^{W_a}).$$

Let us now assume $a\notin\Ac_G^{\el}$. It follows that for a certain $M\in\lc^G$, one has $a\in\Ac_M$ and therefore $W_a\subset W^M$ (cf.\ Proposition~\ref{thm:8.2.1}). In particular, the center of $M$ is contained in $T^{W_a}$. We therefore have $\dim(T^{W_a})>0$ and
$$\delta_a^{\mathrm{aff}}<\delta_a.$$
Let us assume moreover that $a\in\Ac_G^{\mathrm{bon}}$ (cf.\ Definition~\ref{thm:9.1.1}). We therefore have $\delta_a\leq\mathrm{codim}_\Ac(a)$ which, combined with the previous inequality, gives
\begin{equation}\label{eq:10.5.1}
\delta_a^{\mathrm{aff}}<\mathrm{codim}_\Ac(a).
\end{equation}
Suppose that $a\notin\Ac_G^{\el}$ lies in the socle of the restriction of $Rf^\xi_*\Qlb[d_{\mc^\xi}]$ to $\Ac^{\mathrm{bon}}$. By Corollary~\ref{thm:10.3.2} (whose hypotheses we have checked at the beginning of the proof)), we have the inequality
$$\delta_a^{\mathrm{aff}}\geq\mathrm{codim}_\Ac(a).$$
But this manifestly contradicts~(\ref{eq:10.5.1}). Hence every $a\in\Ac_G$ which is in the socle of the restriction of $Rf^\xi_*\Qlb[d_{\mc^\xi}]$ to $\Ac^{\mathrm{bon}}$ belongs to the open subset $\Ac_G^{\el}$.
\end{preuve}
\end{paragr}

\begin{paragr}[The spaces $\Ac_{s,\rho}$.]\label{S:10.6}
In this paragraph, we define the space $\Ac_{s,\rho}$ associated with a ``geometric endoscopic datum'' of $G$ (cf.\ Definition~\ref{thm:10.6.3} below) and we give some of its properties. We do not assume $G$ semi-simple. The notation is that of the previous paragraph. Let $a\in\Ac_G$. We introduced in Section~\ref{S:7.2} the subgroup $W_a$ of $W^G$ which is the stabilizer of the irreducible component of $Y_a$ passing through the point $\infty_{Y_a}$. Let $W_a^0$ be the normal subgroup of $W_a$ generated by the stabilizers in $W_a$ of the points of the irreducible component of the normalization of $Y_a$ passing through the point $\infty_{Y_a}$. If $\bar{a}$ is a geometric point of $\Ac_G$ localized at $a$, the open subset $V_{\bar{a}}$, pointed by $\infty_{Y_{\bar{a}}}$ (with the notation of \S\ref{S:7.2}), is an étale covering of the open subset $U_{\bar{a}}$ pointed by $\infty$. The fiber of this covering at $\infty$ is endowed with an action of the fundamental group $\pi_1(U_{\bar{a}},\infty)$. Using the point $\infty_{Y_{\bar{a}}}$ of this fiber, one obtains a morphism
$$\pi_1(U_{\bar{a}},\infty)\to W_a.$$
The image of the inertia
$$\Ker(\pi_1(U_{\bar{a}},\infty)\to\pi_1(C\times_k k(\bar{a}),\infty))$$
is precisely the subgroup $W_a^0$. One deduces from it a morphism
\begin{equation}\label{eq:10.6.1}
\pi_1(C,\infty)\to W_a/W_a^0.
\end{equation}
We have an exact sequence
$$1\longrightarrow\pi_1(C,\infty)\longrightarrow\pi_1(C_0,\infty)\longrightarrow\Gal(k/\Fq)\longrightarrow 1$$
which is canonically split by the rational point $\infty$. We note that the homomorphism~(\ref{eq:10.6.1}) is stable under $\Gal(k/\Fq)$.

\begin{definition}\label{thm:10.6.1}
For every $s\in\Tc$, we introduce the following groups
\begin{itemize}
\item[--] $W_s$ is the subgroup of $W^G$ defined by
  $$W_s = \{w\in W^G\mid w\cdot s = s\};$$
\item[--] $W_s^0$ is the subgroup of $W^G$ generated by the simple reflections associated with the roots $\al\in\Phi^G$ such that $\al^\vee(s) = 1$; \item[--] $G_s$ is the connected reductive group over $k$ with Langlands dual $Z_{\Gc}(s)^0$, the connected centralizer of $s$ in $\Gc$.
\end{itemize}
\end{definition}

\begin{remark}\label{thm:10.6.2}
The subgroup $W_s^0$ is a normal subgroup of $W_s$. It is moreover identified with the Weyl group of $G_s$.
\end{remark}

\begin{definition}\label{thm:10.6.3}
We call \emph{geometric endoscopic datum} of $G$ the datum of a pair $(s,\rho)$ with
\begin{enumerate}
\item $s\in\Tc$; \item $\rho : \pi_1(C,\infty)\to W_s/W_s^0$ is a continuous homomorphism, fixed by $\Gal(k/\Fq)$.
\end{enumerate}
\end{definition}

Let $(s,\rho)$ be a geometric endoscopic datum of $G$ and let
$$\Ac_{s,\rho}$$
be the part of $\Ac$ formed of the $a$ for which the following conditions are satisfied:
\begin{itemize}
\item[--] $W_a\subset W_s$; \item[--] $W_a^0\subset W_s^0$ \item[--] the composite homomorphism
  $$\pi_1(C,\infty)\to W_a/W_a^0\to W_s/W_s^0$$
is equal to $\rho$.
\end{itemize}

\begin{proposition}\label{thm:10.6.4}
For a geometric endoscopic datum $(s,\rho)$ of $G$, the part $\Ac_{s,\rho}$ is closed in $\Ac_G$ and stable under the Frobenius $\tau$.
\end{proposition}

\begin{preuve}
Let $G_s$ be the connected reductive group over $k$ associated with $s$ (cf.\ Definition~\ref{thm:10.6.1}). The group $W_s$, which acts on $T$, also acts on the set $\Phi_T^{G_s}$ of roots of $T$ in $G_s$. The subgroup $W_s^0$ is identified with the Weyl group of $(G_s,T)$. Let $\Delta_T^{G_s}$ be a set of simple roots of $T$ in $G_s$. Let $W_s^{\mathrm{ext}}$ be the subgroup of $W_s$ which leaves $\Delta_T^{G_s}$ fixed. One then has a decomposition into a semi-direct product $W_s = W_s^0\rtimes W_s^{\mathrm{ext}}$. The set $\Delta_T^{G_s}$ is not unique, but one passes from one choice to another by the action of a unique element of $W_s^0$, so that changing $\Delta_T^{G_s}$ amounts to conjugating $W_s^{\mathrm{ext}}$ by an element $W_s^0$. One completes the datum of $\Delta_T^{G_s}$ into a pinning $(B_s,T,\{X_\al\}_{\al\in\Delta_T^{G_s}})$ of $G_s$. By this choice, one identifies the group $W_s^{\mathrm{ext}}$ with a subgroup, again denoted $W_s^{\mathrm{ext}}$, of the group of automorphisms of $G_s$ which leave the pinning stable. Let $\wc_s^{\mathrm{ext}}$ be the $W_s^{\mathrm{ext}}$-torsor over $C$ endowed with a trivialization of its fiber at $\infty$, deduced from
$$\rho : \pi_1(C,\infty)\to W_s/W_s^0\simeq W_s^{\mathrm{ext}}.$$
Since $W_s^{\mathrm{ext}}$ leaves $W_s^0$ invariant, this group acts on $\car_{G_s} = \tgo//W_s^0$ and one can introduce the $C$-scheme $\car_{G_s,\rho}$ which is the contracted product
$$\car_{G_s,\rho} = \car_{G_s}\times_{\mathrm{k}}^{W_s^{\mathrm{ext}}}\wc_s^{\mathrm{ext}}.$$
One also introduces its variant
$$\car_{G_s,\rho,D} = \lc_D\times_C^{\mathbb{G}_{m,C}}\car_{G_s,\rho}\,.$$
This scheme, fibered over $C$, is endowed with an isomorphism of its fiber at $\infty$ onto $\car_{G_s}$ (deduced from the canonical trivialization of $\lc_D$ and from the fixed trivialization of $\wc_s^{\mathrm{ext}}$). Imitating the construction of \S\ref{S:4.4}, one introduces the quasi-projective scheme $\Ac_{G_s,\rho}^{G\text{-}\reg}$ which classifies the pairs $a = (h_a,t)$, where
\begin{itemize}
\item[--] $h_a\in H^0(C,\car_{G_s,\rho,D})$; \item[--] $t\in\tgo^{G-\reg}$ with $\chi_{G_s}(t) = h_a(\infty)$.
\end{itemize}
One then has an obvious morphism
$$\chi_G^{G_s,\rho} : \Ac_{G_s,\rho}^{G\text{-}\reg}\to\Ac_G.$$
The proposition then follows from the following proposition, the stability under Frobenius being obvious.
\end{preuve}

\begin{proposition}\label{thm:10.6.5}
The morphism $\chi_G^{G_s,\rho}$ defined in the course of the proof of Proposition~\ref{thm:10.6.4} is a closed immersion with image $\Ac_{G_s,\rho}$.
\end{proposition}

\begin{preuve}
We refer the reader to \cite[props.\ 6.3.2 and 6.3.3]{Ngo-lemme}.
\end{preuve}

A particular case of Proposition~\ref{thm:10.6.5} is that of a geometric endoscopic datum $(s,\rho)$ whose homomorphism $\rho$ is the trivial homomorphism, denoted $1$. The closed subset $\Ac_{s,1}$ is then the image of the obvious morphism
$$\chi_G^{G_s} : \Ac_{G_s}^{G\text{-}\reg}\to\Ac_G,$$
where $\Ac_{G_s}^{G\text{-}\reg}$ is the ``$G$-regular'' open subset of the space $\Ac_{G_s}$ for the reductive group $G_s$ of Definition~\ref{thm:10.6.1} (cf.\ \S\ref{S:4.4}). Since this space $\Ac_{G_s}$ will no longer intervene in what follows, we prefer to reserve this notation for the $G$-regular open subset: we set, somewhat abusively,
$$\Ac_{G_s} = \Ac_{G_s}^{G\text{-}\reg}.$$
We introduce, as in Definition~\ref{thm:4.4.1}, an elliptic open subset $\Ac_{G_s}^{\el}$ of $\Ac_{G_s}$, which is nonempty for reasons of dimension.

\begin{proposition}\label{thm:10.6.6}
The locally closed part $\chi_G^{G_s}(\Ac_{G_s})\cap\Ac_G^{\el}$ is nonempty if and only if one has the equality
$$Z_{\Gc}^0 = Z_{\Gc_s}^0$$
between the connected centers of the dual groups $\Gc$ and $\Gc_s$. In case of equality, one has moreover
$$\chi_G^{G_s}(\Ac_{G_s})\cap\Ac_G^{\el} = \chi_G^{G_s}(\Ac_{G_s}^{\el}).$$
\end{proposition}

\begin{remark}\label{thm:10.6.7}
When $G$ is semi-simple, the dual group $\Gc$ is semi-simple too, and the condition on the centers is equivalent to the fact that $\Gc_s$, hence $G_s$, is semi-simple.
\end{remark}

\begin{preuve}
Let $H = G_s$, $a_H\in\Ac_H$ and $a = \chi_G^H(a)$. One then has $W_a = W_{a_H}$ by a variant of Lemma~\ref{thm:8.2.2}. Also, in the notation we shall not distinguish $a$ and $a_H$ and we shall omit $\chi_G^H$. The proof follows from the two lemmas below.
\end{preuve}

\begin{lemma}\label{thm:10.6.8}
For every $a\in\Ac_G^{\el}\cap\Ac_H$, one has $a\in\Ac_H^{\el}$ and $Z_{\Hc}^0 = Z_{\Gc}^0$.
\end{lemma}

\begin{preuve}
Let $M'$ be the smallest element of $\lc^H(T)$ such that $W_a\subset W^{M'}$. Let $M\in\lc^G(T)$ whose dual $\Mc$ is identified with the Levi subgroup $Z_{\Gc}(Z_{\widehat{M'}}^0)$ of $\Gc$ (this is the centralizer in $\Gc$ of the maximal central torus of $\widehat{M'}$). One therefore has $W_a\subset W^{M'}\subset W^M$ and $a\in\Ac_M$ (by Proposition~\ref{thm:8.2.1}). Since $a\in\Ac_G^{\el}$, one has $M = G$ and $Z_{\widehat{M'}}^0\subset Z_{\Gc}$, whence the equality $Z_{\Hc}^0 = Z_{\widehat{M'}}^0 = Z_{\Gc}$ and $a\in\Ac_H^{\el}$ (by an application of Proposition~\ref{thm:8.2.1} to $H$ and its Levi subgroup $M'$).
\end{preuve}

\begin{lemma}\label{thm:10.6.9}
If $Z_{\Hc}^0 = Z_{\Gc}^0$, one has $\Ac_H^{\el}\subset\Ac_G^{\el}$.
\end{lemma}

\begin{preuve}
Let $a\in\Ac_H^{\el}$ and $M$ the smallest element of $\lc^G(T)$ such that $W_a\subset W^M$. By Proposition~\ref{thm:10.6.5}, one has $W_a\subset W^H$. One therefore has $W_a\subset W^{M'}$ where $M'\in\lc^H(T)$ is defined dually by $\widehat{M'} = \Mc\cap\Hc$. Since $a$ is elliptic in $H$, one has $M' = H$, hence $\Hc\subset\Mc$, hence $Z_{\Mc}^0\subset Z_{\Hc}^0$. The condition $Z_{\Hc}^0 = Z_{\Gc}^0$ implies that $M = G$, that is to say $a\in\Ac_G^{\el}$.
\end{preuve}

\begin{proposition}\label{thm:10.6.10}
Assume $G$ semi-simple. The set of geometric endoscopic data $(s,\rho)$ of $G$ such that one has
\begin{equation}\label{eq:10.6.2}
\Ac_{s,\rho}\cap\Ac_G^{\el}\neq\emptyset
\end{equation}
is finite.
\end{proposition}

\begin{preuve}
As $s$ runs through $\Tc$, there are only finitely many possibilities for the groups $\Gc_s$, $W_s$ and $W_s^0$. Since the group $\pi_1(C,\infty)$ is topologically generated by $2g$ generators, there are only finitely many possible morphisms $\rho$. It therefore suffices to prove that there exists a finite subset of $\Tc$ containing all the $s$ which belong to an endoscopic datum $(s,\rho)$ for which the condition~(\ref{eq:10.6.2}) is satisfied. So let such an $s$ be given and let $Z_{\Gc_s}^{W_s}$ be the diagonalizable group of fixed points under $W_s$ of the center of $\Gc_s$. Since this group obviously contains $s$, it suffices to see that it is finite, in other words that its neutral component, denoted $\hat{S}$, is trivial. Let $\Mc$ be the centralizer in $\Gc$ of the torus $\hat{S}$. Its dual $M$ is a Levi subgroup of $G$. Since $W_s$ centralizes $\hat{S}$, one has $W_s\subset W^M$. By the hypothesis~(\ref{eq:10.6.2}), there exists $a\in\Ac_{s,\rho}\cap\Ac_G^{\el}$. One therefore has $W_a\subset W_s\subset W^M$. By Proposition~\ref{thm:8.2.1}, one has $a\in\Ac_M$. Since $a$ is elliptic in $G$, it follows that $M = G$. Hence $\hat{S}$ is a central torus of $\Gc$, which is assumed to be semi-simple. Thus $\hat{S}$ is trivial, which is what we had to prove.
\end{preuve}
\end{paragr}

\begin{paragr}[The cohomological theorem.]\label{S:10.7}
We continue with the same notation and hypotheses. We assume here that $G$ is \emph{semi-simple}. Let $s\in\Tc$ and let $G_s$ be the connected reductive group over $k$ of Definition~\ref{thm:10.6.1}. We endow $G_s$ with the action of the Frobenius $\tau$ which induces the trivial action on the group of characters of $T$.

For a $k$-scheme $X$ endowed with an action of the Frobenius $\tau$, we denote by $\mathrm{Perv}(X)$ the category of perverse sheaves $F$ on $X$ endowed with an isomorphism $\tau^*F\simeq F$. Here is a theorem due to Ngô (cf.\ \cite[Theorems 6.4.1 and 6.4.2]{Ngo-lemme}).

\begin{theorem}\label{thm:10.7.1}
Assume $G$ semi-simple. Let $s\in\Tc$ be such that $G_s$ is semi-simple. Let $\chi_G^{G_s,\el} : \Ac_{G_s}^{\el}\hookrightarrow\Ac_G^{\el}$ be the restriction of $\chi_G^{G_s}$ to the open subset $\Ac_{G_s}^{\el}$.

One then has the following equality in the Grothendieck group of $\mathrm{Perv}(\Ac_{G_s}^{\el})$:
$$(\chi_G^{G_s,\el})^{*\mathrm{p}}\hc^\bullet(Rf_{G,*}^{\el}\Qlb)_s = {}^{\mathrm{p}}\hc^\bullet(Rf_{G_s,*}^{\el}\Qlb)_1[-2r_s](-r_s),$$
where
\begin{itemize}
\item[--] $r_s = \mathrm{codim}_{\Ac_G}(\Ac_{G_s})$; \item[--] ${}^{\mathrm{p}}\hc^\bullet(Rf_{G,*}^{\el}\Qlb)_s$ is the $s$-part of ${}^{\mathrm{p}}\hc^\bullet(Rf_{G,*}^{\el}\Qlb)$ defined in Theorem~\ref{thm:8.5.1}; \item[--] ${}^{\mathrm{p}}\hc^\bullet(Rf_{G_s,*}^{\el}\Qlb)_1$ is the $1$-part of ${}^{\mathrm{p}}\hc^\bullet(Rf_{G_s,*}^{\el}\Qlb)$ defined in Theorem~\ref{thm:8.5.1} for the group $G_s$ $($with $1$ the neutral element of $\Tc)$.
\end{itemize}
\end{theorem}

\begin{remark}\label{thm:10.7.2}
Since $G$ and $G_s$ are semi-simple, so are their duals $\Gc$ and $\Gc_s = Z_{\Gc}(s)^0$. Consequently the centers of the latter are finite. We therefore have $Z_{\Gc}^0 = Z_{\Gc_s}^0 = \{1\}$ and $s$ is a torsion element of $\Tc$.
\end{remark}

Let $M\in\lc^G(T)$. We defined in Section~\ref{S:8.4} an open subset $\uc_M$ of $\Ac_G$. This open subset is clearly $\tau$-stable. The open subset $\Ac_G^{\mathrm{bon}}$ of Definition~\ref{thm:9.1.1} is also $\tau$-stable. For every $s\in\Tc$, the closed subset $\bigcup_\rho\Ac_{s,\rho}$, where $\rho$ runs through the finite set of continuous and nontrivial morphisms
$$\rho : \pi_1(C,\infty)\to W_s/W_s^0,$$
is $\tau$-stable. Let $\bar{s}\in(\Tc/Z_{\Mc}^0)_{\tors}$ and
\begin{equation}\label{eq:10.7.1}
\uc_{\bar{s}} = \uc_M\cap\Ac_G^{\mathrm{bon}}\cap\Big(\Ac_G-\bigcup_{(s,\rho)}\Ac_{s,\rho}\Big),
\end{equation}
where the union is taken over the \emph{finite} set (cf.\ Proposition~\ref{thm:10.6.10}) of geometric endoscopic data $(s,\rho)$ which satisfy
\begin{itemize}
\item[--] $s\in\bar{s}Z_{\Mc}^0$; \item[--] $\Ac_G^{\el}\cap\Ac_{s,\rho}\neq\emptyset$; \item[--] $\rho$ is a nontrivial homomorphism.
\end{itemize}
Then $\uc_{\bar{s}}$ is a $\tau$-stable open subset of $\Ac_G$. Let
$$\uc_{\bar{s}}^{\el} = \uc_{\bar{s}}\cap\Ac_G^{\el}.$$
For every $s\in\bar{s}Z_{\Mc}^0$ such that $G_s$ is semi-simple, we have a commutative diagram, where the arrows denoted $\chi$ are closed immersions and the other arrows are open immersions,
$$\xymatrix{
\uc_{\bar{s}}^{\el}\cap\Ac_{G_s}^{\el} \ar[r]^-{i_{G_s}^{\bar{s}}} \ar[d]_{\chi_{\bar{s}}^{s}} & \Ac_{G_s}^{\el} \ar[d]^{\chi_G^{G_s,\el}}\\
\uc_{\bar{s}}^{\el} \ar[r]^-{i_G^{\bar{s}}} \ar[d]_{j_{\bar{s}}} & \Ac_G^{\el} \ar[d]^{j}\\
\uc_{\bar{s}} \ar[r]^-{i_M^{\bar{s}}} & \uc_M.
}$$

\begin{theorem}\label{thm:10.7.3}
Assume $G$ semi-simple and $\xi\in\ago_T$ in general position. For every $\bar{s}\in(\Tc/Z_{\Mc}^0)_{\tors}$, one has the following equality in the Grothendieck group of $\mathrm{Perv}(\uc_{\bar{s}})$:
$$(i_M^{\bar{s}})^*({}^{\mathrm{p}}\hc^\bullet(Rf_*^{\xi,M\text{-par}}\Qlb)_{\bar{s}}) = \bigoplus_s (j_{\bar{s}})_{!*}(\chi_{\bar{s}}^s)_*(i_{G_s}^{\bar{s}})^{*\mathrm{p}}\hc^\bullet(Rf_{G_s,*}^{\el}\Qlb)_1[-2r_s](-r_s),$$
where
\begin{itemize}
\item[--] the sum runs over the $s\in\bar{s}Z_{\Mc}^0$ such that $G_s$ is semi-simple; \item[--] the $\bar{s}$-part ${}^{\mathrm{p}}\hc^\bullet(Rf_*^{\xi,M\text{-par}}\Qlb)_{\bar{s}}$ is the one defined in Theorem~\ref{thm:8.5.1}; \item[--] the other notation is that used above or in Theorem~\ref{thm:10.7.1}.
\end{itemize}
\end{theorem}

\begin{preuve}
By Theorem~\ref{thm:10.5.1}, we have a canonical isomorphism
$$(i_M^{\bar{s}})^*({}^{\mathrm{p}}\hc^\bullet(Rf_*^{\xi,M\text{-par}}\Qlb)_{\bar{s}})\cong(j_{\bar{s}})_{!*}(j_{\bar{s}})^*(i_M^{\bar{s}})^*({}^{\mathrm{p}}\hc^\bullet(Rf_*^{\xi,M\text{-par}}\Qlb)_{\bar{s}}).$$
By Corollary~\ref{thm:8.5.3}, we have
$$(j_{\bar{s}})^*(i_M^{\bar{s}})^*({}^{\mathrm{p}}\hc^\bullet(Rf_*^{\xi,M\text{-par}}\Qlb)_{\bar{s}}) = \bigoplus_s (i_G^{\bar{s}})^{*\mathrm{p}}\hc^\bullet(Rf_{G,*}^{\el}\Qlb)_s,$$
where the sum is taken over the $s\in\bar{s}Z_{\Mc}^0$ of finite order. For every such $s$, according to Ngô (cf.\ \cite[\S6.4]{Ngo-lemme}), ${}^{\mathrm{p}}\hc^\bullet(Rf_{G,*}^{\el}\Qlb)_s$ is supported in the closed subset $(\bigcup_\rho\Ac_{G_s,\rho})\cap\Ac_G^{\el}\subset\Ac_G^{\el}$ where $\rho$ runs through the continuous homomorphisms from $\pi_1(C,\infty)$ to $W_s/W_s^0$. Consequently, $(i_G^{\bar{s}})^{*\mathrm{p}}\hc^\bullet(Rf_{G,*}^{\el}\Qlb)_s$ is supported in the closed subset $\uc_{\bar{s}}^{\el}\cap\Ac_{G_s}^{\el}$. If this closed subset is nonempty, the intersection $\Ac_{G_s}\cap\Ac_G^{\el}$ is nonempty, which implies that $G_s$ is semi-simple (cf.\ Proposition~\ref{thm:10.6.6} and the remark which follows). For such an $s$, we have
$$(i_G^{\bar{s}})^{*\mathrm{p}}\hc^\bullet(Rf_{G,*}^{\el}\Qlb)_s = (\chi_{\bar{s}}^s)_*(i_{G_s}^{\bar{s}})^*(\chi_G^{G_s,\el})^{*\mathrm{p}}\hc^\bullet(Rf_{G,*}^{\el}\Qlb)_s.$$
We use Ngô's theorem (cf.\ Theorem~\ref{thm:10.7.1}) to conclude.
\end{preuve}
\end{paragr}
\section{The weighted orbital \texorpdfstring{$s$}{s}-integrals}\label{sec:11}

\begin{paragr}\label{S:11.1}
In~\S\ref{S:4.1} we fixed a projective, smooth and connected curve $C$ over $k$, as well as a divisor $D$ and a point $\infty\in C(k)$. From now on we assume that $C$ comes by base change from a curve $C_0$ over $\Fq$. Let $\tau$ be the Frobenius automorphism of $k$ given by $x\mapsto x^q$. We again denote by $\tau$ the automorphism $1\times\tau$ of $C=C_0\times_{\Fq}k$. We assume that $D$ and $\infty$ are fixed under $\tau$.

Let $F$ and $F_0$ be the function fields of $C$ and $C_0$ respectively. We have $F=F_0\otimes_{\Fq}k$, whence an action of $\tau$ on $F$ given by $1\otimes\tau$, whose subfield of fixed points is $F_0$. For every closed point $v_0$ of $C_0$ let $F_{0,v_0}$ be the completion of $F_0$ at $v_0$ and $\oc_{0,v_0}$ its ring of integers. Analogous notations hold for $F$. The completed tensor product
$$
F\hat{\otimes}_{F_0}F_{0,v_0}=\prod_{v\in V,\ v|v_0}F_v
$$
decomposes as the product of the completions $F_v$ for the closed points $v$ of $C$ above $v_0$. We endow this tensor product with the continuous action of $\tau$ given by $\tau\otimes 1$. This action admits $F_{0,v_0}$ as its subring of fixed points. We deduce from it an action of $\tau$ on the ring of adèles $\AAA$ of $F$ whose ring of fixed points is the ring $\AAA_0$ of adèles of $F_0$. This action respects the product
$$
\oc=\prod_v\oc_v
$$
over all the closed points of $C$. We define $\oc_0$ in the same way. Moreover we have $\oc_0=\AAA_0\cap\oc$.

We endow the root datum of the pair $(G,T)$ of~\S\ref{S:2.1} with the trivial action of $\tau$. The pair $(G,T)$ then comes by base change from an analogous pair defined over $\Fq$ where the torus is split over $\Fq$. This endows the pair $(G,T)$ with an action of $\tau$. All the elements of $\pc^G(T)$ or of $\lc^G(T)$, that is to say the parabolic or Levi subgroups of $G$, are stable under the action of $\tau$. For every scheme $S$ defined over $k$ and every $k$-algebra $A$, both endowed with an action of $\tau$, we denote by $S(A)^\tau$ the subset of points fixed under the action of $\tau$. If, moreover, $S$ is a group scheme (and $\tau$ of course acts as an automorphism of group schemes), $S(A)^\tau$ is a group and we denote by $S(A)_\tau$ the set of $\tau$-conjugacy classes (the $\tau$-conjugation being given by $g\cdot x=gx\tau(g)^{-1}$), pointed by the class of the neutral element. If, moreover, $S$ is an abelian group scheme, $S(A)_\tau$ is an abelian group, namely the group of co-invariants.
\end{paragr}

\begin{paragr}\label{S:11.2}
Here is another variant of a result of Kottwitz already mentioned (cf.\ \cite[Lemma 2.2]{Kottwitz-isocrystals}).

\begin{proposition}\label{thm:11.2.1}
Let $U_0$ be an open subset of $C_0$ which contains the geometric point $\infty$. Let $X_*$, $A$ and $B$ be the functors from the category of torus schemes over $U_0$ to that of abelian groups defined respectively, for every torus scheme $\tc$ over $U_0$, by
$$
X_*(\tc)=X_*(\tc_\infty),
$$
where the right-hand side denotes the group of cocharacters of the fiber of $\tc$ at the geometric point $\infty$,
$$
A(\tc)=X_*(\tc_\infty)_{\pi_1(U_0,\infty)},
$$
where the index means that one takes the co-invariants under the fundamental group of $U_0$ pointed by $\infty$, and
$$
B(\tc)=(\tc(F)\back\tc(\AAA))_\tau,
$$
where the index $\tau$ denotes the group of co-invariants.

Let $\varpi_\infty$ be a uniformizer of $F_\infty$. For every $\tc$ as above, the morphism which to $\lambda\in X_*(\tc_\infty)$ associates the image $\varpi_\infty^\lambda$ in the group $(\tc(F)\back\tc(\AAA))_\tau$ does not depend on the choice of $\varpi_\infty$ and defines a morphism of functors
\begin{equation}\label{eq:11.2.1}
X_*\to B.
\end{equation}
There exists a unique isomorphism of functors $A\to B$ which, composed with the canonical morphism $X_*\to A$, gives the morphism~(\ref{eq:11.2.1}).
\end{proposition}

\begin{preuve}
It rests on the methods of Kottwitz (cf.\ \cite[proof of Lemma 2.2]{Kottwitz-isocrystals}). The functor $B$ has the following two fundamental properties: it is right exact and for an induced torus one has $B(\tc)=\ZZ$. Right exactness is shown as in \emph{loc. cit.} \S1.9: the key point is that the group $H^1(F,\tc)$ is trivial for every $F$-torus (and likewise if one replaces $F$ by a completion $F_v$). The second property follows from the following fact: let $E_0$ be a finite extension of $F$ and $E=E_0\otimes_{F_0}F$; one has $(E^\times\back\AAA_E^\times)_\tau\simeq\ZZ$ where $\AAA_E$ is the ring of adèles of $E$. By a lemma of Shapiro, up to replacing $\tau$ by a power, we reduce to the case where $E$ is a field. Let $\AAA_E^1$ be the kernel of the degree homomorphism $\AAA_E^\times\to\ZZ$. The quotient $E^\times\back\AAA_E^1/\oc_E^\times$, where $\oc_E^\times$ is the product of the invertible elements of the completions of the local rings, is the group of $k$-points of the Jacobian of the projective, smooth and connected curve over $k$ whose function field is $E$. By Lang's isogeny theorem, these points lie in the image of $1-\tau$. The same holds for the points of $\oc_E^\times$ (cf.\ \cite[proof of Prop.\ 2.3]{Kottwitz-isocrystals}). It is then easy to see that the degree homomorphism induces a bijection $(E^\times\back\AAA_E^\times)_\tau\simeq\ZZ$.

The morphism $X_*\to B$ does not depend on the choice of $\varpi_\infty$, by a variant of Lang's isogeny theorem. It follows from the proof of \emph{loc. cit.} Lemma 2.2 that the canonical morphism
$$
\Hom(A,B)\to\Hom(X_*,B)
$$
is an isomorphism which sends isomorphism to isomorphism. Moreover, for a morphism $\Hom(X_*,B)$ to be an isomorphism, it is necessary and sufficient that it send a generator of $X_*(\mathbb{G}_{m,U_0})$ to a generator of $(F^\times\back\AAA_F^\times)_\tau$. The proposition follows.
\end{preuve}
\end{paragr}

\begin{paragr}[Characteristic $a_M$ and group scheme $J_{a_M}$.]\label{S:11.3}
The Frobenius automorphism $\tau$ acts on the scheme $\Ac_G$ (cf.\ \S\ref{S:4.4}) and respects the closed subschemes $\Ac_M$ for $M\in\lc^G$. Let $M\in\lc^G(T)$ be a Levi subgroup and $a_M=(h_{a_M},t)\in\Ac_M(k)^\tau$. The open subset
$$
U_{a_M}=h_{a_M}^{-1}(\car_{M,D}^{G\textrm{-}\reg})
$$
is stable under $\tau$ and hence descends to an open subset of $C_0$, denoted $U_{a_M,0}$.

Let $\varepsilon_M$ be a Kostant section of the characteristic morphism $\chi_M$ (cf.\ \S\ref{S:2.3}). Over $\car_M$ we have the smooth commutative group scheme (cf.\ \S\ref{S:3.2})
$$
J^M=\varepsilon_M^*I^M.
$$
By the construction of~\S\ref{S:5.1}, we deduce from it a group scheme $J_D^M$ over $\car_{M,D}$. Let
$$
J_{a_M}=h_{a_M}^*J_D^M.
$$
This is a smooth commutative group scheme over $C$: it is even a torus scheme over the open subset $U_{a_M}$ (and even over a generally larger open subset). One can check that it descends to a smooth commutative group scheme over $C_0$, which is a torus scheme over $U_{a_M,0}$.
\end{paragr}

\begin{paragr}[The invariant $\inv_{a_M}$.]\label{S:11.4}
The short exact sequence
$$
1\longrightarrow J_{a_M}(\AAA)\longrightarrow G(\AAA)\longrightarrow J_{a_M}(\AAA)\back G(\AAA)\longrightarrow 1
$$
gives a long exact sequence in the cohomology of pointed sets
\begin{equation}\label{eq:11.4.1}
\begin{split}
1\longrightarrow J_{a_M}(\AAA_0)\longrightarrow G(\AAA_0)&\longrightarrow
(J_{a_M}(\AAA)\back G(\AAA))^\tau\\
&\longrightarrow \Ker(J_{a_M}(\AAA)_\tau\to G(\AAA)_\tau)\longrightarrow 1,
\end{split}
\end{equation}
where the coboundary morphism
\begin{equation}\label{eq:11.4.2}
(J_{a_M}(\AAA)\back G(\AAA))^\tau\longrightarrow J_{a_M}(\AAA)_\tau
\end{equation}
associates to an element of $G(\AAA)$ which lifts an element of $(J_{a_M}(\AAA)\back G(\AAA))^\tau$ the $\tau$-conjugacy class of $g\tau(g)^{-1}$ in $J_{a_M}(\AAA)$. Let us note that exactness at the third term means that two elements of $(J_{a_M}(\AAA)\back G(\AAA))^\tau$ have the same image in $J_{a_M}(\AAA)_\tau$ if and only if they differ by a right translation by an element of $G(\AAA_0)$.

Proposition~\ref{thm:11.2.1}, applied to the open subset $U_{a_M,0}$ and to the torus scheme over $U_{a_M,0}$ deduced from $J_{a_M}$ by descent, gives an isomorphism
$$
(J_{a_M}(F)\back J_{a_M}(\AAA))_\tau\simeq X_*(J_{a_M,\infty})_{\pi_1(U_{a_M,0},\infty)}.
$$

The notations and the constructions of~\S\ref{S:7.1} are also valid for the group $M$. We thus have a finite, étale covering $V_{a_M}\to U_{a_M}$ with Galois group $W^M$ and a closed point $\infty_{Y_{a_M}}$ of $V_{a_M}$. All these objects are stable under $\tau$ and descend to objects over $\Fq$, denoted by an index $0$.

Reasoning as in the proof of Proposition~\ref{thm:7.5.1}, we obtain identifications $X_*(J_{a_M,\infty})\simeq X_*(T)$ and
$$
X_*(J_{a_M,\infty})_{\pi_1(U_{a_M,0},\infty)}\simeq X_*(T)_{W_\infty},
$$
where $W_\infty$ is the stabilizer in $W^M$ of the connected component of $V_{a_M,0}$ which contains the point $\infty_{Y_{a_M},0}$. The exact sequence
$$
1\longrightarrow\pi_1(U_{a_M},\infty)\longrightarrow\pi_1(U_{a_M,0},\infty) \longrightarrow\Gal(k/\Fq)\longrightarrow 1
$$
is canonically split by the datum of the point $\infty_0\in C_0(\Fq)$ deduced from $\infty$. We thus have a semi-direct product
$$
\pi_1(U_{a_M,0},\infty)=\pi_1(U_{a_M},\infty)\rtimes\Gal(k/\Fq).
$$
The action of $\pi_1(U_{a_M,0},\infty)$ on the fiber of the covering $V_{a_M,0}\to U_{a_M,0}$ above the geometric point $\infty$ is trivial on the factor $\Gal(k/\Fq)$. We deduce that the group $W_\infty$ is none other than the group $W_{a_M}\subset W^M$ (defined as in~\S\ref{S:7.2} but relative to the group $M$). Putting together the preceding isomorphisms, we obtain an isomorphism
\begin{equation}\label{eq:11.4.3}
(J_{a_M}(F)\back J_{a_M}(\AAA))_\tau\simeq X_*(T)_{W_{a_M}}.
\end{equation}

\begin{definition}\label{thm:11.4.1}
Let $a_M\in\Ac_M(k)^\tau$. Let
$$
\inv_{a_M}:(J_{a_M}(\AAA)\back G(\AAA))^\tau\to X_*(T)_{W_{a_M}}
$$
be the morphism defined by composing the coboundary morphism~(\ref{eq:11.4.2}) with the canonical morphism $(J_{a_M}(\AAA))_\tau\to(J_{a_M}(F)\back J_{a_M}(\AAA))_\tau$ and then with the isomorphism~(\ref{eq:11.4.3}).
\end{definition}

Let $M_{\scnx}$ be the simply connected covering of the derived group of $M$. Let $T_{M_{\scnx}}$ be the inverse image of $T$ in $M_{\scnx}$: it is a maximal subtorus of $M_{\scnx}$.

\begin{proposition}\label{thm:11.4.2}
Let $a_M\in\Ac_M(k)^\tau$. The morphism $\inv_{a_M}$ takes its values in the image of the canonical morphism
$$
X_*(T_{M_{\scnx}})_{W_{a_M}}\to X_*(T)_{W_{a_M}}.
$$
\end{proposition}

Before giving the proof, let us state some auxiliary lemmas.

\begin{lemma}\label{thm:11.4.3}
Let $a_M\in\Ac_M(k)^\tau$. We have
\begin{equation}\label{eq:11.4.4}
\Ker(J_{a_M}(\AAA)_\tau\to G(\AAA)_\tau)=\Ker(J_{a_M}(\AAA)_\tau\to M(\AAA)_\tau).
\end{equation}
\end{lemma}

\begin{preuve}
Only the inclusion $\subset$ is not trivial. By Shapiro's lemma, up to replacing $\tau$ by one of its powers, we reduce to a local problem at a closed point of $C$ fixed by $\tau$. Let $g\in G(F_v)$ be such that $g\tau(g)^{-1}\in M(F_v)$. Let $P\in\pc(M)$ and let $B\in\pc(T)$ be a Borel subgroup such that $B\subset P$. A pair $(P,M)$ which contains $(B,T)$ is said to be standard. The pair $\Int(g^{-1})(P,M)$ is $\tau$-stable, hence conjugate by an element of $G(F_v)^\tau$ to a standard pair. Up to translating $g$ on the right by an element of $G(F_v)^\tau$, we are reduced to the case where $\Int(g^{-1})(P,M)$ is standard, hence equal to $(P,M)$. It follows that we may assume $g\in M(F_v)$. If, moreover, $g\tau(g)^{-1}\in M(\oc_v)$, by a variant of Lang's theorem there exists $m\in M(\oc_v)$ such that $g\tau(g)^{-1}=m\tau(m)^{-1}$. This proves the lemma.
\end{preuve}

We define a commutative group scheme over $\mgo$ by
$$
I^{M_{\scnx}}=\{(m,X)\in M_{\scnx}\times_k\mgo\mid\Ad(m)X=X\}.
$$
We deduce from it a smooth commutative group scheme over $\car_M$
$$
J^{M_{\scnx}}=\varepsilon_M^*I^{M_{\scnx}}
$$
and a variant over $\car_{M,D}$, denoted $J_D^{M_{\scnx}}$. Let
$$
J_{a_M}^{M_{\scnx}}=h_{a_M}^*J_D^{M_{\scnx}}.
$$
Let us note that the obvious morphism $I^{M_{\scnx}}\to I^M$ induces morphisms $J^{M_{\scnx}}\to J^M$ and $J_{a_M}^{M_{\scnx}}\to J_{a_M}$.

\begin{lemma}\label{thm:11.4.4}
Let $a_M\in\Ac_M(k)^\tau$. We have the inclusion
$$
\Ker(J_{a_M}(\AAA)_\tau\to M(\AAA)_\tau)\subset \operatorname{Im}(J_{a_M}^{M_{\scnx}}(\AAA)_\tau\to J_{a_M}(\AAA)_\tau),
$$
where $\operatorname{Im}$ denotes the image.
\end{lemma}

\begin{preuve}
As in the proof of Lemma~\ref{thm:11.4.3}, we reduce to a local problem at a closed point $v$ of $C$ stable under $\tau$. By base change, we view $J_{a_M}$ and $M$ as groups over $F_v$. Let $x\in J_{a_M}(F_v)$ and $m\in M(F_v)$ be such that $x=m\tau(m)^{-1}$. We then have
$$
\Coker(M_{\scnx}\to M)\simeq\Coker(J_{a_M}^{M_{\scnx}}\to J_{a_M}).
$$
Since $J_{a_M}^{M_{\scnx}}$ is a torus over $F_v$, we know that, for such a field $F_v$, the group $H^1(F_v,J_{a_M}^{M_{\scnx}})$ is trivial. It follows that there exist $m'\in M_{\scnx}(F_v)$ and $j\in J_{a_M}(F_v)$ such that $m=jm'$. Hence, up to $\tau$-conjugating $x$ by $j$, we may assume that $m\in M_{\scnx}(F_v)$ and therefore $x\in J_{a_M}^{M_{\scnx}}(F_v)$. For almost every $v$, $J_{a_M}^{M_{\scnx}}$ is a torus over $\oc_v$ and an analogous argument shows that if $x\in J_{a_M}(\oc_v)$ and $m\in M(\oc_v)$ satisfy $x=m\tau(m)^{-1}$ then $x$ is $\tau$-conjugate by an element of $J_{a_M}(\oc_v)$ to an element of $J_{a_M}^{M_{\scnx}}(\oc_v)$.
\end{preuve}

\begin{preuve}[of Proposition~\ref{thm:11.4.2}]
As above, we have an identification
$$
X_*(J_{a_M,\infty}^{M_{\scnx}})_{\pi_1(C_0,\infty)}\simeq X_*(T_{M_{\scnx}})_{W_{a_M}}.
$$
By Proposition~\ref{thm:11.2.1}, we obtain a commutative diagram in which the horizontal arrows are isomorphisms
$$
\xymatrix{ (J_{a_M}^{M_{\scnx}}(F)\back J_{a_M}^{M_{\scnx}}(\AAA))_\tau \ar[r]\ar[d] &
X_*(T_{M_{\scnx}})_{W_{a_M}} \ar[d]\\
(J_{a_M}(F)\back J_{a_M}(\AAA))_\tau \ar[r] & X_*(T)_{W_{a_M}}. }
$$
Let $g\in(J_{a_M}(\AAA)\back G(\AAA))^\tau$. By the exact sequence~(\ref{eq:11.4.1}), the class of $g\tau(g)^{-1}$ in $(J_{a_M}(\AAA))_\tau$ belongs to the kernel $\Ker(J_{a_M}(\AAA)_\tau\to G(\AAA)_\tau)$, hence to the image of $(J_{a_M}^{M_{\scnx}}(\AAA))_\tau$ by the combination of Lemmas~\ref{thm:11.4.3} and~\ref{thm:11.4.4}. Proposition~\ref{thm:11.4.2} is then clear.
\end{preuve}
\end{paragr}

\begin{paragr}[Haar measures.]\label{S:11.5}
The group $J_{a_M}(\AAA)_\tau$ is countable: indeed, it appears in the exact sequence
$$
J_{a_M}(F)_\tau\longrightarrow J_{a_M}(\AAA)_\tau\longrightarrow (J_{a_M}(F)\back J_{a_M}(\AAA))_\tau
$$
whose extreme terms are countable (cf.\ Proposition~\ref{thm:11.2.1}). We then endow the kernel $\Ker(J_{a_M}(\AAA)_\tau\to G(\AAA)_\tau)$ with the counting measure. The locally compact groups $G(\AAA_0)$ and $J_{a_M}(\AAA_0)$ are endowed with the Haar measures which give volume $1$ respectively to the compact subgroup $G(\oc_0)$ and to the maximal compact subgroup of $J_{a_M}(\AAA_0)$. By the sequence~(\ref{eq:11.4.1}), we deduce from this a measure on $(J_{a_M}(\AAA)\back G(\AAA))^\tau$ which, by construction, is $G(\AAA_0)$-invariant for the right action of that group. More precisely, the set $(J_{a_M}(\AAA)\back G(\AAA))^\tau$ is a disjoint and countable union of orbits under $G(\AAA_0)$; the orbit of $g\in(J_{a_M}(\AAA)\back G(\AAA))^\tau$ is isomorphic as a measurable set to the quotient $(g^{-1}J_{a_M}(\AAA_0)g)\back G(\AAA_0)$ endowed with the quotient measure (conjugation by $g$ induces a bijection from $g^{-1}J_{a_M}(\AAA_0)g$ onto $J_{a_M}(\AAA_0)$, which determines, by transport, a measure on $g^{-1}J_{a_M}(\AAA_0)g$).
\end{paragr}

\begin{paragr}[Arthur's weight.]\label{S:11.6}
For every $P\in\pc^G(M)$, we have the Levi decomposition $P=MN_P$ and the Iwasawa decomposition $G(\AAA)=P(\AAA)G(\oc)$. We then define a map $H_P$ from $G(\AAA)$ to $X_*(M)$ as follows: for every $\lambda\in X^*(P)=X^*(M)$ and every $g\in G(\AAA)$
$$
\lambda(H_P(g))=-\deg(\lambda(p)),
$$
where $p\in P(\AAA)$ is such that $g\in pG(\oc)$.

\begin{remark}\label{thm:11.6.1}
The function $H_P$ restricted to $M(\AAA)$ no longer depends on $P$: we denote it by $H_M$. Moreover, for all $m\in M(\AAA)$ and $g\in G(\AAA)$ we have the equality $H_P(mg)=H_M(m)+H_p(g)$.
\end{remark}

We recall (cf.\ \S\ref{S:2.4}) that $\ago_T$ is endowed with a $W^G$-invariant scalar product and that all its subspaces are endowed with the Euclidean measure. The following definition obviously depends on the choice of this scalar product.

\begin{definition}\label{thm:11.6.2}
For every $g\in G(\AAA)$, Arthur's weight $v_M^G(g)$ is the volume in $\ago_M^G$ of the convex hull of the projections onto $\ago_M^G$ of the points $-H_P(g)$ for $P\in\pc^G(M)$.
\end{definition}

\begin{remarks}\label{thm:11.6.3}
We have $\ago_M=a_M^G\oplus\ago_G$ and the projection considered above is relative to this decomposition. The weight is by construction right invariant under $G(\oc)$. It is also left invariant under $M(\AAA)$ by virtue of Remark~\ref{thm:11.6.1}. Let us also note that $v_G^G(g)=1$ for every $g\in G(\AAA)$.
\end{remarks}
\end{paragr}

\begin{paragr}[the weighted orbital $s$-integrals.]\label{S:11.7}
We have $D=\sum_{v\in V}d_vv$, where the sum is taken over the closed points of $C$ and the integers $d_v$ are even, non-negative and almost all zero (by hypothesis, cf.\ \S\ref{S:4.1}, $D$ is effective and even). Let $\varpi^{-D}=(\varpi_v^{-d_v})_v\in\AAA$, where the family is taken over all the closed points $v$ of $C$. Let $\mathbf{1}_D$ be the characteristic function of $\varpi^{-D}\ggo(\oc)$. For every $a_M\in\Ac_M(k)$, let $a_{M,\eta}$ be the restriction of $a_M$ to the generic point of $C$, which we view as a point of $\car_M(F)$. Let
$$
X_{a_M}=\varepsilon_M(a_{M,\eta})\in\mgo(F).
$$
Let us note that if $a_M$ is fixed by $\tau$, the same holds for $X_{a_M}$. Let us also note that the fiber of $J_{a_M}$ above $\eta$ is the centralizer of $X_{a_M}$ in $M\times_kF$.

\begin{definition}\label{thm:11.7.1}
Let $M\in\lc^G(T)$ and $a_M\in\Ac_M(k)^\tau$. Let $s\in\Tc^{W_{a_M}}$. The weighted orbital $s$-integral associated to $a_M$ is defined by
$$
J_M^{G,s}(a_M)=q^{-\deg(D)|\Phi^G|/2} \int_{(J_{a_M}(\AAA)\back G(\AAA))^\tau}\bg s,\inv_{a_M}(g)\bd \mathbf{1}_D(\Ad(g^{-1})X_{a_M})v_M^G(g)\,dg.
$$
\end{definition}

\begin{remarks}\label{thm:11.7.2}
One can show that the integrand has compact support in $(J_{a_M}(\AAA)\back G(\AAA))^\tau$. The integral is therefore finite. We have chosen to add the factor $q^{-\deg(D)|\Phi^G|/2}$ in order to obtain uniform local formulas.
\end{remarks}

\begin{proposition}\label{thm:11.7.3}
With the notations above, the map
$$
s\in\Tc^{W_{a_M}}\mapsto J_M^{G,s}(a_M)
$$
is invariant under translation by $Z_{\Mc}$.
\end{proposition}

\begin{preuve}
We have $W_{a_M}\subset W^M$ and hence $Z_{\Mc}\subset\Tc^{W_{a_M}}$. The canonical morphism
$$
\Hom_\ZZ(X_*(T)_{W_{a_M}},\CC^\times)\to \Hom_\ZZ(X_*(T_{M_{\scnx}})_{W_{a_M}},\CC^\times)
$$
is none other than the canonical morphism $\Tc^{W_{a_M}}\to\Tc^{W_{a_M}}/Z_{\Mc}$. The proposition then follows from Proposition~\ref{thm:11.4.2}.
\end{preuve}
\end{paragr}
\section{A first computation of a Frobenius trace}\label{sec:12}

\begin{paragr}[Statement of the theorem.]\label{S:12.1}
We keep the notation and the hypotheses of Section~\ref{sec:11}. We assume moreover that the group $G$ is \emph{semi-simple}. Let $\xi\in\ago_T$ \emph{in general position} in the sense of Remark~\ref{thm:4.9.2}. Let $M\in\lc^G(T)$. On the open subset $\uc_M$ of Section~\ref{S:8.4} and for every integer $i$ and every $s\in(\Tc/Z^0_{\Mc})_{\tors}$, we introduced in Theorem~\ref{thm:8.5.1} the $s$-part of the perverse cohomology of $Rf_*^{\xi,M\text{-par}}\Qlb$, denoted by ${}^{\mathrm p}\hc^i(Rf_*^{\xi,M\text{-par}}\Qlb)_s$. Let $a\in\Ac_G(k)^\tau$. We assume that $a$ is the image, under the closed immersion $\Ac_M\to\Ac_G$, of an element $a_M\in\Ac_M^{\el}(k)^\tau$. We defined in Section~\ref{S:11.3} a group scheme $J_{a_M}$ over $C$. We shall define later a positive real number
$$
c(J_{a_M})
$$
associated with the generic fiber $J_{a_M}$ (cf.\ equality~(\ref{eq:12.2.4}) in Proposition~\ref{thm:12.2.2} below).

Let ${}^{\mathrm p}\hc^i(Rf_*^{\xi,M\text{-par}}\Qlb)_{a,s}$ be the fiber at $a$ of the $s$-part. For every graded endomorphism $\sigma$ (of degree $0$) of ${}^{\mathrm p}\hc^\bullet(Rf_*^{\xi,M\text{-par}}\Qlb)_{a,s}$, let
$$
\trace(\sigma,{}^{\mathrm p}\hc^\bullet(Rf_*^{\xi,M\text{-par}}\Qlb)_{a,s}) =\sum_{i\in\ZZ}(-1)^i \trace(\sigma,{}^{\mathrm p}\hc^i(Rf_*^{\xi,M\text{-par}}\Qlb)_{a,s}).
$$

The goal of this section is to prove the following theorem.

\begin{theorem}\label{thm:12.1.1}
Let $a_M\in\Ac_M^{\el}(k)^\tau$ and $a=\chi_G^M(a_M)\in\Ac_G$. Let $s\in\Tc^{W_a}$ and let $\bar{s}$ be its image in $\Tc/Z^0_{\Mc}$. Then $\bar{s}$ is a torsion element, the group $\Tc^{W_a}/Z^0_{\Mc}$ is finite and we have the following equality:
\begin{align*}
\trace(&\tau^{-1},{}^{\mathrm p}\hc^\bullet(Rf_*^{\xi,M\text{-par}}\Qlb)_{a,\bar{s}})\\
&=\frac{c(J_{a_M})}{\vol(\ago_M/X_*(M))\cdot|\Tc^{W_a}/Z^0_{\Mc}|}
\cdot q^{\deg(D)|\Phi^G|/2}J_M^{G,s}(a_M).
\end{align*}
\end{theorem}

\begin{remark}\label{thm:12.1.2}
Since $a_M$ is elliptic in $\Ac_M$, we have the inclusion $W_a\subset W^M$ (cf.\ Proposition~\ref{thm:8.2.1}) and the group $(\Tc/Z^0_{\Mc})^{W_a}$ is finite since it is the dual of the finite group $X_*(T\cap M_{\der})_{W_a}$ (cf.\ Lemma~\ref{thm:8.3.3}). This explains why $\bar{s}$ is torsion and the group $\Tc^{W_a}/Z^0_{\Mc}$ is finite.
\end{remark}

The proof of Theorem~\ref{thm:12.1.1} will be given in Section~\ref{S:12.3}.

\begin{corollary}\label{thm:12.1.3}
Under the hypotheses of Theorem~\ref{thm:12.1.1}, the trace
$$
\trace(\tau^{-1},{}^{\mathrm p}\hc^\bullet(Rf_*^{\xi,M\text{-par}}\Qlb)_{a,\bar{s}})
$$
depends only on the image of $s$ in $\Tc/Z_{\Mc}$.
\end{corollary}

\begin{preuve}
This is a consequence of Theorem~\ref{thm:12.1.1} above and of Proposition~\ref{thm:11.7.3}.
\end{preuve}

In the following paragraph, we gather some results useful for the proof of Theorem~\ref{thm:12.1.1}.
\end{paragr}

\begin{paragr}\label{S:12.2}
We continue with the notation of the preceding paragraph~\ref{S:12.1}. We recall that $G$ is semi-simple. Let $M\in\lc^G(T)$ and let $X\in\mgo(F_0)$ be a semi-simple and $G$-regular element. Let $J$ be its centralizer in $G$: it is an $F$-subtorus of $M$ stable under $\tau$. We assume that $J$ satisfies the following two properties:
\begin{enumerate}
\item the $F$-torus $J/Z_M$ is anisotropic; \item the torus $J\times_F F_\infty$ is split over $F_\infty$.
\end{enumerate}

Let
$$
K=G(\oc_0).
$$

Let $j\in J(\AAA)$ and
$$
\wc=\{(gK,\delta)\in G(\AAA)/K\times J(F)\mid \delta j\tau(gK)=gK\}.
$$
The group $J(F)$ acts on $\wc$ by
$$
\gamma\cdot(gK,\delta)=(\gamma gK,\gamma\delta\tau(\gamma)^{-1}).
$$
Let $\xi\in\ago_T$ (for the moment, it is not necessary to assume $\xi$ in general position) and let $\wc^\xi$ be the subset of $\wc$ consisting of the pairs $(gK,\delta)$ that satisfy the following two conditions:
\begin{equation}\label{eq:12.2.1}
\Ad(g)^{-1}(X)\in\varpi^{-D}\ggo(\oc)
\end{equation}
and
\begin{equation}\label{eq:12.2.2}
\xi_M\in\operatorname{cvx}(-H_P(g))_{P\in\pc(M)},
\end{equation}
where $\xi_M$ is the orthogonal projection of $\xi$ onto $\ago_M$ and cvx denotes the convex hull.

Since $J(F)$ centralizes $X$ and the function $H_P$ is left invariant under $M(F)$ and \emph{a fortiori} under $J(F)$, the group $J(F)$ leaves the subset $\wc^\xi$ globally stable. Let $[J(F)\back\wc^\xi]$ be the quotient groupoid. The goal of this section is to give a formula for the cardinality of the groupoid $[J(F)\back\wc^\xi]$ defined by
\begin{equation}\label{eq:12.2.3}
|[J(F)\back\wc^\xi]|=\sum_{x\in J(F)\back\wc^\xi}|\Aut_{J(F)}(x)|^{-1},
\end{equation}
where one sums, over a system of representatives of the orbits of $J(F)$ in $\wc^\xi$, the inverses of the orders of the stabilizers (with the convention that the inverse of infinity is $0$).

For every $g\in G(\AAA)$, let $\mathbf{1}_{M,g}$ be the characteristic function on $\ago_M$ of the convex hull of the points $-H_P(g)$ for $P\in\pc(M)$. Let
$$
w_M^\xi(g)=|\{\mu\in X_*(M)\mid \mathbf{1}_{M,g}(\xi_M+\mu)=1\}|\,;
$$
it is an integer.

\begin{lemma}\label{thm:12.2.1}
The function $w_M^\xi$ is left invariant under $M(\AAA)$.
\end{lemma}

\begin{preuve}
We refer the reader to \cite[Lemma~11.7.1]{LFPI}.
\end{preuve}

Let
$$
(J(\AAA)\back G(\AAA))^{\tau,j}
$$
be the open and closed part of $(J(\AAA)\back G(\AAA))^\tau$ consisting of the $g$ such that the class of $g\tau(g)^{-1}$ in $(J(F)\back J(\AAA))_\tau$ is equal to that of $j$. We equip $J(\AAA_0)$ with the Haar measure that gives volume $1$ to the maximal compact subgroup, and $J(F_0)$ with the counting measure. Recall that $G(\AAA_0)$ is equipped with the Haar measure that gives volume $1$ to $K$. As in Section~\ref{S:11.5}, we deduce from this a measure on $(J(\AAA)\back G(\AAA))^\tau$, and hence on $(J(\AAA)\back G(\AAA))^{\tau,j}$.

\begin{proposition}\label{thm:12.2.2}
Let $J(\AAA_0)^1=J(\AAA_0)\cap\Ker(H_M)$ where $H_M$ is the morphism defined in Remark~\ref{thm:11.6.1} and
\begin{equation}\label{eq:12.2.4}
c(J)=\vol(J(F_0)\back J(\AAA_0)^1)\cdot|\Ker(J(F)_\tau\to J(\AAA)_\tau)|.
\end{equation}
Then $c(J)$ is finite and we have the equality
$$
|[J(F)\back\wc^\xi]|=c(J)\cdot\int_{(J(\AAA)\back G(\AAA))^{\tau,j}} \mathbf{1}_D(\Ad(g^{-1})X)\,w_M^\xi(g)\,dg,
$$
where the integral is finite.
\end{proposition}

Before giving the proof let us state two corollaries.

\begin{corollary}\label{thm:12.2.3}
Recall that the set $\wc$ depends on the choice of $j\in J(\AAA)$. The cardinality of the groupoid $[J(F)\back\wc^\xi]$ depends only on the class of $j$ in $(J(F)\back J(\AAA))_\tau$.
\end{corollary}

\begin{preuve}
The set $(J(\AAA)\back G(\AAA))^{\tau,j}$ depends only on the class of $j$ in $(J(F)\back J(\AAA))_\tau$. The result is then an immediate consequence of Proposition~\ref{thm:12.2.2}.
\end{preuve}

\begin{corollary}\label{thm:12.2.4}
Assume, moreover, that $\xi$ is in general position (cf.\ Remark~\ref{thm:4.9.2}). We then have
$$
|[J(F)\back\wc^\xi]|=\frac{c(J)}{\vol(\ago_M/X_*(M))}\cdot \int_{(J(\AAA)\back G(\AAA))^{\tau,j}}\mathbf{1}_D(\Ad(g^{-1})X)\,v_M^G(g)\,dg.
$$
\end{corollary}

\begin{preuve}
For $\xi$ in general position, we have the equality between
$$
\int_{(J(\AAA)\back G(\AAA))^{\tau,j}}\mathbf{1}_D(\Ad(g^{-1})X)\,w_M^\xi(g)\,dg
$$
and
$$
\vol(\ago_M/X_*(M))^{-1}\cdot\int_{(J(\AAA)\back G(\AAA))^{\tau,j}} \mathbf{1}_D(\Ad(g^{-1})X)\,v_M^G(g)\,dg.
$$
The latter rests on the formulas relating the weights $w_M^\xi(g)$ and $v_M^G(g)$, and is proved exactly as Theorem~11.14.2 and its Corollary~11.14.3 of \cite{LFPI}. The result then follows from Proposition~\ref{thm:12.2.2}.
\end{preuve}

\begin{preuve}
Since $J$ is anisotropic modulo $Z_M$ (cf.\ assertion 1), the quotient $J(F_0)\back J(\AAA_0)^1$ is compact, hence of finite volume. The group
$$
\Ker^1(\Ga_{F_0},J(F^{\sep}))
$$
of Lemma~\ref{thm:12.2.5} is finite for every $F_0$-torus $J$. Lemma~\ref{thm:12.2.5} implies the finiteness of $\Ker(J(F)_\tau\to J(\AAA)_\tau)$. Thus the constant $c(J)$ is finite.

Let
$$
\psi:G(\AAA)\to G(\AAA)
$$
be the map defined by $\psi(g)=g\tau(g)^{-1}j^{-1}$. The map $(gK,\delta)\mapsto J(F)gK$ induces a bijection from $J(F)\back\wc$ onto $J(F)\back\psi^{-1}(J(F))/K$. Moreover, the centralizer of $(gK,\delta)$ in $J(F)$ is $J(F_0)\cap gKg^{-1}$. Equality~(\ref{eq:12.2.3}) can therefore be rewritten as
$$
|[J(F)\back\wc^\xi]|=\sum_{g\in J(F)\back\psi^{-1}(J(F))/K} |J(F_0)\cap gKg^{-1}|^{-1}\,\mathbf{1}_D(\Ad(g^{-1})X)\, \mathbf{1}_{M,g}(\xi_M).
$$
The group $G(\AAA_0)$ acts by right translation on $G(\AAA)$ without fixed points. The set $\psi^{-1}(F)$ is therefore a disjoint and countable union of orbits under $G(\AAA_0)$, all isomorphic to $G(\AAA_0)$. By transport, we deduce from this a measure on $\psi^{-1}(F)$ that is right invariant under $G(\AAA_0)$. We equip $J(F)$ with the counting measure and $J(F)\back\psi^{-1}(J(F))$ with the quotient measure. We then have
\begin{equation}\label{eq:12.2.5}
|[J(F)\back\wc^\xi]|=\int_{J(F)\back\psi^{-1}(J(F))}
\mathbf{1}_D(\Ad(g^{-1})X)\,\mathbf{1}_{M,g}(\xi_M)\,dg.
\end{equation}
Since $J_\infty$ is split (cf.\ assertion 2 above), the morphism $H_M$ restricts to a surjective morphism from $J(\AAA_0)$ onto $X_*(M)$ with kernel $J(\AAA_0)^1$. For every $g\in G(\AAA)$, we have the equality (cf.\ \cite[proof of Proposition~11.9.1]{LFPI})
$$
\int_{J(F_0)\back J(\AAA_0)}\mathbf{1}_{M,tg}(\xi_M)\,dt =\vol(J(F_0)\back J(\AAA_0)^1)\cdot w_M^\xi(g).
$$
The exact sequence
$$
0\longrightarrow J(F)\longrightarrow J(\AAA)\longrightarrow J(F)\back J(\AAA)\longrightarrow 0
$$
gives the exact sequence in cohomology
$$
0\longrightarrow J(F_0)\back J(\AAA_0)\longrightarrow (J(F)\back J(\AAA))^\tau\longrightarrow \Ker(J(F)_\tau\to J(\AAA)_\tau)\longrightarrow 0.
$$
We then equip $(J(F)\back J(\AAA))^\tau$ with the measure whose quotient by that of $J(F_0)\back J(\AAA_0)$ gives the counting measure on $\Ker(J(F)_\tau\to J(\AAA)_\tau)$. Using the above exact sequence and the right invariance of $w_M^\xi$ under $M(\AAA)$, we obtain
\begin{equation}\label{eq:12.2.6}
\int_{(J(F)\back J(\AAA))^\tau}\mathbf{1}_{M,tg}(\xi_M)\,dt
=c(J)\cdot w_M^\xi(g).
\end{equation}

The group $(J(F)\back J(\AAA))^\tau$ acts without fixed points on $J(F)\back\psi^{-1}(J(F))$. The map that associates to $g\in\psi^{-1}(J(F))$ its class in $J(\AAA)\back G(\AAA)$ descends to a surjective map $J(F)\back\psi^{-1}(J(F))\to(J(\AAA)\back G(\AAA))^{\tau,j}$ whose fibers are the orbits under $(J(F)\back J(\AAA))^\tau$. We therefore have
$$
(J(F)\back J(\AAA))^\tau\back J(F)\back\psi^{-1}(J(F))\simeq (J(\AAA)\back G(\AAA))^{\tau,j}
$$
and this bijection is compatible with our choices of measures. Combining (\ref{eq:12.2.5}) and (\ref{eq:12.2.6}), we obtain the proposition. The integral is finite because the integrand has compact support.
\end{preuve}

\begin{lemma}\label{thm:12.2.5}
Let $F^{\sep}$ be a separable closure of $F$ and $\Ga_{F_0}$ the Galois group of $F^{\sep}$ over $F_0$. Let
$$
\Ker^1(\Ga_{F_0},J(F^{\sep}))=\Ker(H^1(\Ga_{F_0},J(F^{\sep}))\to H^1(\Ga_{F_0},J(\AAA_0\otimes_{F_0}F^{\sep})).
$$
We have a natural identification
$$
\Ker^1(\Ga_{F_0},J(F^{\sep}))=\Ker(J(F)_\tau\to J(\AAA)_\tau).
$$
\end{lemma}

\begin{preuve}
We have the exact sequence
$$
0\longrightarrow\Ga_F\longrightarrow\Ga_{F_0}\longrightarrow \Gal(F/F_0)\longrightarrow 0,
$$
where $\Ga_F$ is the Galois group of $F^{\sep}$ over $F$. The Galois group $\Gal(F/F_0)$ of $F$ over $F_0$ is identified with the Galois group of $k$ over $\Fq$, hence with $\widehat{\ZZ}$. Let $W_{F_0}$ be the inverse image of $\ZZ$ in $\Ga_{F_0}$. We equip $W_{F_0}$ with the topology that makes $\Ga_F$ an open subgroup of $W_{F_0}$. We therefore have an exact sequence
$$
0\longrightarrow\Ga_F\longrightarrow W_{F_0}\longrightarrow\ZZ \longrightarrow 0.
$$
This exact sequence gives the following commutative diagram with exact rows, of inflation-restriction
$$
\xymatrix@C=1.1pc{ 0\ar[r] & H^1(\ZZ,J(F))\ar[r]\ar[d] & H^1(W_{F_0},J(F^{\sep}))\ar[r]\ar[d] &
H^1(\Ga_F,J(F^{\sep}))\ar[d]\\
0\ar[r] & H^1(\ZZ,J(\AAA))\ar[r] & H^1(W_{F_0},J(\AAA\otimes_F F^{\sep}))\ar[r] & H^1(\Ga_F,J(\AAA\otimes_F F^{\sep})). }
$$
It is well known that for such a field $F$ the terms on the right are zero. Those on the left are identified, respectively from top to bottom, with $J(F)_\tau$ and $J(\AAA)_\tau$. It follows that we have the identification
\begin{equation}\label{eq:12.2.7}
\Ker(J(F)_\tau\to J(\AAA)_\tau)=\Ker(H^1(W_{F_0},J(F^{\sep}))\to
H^1(W_{F_0},J(\AAA\otimes_F F^{\sep})).
\end{equation}
Since $J$ is stable under $\tau$, we may view it as a torus over $F_0$. Let $E_0$ be a finite extension of $F_0$ in $F^{\sep}$ that splits $J$. The exact sequence
$$
0\longrightarrow W_{E_0}\longrightarrow W_{F_0}\longrightarrow \Gal(E_0/F_0)\longrightarrow 0
$$
gives, as before, the following commutative diagram with exact rows, of inflation-restriction, where we set $\Ga_{E_0/F_0}=\Gal(E_0/F_0)$:
$$
\xymatrix@C=1.1pc{ H^1(\Ga_{E_0/F_0},J(E_0))\ar@{^{(}->}[r]\ar[d] & H^1(W_{F_0},J(F^{\sep}))\ar[r]\ar[d] &
H^1(W_{E_0},J(F^{\sep}))\ar[d]\\
H^1(\Ga_{E_0/F_0},J(\AAA_0\otimes_{F_0}E_0))\ar@{^{(}->}[r] & H^1(W_{F_0},J(\AAA\otimes_F F^{\sep}))\ar[r] & H^1(W_{E_0},J(\AAA\otimes_F F^{\sep})). }
$$
Let us show that the right-hand vertical arrow is injective. Since $J$ is split over $E_0$, it suffices to prove it for the torus $\mathbb{G}_{m,E_0}$. By~(\ref{eq:12.2.7}), we are reduced to showing that the arrow $F_\tau^\times\to\AAA_\tau^\times$ is injective. But every $x\in F^\times$ with trivial image in $\AAA_\tau^\times$ is a unit at every place, so such an $x$ belongs to $k^\times$ and is of the form $y\tau(y)^{-1}$ for an element $y\in k^\times$.

This injectivity being established, we deduce from it the equality between
$$
\Ker(H^1(\Gal(E_0/F_0),J(E_0))\to H^1(\Gal(E_0/F_0),J(\AAA_0\otimes_{F_0}E_0)))
$$
and
$$
\Ker(H^1(W_{F_0},J(F^{\sep}))\to H^1(W_{F_0},J(\AAA\otimes_F F^{\sep}))
$$
which, combined with~(\ref{eq:12.2.7}), gives the lemma.
\end{preuve}
\end{paragr}

\begin{paragr}[Proof of Theorem~\ref{thm:12.1.1}.]\label{S:12.3}
The notation is that of Section~\ref{S:12.1}. The finite group $\Ga(\Om_M,\pi_0(\Jc^1))$ of global sections over $\Om_M$ of the sheaf $\pi_0(\Jc^1)$ acts on ${}^{\mathrm p}\hc^\bullet(Rf_*^{\xi,M\text{-par}}\Qlb)$ (cf.\ \ref{S:8.5}). Since this finite group receives the group $X_*(T\cap M_{\der})$, we also have an action of the latter on the perverse cohomology. Let $a_M\in\Ac_M^{\el}(k)^\tau$ and let $a$ be its image in $\Ac_G(k)^\tau$. The decomposition~(\ref{eq:8.5.1}) of Theorem~\ref{thm:8.5.1} gives the following equality of traces for every $\la\in X_*(T\cap M_{\der})$:
\begin{align*}
\trace(&(\la\tau)^{-1},{}^{\mathrm p}\hc^\bullet(Rf_*^{\xi,M\text{-par}}\Qlb)_a)\\
&=\sum_{s\in(\Tc/Z^0_{\Mc})_{\tors}}\bg s,\la\bd^{-1}
\trace(\tau^{-1},{}^{\mathrm p}\hc^\bullet(Rf_*^{\xi,M\text{-par}}\Qlb)_{s,a}),
\end{align*}
where the pairing $\bg\cdot,\cdot\bd$ is the canonical pairing between $\Tc/Z^0_{\Mc}$ and $X_*(T\cap M_{\der})$.

Assertion 1 of Theorem~\ref{thm:8.5.1} shows that, in the above sum, the term associated with $s$ is zero unless $s$ belongs to the finite group $(\Tc/Z^0_{\Mc})^{W_a}$ (for a comment on this group, cf.\ Remark~\ref{thm:12.1.2}). We therefore restrict to such $s$. The trace of $(\la\tau)^{-1}$ on the perverse cohomology depends only on the class of $\la$ in the group of co-invariants $X_*(T\cap M_{\der})_{W_a}$ (cf.\ (\ref{eq:8.3.1}) and Proposition~\ref{thm:8.3.1}). By Fourier inversion on the finite commutative group $(\Tc/Z^0_{\Mc})^{W_a}$, whose dual is $X_*(T\cap M_{\der})_{W_a}$, we obtain the following formula for every $s\in(\Tc/Z^0_{\Mc})^{W_a}$:
\begin{equation}\label{eq:12.3.1}
\begin{split}
\trace(&\tau^{-1},{}^{\mathrm p}\hc^\bullet(Rf_*^{\xi,M\text{-par}}\Qlb)_{s,a})\\
&=|(\Tc/Z^0_{\Mc})^{W_a}|^{-1}\sum_{\la\in X_*(T\cap M_{\der})_{W_a}}
\bg s,\la\bd\,
\trace((\la\tau)^{-1},{}^{\mathrm p}\hc^\bullet(Rf_*^{\xi,M\text{-par}}\Qlb)_a).
\end{split}
\end{equation}

Let $\la\in X_*(T\cap M_{\der})$. We have the following equalities:
\begin{equation}\label{eq:12.3.2}
\begin{split}
\trace((\la\tau)^{-1},{}^{\mathrm p}\hc^\bullet(Rf_*^{\xi,M\text{-par}}\Qlb)_a)
&=\trace((\la\tau)^{-1},H^\bullet(\mc_a^\xi,\Qlb))\\
&=|\mc_a^\xi(k)^{\la\tau}|.
\end{split}
\end{equation}
The first equality relates the trace of $(\la\tau)^{-1}$ on the perverse cohomology to the (alternating) trace of $(\la\tau)^{-1}$ on the cohomology of the fiber at $a$ of $\mc^\xi$ (denoted $\mc_a^\xi$). It is proved as in \cite[Lemma~3.10.1]{Laumon-Ngo}. The second equality follows from a stacky version of the Grothendieck-Lefschetz trace formula. It relates this trace to the cardinality, written between two bars, of the groupoid of $k$-points of $\mc_a^\xi$ that are fixed under $\la\tau$. We are now going to describe this groupoid. We resume the notation of Section~\ref{sec:11}. There we defined a group scheme $J_{a_M}$ over $C$, stable under $\tau$ (cf.\ \S\ref{S:11.3}), as well as an element $X_{a_M}\in\mgo(F_0)$ (cf.\ \S\ref{S:11.7}). We have the following result.

\begin{lemma}\label{thm:12.3.1}
The groupoid $\mc_a^\xi(k)$ of $k$-points of the fiber at $a$ of $\mc^\xi$ is equivalent to the quotient groupoid of the set of $g\in G(\AAA)/G(\oc)$ such that
\begin{equation}\label{eq:12.3.3}
\Ad(g)^{-1}(X)\in\varpi^{-D}\ggo(\oc)
\end{equation}
and
\begin{equation}\label{eq:12.3.4}
\xi_M\in\operatorname{cvx}(-H_P(g))_{P\in\pc(M)}
\end{equation}
by the left translation action of the group $J_{a_M}(F)$.
\end{lemma}

\begin{preuve}
This is an easy consequence of Proposition~7.6.3 of \cite{LFPI}.
\end{preuve}

The group $X_*(T\cap M_{\der})$ acts on the fiber $\mc_a^\xi$ via the morphism $X_*(T\cap M_{\der})\to\Jc_a^1$ considered in Section~\ref{S:8.3} (cf.\ also \S\ref{S:7.3}) and the action of $\Jc_a^1$ on $\mc_a^\xi$ (cf.\ \ref{S:6.6}). Let $\varpi_\infty$ be a uniformizer of $\oc_\infty$. Recall that in order to construct this morphism, we identify $J_a$ canonically with $T$ over $\oc_\infty$ (cf.\ Proposition~\ref{thm:7.3.1}) and the element $\varpi_\infty^\la\in T(\oc_\infty)\simeq J_a(\oc_\infty)$ is interpreted as gluing data defining a $J_a$-torsor over $C$. In fact, we also have a canonical identification of $J_{a_M}$ with $T$ over $\oc_\infty$, and we denote by $j_\la\in J_{a_M}(\oc_\infty)$ the element corresponding to $\varpi_\infty^\la$ under this identification.

\begin{lemma}\label{thm:12.3.2}
The groupoid $\mc_a^\xi(k)^{\la\tau}$ is equivalent to the quotient groupoid of the set of pairs $(g,\delta)\in G(\AAA)/G(\oc)\times J_{a_M}(F)$ that satisfy conditions~(\ref{eq:12.3.3}) and~(\ref{eq:12.3.4}) of Lemma~\ref{thm:12.3.1} for $g$ and the condition
$$
j_\la\delta\tau(g)=g
$$
by the action of $J_{a_M}(F)$ given by left translation on $g$ and $\tau$-conjugation on $\delta$.
\end{lemma}

\begin{preuve}
It follows from the description of Lemma~\ref{thm:12.3.1}, from the definition of the fixed points of a groupoid under an automorphism (cf.\ for example \cite[\S11.2]{LFPI}) and from the adelic description of the action of $\Jc_a^1$ (left to the reader).
\end{preuve}

Here is a variant of Lemma~\ref{thm:12.3.2} that will allow us to use the results of Section~\ref{S:12.2}.

\begin{lemma}\label{thm:12.3.3}
The groupoid $\mc_a^\xi(k)^{\la\tau}$ is equivalent to the quotient groupoid $[J_{a_M}(F)\back\wc^\xi]$ defined in Section~\ref{S:12.2} relative to the data $X_{a_M}$ and $j_\la$.
\end{lemma}

\begin{preuve}
The only difficulty is to see that the obvious morphism from $[J_{a_M}(F)\back\wc^\xi]$ to the quotient groupoid considered in Lemma~\ref{thm:12.3.2} is essentially surjective. For this one uses the variant of Lang's theorem showing that $1-\tau$ is a surjective map from $G(\oc)$ to $G(\oc)$.
\end{preuve}

We deduce the following lemma.

\begin{lemma}\label{thm:12.3.4}
For every $\la\in X_*(T\cap M_{\der})$, we have
\begin{align*}
\trace(&(\la\tau)^{-1},{}^{\mathrm p}\hc^\bullet(Rf_*^{\xi,M\text{-par}}\Qlb)_a)\\
&=\frac{c(J_{a_M})}{\vol(\ago_M/X_*(M))}\cdot
\int_{(J_{a_M}(\AAA)\back G(\AAA))^{\tau,j_\la}}
\mathbf{1}_D(\Ad(g^{-1})X)\,v_M^G(g)\,dg.
\end{align*}
\end{lemma}

\begin{preuve}
The desired equality is a simple combination of equality~(\ref{eq:12.3.2}) and of Corollary~\ref{thm:12.2.4} (recall that we assume that $\xi$ is in general position).
\end{preuve}

\begin{lemma}\label{thm:12.3.5}
For every $s\in\Tc^{W_a}/Z^0_{\Mc}$, we have the equality between
$$
\trace(\tau^{-1},{}^{\mathrm p}\hc^\bullet(Rf_*^{\xi,M\text{-par}}\Qlb)_{s,a})
$$
and
$$
\frac{c(J_{a_M})\cdot|\Ker(X_*(T\cap M_{\der})_{W_a}\to X_*(T)_{W_a})|} {\vol(\ago_M/X_*(M))\cdot|(\Tc/Z^0_{\Mc})^{W_a}|} \cdot q^{\deg(D)|\Phi^G|/2}\cdot J_M^{G,s}(a_M).
$$
\end{lemma}

\begin{preuve}
Let $s\in\Tc^{W_a}/Z^0_{\Mc}$. Equality~(\ref{eq:12.3.1}) combined with Lemma~\ref{thm:12.3.4} gives the following formula for $\trace(\tau^{-1},{}^{\mathrm p}\hc^\bullet(Rf_*^{\xi,M\text{-par}}\Qlb)_{s,a})$:
\begin{align*}
&\frac{c(J_{a_M})}{\vol(\ago_M/X_*(M))\cdot|(\Tc/Z^0_{\Mc})^{W_a}|}\\
&\qquad\times\sum_{\la\in X_*(T\cap M_{\der})_{W_a}}\bg s,\la\bd
\int_{(J_{a_M}(\AAA)\back G(\AAA))^{\tau,j_\la}}
\mathbf{1}_D(\Ad(g^{-1})X)\,v_M^G(g)\,dg.
\end{align*}
The quotient
$$
(J_{a_M}(\AAA)\back G(\AAA))^{\tau,j_\la}
$$
is none other than the set of $g\in(J_{a_M}(\AAA)\back G(\AAA))^\tau$ such that $\inv_{a_M}(g)$ is equal to the image of $\la$ in $X_*(T)_{W_a}$. By Proposition~\ref{thm:11.4.2}, the image of the invariant $\inv_{a_M}$ is contained in the image of the morphism $X_*(T\cap M_{\der})\to X_*(T)_{W_a}$. We deduce from this that the sum
$$
\sum_{\la\in X_*(T\cap M_{\der})_{W_a}}\bg s,\la\bd \int_{(J_{a_M}(\AAA)\back G(\AAA))^{\tau,j_\la}} \mathbf{1}_D(\Ad(g^{-1})X)\,v_M^G(g)\,dg
$$
is equal to
$$
|\Ker(X_*(T\cap M_{\der})_{W_a}\to X_*(T)_{W_a})|\cdot q^{\deg(D)|\Phi^G|/2}\cdot J_M^{G,s}(a_M)
$$
which gives the result.
\end{preuve}

Lemma~\ref{thm:12.3.5} gives the theorem by virtue of the following Lemma~\ref{thm:12.3.6}.

\begin{lemma}\label{thm:12.3.6}
We have the equality
$$
\frac{|\Ker(X_*(T\cap M_{\der})_{W_a}\to X_*(T)_{W_a})|} {|(\Tc/Z^0_{\Mc})^{W_a}|}=|\Tc^{W_a}/Z^0_{\Mc}|^{-1}.
$$
\end{lemma}

\begin{preuve}
The exact sequence
$$
0\longrightarrow X_*(T\cap M_{\der})\longrightarrow X_*(T)\longrightarrow X_*(M)\longrightarrow 0
$$
gives the exact sequence of co-invariants (recall that $W_a\subset W^M$ acts trivially on $X_*(M)$)
$$
0\longrightarrow I\longrightarrow X_*(T\cap M_{\der})_{W_a}\longrightarrow X_*(T)_{W_a}\longrightarrow X_*(M)\longrightarrow 0
$$
with $I=\Ker(X_*(T\cap M_{\der})_{W_a}\to X_*(T)_{W_a})$. Applying the exact functor $\Hom_{\ZZ}(\cdot,\CC^\times)$ to this sequence, we obtain the exact sequence of groups
$$
1\longrightarrow Z^0_{\Mc}\longrightarrow\Tc^{W_a}\longrightarrow (\Tc/Z^0_{\Mc})^{W_a}\longrightarrow\Hom_{\ZZ}(I,\CC^\times)\longrightarrow 0.
$$
We deduce from this that the order of $\Hom_{\ZZ}(I,\CC^\times)$ is equal to the ratio between the orders of the finite groups $(\Tc/Z^0_{\Mc})^{W_a}$ and $\Tc^{W_a}/Z^0_{\Mc}$, which gives the desired equality since this order is also that of $I$.
\end{preuve}
\end{paragr}
\section{Stable avatars of weighted orbital integrals}\label{sec:13}

\begin{paragr}\label{S:13.1}
In this section we take up, in our context, the formalism introduced by Arthur in order to stabilize weighted orbital integrals. As throughout the article, $G$ is a connected reductive group defined over $k$ and $T$ a maximal subtorus. We have a natural bijection $M\mapsto\Mc$ (cf.\ \ref{S:2.5}) between $\lc^G(T)$ and the Levi subgroups of $\Gc$ that contain $\Tc$. For every $s\in\Tc$ and every subgroup $\Hc$ of $\Gc$, we denote by $\Hc_s$ the neutral component of the centralizer of $s$ in $\Hc$, and by $H_s$ its dual over $k$.
\end{paragr}

\begin{paragr}\label{S:13.2}
Let $s_M\in\Tc$ and $\Mc'=\Mc_{s_M}$. We have $\Mc'=\Mc$ if and only if $s_M$ is central in $\Mc$. Let $\ec_M^G(s_M)$ be the (finite) set of subgroups $\Hc$ of $\Gc$ for which there exist an integer $n\geq 1$ and a finite family $\underline{s}=(s_1,\ldots,s_n)$ of elements of $\Tc$ satisfying
\begin{itemize}
\item[--] $s_1\in s_M Z_{\Mc}$ and $s_i\in Z_{\Mc'}$ for $i\geq 2$; \item[--] $\Hc$ is the neutral component of the group
$$
\Gc_{s_1}\cap\ldots\cap\Gc_{s_n}.
$$
\end{itemize}

\begin{remarks}\label{thm:13.2.1}
All the groups of $\ec_M^G(s_M)$ are reductive, connected and of the same rank, namely the dimension of $\Tc$.

For every $\Hc\in\ec_M^G(s_M)$, we have $\Mc\cap\Hc=\Mc_{s_M}$. In particular, $\Gc\in\ec_M^G(s_M)$ if and only if $s_M\in Z_{\Mc}$.
\end{remarks}

In the same way one defines a set $\ec_M^L(s_M)$ for every $L\in\lc^G(M)$. We have an obvious inclusion $\ec_M^L(s_M)\subset\ec_M^G(s_M)$. More generally, for every $s\in s_M Z_{\Mc}$, one defines $\ec_{M'}^{G_s}(s_M)$. Moreover we have $\ec_{M'}^{G_s}(s_M)=\ec_{M'}^{G_s}(1)$, where $1$ is the neutral element of $\Tc$.

\begin{definition}\label{thm:13.2.2}
Let $\Lc_1\subset\Lc_2$ be two connected reductive subgroups of $\Gc$ that contain $\Tc$. We say that $\Lc_1$ is \emph{elliptic} in $\Lc_2$ if one of the two following equivalent conditions is satisfied:
\begin{enumerate}
\item the connected centers of $\Lc_1$ and $\Lc_2$ are equal; \item the quotient $Z_{\Lc_1}/Z_{\Lc_2}$ is finite.
\end{enumerate}
\end{definition}

For every $L\in\lc^G(M)$, let $\ec_{M,\el}^L(s_M)\subset\ec_M^L(s_M)$ be the subset of the subgroups $\Hc\in\ec_M^L(s_M)$ that are elliptic in $\Lc$.

\begin{lemma}\label{thm:13.2.3}
We have the following equality, where the union on the right-hand side is disjoint:
$$
\ec_M^G(s_M)=\bigcup_{L\in\lc(M)}\ec_{M,\el}^L(s_M).
$$
\end{lemma}

\begin{preuve}
Let $\Hc\in\ec_M^G(s_M)$. The centralizer of the connected center of $\Hc$ is a Levi subgroup $\Lc$ of $\Gc$ that contains $\Hc$. Consequently, we have $Z_{\Hc}^0\subset Z_{\Lc}\subset Z_{\Hc}$. We therefore have $Z_{\Lc}^0=Z_{\Hc}^0$ and $\Hc$ is elliptic in $\Lc$. Any other Levi subgroup $\Lc'$ which contains $\Hc$ as an elliptic subgroup satisfies $Z_{\Lc'}^0=Z_{\Hc}^0$ and therefore $\Lc'=\Lc$.
\end{preuve}

\begin{proposition}\label{thm:13.2.4}
Let $s_M\in\Tc$ and $\Mc'=\Mc_{s_M}$. For every map
$$
J\ :\ \ec_{M,\el}^G(s_M)\to\CC,
$$
there exists a unique map
$$
S\ :\ \ec_{M,\el}^G(s_M)\to\CC
$$
such that the following relation
\begin{equation}\label{eq:13.2.1}
S(\Hc)=J(\Hc)-\sum_{s\in Z_{\Mc'}/Z_{\Hc},\ s\neq 1}|Z_{\Hc_s}/Z_{\Hc}|^{-1}S(\Hc_s)
\end{equation}
is satisfied for every $\Hc\in\ec_{M,\el}^G(s_M)$.
\end{proposition}

\begin{remark}\label{thm:13.2.5}
Let $\Hc\in\ec_{\el}^G(s_M)$. For every $s\in Z_{\Mc'}$, the group $\Hc_s$ belongs to $\ec_M^G(s_M)$ and depends only on the class of $s$ in $Z_{\Mc'}/Z_{\Hc}$. If $\Hc_s$ is not elliptic in $\Hc$, the quotient $Z_{\Hc_s}/Z_{\Hc}$ is not finite and, by convention, we set $|Z_{\Hc_s}/Z_{\Hc}|^{-1}=0$. Thus the sum over $s$ only involves those $S(\Hc_s)$ for which $\Hc_s$ belongs to $\ec_{M,\el}^G(s_M)$. As there are only finitely many such $\Hc_s$, the support of the sum over $s$ is finite, since it is contained in a finite union of finite quotients $Z_{\Hc_s}/Z_{\Hc}$.
\end{remark}

\begin{preuve}
Let $\Hc\in\ec_{\el}^G(s_M)$. Let us observe that if $s\in Z_{\Mc'}$ and $s\notin Z_{\Hc}$ then $\dim(\Hc_s)<\dim(\Hc)$. In particular, if $\Hc$ is of minimal dimension, the relation~(\ref{eq:13.2.1}) reduces to
$$
S(\Hc)=J(\Hc).
$$
An induction on the dimension of the groups $\Hc$ gives the existence and the uniqueness of $S$.
\end{preuve}

\begin{corollary}\label{thm:13.2.6}
Let $s_M\in\Tc$. For every map
$$
J\ :\ \ec_M^G(s_M)\to\CC,
$$
there exists a unique map
$$
S\ :\ \ec_M^G(s_M)\to\CC
$$
such that for every $L\in\lc(M)$, the restriction of $S$ to $\ec_{M,\el}^L(s_M)$ is the map deduced from the restriction of $J$ to $\ec_M^L(s_M)$ by Lemma~\ref{thm:13.2.4}.
\end{corollary}

\begin{preuve}
It follows from Lemma~\ref{thm:13.2.3} and Proposition~\ref{thm:13.2.4}.
\end{preuve}
\end{paragr}

\begin{paragr}[Stable avatar of $J_M^{G,1}(a_M)$.]\label{S:13.3}
We take up the notation of Section~\ref{sec:11}. We introduce the following definition. We denote by $1$ the neutral element of $\Tc$.

\begin{definition}\label{thm:13.3.1}
Let $M\in\lc^G(T)$ and $a_M\in\Ac_M(k)^{\tau}$. The map
$$
H\in\ec_M^G(1)\mapsto S_M^H(a_M)
$$
is the map deduced from the map of Definition~\ref{thm:11.7.1}
$$
H\in\ec_M^G(1)\mapsto J_M^{H,1}(a_M)
$$
by Corollary~\ref{thm:13.2.6}.
\end{definition}
\end{paragr}

\section{A second computation of a Frobenius trace}\label{sec:14}

\begin{paragr}\label{S:14.1}
We take up the hypotheses and the notation of Sections~\ref{sec:10}, \ref{sec:11}, \ref{sec:12}. The main result of this section is the following theorem, which is a computation of a Frobenius trace on an intermediate extension in terms of the stable variant of the weighted orbital integrals $J_M^{G,1}(a_M)$ of Definition~\ref{thm:11.7.1}.

\begin{theorem}\label{thm:14.1.1}
Let $G$ be a semisimple group, $M\in\lc^G$ and $a_M\in\Ac_M^{\el}(k)^{\tau}$. Let $a=\chi_G^M(a_M)$. Assume that $a\in\Ac_G^{\mathrm{bon}}$. Let $j$ be the open immersion
$$
j:\Ac_G^{\el}\to\Ac_G.
$$
Then we have the equality between
$$
\trace(\tau^{-1},(j_{!*}({}^{\mathrm{p}}\hc^{\bullet}(Rf_{G,*}^{\el}\Qlb)_1))_a)
$$
and
$$
\frac{c(J_{a_M})\cdot q^{\deg(D)|\Phi^G|/2}} {\vol(\ago_M/X_*(M))\cdot|\Tc^{W_a}/Z_{\Mc}^0|\cdot|Z_{\Mc}^0\cap Z_{\Gc}|} \cdot S_M^G(a_M),
$$
where
\begin{enumerate}
\item ${}^{\mathrm{p}}\hc^{\bullet}(Rf_{G,*}^{\el}\Qlb)_1$ is the $1$-part of ${}^{\mathrm{p}}\hc^{\bullet}(Rf_{G,*}^{\el}\Qlb)$ (cf.\ Theorem~\ref{thm:8.5.1}); \item its intermediate extension to $\Ac_G$ is denoted $j_{!*}$; \item the notation on the right-hand side, with the exception of $S_M^G(a_M)$, is that of Theorem~\ref{thm:12.1.1}; \item $S_M^G(a_M)$ is the variant of the integral $J_M^{G,1}(a_M)$ introduced in Definition~\ref{thm:13.3.1}.
\end{enumerate}
\end{theorem}

This theorem is nontrivial. Although it rests in part on counting arguments (that is to say, on Theorem~\ref{thm:12.1.1}), its proof also requires the cohomological Theorem~\ref{thm:10.7.3}. This last point explains the condition $a\in\Ac_G^{\mathrm{bon}}$. The proof is to be found in the following paragraph.
\end{paragr}

\begin{paragr}[Proof of Theorem~\ref{thm:14.1.1}.]\label{S:14.2}
We take up the hypotheses of Theorem~\ref{thm:14.1.1}. For every $H\in\ec_{M,\el}^G(1)$, we define $\tilde{S}_M^H(a_M)$ by the following equality:
$$
\trace(\tau^{-1},(j_{!*}({}^{\mathrm{p}}\hc^{\bullet}(Rf_{H,*}^{\el}\Qlb)_1))_a)=
$$
$$
\frac{c(J_{a_M})\cdot q^{\deg(D)|\Phi^G|/2}} {\vol(\ago_M/X_*(M))\cdot|\Tc^{W_a}/Z_{\Mc}^0|\cdot|Z_{\Mc}^0\cap Z_{\Gc}|} \cdot\tilde{S}_M^H(a_M).
$$

It is a matter of seeing that we have
\begin{equation}\label{eq:14.2.1}
\tilde{S}_M^H(a_M)=S_M^H(a_M).
\end{equation}

Since $\Ac_G^{\mathrm{bon}}\cap\Ac_H\subset\Ac_H^{\mathrm{bon}}$ (cf.\ Lemma~\ref{thm:14.3.1}), we may assume, arguing by induction on the dimension of the groups, that the equality~(\ref{eq:14.2.1}) holds for every $H\in\ec_{M,\el}^G(1)$ such that $H\neq G$. Let us show that the point $a=\chi_G^M(a_M)$ belongs to the open subset (notation of~(\ref{eq:10.7.1}) for $\bar{s}=1$)
$$
\uc_1=\uc_M\cap\Ac_G^{\mathrm{bon}}\cap\Bigl(\Ac_G-\bigcup_{(s,\rho)}\Ac_{s,\rho}\Bigr),
$$
where the union is taken over the finite set (cf.\ Proposition~\ref{thm:10.6.10}) of geometric endoscopic data $(s,\rho)$ of $G$ that satisfy
\begin{itemize}
\item[--] $s\in Z_{\Mc}^0$; \item[--] $\Ac_G^{\el}\cap\Ac_{s,\rho}\neq\emptyset$; \item[--] $\rho$ is a nontrivial homomorphism.
\end{itemize}

Since $\Ac_M^{\el}\subset\uc_M$ (cf.\ the definition of $\uc_M$ given in~\S\ref{S:8.4}), it remains to see that we have $a\notin\Ac_{s,\rho}$ for every pair $(s,\rho)$ that satisfies the three conditions above. Let $(s,\rho)$ be a geometric endoscopic datum of $G$ such that $s\in Z_{\Mc}^0$ and $a\in\Ac_{s,\rho}$. Since $a=\chi_G^M(a_M)$, we have $W_a\subset W^M$ (cf.\ Proposition~\ref{thm:8.2.1}). But the condition $s\in Z_{\Mc}^0$ implies $W^M\subset W_s^0$. Hence $\rho$ is necessarily trivial.

We may therefore take the trace of $\tau^{-1}$ on the fiber at $a$ of the equality of Theorem~\ref{thm:10.7.3}. After simplifying by the factor $\frac{c(J_{a_M})\cdot q^{\deg(D)|\Phi^G|/2}}{\vol(\ago_M/X_*(M))\cdot|\Tc^{W_a}/Z_{\Mc}^0|}$, we find the equality
$$
J_M^{G,1}(a_M)=\sum_s|Z_{\Mc}^0\cap Z_{\Gc_s}|^{-1}\cdot\tilde{S}_M^{G_s}(a_M),
$$
where the sum is taken over the $s\in Z_{\Mc}^0$ such that $\Gc_s$ is semisimple. As the terms in the sum on the right are invariant under translation by $Z_{\Mc}^0\cap Z_{\Gc}$, we also have
$$
J_M^{G,1}(a_M)=\sum_s|Z_{\Mc}^0\cap Z_{\Gc}|\cdot|Z_{\Mc}^0\cap Z_{\Gc_s}|^{-1} \cdot\tilde{S}_M^{G_s}(a_M),
$$
where, this time, the sum is taken over the $s\in Z_{\Mc}^0/(Z_{\Mc}^0\cap Z_{\Gc})$ such that $\Gc_s$ is semisimple. Since $\Mc$ is a Levi subgroup of $\Gc$ and of $\Gc_s$, we have $Z_{\Mc}=Z_{\Gc}Z_{\Mc}^0$ and $Z_{\Mc}=Z_{\Gc_s}Z_{\Mc}^0$. In particular, we have $Z_{\Gc_s}\subset Z_{\Gc}(Z_{\Mc}^0\cap Z_{\Gc_s})$. One deduces from this the isomorphisms
$$
Z_{\Mc}^0/(Z_{\Mc}^0\cap Z_{\Gc})\cong Z_{\Mc}/Z_{\Gc}
$$
and
$$
(Z_{\Mc}^0\cap Z_{\Gc_s})/(Z_{\Mc}^0\cap Z_{\Gc})\cong Z_{\Gc_s}/Z_{\Gc},
$$
and one obtains
$$
J_M^{G,1}(a_M)=\sum_{s\in Z_{\Mc}/Z_{\Gc}}|Z_{\Gc_s}/Z_{\Gc}|^{-1} \tilde{S}_M^{G_s}(a_M),
$$
where, by the usual convention, the inverse of the cardinality of an infinite set is zero. By the induction hypothesis, we have $\tilde{S}_M^{G_s}(a_M)=S_M^{G_s}(a_M)$ for $s\neq 1$. It then follows that
\begin{align*}
\tilde{S}_M^G(a_M)&=J_M^{G,1}(a_M)-\sum_{s\in Z_{\Mc}/Z_{\Gc},\,s\neq 1}
|Z_{\Gc_s}/Z_{\Gc}|^{-1}\tilde{S}_M^{G_s}(a_M)\\
&=J_M^{G,1}(a_M)-\sum_{s\in Z_{\Mc}/Z_{\Gc},\,s\neq 1}
|Z_{\Gc_s}/Z_{\Gc}|^{-1}S_M^{G_s}(a_M).
\end{align*}

But this last expression is precisely $S_M^G(a_M)$, which completes the proof.
\end{paragr}

\begin{paragr}\label{S:14.3}
The following lemma was used in the proof of Theorem~\ref{thm:14.1.1}.

\begin{lemma}\label{thm:14.3.1}
Let $H$ be a dual group of a reductive subgroup $\Hc$ of $\Gc$ that contains $\Tc$. Let $\Ac_H\hookrightarrow\Ac_G$ be the canonical closed immersion (induced by the inclusion $W^H\subset W^G$). Then we have
$$
\Ac_G^{\mathrm{bon}}\cap\Ac_H\subset\Ac_H^{\mathrm{bon}}.
$$
\end{lemma}

\begin{preuve}
By Section~\ref{sec:9}, we have at our disposal a function $\delta^H$ on $\Ac_H$ and a function $\delta^G$ on $\Ac_G$. Let $a\in\Ac_G^{\mathrm{bon}}$. We therefore have $\mathrm{codim}_{\Ac_G}(a)\geq\delta_a^G$. Assume moreover that $a\in\Ac_H$. One checks the following formula:
$$
\mathrm{codim}_{\Ac_G}(\Ac_H)=(\dim(G)-\dim(H))\deg(D)/2=\delta_a^G-\delta_a^H.
$$
The first equality comes from the computation of the dimension of $\Ac_G$ by the Riemann-Roch formula, and the second follows from the Ngô-Bezrukavnikov formula, cf.\ Lemma~\ref{thm:9.12.1} (cf.\ also \cite[\S4.17]{Ngo-lemme}). One immediately deduces from this the inequality $\mathrm{codim}_{\Ac_H}(a)\geq\delta_a^H$, which is what we had to prove.
\end{preuve}
\end{paragr}
\section{A global form of the weighted fundamental lemma}\label{sec:15}

\begin{paragr}\label{S:15.1}
The following theorem, which is the main result of this section, is a global form of Arthur's weighted fundamental lemma.

\begin{theorem}\label{thm:15.1.1}
Let $M\in\lc^G$ and $s_M\in\Tc$. Let $M'$ be of dual $\Mc'=\Mc_{s_M}$. Let $a_M\in\Ac_{M'}^{\el}(k)^{\tau}$. By abuse of notation, we still denote by $a_M$ the element $\chi_M^{M'}(a_M)$. We make the following hypotheses:
\begin{itemize}
\item[--] $\Mc'\subsetneq\Mc$ (that is to say $s_M\notin Z_{\Mc}$); \item[--] $\Mc'$ is elliptic in $M$ (which amounts to $a_M\in\Ac_M^{\el}(k)$); \item[--] the element $\chi_G^{M'}(a_M)$ belongs to $\Ac_G^{\mathrm{bon}}$.
\end{itemize}
Then we have the following equality:
$$
J_M^{G,s_M}(a_M)=|Z_{\Mc'}/Z_{\Mc}|\sum_{s\in s_M Z_{\Mc}/Z_{\Gc}} |Z_{\Gc_s}/Z_{\Gc}|^{-1}S_{M'}^{G_s}(a_M).
$$
\end{theorem}
\end{paragr}

\begin{paragr}[Proof of Theorem~\ref{thm:15.1.1}.]\label{S:15.2}
Let us begin with the case of a semisimple group. Let us prove the following lemma.

\begin{lemma}\label{thm:15.2.1}
Let $a=\chi_G^{M'}(a_M)$. Then $a$ belongs to the open subset $\uc_{s_M}$ defined in~(\ref{eq:10.7.1}).
\end{lemma}

\begin{preuve}
By hypothesis, we know that $a$ belongs to $\Ac_G^{\mathrm{bon}}$. The hypothesis that $\Mc'$ is elliptic in $M$ implies that $a\in\Ac_M^{\el}$ (cf.\ Proposition~\ref{thm:10.6.6}) and therefore $a$ belongs to the open subset $\uc_M$ of~\S\ref{S:8.4}. Let $(s,\rho)$ be a geometric endoscopic datum (cf.\ Definition~\ref{thm:10.6.3}) such that the three following conditions are satisfied:
\begin{itemize}
\item[--] $s\in s_M Z_{\Mc}^0$; \item[--] $\Ac_G^{\el}\cap\Ac_{s,\rho}\neq\emptyset$; \item[--] $a\in\Ac_{s,\rho}$.
\end{itemize}
Since $a_M\in\Ac_{M'}$, we have $W_a\subset W^{M'}$ (by a variant of Proposition~\ref{thm:8.2.1}). As we have $W^{M'}\subset W_s^0$, we also have $W_a\subset W_s^0$ and $\rho$ is necessarily trivial. This completes the proof.
\end{preuve}

Let us take $\xi$ in general position. Theorem~\ref{thm:10.7.3} specialized at the point $a\in\uc_{s_M}$ entails an identity of Frobenius traces which is made explicit as follows by Theorems~\ref{thm:12.1.1} and~\ref{thm:14.1.1} (cf.\ also Lemma~\ref{thm:14.3.1}):
\begin{multline*}
\frac{c(J_{a_M})\cdot q^{\deg(D)|\Phi^G|/2}}
{\vol(\ago_M/X_*(M))\cdot|\Tc^{W_a}/Z_{\Mc}^0|}J_M^{G,s_M}(a_M)\\
=\sum_s\frac{c(J_{a_M})\cdot q^{\deg(D)|\Phi^G|/2}}
{\vol(\ago_{M'}/X_*(M'))\cdot|\Tc^{W_a}/Z_{\Mc'}^0|\cdot
|Z_{\Mc'}^0\cap Z_{\Gc_s}|}\cdot S_{M'}^{G_s}(a_M),
\end{multline*}
where the sum is taken over the $s\in s_M Z_{\Mc}^0$ such that $G_s$ is semisimple. The hypothesis that $M'$ is elliptic entails the equalities $Z_{\Mc}^0=Z_{\Mc'}^0$, $\ago_M=\ago_{M'}$ (as Euclidean spaces) and $X_*(M)=X_*(M')$. After simplification, we obtain the equality
$$
J_M^{G,s_M}(a_M)=\sum_s|Z_{\Mc}^0\cap Z_{\Gc_s}|^{-1}S_{M'}^{G_s}(a_M)
$$
for the same summation set, which one may rewrite (as in the proof of Theorem~\ref{thm:14.1.1})
$$
J_M^{G,s_M}(a_M)=\sum_s\frac{[Z_{\Mc}^0\cap Z_{\Gc}|} {|Z_{\Mc}^0\cap Z_{\Gc_s}|}S_{M'}^{G_s}(a_M),
$$
where this time the sum is over the $s\in s_M Z_{\Mc}^0/Z_{\Mc}^0\cap Z_{\Gc}$ such that $G_s$ is elliptic. The exact sequence
$$
1\longrightarrow(Z_{\Mc}\cap Z_{\Gc_s})/Z_{\Gc}\longrightarrow Z_{\Gc_s}/Z_{\Gc}\longrightarrow Z_{\Mc'}/Z_{\Mc}\longrightarrow 1
$$
(where the surjectivity comes from the equalities $Z_{\Mc'}=Z_{\Gc_s}Z_{\Mc'}^0$ and $Z_{\Mc'}^0=Z_{\Mc}^0$) and the isomorphism
$$
(Z_{\Gc_s}\cap Z_{\Mc})/Z_{\Gc}\cong (Z_{\Mc}^0\cap Z_{\Gc_s})/(Z_{\Mc}^0\cap Z_{\Gc})
$$
entail the formula
$$
\frac{[Z_{\Mc}^0\cap Z_{\Gc}|}{|Z_{\Mc}^0\cap Z_{\Gc_s}|} =\frac{|Z_{\Mc'}/Z_{\Mc}|}{|Z_{\Gc_s}/Z_{\Gc}|}.
$$
All the orders in the formula above are finite if $G_s$ is elliptic. In the right-hand side of this formula, the denominator is infinite if $G_s$ is not elliptic, in which case its inverse is by convention $0$. Using the isomorphism $Z_{\Mc}^0/Z_{\Mc}^0\cap Z_{\Gc}\cong Z_{\Mc}/Z_{\Gc}$, we obtain the equality
$$
J_M^{G,s_M}(a_M)=|Z_{\Mc'}/Z_{\Mc}|\sum_{s\in s_M Z_{\Mc}/Z_{\Gc}} |Z_{\Gc_s}/Z_{\Gc}|^{-1}S_{M'}^{G_s}(a_M)
$$
which is what we had to prove.

The passage from a semisimple group to an arbitrary reductive group is left to the reader.
\end{paragr}

\section{Arthur's formalism for the stabilization of weighted orbital integrals}\label{sec:16}

\begin{paragr}\label{S:16.1}
This section takes up the notation and the definitions of \S\S\ref{S:13.1} and~\ref{S:13.2}. It is independent of the other parts of the article. We develop, in an abstract context, Arthur's formalism for the stabilization of weighted orbital integrals and of their descent and splitting formulas (cf.\ \cite{Arthur-transfer} for example). We give several propositions which will then allow us to deduce the weighted fundamental lemma from the global variant given by Theorem~\ref{thm:15.1.1}. All the results rest on Arthur's methods.
\end{paragr}

\begin{paragr}\label{S:16.2}
Let $M\in\lc^G$. Recall that for every $s_M\in\Tc$ we have defined in \S\ref{S:13.2} finite sets $\ec_{M,\el}^G(s_M)\subset\ec_M^G(s_M)$ of reductive subgroups of $\Gc$ which contain $\Tc$. For every Levi subgroup of $\Gc$ containing $\Tc$, let $\lc^{\Gc}(\Mc)$ be the finite set of Levi subgroups of $\Gc$ which contain $\Mc$. This set is in natural bijection with $\lc^G(M)$. We set $\lc^{\Gc}=\lc^{\Gc}(\Tc)$. This definition holds more generally for a reductive subgroup of $\Gc$ and one of its Levi subgroups.

\begin{definition}\label{thm:16.2.1}
Let $s_M\in\Tc$ and $\Mc'=\Mc_{s_M}$. For every map
$$
J\ :\ \ec_{M,\el}^G(s_M)\to\CC.
$$
Let
$$
J^{\ec}(\Gc)=|Z_{\Mc'}/Z_{\Mc}|\sum_{s\in s_M Z_{\Mc}/Z_{\Gc}} |Z_{\Gc_s}/Z_{\Gc}|^{-1}S(\Gc_s),
$$
where $S$ is the map deduced from $J$ by Proposition~\ref{thm:13.2.4}.
\end{definition}

\begin{remark}\label{thm:16.2.2}
If $s_M\in Z_{\Mc}$, we have $\Mc=\Mc'$ and $\Gc\in\ec_{M,\el}^G(s_M)$. In this case, we have $J^{\ec}(\Gc)=J(\Gc)$. Indeed the formula above may be rewritten
$$
J^{\ec}(\Gc)=S(\Gc)+\sum_{s\in s_M Z_{\Mc}/Z_{\Gc},\,s\neq 1} |Z_{\Gc_s}/Z_{\Gc}|^{-1}S(\Gc_s)
$$
and the sum over $s$ above is equal to $J(\Gc)-S(\Gc)$ by the formula~(\ref{eq:13.2.1}).
\end{remark}

\begin{remark}\label{thm:16.2.3}
By Lemma~\ref{thm:13.2.3}, we have $\ec_M^G(s_M)=\bigcup_{L\in\lc(M)}\ec_{M,\el}^L(s_M)$ where the union is disjoint. For every map $J\ :\ \ec_M^G(s_M)\to\CC$, we therefore obtain a map
$$
J^{\ec}:\lc^{\Gc}(\Mc)\to\CC
$$
given by Definition~\ref{thm:16.2.1}.
\end{remark}

\begin{definition}\label{thm:16.2.4}
Let $M\in\lc^G$ and $s_M\in\Tc$. We say that a map
$$
J\ :\ \ec_{M,\el}^G(s_M)\cup\{\Gc\}\to\CC
$$
is \emph{stabilizing} if it satisfies the equality $J(\Gc)=J^{\ec}(\Gc)$. We say that a map
$$
J\ :\ \ec_M^G(s_M)\cup\lc^{\Gc}(\Mc)\to\CC
$$
is \emph{stabilizing} if for every $L\in\lc^G(M)$ the map restricted to $\ec_{M,\el}^L(s_M)\cup\{\Lc\}$ is stabilizing.
\end{definition}

Here is a first, trivial example.

\begin{example}\label{thm:16.2.5}
If $s_M\in Z_{\Mc}$, every map $J\ :\ \ec_{M,\el}^G(s_M)\cup\{\Gc\}\to\CC$ is stabilizing by virtue of Remark~\ref{thm:16.2.2}.
\end{example}

Here is a second one which is not.

\begin{example}\label{thm:16.2.6}
Let $s_M\in\Tc$ and $\Mc'=\Mc_{s_M}$. We assume that $s_M\notin Z_{\Mc}$. For $a_M\in\Ac_{M'}^{\el}(k)^{\tau}$ which satisfies the hypotheses of Theorem~\ref{thm:15.1.1}, the map $J$ defined by $J(\Gc)=J_M^{G,s_M}(a_M)$ and for $\Hc\in\ec_{M,\el}^G(s_M)$ by $J(\Hc)=J_{M'}^{H,1}(a_M)$ is stabilizing: this is nothing but a reformulation of Theorem~\ref{thm:15.1.1}.
\end{example}
\end{paragr}

\begin{paragr}[Splitting and descent.]\label{S:16.3}
In all the rest of this section, we fix $M\in\lc^G$ and $s_M\in\Tc$. Weighted orbital integrals admit descent and splitting formulas brought to light by Arthur (cf.\ \cite{Arthur-invariant-I}). The coefficient $d$ of Definition~\ref{thm:16.3.1} below plays a key role in these formulas. Beforehand, for every $\Hc\in\ec_M^G(s_M)$, we set
$$
\ago_{\Hc}=X^*(\Hc)\otimes_{\ZZ}\RR=X^*(Z_{\Hc}^0)\otimes_{\ZZ}\RR.
$$
We have moreover $\ago_H=\ago_{\Hc}$ if $H$ is the dual of $\Hc$. We then have decompositions $\ago_{\Tc}=\ago_{\Tc}^{\Hc}\oplus\ago_{\Hc}$ as in \S\ref{S:2.4}.

\begin{definition}\label{thm:16.3.1}
Let $\Hc\in\ec_M^G(s_M)$ and $\Rc\in\lc^{\Hc}(\Tc)$. For all Levi subgroups $\Lc_1$ and $\Lc_2$ in $\lc^{\Hc}(\Rc)$, we set
$$
d_{\Rc}^{\Hc}(\Lc_1,\Lc_2)=0
$$
except if we have a direct sum $\ago_{\Rc}^{\Lc_1}\oplus\ago_{\Rc}^{\Lc_2}=\ago_{\Rc}^{\Hc}$, in which case $d$ is the absolute value of the determinant of the canonical isomorphism $\ago_{\Rc}^{\Lc_1}\times\ago_{\Rc}^{\Lc_2}\to\ago_{\Rc}^{\Hc}$ written in orthonormal bases.
\end{definition}

\begin{remark}\label{thm:16.3.2}
We have $d_{\Rc}^{\Hc}(\Lc_1,\Lc_2)=d_{\Rc}^{\Hc}(\Lc_2,\Lc_1)$ and this coefficient is nonzero if and only if $\ago_{\Lc_1}^{\Gc}\oplus\ago_{\Lc_2}^{\Gc}=\ago_{\Rc}^{\Hc}$. We also have $d_{\Rc}^{\Hc}(\Rc,\Hc)=1$.
\end{remark}

We collect in the following lemma some well-known elementary facts whose proof is left to the reader.

\begin{lemma}\label{thm:16.3.3}
\mbox{}
\begin{enumerate}
\item We have $a_{\Lc_1}^{\Hc}\cap a_{\Lc_2}^{\Hc}=0$ if and only if the quotient
$$
(Z_{\Lc_1}\cap Z_{\Lc_2})/Z_{\Hc}
$$
is finite; \item We have $a_{\Rc}=a_{\Lc_1}+a_{\Lc_2}$ if and only if $Z_{\Rc}=Z_{\Lc_1}Z_{\Lc_2}$; \item Let $s\in\Tc$ be such that $\Gc_s$ is elliptic in $\Gc$ and $\Mc_s$ is elliptic in $\Mc$. The map $\Lc\to\Lc_s$ induces a bijection from the set of the $\Lc\in\lc(\Mc)$ such that $\Lc_s$ is elliptic in $\Lc$ onto $\lc^{\Hc_s}(\Mc_s)$.
\end{enumerate}
\end{lemma}

\begin{definition}\label{thm:16.3.4}
Let $\Hc\in\ec_M^G(s_M)$ and $\Rc\in\lc^{\Hc}(\Tc)$. For all Levi subgroups $\Lc_1$ and $\Lc_2$ in $\lc^{\Hc}(\Rc)$, we set
$$
e_{\Rc}^{\Hc}(\Lc_1,\Lc_2)=|(Z_{\Lc_1}\cap Z_{\Lc_2})/Z_{\Hc}|^{-1} d_{\Rc}^{\Hc}(\Lc_1,\Lc_2),
$$
with the usual convention that the inverse of the cardinality of an infinite set is zero.
\end{definition}

\begin{remark}\label{thm:16.3.5}
By assertion 1 of Lemma~\ref{thm:16.3.3}, we have $e_{\Rc}^{\Hc}(\Lc_1,\Lc_2)=0$ if and only if $d_{\Rc}^{\Hc}(\Lc_1,\Lc_2)=0$.
\end{remark}

We can now state the two main results of this section.

\begin{proposition}\label{thm:16.3.6}
(Splitting) Let $s_M\in\Tc$ be such that $s_M\notin Z_{\Mc}$ and $\Mc'=\Mc_{s_M}$ is an elliptic subgroup of $\Mc$.

Let $J_1$ and $J_2$ be two maps from $\ec_M^G(s_M)\cup\lc^{\Gc}(\Mc)$ to $\CC$. Let $J$ be the map from $\ec_M^G(s_M)\cup\lc^{\Gc}(\Mc)$ to $\CC$ defined for every $\Hc\in\ec_M^G(s_M)$ by
$$
J(\Hc)=\sum_{\Lc_1,\Lc_2\in\lc^{\Hc}(\Mc')} d_{\Mc'}^{\Hc}(\Lc_1,\Lc_2)\cdot J_1(\Lc_1)\cdot J_2(\Lc_2)
$$
and for every $\Lc\in\lc^{\Gc}(\Mc)$ by
$$
J(\Lc)=\sum_{\Lc_1,\Lc_2\in\lc^{\Lc}(\Mc)} d_{\Mc}^{\Lc}(\Lc_1,\Lc_2)\cdot J_1(\Lc_1)\cdot J_2(\Lc_2).
$$

We have
\begin{enumerate}
\item If $J_1$ and $J_2$ are stabilizing then $J$ is stabilizing; \item If $J$ and $J_2$ are stabilizing and $J_2(\Mc)\neq 0$ then $J_1$ is stabilizing.
\end{enumerate}
\end{proposition}

\begin{proposition}\label{thm:16.3.7}
(Descent) Let $s_M\in\Tc$ be such that $s_M\notin Z_{\Mc}$ and $\Mc'=\Mc_{s_M}$ is an elliptic subgroup of $\Mc$. Let $\Rc\in\lc^{\Mc}$ be such that $\Rc_{s_M}$ is elliptic in $\Rc$. Let $\Rc'=\Rc_{s_M}$.

Let $J_R$ be a map from $\ec_R^G(s_M)\cup\lc^{\Gc}(\Rc)$ to $\CC$. Let $J_M$ be the map from $\ec_M^G(s_M)\cup\lc^{\Gc}(\Mc)$ to $\CC$ defined for every $\Hc\in\ec_M^G(s_M)$ by
$$
J_M(\Hc)=\sum_{\Lc\in\lc^{\Hc}(\Rc')}d_{\Rc'}^{\Hc}(\Mc',\Lc)J_R(\Lc)
$$
and for every $\Lc\in\lc^{\Gc}(\Mc)$ by
$$
J_M(\Lc)=\sum_{\Lc_1\in\lc^{\Lc}(\Rc)}d_{\Rc}^{\Lc}(\Mc,\Lc_1)J_R(\Lc_1).
$$
If $J_R$ is stabilizing, then $J_M$ is stabilizing.
\end{proposition}

\begin{remark}\label{thm:16.3.8}
Using assertion 3 of Lemma~\ref{thm:16.3.3}, one sees that for every $\Hc\in\ec_M^G(s_M)$ and every $\Lc\in\lc^{\Hc}(\Rc)$, we have $\Lc\in\ec_R^G(s_M)$.
\end{remark}

The rest of the section (in particular \S\S\ref{S:16.4} and~\ref{S:16.6}) is devoted to the proofs of Propositions~\ref{thm:16.3.6} and~\ref{thm:16.3.7}.
\end{paragr}

\begin{paragr}[Proof of Proposition~\ref{thm:16.3.6}.]\label{S:16.4}
It rests on the following proposition, of which the reader will find, for his convenience, a proof in \S\ref{S:16.5}.

\begin{proposition}\label{thm:16.4.1}
(Splitting formulas) We take up the hypotheses of Proposition~\ref{thm:16.3.6}
\begin{enumerate}
\item Let $S$, resp.\ $S_i$, be the map deduced from $J$, resp.\ $J_i$, by Corollary~\ref{thm:13.2.6}. For every $\Hc\in\ec_M^G(s_M)$, we have
$$
S(\Hc)=\sum_{\Lc_1,\Lc_2\in\lc^{\Hc}(\Mc')} e_{\Mc'}^{\Hc}(\Lc_1,\Lc_2)\cdot S_1(\Lc_1)\cdot S_2(\Lc_2)\ ;
$$
\item Let $J^{\ec}$, resp.\ $J_i^{\ec}$, be the map deduced from $J$ (cf.\ Definition~\ref{thm:16.2.1} and Remark~\ref{thm:16.2.3}). For every $\Lc\in\lc^{\Gc}(\Mc)$, we have
$$
J^{\ec}(\Lc)=\sum_{\Lc_1,\Lc_2\in\lc^{\Lc}(\Mc)} d_{\Mc}^{\Lc}(\Lc_1,\Lc_2)\cdot J_1^{\ec}(\Lc_1)\cdot J_2^{\ec}(\Lc_2).
$$
\end{enumerate}
\end{proposition}

We can now give the proof of Proposition~\ref{thm:16.3.6}. Assertion 1 is an obvious consequence of assertion 2 of Proposition~\ref{thm:16.4.1} above. Let us prove assertion 2 of Proposition~\ref{thm:16.3.6}. By hypothesis, $J$ and $J_2$ are stabilizing. We therefore have
$$
J(\Gc)=J^{\ec}(\Gc).
$$
One may rewrite this equality by using, for the left-hand side, the definition of $J$ in terms of $J_1$ and $J_2$ and, for the right-hand side, assertion 2 of Proposition~\ref{thm:16.4.1}. It follows that
$$
\sum_{\Lc_1,\Lc_2\in\lc^{\Gc}(\Mc)}d_{\Mc}^{\Gc}(\Lc_1,\Lc_2)\cdot J_1(\Lc_1)\cdot J_2(\Lc_2)= \sum_{\Lc_1,\Lc_2\in\lc^{\Gc}(\Mc)}d_{\Mc}^{\Gc}(\Lc_1,\Lc_2)\cdot J_1^{\ec}(\Lc_1)\cdot J_2^{\ec}(\Lc_2).
$$
Since $J_2$ is stabilizing, we have $J_2^{\ec}=J_2$. Arguing by induction, we may indeed assume that we have $J_1^{\ec}(\Lc_1)=J_1(\Lc_1)$ for every $\Lc_1\in\lc^{\Gc}(\Mc)-\{\Gc\}$. Subtracting the equal terms from each side, one sees that the preceding equality is equivalent to
$$
J_1(\Gc)J_2(\Mc)=J_1^{\ec}(\Gc)J_2(\Mc)
$$
whence $J_1(\Gc)=J_1^{\ec}(\Gc)$ since we have assumed $J_2(\Mc)\neq 0$. Hence $J_1$ is indeed stabilizing.
\end{paragr}

\begin{paragr}[Proof of Proposition~\ref{thm:16.3.6}.]\label{S:16.5}
Let us begin with the following lemma.

\begin{lemma}\label{thm:16.5.1}
Let $s_M\in\Tc$ be such that $\Mc'=\Mc_{s_M}$ is elliptic in $\Mc$. Let $\tilde{\ec}$ be the set of triples
$$
(s,\Lc'_1,\Lc'_2)
$$
formed of $s\in s_M Z_{\Mc}/Z_{\Gc}$ and $\Lc'_1,\Lc'_2\in\lc^{\Gc_s}(\Mc')$ such that
\begin{itemize}
\item[--] $\Gc_s$ is elliptic in $\Gc$; \item[--] $d_{\Mc'}^{\Gc_s}(\Lc'_1,\Lc'_2)\neq 0$.
\end{itemize}
Let $\tilde{\lc}$ be the set of quadruples
$$
(\Lc_1,\Lc_2,s_1,s_2)
$$
formed of $\Lc_1,\Lc_2\in\lc^{\Gc}(\Mc)$ and $s_i\in Z_{\Mc}/Z_{\Lc_i}$ such that
\begin{itemize}
\item[--] $\Lc_{i,s_i}$ is elliptic in $\Lc_i$; \item[--] $d_{\Mc}^{\Gc}(\Lc_1,\Lc_2)\neq 0$.
\end{itemize}
Let $\tilde{\ec}\to\tilde{\lc}$ be the map which to $(s,\Lc'_1,\Lc'_2)$ associates $(\Lc_1,\Lc_2,s_1,s_2)$ defined by
\begin{itemize}
\item[--] $\Lc_i$ is the unique element of $\lc^{\Gc}(\Mc)$ such that $\Lc'_i$ is elliptic in $\Lc_i$ and $\Lc'_i=\Lc_{i,s}$ (cf.\ Lemma~\ref{thm:16.3.3}, assertion 3); \item[--] $s_i$ is the image of $s$ under the projection $s_M Z_{\Mc}/Z_{\Gc}\to s_M Z_{\Mc}/Z_{\Lc_i}$.
\end{itemize}
Then this map is surjective and its fibers all have the same cardinality, namely
$$
|(Z_{\Lc_1}\cap Z_{\Lc_2})/Z_{\Gc}|.
$$
Moreover, if $(s,\Lc'_1,\Lc'_2)$ has image $(\Lc_1,\Lc_2,s_1,s_2)$, we have the equality between
$$
|(Z_{\Lc_1}\cap Z_{\Lc_2})/Z_{\Gc}|\cdot|Z_{\Mc'}/Z_{\Mc}|\cdot |Z_{\Gc_s}/Z_{\Gc}|^{-1}\cdot e_{\Mc'}^{\Gc_s}(\Lc'_1,\Lc'_2)
$$
and
$$
|Z_{\Mc'}/Z_{\Mc}|^2\cdot|Z_{\Lc_{1,s}}/Z_{\Lc_1}|^{-1}\cdot |Z_{\Lc_{2,s}}/Z_{\Lc_2}|^{-1}\cdot d_{\Mc}^{\Gc}(\Lc_1,\Lc_2).
$$
\end{lemma}

\begin{preuve}
Let us first observe that if $(s,\Lc'_1,\Lc'_2)$ has image $(\Lc_1,\Lc_2,s_1,s_2)$ then $d_{\Mc'}^{\Gc_s}(\Lc'_1,\Lc'_2)=d_{\Mc}^{\Gc}(\Lc_1,\Lc_2)$ (for if $\Lc'$ is elliptic in $\Lc$, we have $\ago_{\Lc}=\ago_{\Lc'}$).

The map $\tilde{\ec}\to\tilde{\lc}$ is \emph{surjective}. Indeed, if $(\Lc_1,\Lc_2,s_1,s_2)\in\tilde{\lc}$, we have $d_{\Mc}^{\Gc}(\Lc_1,\Lc_2)\neq 0$ and by Lemma~\ref{thm:16.3.3} we have $Z_{\Mc}=Z_{\Lc_1}Z_{\Lc_2}$. Hence $(s_1,s_2)$ lifts to $s\in Z_{\Mc}/Z_{\Gc}$. The triple $(s,\Lc_{1,s},\Lc_{2,s})$ belongs to $\tilde{\ec}$: the only condition which is not obvious is that $\Gc_s$ is elliptic in $\Gc$. By hypothesis, $\Lc_{i,s}=\Lc_{i,s_i}$ is elliptic in $\Lc_i$. We therefore have
$$
\ago_{\Mc}^{\Gc}=\ago_{\Mc}^{\Lc_1}\oplus\ago_{\Mc}^{\Lc_2} =\ago_{\Mc}^{\Lc_{1,s}}\oplus\ago_{\Mc}^{\Lc_{2,s}}
$$
and, passing to the orthogonal complements, we have
$$
\ago_{\Gc}=\ago_{\Lc_{1,s}}\cap\ago_{\Lc_{2,s}}.
$$
As the intersection above contains $\ago_{\Gc_s}$, we have $\ago_{\Gc}=\ago_{\Gc_s}$ and $\Gc_s$ is elliptic in $\Gc$. It is clear that the image of $(s,\Lc_{1,s},\Lc_{2,s})$ is the quadruple $(\Lc_1,\Lc_2,s_1,s_2)\in\tilde{\lc}$ we started from.

The assertion on the cardinality of the fibers is obvious.

Let $(s,\Lc'_1,\Lc'_2)$ have image $(\Lc_1,\Lc_2,s_1,s_2)$. Let us prove the last assertion. After simplification by $d_{\Mc'}^{\Gc_s}(\Lc'_1,\Lc'_2)=d_{\Mc}^{\Gc}(\Lc_1,\Lc_2)\neq 0$, one arrives at the equality between
$$
|(Z_{\Lc_1}\cap Z_{\Lc_2})/Z_{\Gc}|\cdot|Z_{\Mc'}/Z_{\Mc}|\cdot |Z_{\Gc_s}/Z_{\Gc}|^{-1}|(Z_{\Lc'_1}\cap Z_{\Lc'_2})/Z_{\Gc_s}|^{-1}
$$
and
$$
|(Z_{\Lc_1}\cap Z_{\Lc_2})/Z_{\Gc}|^{-1}\cdot|Z_{\Mc'}/Z_{\Mc}|^2\cdot |Z_{\Lc_{1,s}}/Z_{\Lc_1}|^{-1}\cdot|Z_{\Lc_{2,s}}/Z_{\Lc_2}|^{-1}.
$$
This one results from the commutative diagram with exact rows
{\small
\[
\xymatrix@C=1pc@R=2.6pc{ 1\ar[r] & (Z_{\Lc_1}\cap Z_{\Lc_2})/Z_{\Gc}\ar[r]\ar[d]^{\varphi} & Z_{\Mc}/Z_{\Gc}\ar[r]\ar[d] &
Z_{\Mc}/Z_{\Lc_1}\times Z_{\Mc}/Z_{\Lc_2}\ar[r]\ar[d] & 1\\
1\ar[r] & (Z_{\Lc_{1,s}}\cap Z_{\Lc_{2,s}})/Z_{\Gc_s}\ar[r] & Z_{\Mc'}/Z_{\Gc_s}\ar[r] & Z_{\Mc'}/Z_{\Lc_{1,s}}\times Z_{\Mc'}/Z_{\Lc_{2,s}}\ar[r] & 1 }
\]
}
and from the following lemma.
\end{preuve}

\begin{lemma}\label{thm:16.5.2}
Let $s$ be such that $\Mc=\Mc_s$ is elliptic in $\Mc$ and $\Gc_s$ is elliptic in $\Gc$. For every $\Lc\in\lc^{\Gc}(\Mc)$ such that $\Lc_s$ is elliptic in $\Lc$, the morphism
$$
Z_{\Mc}/Z_{\Lc}\to Z_{\Mc'}/Z_{\Lc_s}
$$
is surjective and its kernel has cardinality
$$
|Z_{\Lc_s}/Z_{\Lc}|\cdot|Z_{\Mc'}/Z_{\Mc}|^{-1}.
$$
\end{lemma}

\begin{preuve}
The surjectivity is a consequence of the equalities
$$
Z_{\Mc'}=Z_{\Gc_s}(Z_{\Mc'})^0=Z_{\Gc_s}(Z_{\Mc})^0.
$$
As $Z_{\Lc_s}=Z_{\Gc_s}(Z_{\Lc_s}^0)=Z_{\Gc_s}(Z_{\Lc}^0)$, the kernel is written
$$
(Z_{\Mc}\cap Z_{\Lc_s})/Z_{\Lc}=Z_{\Lc}(Z_{\Gc_s}\cap Z_{\Mc})/Z_{\Lc}\simeq (Z_{\Gc_s}\cap Z_{\Mc})/(Z_{\Gc_s}\cap Z_{\Lc}).
$$
One concludes by remarking that $(Z_{\Gc_s}\cap Z_{\Lc})/Z_{\Gc}$ is the kernel of the surjective morphism
$$
Z_{\Gc_s}/Z_{\Gc}\to Z_{\Lc_s}/Z_{\Lc}.
$$
\end{preuve}

We can now give the proof of Proposition~\ref{thm:16.4.1}. Let us prove assertion 1. Let $\Hc\in\ec_M^G(s_M)$. Let $\tilde{\ec}$ and $\tilde{\lc}$ be the sets defined in Lemma~\ref{thm:16.5.1} relative to the group $\Hc$, to its Levi subgroup $\Mc'$ and to the element $s_M=1$. By the definition of the maps $S$ and $S_i$, we have
\begin{align*}
J(\Hc)&=\sum_{\Lc_1,\Lc_2\in\lc^{\Hc}(\Mc')}
d_{\Mc'}^{\Hc}(\Lc_1,\Lc_2)\cdot J_1(\Lc_1)\cdot J_2(\Lc_2)\\
&=\sum_{\Lc_1,\Lc_2\in\lc^{\Hc}(\Mc')}d_{\Mc'}^{\Hc}(\Lc_1,\Lc_2)
\Bigl[\sum_{s_1\in Z_{\Mc'}/Z_{\Lc_1}}
|Z_{\Lc_{1,s_1}}/Z_{\Lc_1}|^{-1}S_1(\Lc_{1,s_1})\Bigr]\\
&\qquad\cdot\Bigl[\sum_{s_2\in Z_{\Mc'}/Z_{\Lc_2}}
|Z_{\Lc_{2,s_2}}/Z_{\Lc_2}|^{-1}S_2(\Lc_{2,s_2})\Bigr]\\
&=\sum_{(\Lc_1,\Lc_2,s_1,s_2)\in\tilde{\lc}}d_{\Mc'}^{\Hc}(\Lc_1,\Lc_2)\cdot
|Z_{\Lc_{1,s_1}}/Z_{\Lc_1}|^{-1}\cdot|Z_{\Lc_{2,s_2}}/Z_{\Lc_2}|^{-1}\\
&\qquad\cdot S_1(\Lc_{1,s_1})\cdot S_2(\Lc_{2,s_2})
\end{align*}

Using Lemma~\ref{thm:16.5.1}, we obtain
\begin{align*}
J(\Hc)&=\sum_{(s,\Lc'_1,\Lc'_2)\in\tilde{\ec}}|Z_{\Hc_s}/Z_{\Hc}|^{-1}
e_{\Mc'}^{\Hc_s}(\Lc'_1,\Lc'_2)\cdot S_1(\Lc_1)\cdot S_2(\Lc_2)\\
&=\sum_{s\in Z_{\Mc'}/Z_{\Hc}}|Z_{\Hc_s}/Z_{\Hc}|^{-1}
\Bigl(\sum_{\Lc_1,\Lc_2\in\lc^{\Hc_s}(\Mc')}
e_{\Mc'}^{\Hc_s}(\Lc_1,\Lc_2)\cdot S_1(\Lc_1)\cdot S_2(\Lc_2)\Bigr).
\end{align*}

By the uniqueness in Corollary~\ref{thm:13.2.6}, we obtain the desired expression for $S(\Hc)$.

Let us next prove assertion 2. Let $\tilde{\ec}$ and $\tilde{\lc}$ be the sets defined in Lemma~\ref{thm:16.5.1} relative to the group $\Gc$, to its Levi subgroup $\Mc$ and to the element $s_M$. By the definition of $J^{\ec}$ and assertion 1 which we have just proved, we have
\begin{align*}
J^{\ec}(\Gc)&=|Z_{\Mc'}/Z_{\Mc}|\sum_{s\in s_M Z_{\Mc}/Z_{\Gc}}
|Z_{\Gc_s}/Z_{\Gc}|^{-1}S(\Gc_s)\\
&=|Z_{\Mc'}/Z_{\Mc}|\sum_{s\in s_M Z_{\Mc}/Z_{\Gc}}|Z_{\Gc_s}/Z_{\Gc}|^{-1}\\
&\qquad\cdot\sum_{\Lc'_1,\Lc'_2\in\lc^{\Gc_s}(\Mc')}
e_{\Mc'}^{\Gc_s}(\Lc'_1,\Lc'_2)\cdot S_1(\Lc'_1)\cdot S_2(\Lc'_2)\\
&=|Z_{\Mc'}/Z_{\Mc}|\sum_{(s,\Lc'_1,\Lc'_2)\in\tilde{\ec}}
|Z_{\Gc_s}/Z_{\Gc}|^{-1}\cdot e_{\Mc'}^{\Gc_s}(\Lc'_1,\Lc'_2)\cdot
S_1(\Lc'_1)\cdot S_2(\Lc'_2).
\end{align*}

Lemma~\ref{thm:16.5.1} implies that we have
\begin{align*}
J^{\ec}(\Gc)&=\sum_{(\Lc_1,\Lc_2,s_1,s_2)\in\tilde{\lc}}
|Z_{\Mc'}/Z_{\Mc}|^2\cdot|Z_{\Lc_{1,s}}/Z_{\Lc_1}|^{-1}\cdot
|Z_{\Lc_{2,s}}/Z_{\Lc_2}|^{-1}\\
&\qquad\cdot d_{\Mc}^{\Gc}(\Lc_1,\Lc_2)\cdot S_1(\Lc'_1)\cdot S_2(\Lc'_2)\\
&=\sum_{\Lc_1,\Lc_2\in\lc^{\Lc}(\Mc)}d_{\Mc}^{\Lc}(\Lc_1,\Lc_2)\cdot
J_1^{\ec}(\Lc_1)\cdot J_2^{\ec}(\Lc_2).
\end{align*}

The last equality results from the definition of $J_1^{\ec}$ and $J_2^{\ec}$. This therefore proves assertion 2.
\end{paragr}

\begin{paragr}[Proof of Proposition~\ref{thm:16.3.7}.]\label{S:16.6}
Let us state the following proposition, which will be proved in \S\ref{S:16.7}.

\begin{proposition}\label{thm:16.6.1}
The hypotheses are those of Proposition~\ref{thm:16.3.7}.
\begin{enumerate}
\item Let $S_M$ and $S_R$ be the maps deduced from $J_M$ and $J_R$ by Corollary~\ref{thm:13.2.6}. For every $\Hc\in\ec_M^G(s_M)$, we have
$$
S_M(\Hc)=\sum_{\Lc\in\lc^{\Hc}(\Rc')}e_{\Rc'}^{\Hc}(\Mc',\Lc)S_R(\Lc)\ ;
$$
\item Let $J_M^{\ec}$ and $J_R^{\ec}$ be the maps deduced from $J_M$ and $J_R$ (cf.\ Definition~\ref{thm:16.2.1} and Remark~\ref{thm:16.2.3}). For every $\Lc\in\lc^{\Gc}(\Mc)$, we have
$$
J_M^{\ec}(\Lc)=\sum_{\Lc_1\in\lc^{\Lc}(\Mc)} d_{\Rc}^{\Lc}(\Mc,\Lc_1)J_R^{\ec}(\Lc_1).
$$
\end{enumerate}
\end{proposition}

One may then give the proof of Proposition~\ref{thm:16.3.7}. Using assertion 2 of Proposition~\ref{thm:16.6.1} and the fact that $J_R$ is stabilizing, we have
\begin{align*}
J_M^{\ec}(\Lc)&=\sum_{\Lc_1\in\lc^{\Lc}(\Rc)}
d_{\Rc}^{\Lc}(\Mc,\Lc_1)J_R^{\ec}(\Lc_1)\\
&=\sum_{\Lc_1\in\lc^{\Lc}(\Rc)}d_{\Rc}^{\Lc}(\Mc,\Lc_1)J_R(\Lc_1)\\
&=J_M(\Lc).
\end{align*}
Hence $J_M$ is stabilizing.
\end{paragr}

\begin{paragr}[Proof of Proposition~\ref{thm:16.6.1}.]\label{S:16.7}
It is deduced from Lemma~\ref{thm:16.7.1} below (cf.\ the proof of Proposition~\ref{thm:16.4.1}).

\begin{lemma}\label{thm:16.7.1}
Let $\Rc\subset\Mc\subset\Gc$ be two Levi subgroups of $\Gc$. Let $s_M\in\Tc$, $\Mc'=\Mc_{s_M}$ and $\Rc'=\Rc_{s_M}$. We assume that $\Mc'$ and $\Rc'$ are elliptic in $\Mc$ and $\Rc$ respectively.

Let $\tilde{\ec}$ be the set of pairs $(s',\Lc')$ formed of $s'\in s_M Z_{\Mc}/Z_{\Gc}$ and $\Lc'\in\lc^{\Gc_s}(\Rc')$ such that
\begin{itemize}
\item[--] $\Gc_{s'}$ is elliptic in $\Gc$; \item[--] $d_{\Rc'}^{\Gc_{s'}}(\Mc',\Lc')\neq 0$.
\end{itemize}
Let $\tilde{\lc}$ be the set of pairs $(\Lc,s)$ formed of $\Lc\in\lc^{\Gc}(\Rc)$ and $s\in Z_{\Mc}/Z_{\Lc}$ such that
\begin{itemize}
\item[--] $\Lc_s$ is elliptic in $\Lc$; \item[--] $d_{\Rc}^{\Gc}(\Mc,\Lc)\neq 0$.
\end{itemize}
Let $\tilde{\ec}\to\tilde{\lc}$ be the map which to $(s',\Lc')$ associates $(\Lc,s)$ defined by
\begin{itemize}
\item[--] $\Lc$ is the unique element of $\lc^{\Gc}(\Rc)$ such that $\Lc'$ is elliptic in $\Lc$ and $\Lc'=\Lc_{s'}$ (cf.\ Lemma~\ref{thm:16.3.3}, assertion 3); \item[--] $s$ is the image of $s'$ under the projection $s_M Z_{\Mc}/Z_{\Gc}\to s_M Z_{\Rc}/Z_{\Lc}$.
\end{itemize}
This map is surjective and its fibers all have the same cardinality, namely
$$
|(Z_{\Mc}\cap Z_{\Rc})/Z_{\Gc}|.
$$
Moreover, if $(s',\Lc')$ has image $(\Lc,s)$, we have the equality between
$$
|(Z_{\Mc}\cap Z_{\Lc})/Z_{\Gc}|\cdot|Z_{\Mc'}/Z_{\Mc}|\cdot |Z_{\Gc_{s'}}/Z_{\Gc}|^{-1}\cdot e_{\Rc'}^{\Gc_{s'}}(\Mc',\Lc')
$$
and
$$
|Z_{\Rc'}/Z_{\Rc}|\cdot|Z_{\Lc_s}/Z_{\Lc}|^{-1}\cdot d_{\Rc}^{\Gc}(\Mc,\Lc).
$$
\end{lemma}

\begin{preuve}
It is analogous to that of Lemma~\ref{thm:16.5.1}. The details are left to the reader.
\end{preuve}
\end{paragr}
\section{From the global to the local}\label{sec:17}

\begin{paragr}\label{S:17.1}
The main result is Theorem~\ref{thm:17.3.1}, which is a variant of a ``local'' nature of Theorem~\ref{thm:15.1.1}. The splitting and descent results of Section~\ref{sec:16} play a key role, as do certain local computations of weighted orbital integrals. The notations are those used throughout the article, in particular those of Sections~\ref{sec:11}, \ref{sec:13} and \ref{sec:16}.
\end{paragr}

\begin{paragr}[Notations.]\label{S:17.2}
Let $V$ be a finite or infinite but non-empty set of closed points of the curve $C$ which is stable under the Frobenius $\tau$. Let $\AAA_V$ be the subgroup of adeles of $\AAA$ whose components outside $V$ are zero. The group $\AAA_V$ has an obvious ring structure. If $V$ is finite (resp.\ infinite), $\AAA_V$ is the product (resp.\ restricted product) of the completions of the function field $F$ at the points of $V$. Let $\oc_V=\prod_{v\in V}\oc_v$ where $\oc_v$ is the completion of the local ring of $C$ at $v$.

Let $M\in\lc^G$, $a_M\in\Ac_M(k)^\tau$ and $X_{a_M}=\eps_M(a_{M,\eta})\in \mgo(F)$ (cf.\ \S\ref{S:11.7}). Let $J_{a_M}$ be the group scheme over $C$ defined in Section~\ref{S:11.3} whose fibre at the generic point of $C$ is the centralizer of $X_{a_M}$ in $M\times_k F$ (it is a maximal $F$-subtorus of $M\times_k F$).

We complete the choices of Haar measures made in Section~\ref{S:11.5}. We endow $J_{a_M}(\AAA_V)^\tau$ and $G(\AAA_V)^\tau$ with the Haar measures which give volume $1$ respectively to the maximal compact subgroup of $J_{a_M}(\AAA_V)^\tau$ and to $G(\oc_V)^\tau$. As in Section~\ref{S:11.5}, we deduce from this a measure on $(J_{a_M}(\AAA_V)\back G(\AAA_V))^\tau$ which is invariant on the right under $G(\AAA_V)^\tau$. If $V'$ is the complement of $V$ in the set of closed points of $C$, we have a canonical bijection
$$
(J_{a_M}(\AAA_V)\back G(\AAA_V))^\tau\times (J_{a_M}(\AAA_{V'})\back G(\AAA_{V'}))^\tau\to (J_{a_M}(\AAA)\back G(\AAA))^\tau
$$
which preserves the measures.

For every $s\in\Tc^{W_{a_M}}$, we define the weighted $s$-orbital integral on $V$ associated with $a_M$ by
\begin{equation}\label{eq:17.2.1}
\begin{split}
J_{M,V}^{G,s}(a_M)&=q^{-\deg(D_V)|\Phi^G|/2}|D^G(X_{a_M})|_V^{1/2}\\
&\quad\times\int_{(J_{a_M}(\AAA_V)\back G(\AAA_V))^\tau}
\bg s,\inv_{a_M}(g)\bd\mathbf{1}_D(\Ad(g^{-1})X_{a_M})v_M^G(g)\,dg,
\end{split}
\end{equation}
where $D_V$ is the restriction of $D$ to $V$ and where $|D^G(X_a)|_V$ is the product over $V$ of the usual absolute values of the discriminant $D^G$ of $X_a$ (cf.\ \S\ref{S:2.7}).

\begin{remark}\label{thm:17.2.1}
When $V$ is the set of all closed points, the weighted orbital integral defined by~(\ref{eq:17.2.1}) is nothing but the integral of Definition~\ref{thm:11.7.1}.
\end{remark}
\end{paragr}

\begin{paragr}[The main statement.]\label{S:17.3}
We can now state the main result of this section. We use there the divisors $\mathfrak{R}^G_M$ and $\mathfrak{D}^M$ of $\car_{M,D}$ defined in Section~\ref{S:9.1}.

\begin{theorem}\label{thm:17.3.1}
We place ourselves under the hypotheses of Theorem~\ref{thm:15.1.1}. Let
$$
h_{a_M}:C\to\car_{M',D}
$$
be the section associated with $a_M\in\Ac_{M'}(k)$. Let $V$ be a non-empty $\tau$-stable set of closed points of $C$ such that every closed point $v\notin V$ satisfies one of the two following conditions{\rm :}
\begin{enumerate}
\item $h_{a_M}(v)$ is a smooth point of $\mathfrak{R}^G_{M'}$ and does not belong to $\mathfrak{D}^{M'}${\rm ;} moreover, the intersection of $h_a(C)$ and $\mathfrak{R}^G_{M'}$ is transverse at $h_a(v)${\rm ;} \item $h_{a_M}(v)$ is a smooth point of $\mathfrak{D}^{M'}$ and does not belong to $\mathfrak{R}^G_{M'}${\rm ;} moreover, the intersection of $h_a(C)$ and $\mathfrak{D}^{M'}$ is transverse at $h_a(v)$.
\end{enumerate}
Then the map $J_V:\ec^G_M(s_M)\cup\lc^{\Gc}(\Mc)\to\CC$ defined for every $\Hc\in\ec^G_M(s_M)$ by $J_V(\Hc)=J^{H,1}_{M',V}(a_M)$ and for every $\Lc\in\lc^{\Gc}(\Mc)$ by $J_V(\Lc)=J^{L,s_M}_{M,V}(a_M)$ is stabilizing in the sense of Definition~\ref{thm:16.2.4}.
\end{theorem}

\begin{remark}\label{thm:17.3.2}
If $V$ is the set of all closed points of $C$, the theorem is merely a paraphrase of Theorem~\ref{thm:15.1.1} (cf.\ Example~\ref{thm:16.2.6}). The set of the $v$ which satisfy neither assertion 1 nor assertion 2 is finite. In other words there always exists a finite set $V$ which satisfies the hypotheses of the theorem. In this sense, one may say that Theorem~\ref{thm:17.3.1} is a local variant of Theorem~\ref{thm:15.1.1}.
\end{remark}

\begin{preuve}
Let $V$ be the set of closed points of $C$. Let $V_1$ be a $\tau$-stable set of closed points and $V_2$ its complement in $V$. We suppose that every $v\in V_2$ satisfies one of the assertions 1 and 2 of Theorem~\ref{thm:17.3.1}. We have to see that $J_{V_1}$ is stabilizing. Now $J_V$, $J_{V_1}$ and $J_{V_2}$ satisfy the splitting relations of Proposition~\ref{thm:16.3.6}, as Proposition~\ref{thm:17.4.1} below shows. It therefore suffices, by assertion 2 of Proposition~\ref{thm:16.3.6}, to show that $J_V$ and $J_{V_2}$ are stabilizing and that $J_2(\Mc)\neq 0$. We have already said in Remark~\ref{thm:17.3.2} above that $J_V$ is stable. There exists a (finite) partition of $V_2$ into parts $V'_i$ which satisfy the hypotheses of Lemma~\ref{thm:17.8.1}. We thus have $J_2(\Mc)=\prod_i J_{V'_i}(\Mc)\neq 0$ (with the notations of Lemma~\ref{thm:17.8.1}) and, by a repeated application of assertion 1 of Proposition~\ref{thm:16.3.6}, we see that $J_2$ is stabilizing and that $J_2(\Mc)\neq 0$.
\end{preuve}
\end{paragr}

\begin{paragr}[Splitting formulas for weighted orbital integrals.]\label{S:17.4}
For all $L_1,L_2\in\lc^G(M)$, we set
$$
d^G_M(L_1,L_2)=d^{\Gc}_{\Mc}(\Lc_1,\Lc_2),
$$
where the second one is that of Definition~\ref{thm:16.3.1}. Here are Arthur's splitting formulas.

\begin{proposition}\label{thm:17.4.1}
Let $V_1$ and $V_2$ be two disjoint non-empty sets of closed points of $C$. Let $V=V_1\cup V_2$. For every $a\in\Ac_M(k)^\tau$ and every $s\in\Tc^{W_a}$, we have the equality
$$
J^{G,s}_{M,V}(a_M)=\sum_{L_1,L_2\in\lc^G(M)}d^G_M(L_1,L_2)\cdot J^{L_1,s}_{M,V}(a_M)\cdot J^{L_2,s}_{M,V}(a_M).
$$
\end{proposition}

\begin{preuve}
Arthur has defined a section $(L_1,L_2)\mapsto(Q_{L_1},Q_{L_2})$ of the map $\fc(M)\times\fc(M)\to\lc(M)\times\lc(M)$ given by $(Q_1,Q_2)\mapsto(M_{Q_1},M_{Q_2})$, defined on the set of the $(L_1,L_2)$ such that $d^G_M(L_1,L_2)\neq 0$. This section has the following property: for every $g\in(J_{a_M}(\AAA_V)\back G(\AAA_V))^\tau$ with inverse image $(g_1,g_2)$ under the canonical bijection
\begin{equation}\label{eq:17.4.1}
(J_{a_M}(\AAA_{V_1})\back G(\AAA_{V_1}))^\tau\times
(J_{a_M}(\AAA_{V_2})\back G(\AAA_{V_2}))^\tau\to
(J_{a_M}(\AAA_V)\back G(\AAA_V))^\tau,
\end{equation}
we have
\begin{equation}\label{eq:17.4.2}
v^G_M(g)=\sum_{L_1,L_2\in\lc^G(M)}d^G_M(L_1,L_2)\cdot v^{Q_{L_1}}_M(g_1)\cdot
v^{Q_{L_2}}_M(g_2),
\end{equation}
where for every $Q\in\fc(M)$ and every $h\in G(\AAA)$ we set
$$
v^Q_M(h)=v^{M_Q}_M(l),
$$
where $l$ is the unique element of $M(\AAA)\back M(\oc)$ such that $g\in lN(\AAA)G(\oc)$ (Iwasawa decomposition). Let us introduce the integral $J^{Q,s}_{M,V}(a_M)$ given by the equality~(\ref{eq:17.2.1}) where one has replaced the weight $v^G_M$ by the weight $v^Q_M$. Since the bijection~(\ref{eq:17.4.1}) preserves the measures, formula~(\ref{eq:17.4.2}) implies that we have
$$
J^{G,s}_{M,V}(a_M)=\sum_{L_1,L_2\in\lc^G(M)}d^G_M(L_1,L_2)\cdot J^{Q_{L_1},s}_{M,V}(a_M)\cdot J^{Q_{L_2},s}_{M,V}(a_M).
$$
To conclude, it suffices to show that for every $L\in\lc(M)$ and every $Q\in\pc(L)$, we have
$$
J^{Q,s}_{M,V}(a_M)=J^{L,s}_{M,V}(a_M).
$$
One shows with the help of Lang's isogeny theorem that one has a decomposition
\begin{equation}\label{eq:17.4.3}
(J_{a_M}(\AAA_V)\back G(\AAA_V))^\tau=(J_{a_M}(\AAA_V)\back L(\AAA_V))^\tau
N_Q(\AAA_V)^\tau G(\oc_V)^\tau
\end{equation}
which is compatible with the measures provided that $N_Q(\AAA_V)^\tau$ is endowed with the Haar measure which gives volume $1$ to the subgroup $(N_Q(\oc_V))^\tau$. By the change of variable $g=lnk$, with notations that we hope are obvious, we see that $J^{Q,s}_{M,V}(a_M)$ is equal to
\[
\begin{aligned}
q^{-\deg(D_V)|\Phi^G|/2}&|D^G(X_a)|_V^{1/2}
\int_{(J_{a_M}(\AAA_V)\back L(\AAA_V))^\tau}\int_{N_Q(\AAA_V)^\tau}\\
&\times\int_{G(\oc_V)^\tau}\bg s,\inv_{a_M}(lnk)\bd
\mathbf{1}_D(\Ad((lnk)^{-1})X_{a_M})v^L_M(l)\,dl\,dn\,dk.
\end{aligned}
\]
By Definition~\ref{thm:11.4.1}, the invariant $\inv_{a_M}(lnk)$ depends only on $lnk\tau(lnk)^{-1}=l\tau(l)^{-1}$. We thus have $\inv_{a_M}(lnk)=\inv_{a_M}(l)$. The integrand is then $G(\oc_V)^\tau$-invariant and the integral over $G(\oc_V)^\tau$, which is equal to $1$, may be omitted. We then use the change of variable $\Ad((ln)^{-1})X_{a_M}=\Ad(l^{-1})X_{a_M}+U$ (where $U$ belongs to $\ngo_Q(\AAA_V)^\tau$, which is endowed with the Haar measure that gives volume $1$ to $\ngo_Q(\oc_V)^\tau$), of Jacobian $|D^G(X_a)|_V^{-1/2}|D^L(X_a)|_V^{1/2}$, to obtain that $J^{Q,s}_{M,V}(a_M)$ is equal to
\[
\begin{aligned}
q^{-\deg(D_V)|\Phi^G|/2}&|D^L(X_a)|_V^{1/2}
\int_{(J_{a_M}(\AAA_V)\back L(\AAA_V))^\tau}\\
&\times\int_{\ngo_Q(\AAA_V)^\tau}\bg s,\inv_{a_M}(l)\bd
\mathbf{1}_D(\Ad(l^{-1})X_{a_M}+U)\,v^L_M(l)\,dl\,dU.
\end{aligned}
\]
Since we have
\[
\begin{aligned}
\int_{\ngo_Q(\AAA_V)^\tau}&\mathbf{1}_D(\Ad(l^{-1})X_{a_M}+U)\,dU\\
&=\mathbf{1}_D(\Ad(l^{-1})X_{a_M})\int_{\ngo_Q(\AAA_V)^\tau}
\mathbf{1}_D(\Ad(l^{-1})X_{a_M}+U)\,dU\\
&=q^{\deg(D_V)\dim(\ngo_Q)}\mathbf{1}_D(\Ad(l^{-1})X_{a_M})
\end{aligned}
\]
and $q^{-\deg(D_V)|\Phi^G|/2}q^{\deg(D_V)\dim(\ngo_Q)}= q^{-\deg(D_V)|\Phi^L|/2}$, we indeed obtain $J^{Q,s}_{M,V}(a_M)= J^{L,s}_{M,V}(a_M)$, which concludes the proof.
\end{preuve}
\end{paragr}

\begin{paragr}\label{S:17.5}
Let $a_M=(h_{a_M},t)\in\Ac_M(k)^\tau$. Let $h_{a_M,\eta}$ be the restriction of $h_a$ to the generic point of the curve $C$. One then views $h_{a_M,\eta}$ as a point in $\car_M(F)$ where $F$ is the field of functions of $F$. Let $V$ be a $\tau$-stable set of closed points of $C$. Let $\varpi^{D_V}=\prod_{v\in v}\varpi_v^{d_v}\in\oc_V$. We have $\varpi^{D_V}h_{a_M,\eta}\in\car_M(\oc_V)$. Let $X_{a_M,V}=\eps_M(\varpi^{D_V}h_{a_M,\eta})\in\mgo(\oc_V)$ be the image under the Kostant section $\eps_M$. Let us recall that we have $X_{a_M}=\eps_M(h_{a_M,\eta})\in\mgo(F)$. Outside the support of $D$, these two elements are in general not equal. Nevertheless, they always satisfy the following relation
$$
X_{a_M,V}=\varpi^{D_V}\Ad(\rho_M(\varpi^{D_V/2}))X_{a_M},
$$
where $\rho_M$ is the sum of the coroots in $M$ that are positive for the choice of the pinning which determines the Kostant section $\eps_M$. We recall that the divisor $D$ is even. The element $\rho_M(\varpi^{D_V/2})$ is $\tau$-stable. It conjugates the centralizers $J_{a_M}$ and $J_{X_{a_M,V}}$. The translation by $\rho_M(\varpi^{D_V/2})$ sends $(J_{a_M}(\AAA_V)\back G(\AAA_V))^\tau$ bijectively onto $(J_{a_M,V}(\AAA_V)\back G(\AAA_V))^\tau$, which we endow with the measure obtained by transport. In the same way, we denote by $\inv_{a_M,V}$ the map defined on $(J_{a_M,V}(\AAA_V)\back G(\AAA_V))^\tau$ which, composed with the previous bijection, gives back $\inv_{a_M}$. Let $\mathbf{1}_{\ggo(\oc)}$ be the characteristic function of $\ggo(\oc)$. By a change of variables, we obtain the following lemma.

\begin{lemma}\label{thm:17.5.1}
The notations are those used above. For every $s\in\Tc^{W_{a_M}}$, we have the equality
\[
\begin{aligned}
J^{G,s}_{M,V}(a_M)&=|D^G(X_{a_M,V})|_V^{1/2}\\
&\quad\times\int_{(J_{a_M,V}(\AAA_V)\back G(\AAA_V))^\tau}
\bg s,\inv_{a_M,V}(g)\bd\mathbf{1}_{\ggo(\oc)}(\Ad(g^{-1})X_{a_M,v})
v^G_M(g)\,dg.
\end{aligned}
\]
\end{lemma}
\end{paragr}

\begin{paragr}[Local computations of weighted orbital integrals.]\label{S:17.6}
Throughout this paragraph, we suppose that $M\in\lc^G$ satisfies $M\neq G$. Let $a_M=(h_{a_M},t)\in\Ac_M(k)$ and $s\in\Tc^{W_{a_M}}$. We are going to compute in a few simple cases the local weighted orbital integrals.

\begin{lemma}\label{thm:17.6.1}
Let $V$ be a $\tau$-stable set of closed points of $C$. We suppose that for every $v\in V$, the section $h_{a_M}$ is such that $h_{a_M}(v)$ does not belong to $\mathfrak{R}^G_M$. Then $J^{G,s}_{M,V}(a_M)=0$.
\end{lemma}

\begin{preuve}
We use the notations of Section~\ref{S:17.5}. By Lemma~\ref{thm:17.5.1}, we have
\[
\begin{aligned}
J^{G,s}_{M,V}(a_M)&=|D^G(X_{a_M,V})|_V^{1/2}\\
&\quad\times\int_{(J_{a_M,V}(\AAA_V)\back G(\AAA_V))^\tau}
\bg s,\inv_{a_M,V}(g)\bd\mathbf{1}_{\ggo(\oc)}(\Ad(g^{-1})X_{a_M,V})
v^G_M(g)\,dg.
\end{aligned}
\]
Let $g\in G(\AAA_V)$ be such that the integrand is non-zero. We thus have $\Ad(g^{-1})X_{a_M,V}\in\ggo(\oc_V)$. Let $P\in\pc(M)$. We write $g=mnk$ with $m\in M(\AAA_V)$, $n\in N_P(\AAA_V)$ and $k\in G(\oc_V)$ (Iwasawa decomposition). It follows that
$$
\Ad(m^{-1})X_{a_M,V}\in\mgo(\oc_V)
$$
and
\begin{equation}\label{eq:17.6.1}
\Ad(n^{-1})\Ad(m^{-1})X_{a_M,V}-\Ad(m^{-1})X_{a_M,V}\in\ngo_P(\oc_V).
\end{equation}
The hypothesis on $h_{a_M}$ implies that the adjoint action of $\Ad(m^{-1})X_{a_M,V}$ on $\ngo_P(\oc_V)$ is an $\oc_V$-isomorphism. By a standard dévissage argument, we deduce from this that condition~(\ref{eq:17.6.1}) implies $n\in N_P(\oc)$. But then we have $v^G_M(g)=v^G_M(n)=v^G_M(1)=0$ (since $M\neq G$). Hence the integrand is always zero and the integral as well, a fortiori.
\end{preuve}

\begin{lemma}\label{thm:17.6.2}
Let $v$ be a closed point of $C$ and $n\geq 1$ the smallest integer such that $\tau^n(v)=v$. Let $V$ be the orbit of $v$ under the action of $\tau$. We suppose that $v$ satisfies the following conditions{\rm :}
\begin{itemize}
\item[--] $h_{a_M}(c)$ is a smooth point of $\mathfrak{R}^G_M${\rm ;} \item[--] $h_{a_M}(c)$ does not belong to $\mathfrak{D}^M${\rm ;} \item[--] the intersection of $h_{a_M}(C)$ with $\mathfrak{R}^G_M$ is transverse at $h_{a_M}(c)$.
\end{itemize}
We then have the following conclusions{\rm :}
\begin{enumerate}
\item There exists $Y_v\in\tgo(\oc_v)$ such that $\chi^T_M(Y_v)=\chi_M(X_{a_M,v})$. For every such $Y_v$, there exists a unique pair $\{\alpha,-\alpha\}\subset\Phi^G_T-\Phi^M_T$ such that $\val_v(\pm\alpha(Y_v))=1$. Let $w$ be the unique element of $W^M$ such that $w\tau(Y_v)=Y_v$. Then $w(\alpha)=\alpha$. \item With the notations above, let $L\in\lc^G$ be defined by $\Phi^L_T=\{\alpha,-\alpha\}$ and let $(B_L,T,X_\alpha)$ be a pinning of $L$. Let $w$ be the unique automorphism of $L$ which fixes the pinning and which is of the form $\Ad(m)$ where $m\in\Norm_{M(\oc_v)}(T)$ is a lift of $w$. We have
\[
\begin{aligned}
J^{G,s}_{M,V}(a_M)&=d^G_T(M,L)|D^L(Y_v)|^{|V|/2}|V|^{\dim(\ago^G_M)}\\
&\quad\times\int_{T(F_v)^{w\tau^n}\back L(F_v)^{w\tau^n}}
\mathbf{1}_{\lgo(\oc_v)}(\Ad(l^{-1})Y_v)v^L_T(l)\,dl,
\end{aligned}
\]
where the groups $T(F_v)^{w\tau^n}$ and $L(F_v)^{w\tau^n}$ are endowed with the Haar measures which give volume $1$ respectively to the groups $T(\oc_v)^{w\tau^n}$ and $L(\oc_v)^{w\tau^n}$.
\end{enumerate}
\end{lemma}

\begin{preuve}
Since $h_{a_M}(c)$ does not belong to $\mathfrak{D}^M$, we have $\chi_M(X_{a_M,v})\in\car_M^{M-\reg}(\oc_v)$. Its reduction mod $\varpi_v$ therefore defines an element of $\car_M^{M-\reg}(k)$ which obviously lifts to an element $y$ of $\tgo^{M-\reg}(k)$. Now over the open subset $\tgo^{M-\reg}$, we know that $\chi^T_M$ is étale. It follows that $y$ lifts to a unique element $Y_v\in\tgo^{M-\reg}(\oc_v)$ such that $\chi^T_M(Y_v)=\chi_M(X_{a_M,v})$.

Let $P\in\pc(M)$. The hypotheses on the intersection of $h_{a_M}(C)$ with $\mathfrak{R}^G_M$ imply that there exists a unique root $\alpha$ of $T$ in $N_P$ such that $\val_v(\alpha(Y_v))=1$, the other roots of $T$ in $N_P$ being of zero valuation at $Y_v$. Reasoning with the parabolic subgroup opposite to $P$, we see that $\alpha$ and $-\alpha$ are the only roots outside $M$ such that $\val_v(\pm\alpha(Y_v))=1$. All the other roots, including those in $M$ (given that we have $h_{a_M}(c)\notin\mathfrak{D}^M$), have valuation $0$ at $Y_v$. Let $Y_v$ and $\alpha\in\Phi^G_T$ satisfy assertion 1. There exists $w\in W^M$ such that $w\tau(Y_v)=Y_v$ since $\chi^T_M(Y_v)\in\car_M(\oc_v)^\tau$. We have $\val_v(\alpha(Y_v))=\val_v(\tau(\alpha(Y_v)))=\val_v(\tau(\alpha)(Y_v))$. We therefore necessarily have $w(\alpha)=\pm\alpha$, but since $\alpha$ is a root in the unipotent radical of a parabolic subgroup $P\in\pc(M)$, the same holds for $w(\alpha)$ and hence $w(\alpha)=\alpha$. This therefore proves assertion 1.

Let us next prove assertion 2. Let us start from the equality (cf.\ Lemma~\ref{thm:17.5.1})
\[
\begin{aligned}
J^{G,s}_{M,V}(a_M)&=|D^G(X_{a_M,V})|_V^{1/2}\\
&\quad\int_{(J_{a_M,V}(\AAA_V)\back G(\AAA_V))^\tau}
\bg s,\inv_{a_M,V}(g)\bd\mathbf{1}_{\ggo(\oc)}(\Ad(g^{-1})X_{a_M,V})
v^G_M(g)\,dg.
\end{aligned}
\]
Let $g\in G(\AAA_V)$ and $j\in J_{a_M,V}(\AAA_V)$ be such that $\tau(g)=jg$. We shall first prove that if $\mathbf{1}_{\ggo(\oc_V)}(\Ad(g^{-1})X_{a_M,V})\neq 0$ then $\inv_{a_M,V}(g)=0$. For this, it suffices to prove that the image of $j$ in the coinvariants $J(\AAA_V)_\tau$ is trivial. Let $P\in\pc(M)$ and $g=mnk$ be the Iwasawa decomposition of $g$ associated with $P$. We thus have $\tau(m)\in jmM(\oc_V)$. By an application of Lang's isogeny theorem, we see that, after translating $m$ by an element of $M(\oc_V)$, we have $\tau(m)=jm$, which we suppose from now on. As in the proof of Lemma~\ref{thm:17.5.1}, the condition $\Ad(g^{-1})X_{a_M,V}\in\ggo(\oc_V)$ implies
$$
\Ad(m^{-1})X_{a_M,V}\in\mgo(\oc_V).
$$
Since $h_{a_M}(c)$ does not belong to $\mathfrak{D}^M$, we have $\chi_M(X_{a_M,V})\in\car_M^{M-\reg}(\oc_V)$. Consequently, we have $m\in J_{a_M,V}(\AAA_V)M(\oc_V)$. Since $\inv_{a_M,V}$ is invariant under left translation by $J_{a_M,V}(\AAA_V)$, we may and shall suppose that $m\in M(\oc_V)$. But then we have $j\in M(\oc_V)\cap J_{a_M,V}(\AAA_V)$. Let $J$ be the $\oc_V$-group scheme defined as the centralizer of $X_{a_M,V}$ in $M\times_k\AAA_V$. Since $\chi_M(X_{a_M,V})\in\car_M^{M-\reg}(\oc_V)$, it is in fact a scheme in tori. We have $M(\oc_V)\cap J_{a_M,V}(\AAA_V)=J(\oc_V)$ and a new application of Lang's theorem shows that $j$ is trivial in the coinvariants. It follows that one may restrict the integral to the $g\in(J_{a_M,V}(\AAA_V)\back G(\AAA_V))^\tau$ with trivial invariant $\inv_{a_M,V}(g)$, that is to say that one may restrict the integral to $J_{a_M,V}(\AAA_V)^\tau\back G(\AAA_V)^\tau$. We thus have
\[
\begin{aligned}
J^{G,s}_{M,V}(a_M)&=|D^G(X_{a_M,V})|_V^{1/2}\\
&\quad\times\int_{J_{a_M,V}(\AAA_V)^\tau\back G(\AAA_V)^\tau}
\mathbf{1}_{\ggo(\oc_V)}(\Ad(g^{-1})X_{a_M,V})v^G_M(g)\,dg\\
&=|D^G(X_{a_M,v})|_v^{|V|/2}|V|^{\dim(\ago^G_M)}\\
&\quad\times\int_{J_{a_M,v}(F_v)^{\tau^n}\back G(F_v)^{\tau^n}}
\mathbf{1}_{\ggo(\oc_v)}(\Ad(g^{-1})X_{a_M,v})v^G_M(g)\,dg,
\end{aligned}
\]
where one uses the obvious bijection
$$
J_{a_M,V}(\AAA_V)^\tau\back G(\AAA_V)^\tau\to J_{a_M,v}(\AAA_v)^{\tau^n}\back G(\AAA_v)^{\tau^n}.
$$
The factor $|V|^{\dim(\ago^G_M)}$ comes from the comparison of the weights.

It is also clear that we have
$$
|D^G(X_{a_M,v})|_v=|D^G(Y_v)|_v=|D^L(Y_v)|_v.
$$

From now on, in order to lighten the notations, we write $\tau$ instead of $\tau^n$. Since $Y_v$ and $X_{a_M,v}$ have the same image under $\chi_M$ and since this image belongs to the open subset $\car_M^{M-\reg}(\oc_v)$, there exists $m\in M(\oc_v)$ such that $X_{a_M,v}=\Ad(m^{-1})Y_v$. It follows that $m\tau(m)^{-1}$ is an element of the normalizer of $T$ in $M(\oc_v)$ and that its image in $W^M$ is $w$. We deduce from this that $m\tau(m)^{-1}$ normalizes $L$. The automorphism $\Int(m\tau(m)^{-1})$ preserves the pair $(B_L,T)$ of the pinning $(N_L,T,X_\alpha)$ and acts on $X_\alpha$ by a homothety by an element of $\oc_v^\times$. If one changes $m$ into $tm$ with $t\in T(\oc_v)$, the ratio of this homothety is multiplied by $\alpha(t\tau(t)^{-1})$. We therefore see that by taking a judicious element $t$, one may suppose that $\Int(m\tau(m)^{-1})$ preserves the pinning.

It is clear that we have
$$
|D^G(X_{a_M,v})|_v=|D^G(Y_v)|_v=|D^L(Y_v)|_v.
$$
By abuse of notation, we still denote by $w$ the inner automorphism of $G$ given by $\Int(m\tau(m)^{-1})$. The map $g\mapsto mgm^{-1}$ induces a bijection from $J_{a_M,v}(F_v)^\tau\back G(F_v)^\tau$ onto $T(F_v)^{w\tau}\back G(F_v)^{w\tau}$ and one has
\[
\begin{aligned}
\int_{J_{a_M,v}(F_v)^\tau\back G(F_v)^\tau}&
\mathbf{1}_{\ggo(\oc_v)}(\Ad(g^{-1})X_{a_M,v})v^G_M(g)\,dg\\
&=\int_{T(F_v)^{w\tau}\back G(F_v)^{w\tau}}
\mathbf{1}_{\ggo(\oc_v)}(\Ad(g^{-1})Y_v)v^G_M(g)\,dg,
\end{aligned}
\]
where the measure on $T(F_v)^{w\tau}\back G(F_v)^{w\tau}$ is deduced from that on $J_{a_M,v}(F_v)^\tau\back G(F_v)^\tau$ by transport by the previous bijection. Using as before the Iwasawa decomposition for $Q\in\pc(L)$, we see that if $g\in G(F_v)$ satisfies $\Ad(g^{-1})Y_v\in\ggo(\oc_v)$ then $g=ls$ with $l\in L(F_v)$ and $s\in G(\oc_v)$. Suppose moreover that we have $w\tau(g)=g$. Then $l^{-1}w\tau(l)\in L(\oc_v)$ and, by a new application of Lang's theorem, we see that one may suppose $w\tau(l)=l$ and $s\in G(\oc)^{w\tau}$. We thus have
\[
\begin{aligned}
\int_{T(F_v)^{w\tau}\back G(F_v)^{w\tau}}&
\mathbf{1}_{\ggo(\oc_v)}(\Ad(g^{-1})Y_v)v^G_M(g)\,dg\\
&=\int_{T(F_v)^{w\tau}\back L(F_v)^{w\tau}}
\mathbf{1}_{\lgo(\oc_v)}(\Ad(l^{-1})Y_v)v^G_M(l)\,dl.
\end{aligned}
\]
We then use Arthur's descent formula
$$
v^G_M(l)=\sum_{R\in\lc^G(T)}d^G_T(M,R)v^{Q_R}_T(l),
$$
where $R\in\lc^G(T)\mapsto Q_R\in\pc(R)$ is a certain section (defined on the $R$ such that $d^G_T(M,R)\neq 0$) of the map $\fc^G(T)\to\lc(T)$ defined by $Q\mapsto M_Q$. In what follows we restrict ourselves to the $R$ such that $d^G_T(M,R)\neq 0$. This therefore implies $R\neq T$. The weight $v^{Q_R}_T$ was defined in the proof of Proposition~\ref{thm:17.4.1}. Here is how it is computed here. The group $Q_R\cap L$ is a parabolic subgroup of $L$. Since $L$ is of semi-simple rank 1, one of two things holds
\begin{itemize}
\item[--] either $Q_R\cap L$ is a Borel subgroup of $L$; one then has $l\in tN_{Q_R\cap L}(F_v)L(\oc_v)$ with $t\in T(F_v)$ and hence $v^{Q_R}_T(l)=v^R_T(t)=0$ (since $T\neq R$); \item[--] or else $L\subset R$ and then $v^{Q_R}_T(l)=v^R_T(l)$.
\end{itemize}
We now place ourselves in the second case. Let us recall that the weight $v^R_T(l)$ is the volume of the projection onto $\ago^R_T$ of the convex hull of the points $-H_B(l)$ for $B\in\pc^R(T)$. Now this convex hull is here contained in a line, so the volume is therefore zero unless $\dim(\ago^R_T)=1$, in which case $R=L$. Arthur's descent formula therefore simplifies into
$$
v^G_M(l)=d^G_T(M,L)v^L_T(l)
$$
which gives the desired result.
\end{preuve}
\end{paragr}

\begin{paragr}[Local computations of ordinary orbital integrals.]\label{S:17.7}
We give some local computations of ordinary orbital integrals. Let $v\in V$ be a closed point of $C$ and $V$ the orbit of $v$ under the Frobenius $\tau$. Let $s\in\Tc$ and $\Mc'=\Mc_s$. Let $a_M=(h_{a_M},t)\in\Ac_{M'}(k)^\tau$.

\begin{lemma}\label{thm:17.7.1}
We suppose that $h_{a_M}(c)$ does not belong to $\mathfrak{D}^M$. Then
$$
J^{M,s}_{M,V}(a_M)=1.
$$
\end{lemma}

\begin{preuve}
Let us start from the equality
\[
\begin{aligned}
J^{M,s}_{M,V}(a_M)&=|D^M(X_{a_M,V})|_V^{1/2}\\
&\quad\times\int_{(J_{a_M,V}(\AAA_V)\back M(\AAA_V))^\tau}
\bg s,\inv_{a_M,V}(m)\bd\mathbf{1}_{\mgo(\oc_V)}(\Ad(m^{-1})X_{a_M,V})\,dm.
\end{aligned}
\]
Let $m\in M(\AAA_V)$ be such that $\Ad(m^{-1})X_{a_M,V}\in\mgo(\oc_V)$. Since $X_{a_M,V}$ is still regular semi-simple in reduction (given that $h_{a_M}(c)$ does not belong to $\mathfrak{D}^M$), we see that $m\in J_{a_M,V}(\AAA_V)M(\oc_V)$. As in the proof of the lemma, we deduce from this that $\inv_{a_M,V}(m)=0$. Since moreover $|D^M(X_{a_M,V})|_V^{1/2}=1$, we have
\[
\begin{aligned}
J^{M,s}_{M,V}(a_M)&=\int_{J_{a_M,V}(\oc_V)^\tau\back M(\oc_V)^\tau}
\mathbf{1}_{\mgo(\oc_V)}(\Ad(m^{-1})X_{a_M,V})\,dm.\\
&=\vol(J_{a_M,V}(\oc_V)^\tau)^{-1}\vol(M(\oc_V)^\tau).
\end{aligned}
\]
Now $\vol(M(\oc_V)^\tau)=1$ and $\vol(J_{a_M,V}(\oc_V)^\tau)=1$ since $J_{a_M,V}(\oc_V)^\tau$ is the maximal compact subgroup of $J_{a_M,V}(\AAA_V)^\tau$.
\end{preuve}

\begin{lemma}\label{thm:17.7.2}
We suppose that
\begin{itemize}
\item[--] $h_{a_M}(c)$ is a smooth point of $\mathfrak{R}^M_{M'}${\rm ;} \item[--] $h_{a_M}(c)$ does not belong to $\mathfrak{D}^{M'}${\rm ;} \item[--] the intersection of $h_{a_M}(C)$ with $\mathfrak{R}^M_{M'}$ is transverse at $h_{a_M}(c)$.
\end{itemize}
Then
$$
J^{M,s}_{M,V}(a_M)=J^{M',1}_{M',V}=1.
$$
\end{lemma}

\begin{preuve}
The first equality is merely a particular case of the Langlands-Shelstad fundamental lemma. The second follows from Lemma~\ref{thm:17.7.1} above. The reader who would balk at using the fundamental lemma may check the equality $J^{M,s}_{M,V}(a_M)=1$ directly: one reduces to a computation on a group of semi-simple rank 1 which one does by hand (cf.\ also \cite[\S8.5.9]{Ngo-lemme}).
\end{preuve}

\begin{lemma}\label{thm:17.7.3}
We suppose that
\begin{itemize}
\item[--] $h_{a_M}(c)$ does not belong to $\mathfrak{R}^M_{M'}${\rm ;} \item[--] $h_{a_M}(c)$ is a smooth point of $\mathfrak{D}^{M'}${\rm ;} \item[--] the intersection of $h_{a_M}(C)$ with $\mathfrak{D}^{M'}$ is transverse at $h_{a_M}(c)$.
\end{itemize}
Then
$$
J^{M,s}_{M,V}(a_M)=J^{M',1}_{M',V}(a_M)\neq 0.
$$
\end{lemma}

\begin{preuve}
Here again, the first equality is a particular case of the Langlands-Shelstad fundamental lemma. The integrand in
$$
J^{M',1}_{M',V}=|D^{M'}(X_{a_M,V})|_V^{1/2} \int_{(J_{a_M,V}(\AAA_V)\back M'(\AAA_V))^\tau} \mathbf{1}_{\mgo'(\oc_V)}(\Ad(m^{-1})X_{a_M,V})\,dm
$$
is positive and the integral is clearly bounded below by $\vol((J_{a_M,V}(\oc_V)\back M'(\oc_V))^\tau)>0$. Of course, it is also possible to carry out directly the computations of $J^{M,s}_{M,V}(a_M)$ and $J^{M',1}_{M',V}$ (one reduces to a computation on a group of semi-simple rank 1, cf.\ also \cite[\S8.5.8]{Ngo-lemme}).
\end{preuve}
\end{paragr}

\begin{paragr}\label{S:17.8}
Let $V$ be a non-empty $\tau$-stable set of closed points of $C$. Let $s\in\Tc$ be such that $\Mc'=\Mc_s\neq\Mc$. Let $a_M=(h_{a_M},t)\in\Ac_{M'}(k)^\tau$. Let $J_V:\ec^G_M(s)\cup\lc^{\Gc}(\Mc)\to\CC$ be the map defined for every $\Hc\in\ec^G_M(s)$ by $J_V(\Hc)=J^{H,1}_{M',V}(a_M)$ and for every $\Lc\in\lc^{\Gc}(\Mc)$ by $J_V(\Lc)=J^{L,s}_{M,V}(a_M)$.

\begin{lemma}\label{thm:17.8.1}
We suppose that $V$ satisfies one of the three following conditions{\rm :}
\begin{enumerate}
\item For every $v\in V$, the point $h_{a_M}(v)$ does not belong to $\mathfrak{D}^G${\rm ;} \item $V$ is the orbit under $\tau$ of a point $v$ which satisfies assertion 1 of Theorem~\ref{thm:17.3.1}{\rm ;} \item $V$ is the orbit under $\tau$ of a point $v$ which satisfies assertion 2 of Theorem~\ref{thm:17.3.1}.
\end{enumerate}
Then $J_V$ is stabilizing and $J_V(\Mc)\neq 0$.
\end{lemma}

\begin{preuve}
Suppose assertion 1 holds. By Lemma~\ref{thm:17.6.1}, we have $J_V(\Hc)=0$ unless $\Hc=\Mc$ or $\Hc=\Mc'$. Moreover, $J_V(\Mc)=J_V(\Mc')=1$ by Lemma~\ref{thm:17.7.1}. Hence $J_V$ is stabilizing.

Suppose assertion 2 holds. In this case, $h_{a_M}(v)$ is a smooth point of $\mathfrak{R}^G_{M'}$ and does not belong to $\mathfrak{D}^{M'}$, and the intersection of $h_{a_M}(v)$ with $\mathfrak{R}^G_{M'}$ is transverse at $h_{a_M}(v)$. Suppose moreover, to begin with, that $h_{a_M}(v)$ does not belong to $\mathfrak{R}^G_M$. Consequently, the intersection of $h_a(C)$ with $\mathfrak{R}^M_{M'}$ is transverse at $h_{a_M}(v)$. We thus have $J_V(\Mc)=J_V(\Mc')=1$ by Lemma~\ref{thm:17.7.2}. By Lemma~\ref{thm:17.6.1}, we have $J_V(\Lc)=0$ for every $\Lc\in\lc^{\Gc}(\Mc)$ with $\Lc\neq\Mc$. For every $\Hc\in\ec^G_M(s)$, we have $\Hc\cap\Mc=\Mc'$; consequently, $h_{a_M}(v)$ does not belong to $\mathfrak{R}^H_{M'}$ and we have $J_V(\Hc)=0$ by Lemma~\ref{thm:17.6.1}. Thus $J_V$ is stabilizing.

Suppose still that assertion 2 holds but this time suppose that $h_{a_M}(v)$ belongs to $\mathfrak{R}^G_M$: under our hypotheses, $h_{a_M}(v)$ is a point of transverse intersection of $h_{a_M}(C)$ and $\mathfrak{R}^G_M$ and it does not belong to $\mathfrak{D}^M$.

Let $Y_v\in\tgo(\oc_v)$, $L\in\lc^G$, $n$ an integer and $w$ an automorphism of $L$ which preserves $T$, such that these elements satisfy the conclusions of Lemma~\ref{thm:17.6.2}. Let us introduce the map $J_T:\ec^G_T(1)\to\CC$ defined by $J_T(\Tc)=1$,
$$
J_T(\Lc)=|D^L(Y_v)|^{|V|/2}|V|^{\dim(\ago^G_M)} \int_{T(F_v)^{w\tau^n}\back L(F_v)^{w\tau^n}} \mathbf{1}_{\lgo(\oc_v)}(\Ad(l^{-1})Y_v)v^L_T(l)\,dl,
$$
and $J_T(\Hc)=0$ if $\Hc\notin\{\Tc,\Lc\}$. For every $\Hc\in\lc^{\Gc}(\Mc)$, the expression
\begin{equation}\label{eq:17.8.1}
\sum_{\Rc\in\lc^{\Hc}(\Tc)}d^{\Hc}_{\Tc}(\Mc,\Rc)J_T(\Rc)
\end{equation}
equals
\begin{itemize}
\item[--] $J_T(\Tc)=1=J_V(\Mc)$ if $\Hc=\Mc$ (cf.\ Lemma~\ref{thm:17.7.1}); \item[--] $d^{\Hc}_{\Tc}(\Mc,\Lc)J_T(\Lc)=J_V(\Hc)$ if $\Hc\neq\Mc$ (cf.\ Lemma~\ref{thm:17.6.2}).
\end{itemize}
Hence the expression~(\ref{eq:17.8.1}) is equal to $J_V(\Hc)$. For every $\Hc\in\ec^G_M(s)$, we have $h_{a_M}(v)\in\mathfrak{R}^H_{M'}$ if and only if $\Lc\subset\Hc$, and the expression
\begin{equation}\label{eq:17.8.2}
\sum_{\Rc\in\lc^{\Hc}(\Tc)}d^{\Hc}_{\Tc}(\Mc',\Rc)J_T(\Rc)=J_V(\Hc)
\end{equation}
equals
\begin{itemize}
\item[--] $J_T(\Tc)=1=J_V(\Mc')$ if $\Hc=\Mc'$ (cf.\ Lemma~\ref{thm:17.7.1}); \item[--] $d^{\Hc}_{\Tc}(\Mc,\Lc)J_T(\Lc)=J_V(\Hc)$ if $\Lc\subset\Hc$ (cf.\ Lemma~\ref{thm:17.6.2}); \item[--] $0$ if $\Lc\not\subset\Hc$ (Lemma~\ref{thm:17.6.1}).
\end{itemize}
In any case the expression~(\ref{eq:17.8.2}) is equal to $J_V$. Since $J_T$ is (trivially) stabilizing, $J_V$ is stabilizing (cf.\ Proposition~\ref{thm:16.3.7}). Lemma~\ref{thm:17.7.1} gives $J_V(\Mc)=J_V(\Mc')=$.

Suppose now that assertion 3 holds. In this case, Lemma~\ref{thm:17.6.1} indicates that $J_V(\Hc)=0$ unless $\Hc=\Mc$ or $\Mc'$. We have $J_V(\Mc)=J_V(\Mc')\neq 0$ by Lemma~\ref{thm:17.7.3}. Hence $J_V$ is stabilizing.
\end{preuve}
\end{paragr}
\section{Truncated affine Springer fibers}\label{sec:18}

\begin{paragr}\label{S:18.1}
In this section, the notation is that used throughout the article. We add the following specifications. Let $F=k((\eps))$ and $\oc=k[[\eps]]$. Let $\val$ be the usual valuation of $F$. Let $F_0=\Fq((\eps))$. We recall that we denote by $\tau$ the Frobenius automorphism of $k$ given by raising to the power $q$. We again denote by $\tau$ the Frobenius automorphism of $F$, so that we have $F^\tau=F_0$. Let $\Fcl$ be a separable closure of $F$ and $I=\Gal(\Fcl/F)$. Let $\Fcl_0$ be the algebraic closure of $F_0$ in $\Fcl$. Let $\Ga=\Gal(\Fcl_0/F_0)$.

We assume that the pair $(G,T)$ is defined over $\Fq$. For every $P\in\pc(M)$, we define, as in~\S\ref{S:11.6}, a map
$$
H_P:G(F)\to X_*(M)
$$
defined for every $\la\in X^*(P)$ by $\la(H_P(g))=-\val(\la(p))$ where $p\in P(F)$ satisfies $g\in pG(\oc)$ (by the Iwasawa decomposition).

We recall that $\ago_T$ is equipped with a $W^G$-invariant scalar product and that all its subspaces are equipped with the Euclidean measure deduced from it (cf.~\S\ref{S:2.4}). Let
$$
H_P^G:G(F)\to\ago_M^G
$$
be the composite of the map $H_P$ with the orthogonal projection $\ago_M\to\ago_M^G$. The restrictions to $M(F)$ of these maps do not depend on the choice of $P\in\pc(M)$: we denote them respectively by $H_M$ and $H_M^G$.
\end{paragr}

\begin{paragr}[Truncated affine Grassmannian.]\label{S:18.2}
Let $\Xgo^G$ be the affine Grassmannian: it is a projective ind-$k$-scheme whose set of $k$-points is the quotient $G(F)/G(\oc)$. Let $M\in\lc$. For every $P\in\pc(M)$, the map $H_P^G$ on $G(F)$ induces a map $H_P^G:\Xgo^G\to\ago_M^G$. Let $\xi\in\ago_M^G$. We define the truncated affine Grassmannian
$$
\Xgo_M^{G,\xi}
$$
by the following condition: for every $x\in\Xgo^G$, we have $x\in\Xgo_M^{G,\xi}$ if and only if $\xi$ belongs to the convex hull of the projections onto $\ago_M^G$ of the points $-H_P^G(x)$ for $P\in\pc(M)$.

\begin{proposition}\label{thm:18.2.1}
The truncated affine Grassmannian $\Xgo_M^{G,\xi}$ is an open sub-ind-scheme of $\Xgo^G$.
\end{proposition}

\begin{preuve}
It is well known that the convex hull of the points $-H_P^G(x)$ is the intersection of the cones $-H_P^G(x)-{}^+\overline{\ago_P}$, where ${}^+\overline{\ago_P}$ is the positive obtuse Weyl chamber in $\ago_M$ (this is a result of Arthur, cf.~\cite{Arthur-characters}). It therefore suffices to show that the condition $H_P^G(x)\in-\xi-{}^+\overline{\ago_P}$ defines an open subset of $\Xgo^G$. For that, we proceed as in \cite[proof of Proposition 6.4]{CL-homologie}.
\end{preuve}
\end{paragr}

\begin{paragr}[Truncated affine Springer fiber.]\label{S:18.3}
Let $Y\in\mgo(\oc)$ be a semi-simple and $G$-regular element. Following Kazhdan-Lusztig (cf.~\cite{Kazhdan-Lusztig}), we introduce the affine Springer fiber
$$
\Xgo_Y^G=\{x\in\Xgo^G\mid \Ad(x^{-1})Y\in\ggo(\oc)\}.
$$
It is a closed sub-ind-scheme of $\Xgo^G$. In fact, $\Xgo_Y^G$ is a $k$-scheme locally of finite type and of finite dimension (cf.~\cite{Kazhdan-Lusztig}). The truncated affine Springer fiber $\Xgo_{M,Y}^{G,\xi}$ is the open subset of $\Xgo_Y^G$ defined by
$$
\Xgo_{M,Y}^{G,\xi}=\Xgo_Y^G\cap\Xgo_M^{G,\xi}.
$$
Let $J$ be the centralizer of $Y$ in $G\times_k F$. It is a maximal $F$-subtorus of $G\times_k F$. The group $J(F)$ acts on the set of $k$-points of $\Xgo_Y^G$. It does not preserve the truncation $\Xgo_{M,Y}^{G,\xi}$. Let $J'$ be the $F$-subtorus of $J$ defined by
$$
X_*(J')=\Ker(\la\in X_*(J)\mapsto H_M^G(\eps^\la)).
$$
The group $J'(F)$ acts on $\Xgo_Y^G$ and preserves the truncation $\Xgo_{M,Y}^{G,\xi}$. Let
$$
\Lambda=X_*(J')^I
$$
be the group of cocharacters of the maximal split subtorus of $J'$. Via the morphism $\la\mapsto\eps^\la$, we identify the discrete group $\Lambda$ with a subgroup of $J'(F)$.

\begin{proposition}\label{thm:18.3.1}
The quotient $\Lambda\back\Xgo_M^{G,\xi}(Y)$ is a quasi-projective variety defined over $k$.
\end{proposition}

\begin{preuve}
This is proved as in \cite{Kazhdan-Lusztig}.
\end{preuve}
\end{paragr}

\begin{paragr}[A local weighted orbital $s$-integral.]\label{S:18.4}
We continue with the preceding notation, assuming henceforth that $\tau(Y)=Y$. In this case, $J$ is equipped with an action of $\tau$ and we have an exact sequence in pointed set cohomology
\begin{equation}\label{eq:18.4.1}
1\longrightarrow J(F)^\tau\longrightarrow G(F)^\tau\longrightarrow
(J(F)\back G(F))^\tau\longrightarrow \Ker(J(F)_\tau\to G(F)_\tau)\longrightarrow 1,
\end{equation}
where, as in~\S\ref{S:11.1}, the index $\tau$ denotes the set of $\tau$-conjugacy classes. We equip the groups $J(F)^\tau$ and $G(F)^\tau$ with the Haar measures which give measure $1$ respectively to the maximal compact subgroup of $J(F)^\tau$ and to the subgroup $G(\oc)^\tau$ of $G(F)^\tau$. By Kottwitz (\cite[\S7.6]{Kottwitz-isocrystals-II}), we have a canonical exact sequence
\begin{equation}\label{eq:18.4.2}
0\longrightarrow J(F)_1\longrightarrow J(F)\longrightarrow X_*(J)_I\longrightarrow 0
\end{equation}
which induces, at the level of the co-invariants under $\tau$, an isomorphism
\begin{equation}\label{eq:18.4.3}
J(F)_\tau\simeq X_*(J)_\Ga.
\end{equation}
In particular, the group $J(F)_\tau$ is discrete. We equip it with the counting measure. We deduce from this a right $G(F)^\tau$-invariant measure on the quotient $(J(F)\back G(F))^\tau$.

Let
$$
\inv_Y:(J(F)\back G(F))^\tau\to X_*(J)_\Ga
$$
be the composite of the coboundary map of~(\ref{eq:18.4.1}) with Kottwitz's isomorphism~(\ref{eq:18.4.3}).

\begin{definition}\label{thm:18.4.1}
For $Y$ and $M$ as above and $s\in\hat{J}^\Ga$, the weighted orbital $s$-integral $J_M^{G,s}(Y)$ is defined by
\begin{equation}\label{eq:18.4.4}
J_M^{G,s}(Y)=|D^G(Y)|_F^{1/2}\int_{(J(F)\back G(F))^\tau}
\bg s,\inv_Y(g)\bd\,\mathbf{1}_{\ggo(\oc)}(\Ad^{-1}(Y))\,v_M^G(g)\,dg
\end{equation}
where
\begin{itemize}
\item[--] the absolute value is $|\cdot|_F=q^{-\val(\cdot)}$; \item[--] the pairing is the one between $\hat{J}^\Ga$ and $X_*(J)_\Ga$; \item[--] $\mathbf{1}_{\ggo(\oc)}$ is the characteristic function of $\ggo(\oc)$; \item[--] Arthur's weight $v_M^G(g)$ is the volume in $\ago_M^G$ of the convex hull of the points $-H_P^G(g)$ for $P\in\pc(M)$.
\end{itemize}
\end{definition}
\end{paragr}

\begin{paragr}[Geometric interpretation of the weighted orbital integrals.]\label{S:18.5}
Since $Y$ is $\tau$-stable, the quotient $\Lambda\back\Xgo_{M,Y}^{G,\xi}$ is defined over $\Fq$. Let $\eta\in\Hom_\ZZ(\Lambda,\Qb^\times)$ be a character of finite order. We assume that this character is fixed by $\tau$ or, what amounts to the same, fixed by $\Ga$. By the choice of an embedding of $\Qb$ into $\Ql$, we interpret $\eta$ as a continuous $\ell$-adic character of $\Lambda\rtimes\Gal(k/\Fq)$ trivial on the second factor. We deduce from this a local system $\lc_\eta$ on $\Lambda\back\Xgo_{M,Y}^{G,\xi}$ defined over $\Fq$ which becomes trivial when it is pulled back by the covering $\Xgo_{M,Y}^{G,\xi}\to\Lambda\back\Xgo_{M,Y}^{G,\xi}$. The action of $J'(F)$ on $\Xgo_{M,Y}^{G,\xi}$ commutes with that of $\Lambda$. We deduce from this an action of $J'(F)$ on the cohomology $H^\bullet(\Xgo_{M,Y}^{G,\xi},\lc_\eta)$. By a homotopy lemma, this action factors through its group of connected components. The latter is identified with the group $X_*(J')_I$ by the sequence~(\ref{eq:18.4.2}) applied to the $F$-torus $J'$. Now the discrete subgroup $\Lambda$ of $J'(F)$ acts by the character $\eta^{-1}$ on the cohomology $H^\bullet(\Xgo_{M,Y}^{G,\xi},\lc_\eta)$. Since the image of $\Lambda$ in $X_*(J')_I$ is of finite index, we deduce from this the following decomposition:
\begin{equation}\label{eq:18.5.1}
H^\bullet(\Xgo_{M,Y}^{G,\xi},\lc_\eta)=\bigoplus_{s'}
H^\bullet(\Xgo_{M,Y}^{G,\xi},\lc_\eta)_{s'},
\end{equation}
where the sum is taken over the finite set of $s'\in(\hat{J}')^I= \Hom(X_*(J')_I,\CC^\times)$ which induce the character $\eta^{-1}$ on $\Lambda$ and where $H^\bullet(\Xgo_{M,Y}^{G,\xi},\lc_\eta)_{s'}$ denotes the direct factor of $H^\bullet(\Xgo_{M,Y}^{G,\xi},\lc_\eta)$ on which $J'(F)$ acts via the character $s'$ of $X_*(J')_I$. Let $p_M^G$ be the orthogonal projection of $\ago_M$ onto $\ago_M^G$. The lattice $H_M^G(J(F)^\tau)$ is of finite index in $p_M^G(X_*(M))$. Let $\mu_1,\ldots,\mu_n$ be a system of representatives of the quotient. Let
$$
\psi:\Lambda_\tau\to X_*(J')_\Ga\simeq J'(F)_\tau
$$
and
$$
\varphi:X_*(J')_\Ga\simeq J'(F)_\tau\to X_*(J)_\Ga\simeq J(F)_\tau
$$
be the canonical maps, the isomorphisms being those given by~(\ref{eq:18.4.3}). We note that the kernels as well as the cokernel of $\psi$ are finite. Let
$$
J(F)^0=J(F)\cap\Ker(H_M^G).
$$
The quotient $J(F)^{0,\tau}\back J(F)^\tau$ is discrete and it is equipped with the counting measure. We deduce from this a Haar measure on $J(F)^{0,\tau}$.

\begin{theorem}\label{thm:18.5.1}
Let $s\in\hat{J}^\Ga$ be an element of finite order. Let $s'\in(\hat{J}')^\Ga$ be the image of $s$ under the canonical morphism $\hat{J}^\Ga\to(\hat{J}')^\Ga$. Let $\eta\in\Hom_\ZZ(\Lambda,\Qb^\times)$ be the character induced by the inverse of $s'$.

For every $\xi$ in general position, we have the equality
$$
J_M^{G,s}(Y)=c_M^G(Y)\sum_{i=1}^n \trace(\tau^{-1},H^\bullet(\Xgo_{M,Y}^{G,\xi+\mu_i},\lc_\eta)_{s'}),
$$
where
$$
c_M^G(Y)=\frac{|\Ker(\psi)|}{|\Coker(\psi)|}\cdot|\Ker(\varphi)|\cdot \frac{\vol(\Lambda^\tau\back J(F)^{0,\tau})} {\vol\bigl(\ago_M^G/p_M^G(X_*(M))\bigr)}\cdot|D^G(Y)|_F^{1/2}.
$$
\end{theorem}
\end{paragr}

\begin{paragr}[Proof of Theorem~\ref{thm:18.5.1}.]\label{S:18.6}
We continue with the preceding notation. For $g\in G(F)$, let $\mathbf{1}_{M,g}$ be the function on $\ago_M^G$ which is the characteristic function of the convex hull of the points $-H_P(g)$, $P\in\pc(M)$. For every $\xi\in\ago_M^G$, let
$$
w_M^{G,\xi}(g)=|\{\mu\in p_M^G(X_*(M))\mid \mathbf{1}_{M,g}(\xi+\mu)=1\}|.
$$
The weight $w_M^{G,\xi}$ is right invariant under $M(F)$, hence under $J(F)$. We then define the weighted orbital integral $J_M^{G,\xi,s}$ by formula~(\ref{eq:18.4.4}) in which we replace the weight $v_M^G$ by $w_M^{G,\xi}$.

\begin{proposition}\label{thm:18.6.1}
Under the hypotheses of Theorem~\ref{thm:18.5.1}, for every $\xi\in\ago_M^G$, we have
$$
J_M^{G,\xi,s}(Y)=c_M^G(Y)\cdot\vol\bigl(\ago_M^G/p_M^G(X_*(M))\bigr)\cdot \sum_{i=1}^n\trace(\tau^{-1},H^\bullet(\Xgo_{M,Y}^{G,\xi+\mu_i},\lc_\eta)_{s'}).
$$
\end{proposition}

\begin{preuve}
This proposition generalizes to the truncated case Theorem 15.8 of the article~\cite{GKM} of Goresky-Kottwitz-MacPherson. Our proof is inspired by theirs. Let $j\in J'(F)$. We do not distinguish in the notation between $j$ and its image in $J'(F)_\tau$. Let $\xi\in\ago_M^G$. The Grothendieck-Lefschetz trace formula gives the equality
$$
\bg s',j\bd\trace((j\tau)^{-1},H^\bullet(\Xgo_{M,Y}^{G,\xi},\lc_\eta))= \sum_{x\in(\Lambda\back\Xgo_{M,Y}^{G,\xi})^{j\tau}}\bg s',\psi(\la_x)j\bd
$$
where $\la_x$ is the class in $\Lambda_\tau$ of an element $\la\in\Lambda$ which satisfies $\la j\tau(x)=x$ and where the pairing is the canonical pairing between . The above expression depends only on the image of $j$ in the cokernel $\Coker(\psi)$ of $\psi$. We thus have
\begin{align*}
\sum_{j\in\Coker(\psi)}&\bg s',j\bd
\trace((j\tau)^{-1},H^\bullet(\Xgo_{M,Y}^{G,\xi},\lc_\eta))\\
&=|\Ker(\psi)|\sum_{j\in J'(F)_\tau}\bg s',j\bd\,
|\Lambda^\tau\back(\Xgo_{M,Y}^{G,\xi})^{j\tau}|.
\end{align*}
By the decomposition~(\ref{eq:18.5.1}), the left-hand side in the above expression equals
\begin{equation}\label{eq:18.6.1}
|\Coker(\psi)|\trace(\tau^{-1},H^\bullet(\Xgo_{M,Y}^{G,\xi},\lc_\eta)_{s'}).
\end{equation}
With the help of Lang's isogeny theorem, we see that we have
\begin{align*}
\sum_{j\in J'(F)_\tau}&\bg s',j\bd\,
|\Lambda^\tau\back(\Xgo_{M,Y}^{G,\xi})^{j\tau}|\\
&=\sum_{j\in J'(F)_\tau}\bg s',j\bd\int_{\Lambda^\tau\back G(F)^{j\tau}}
\mathbf{1}_{\ggo(\oc)}(\Ad(g^{-1})Y)\mathbf{1}_{M,g}(\xi)\,dg\\
&=\vol(\Lambda^\tau\back J(F)^{0,\tau})\sum_{j\in J'(F)_\tau}\bg s',j\bd
\int_{J(F)^\tau\back G(F)^{j\tau}}
\mathbf{1}_{\ggo(\oc)}(\Ad(g^{-1})Y)w_J^\xi(g)\,dg,
\end{align*}
where
$$
w_J^\xi(g)=\sum_{h\in J(F)^{0,\tau}\back J(F)^\tau}\mathbf{1}_{M,g}(\xi)
$$
is the number of points of the lattice $\xi+H_M^G(J(F)^\tau)$ in the convex hull of the $-H_P(g)$, $P\in\pc(M)$. We thus have
$$
\sum_{i=1}^n w_J^{\xi+\mu_i}(g)=w_M^{G,\xi}(g).
$$
Consequently, when we sum the expression~(\ref{eq:18.6.1}) over $\xi+\mu_1,\ldots,\xi+\mu_n$, we find
\begin{equation}\label{eq:18.6.2}
\begin{split}
|\Ker(\psi)|\,&\vol(\Lambda^\tau\back J(F)^{0,\tau})\\
&\times\sum_{j\in J'(F)_\tau}\bg s',j\bd\int_{J(F)^\tau\back G(F)^{j\tau}}
\mathbf{1}_{\ggo(\oc)}(\Ad(g^{-1})Y)w_M^{G,\xi}(g)\,dg.
\end{split}
\end{equation}
We have $\bg s',j\bd=\bg s,\varphi(j)\bd$ and $J(F)^\tau\back G(F)^{j\tau}$ is identified with the set of $g\in(J(F)\back G(F))^\tau$ such that $\inv_Y(g)=\varphi(j)$. In fact, it follows from Lemma~\ref{thm:11.4.4} that the map $\inv_Y$ takes its values in the image of $\varphi$. The expression~(\ref{eq:18.6.2}) is therefore equal to
$$
|\Ker(\psi)|\cdot|\Ker(\varphi)|\cdot\vol(\Lambda^\tau\back J(F)^{0,\tau}) |D^G(Y)|^{-1/2}J_M^{G,\xi,s}(Y).
$$
The proposition follows.
\end{preuve}

\begin{proposition}\label{thm:18.6.2}
With the notation of Theorem~\ref{thm:18.5.1}, for every $\xi$ in general position, we have
$$
J_M^{G,\xi,s}(Y)=\vol\bigl(\ago_M^G/p_M^G(X_*(M))\bigr)J_M^{G,s}(Y).
$$
\end{proposition}

\begin{preuve}
It is identical to the proof of Theorem 11.14.2 of~\cite{LFPI}.
\end{preuve}
\end{paragr}

\begin{paragr}[Local constancy of the truncated affine Springer fibers.]\label{S:18.7}
Let $M\in\lc$ and $a\in\car_M(\oc)$ a generically $G$-regular element. We fix a Kostant section in $\mgo$, denoted by $\eps^M$. We set $X_a=\eps^M(a)$. We denote by $J_a$ the centralizer of $X_a$ in $M\times_k\oc$: it is a smooth $\oc$-group scheme. Let $\xi\in\ago_M^G$. We set
$$
\Xgo_M^{G,\xi}(a)=\Xgo_{M,X_a}^{G,\xi}.
$$
The aim of this paragraph is to prove the following proposition, which generalizes to the truncated case Proposition 3.5.1 of~\cite{Ngo-lemme}.

\begin{proposition}\label{thm:18.7.1}
Let $a\in\car_M(\oc)\cap\car_M^{\reg}(F)$. There exists an integer $n$ such that for every $a'\in\car_M(\oc)$ such that
$$
a'=a\quad\mod\eps^n
$$
there exists $h\in M(\oc)$ such that
\begin{enumerate}
\item $\Ad(h)$ induces an isomorphism of group schemes from $J_a$ onto $J_{a'}$; \item left translation by $h$ induces an isomorphism from $\Xgo_M^{G,\xi}(a)$ onto $\Xgo_M^{G,\xi}(a')$.
\end{enumerate}
Moreover, if $a$ and $a'$ are fixed by a power of $\tau$, we may further assume that $h$ is also fixed by the same power of $\tau$.
\end{proposition}

\begin{preuve}
Let $r$ and $n_0$ be integers for which the conclusions of Lemma~\ref{thm:18.7.3} hold for the element $X_a$ and the group $M$. Let
$$
v=\val(\det(\ad(X_a)_{|\ngo});
$$
it is an integer. Let $\zeta:\ngo\to N$ be an isomorphism of affine varieties and $i$ an integer which satisfy the conclusions of Lemma~\ref{thm:18.7.2}. Let $c\geq 0$ be an integer such that
\begin{equation}\label{eq:18.7.1}
\forall Y\in\eps^c\mgo(\oc)\ \text{and}\ \forall n\in\zeta(\eps^{-iv}\oc)
\ \text{we have}\ \Ad(n^{-1})Y-Y\in\ngo(\oc).
\end{equation}
Let $c_1\geq\max(1,n_0,c+r)$. We further assume that $c_1$ is large enough so that for every $a'\in\car_M(\oc)$ congruent to $a$ modulo $\eps^{c_1}$ we have $a'\in\car^{\reg}(F)$ and $\val(\det(\ad(X_a)_{|\ngo})=v$. By the approximation Lemma~\ref{thm:18.7.3}, for every $a'\in\car_M(\oc)$ such that $a'=a \mod\eps^{c_1}$, there exists $h\in M(\oc)$ such that
\begin{equation}\label{eq:18.7.2}
h=1\quad\mod\eps^c
\end{equation}
and
\begin{equation}\label{eq:18.7.3}
[\Ad(h)X_a,X_{a'}]=0.
\end{equation}

Conjugation by $h$ induces a morphism from $J_a$ onto the centralizer $I_{\Ad(h)X_a}$ of $\Ad(h)X_a$, which is none other than $J_{a'}$ since $X_a$ and $X_{a'}$ are generically semi-simple regular. By Proposition~\ref{thm:3.2.1}, this morphism extends. Arguing with the inverse morphism $\Ad(h^{-1})$, we see that $\Ad(h)$ induces an isomorphism from $J_a$ onto $J_{a'}$. Let $gG(\oc)\in\Xgo_M^{G,\xi}(a)$. Let us show that $hgG(\oc)\in\Xgo_M^{G,\xi}(a)$. Let us note right away that we have $H_M(h)=0$ since $h=1\mod\eps$. In particular, translation by $h$ preserves the truncation. It therefore suffices to prove that we have
$$
\Ad(hg)^{-1}X_{a'}\in\ggo(\oc).
$$
Let $P\in\pc(M)$. By the Iwasawa decomposition, we may well assume that we have $g=mn$ with $m\in M(F)$ and $n\in N(F)$. Since $gG(\oc)\in\Xgo_M^{G,\xi}(a)$, we have $\Ad(mn)^{-1}X_a\in\ggo(\oc)$. We write
$$
\Ad(mn)^{-1}X_a=[\Ad(n^{-1})(\Ad(m)^{-1}X_a)-\Ad(m)^{-1}X_a]+\Ad(m^{-1})X_a,
$$
where the bracket belongs to $\ngo(F)$ and where $\Ad(m^{-1})X_a\in\mgo(F)$. We thus have the two conditions
\begin{equation}\label{eq:18.7.4}
\Ad(m^{-1})X_a\in\mgo(\oc)
\end{equation}
and
\begin{equation}\label{eq:18.7.5}
\Ad(n^{-1})(\Ad(m)^{-1}X_a)-\Ad(m)^{-1}X_a\in\ngo(\oc).
\end{equation}
Let $Y=\Ad(m^{-1})X_a$. Condition~(\ref{eq:18.7.4}) implies that $Y\in\mgo(\oc)$. Let $I_Y$ be the centralizer of $Y$ in $M$. The inner automorphism $\Ad(m^{-1})$ induces an isomorphism on the generic fiber from $J_a$ onto $I_Y$. By Proposition~\ref{thm:3.2.1}, we know that this isomorphism extends to a morphism of group schemes over $\oc$ from $J_a$ into $I_Y$. In particular, we have
\begin{equation}\label{eq:18.7.6}
\Ad(m^{-1})(\Lie(J_a)(\oc))\subset I_Y(\oc)\subset\mgo(\oc).
\end{equation}
Since $\Ad(h^{-1})(X_{a'})$ commutes with $X_a$, we have $\Ad(h^{-1})(X_{a'})\in\Lie(J_a)(\oc)$. In view of~(\ref{eq:18.7.6}), we obtain
\begin{equation}\label{eq:18.7.7}
\Ad(m^{-1})(\Ad(h^{-1})(X_{a'}))\in\mgo(\oc).
\end{equation}
Conditions~(\ref{eq:18.7.2}) and~(\ref{eq:18.7.3}) imply that we have $X_a-\Ad(h^{-1})X_{a'}\in\eps^c\Lie(J_a)(\oc)$. By~(\ref{eq:18.7.6}), we have
\begin{equation}\label{eq:18.7.8}
\Ad(m^{-1})(X_a-\Ad(h^{-1})X_{a'})\in\eps^c\mgo(\oc).
\end{equation}
By Lemma~\ref{thm:18.7.2}, conditions~(\ref{eq:18.7.4}) and~(\ref{eq:18.7.5}) and the fact that $v=\val(\det(\ad(X_a)_\ngo))$ imply that $n\in\zeta(\eps^{-iv}\ngo(\oc))$. By~(\ref{eq:18.7.8}) and~(\ref{eq:18.7.1}), we have
\begin{equation}\label{eq:18.7.9}
\Ad(n^{-1})(\Ad(m^{-1})(X_a-\Ad(h^{-1})X_{a'}))-\Ad(m^{-1})(X_a-\Ad(h^{-1})X_{a'})
\in\ngo(\oc).
\end{equation}
Hence condition~(\ref{eq:18.7.5}) implies the assertion
\begin{equation}\label{eq:18.7.10}
\Ad(n^{-1})(\Ad(m^{-1})(\Ad(h^{-1})X_{a'}))-\Ad(m^{-1})(\Ad(h^{-1})X_{a'})
\in\ngo(\oc)
\end{equation}
which, combined with~(\ref{eq:18.7.7}), gives
$$
\Ad(hmn)^{-1}(X_{a'})\in\ggo(\oc)
$$
which is what we were seeking to prove. The same reasoning shows the converse implication. The proposition follows, the last assertion resulting from the last assertion of~\ref{thm:18.7.3}.
\end{preuve}

The following two lemmas were used in the proof of Proposition~\ref{thm:18.7.1}.

\begin{lemma}\label{thm:18.7.2}
Let $M\in\lc$ and $P\in\pc(M)$. There exist an isomorphism of affine varieties $\zeta:\ngo\to N$, a morphism $\varphi:\ngo\times\mgo\to\ngo$ and an integer $i$ such that for every $Y\in\mgo$ with $\det(\ad(Y)_{|\ngo})\neq 0$ the morphism $N\to\ngo$ induced by
$$
n\mapsto\Ad(n^{-1})Y-Y
$$
is an isomorphism with inverse $U\in\ngo\mapsto \zeta((\det(\ad(Y)_{|\ngo}))^{-i}\varphi(U,Y))$.
\end{lemma}

\begin{preuve}
This statement is probably well known, but for lack of a reference we sketch a proof. Let $B\subset P$ be a Borel subgroup which contains $T$. Let $\Delta_B$ be the set of simple roots of $T$ in $B$ and $\Delta_M\subset\Delta_B$ the subset of the roots outside $M$. The restriction of the roots in $\Delta_M$ to the neutral component, denoted by $A_M$, of the center of $M$ is injective. We identify the elements of $\Delta_M$ with their restriction to $A_M$. Let $\Phi^N$ be the set of roots of $T$ in $N$. The restriction of a root in $\Phi^N$ to $A_M$ is a linear combination with non-negative integer coefficients, not all zero, of elements of $\Delta_M$. For every integer $l$, let $\Phi_l^N$ be the subset of the roots for which the sum of the coefficients on $\Delta_M$ is equal to $l$. For $l$ large enough the set $\Phi_l^N$ is empty. For every $\al\in\Phi^N$, we fix an additive $1$-parameter subgroup $\zeta_\al:\mathbb{G}_a\to N$ such that $t\zeta_\al(u)t^{-1}=\zeta_\al(\al(t)u)$ for $t\in T$ and $u\in\mathbb{G}_a$.

For every $\al\in\Phi^N$, let $\ngo_\al\subset\ngo$ be the weight subspace of weight $\al$. For every integer $l$ we set
$$
\ngo_l=\bigoplus_{\al\in\Phi_l^N}\ngo_\al
$$
and
$$
\ngo^l=\bigoplus_{i>l}\ngo_i.
$$
By the choice of a basis, we identify $\ngo_l$ with $\mathbb{G}_a^{\Phi_l^N}$. The morphism $\zeta_l=\prod_{\al\in\Phi_l^N}\zeta_\al$ (the product is the product in $N$ according to an order fixed once and for all) sends $\ngo_l$ onto a subvariety of $N$ denoted by $N_l$. Let $N^l=N_{l+1}\cdot N_{l+2}\cdot\ldots$\,. (the $\cdot$ denotes the product in $N$). We then have the following properties for every $l\geq 0$:
\begin{enumerate}
\item $\ngo_l$ is stable under the adjoint action of $M$; \item $[\ngo,\ngo_l]\subset\ngo^l$; \item $N^l$ is a subgroup of $N$ stable under conjugation by $M$ and $[N,N_l]\subset N^l$; \item the quotient $N^l/N^{l+1}$ is commutative and the morphism $\zeta_{l+1}$ induces a group isomorphism from $\ngo_{l+1}$ onto $N^l/N^{l+1}$. We may furthermore assume that the differential of this morphism is the canonical morphism $\ngo_{l+1}\to\ngo^l/\ngo^{l+1}$.
\end{enumerate}
Let $l\geq 0$, $n\in N^l$ and $y\in M$. We thus have $n^{-1}yny^{-1}\in N^l$. There exists $U\in\ngo_{l+1}$ such that $n\in\zeta_{l+1}(U)N^{l+1}$. Using the above properties, we see that there exists a morphism $\psi:\ngo_{l+1}\times M\to\ngo_{l+1}$ such that $n^{-1}yny^{-1}\in\zeta_{l+1}(\psi(U,y))N^{l+1}$. We note that the morphism $\psi$ is linear in the first variable. Taking a differential, we obtain the following relation for every $Y\in\mgo$:
\begin{equation}\label{eq:18.7.11}
\Ad(n^{-1})Y-Y=\psi'(U,Y)\quad\mod(\ngo^{l+1})
\end{equation}
for a certain bilinear map $\psi'$. Differentiating once more, we find
\begin{equation}\label{eq:18.7.12}
\psi'(U,Y)=[Y,U]\quad\mod(\ngo^{l+1}).
\end{equation}

This said, we are in a position to prove the lemma. Since the morphism $n\mapsto\Ad(n^{-1})Y-Y$ sends $N^l$ into $\ngo^l$, we may refine the statement for the subgroups $N^l$. By induction, we assume the refined statement true for $N^1$. It then remains to prove it for $N^0=N$. Let us start from $V\in\ngo$. Let $V_1$ be the projection of $V$ onto $\ngo_1$ along $\ngo=\ngo_1\oplus\ngo^1$. Let $n=\zeta_1(U)$ with $U\in\ngo_1$. From~(\ref{eq:18.7.11}) and~(\ref{eq:18.7.12}), we deduce
$$
\Ad(n^{-1})Y-Y=[Y,U]\quad\mod(\ngo^1).
$$

Let us further assume that $\det(\ad(Y)_{|\ngo})\neq 0$. There then exists a unique $U\in\ngo_1$ such that $[Y,U]=V_1$. Moreover $U$ is a polynomial $\varphi_1$ in $V_1$ and $Y$ multiplied by the inverse of $\det(\ad(Y)_{|\ngo_1})\neq 0$. For such a $U$, we thus have
$$
\Ad(n^{-1})Y-Y-V\in\ngo^1.
$$
Applying the induction hypothesis, an element of $\ngo_1$ is of the form $-(\Ad(n_1^{-1})Y-Y)$ with $n_1\in N^1$ which is given by a certain morphism of the type described in the statement. We then have
\begin{align*}
\Ad((nn_1)^{-1})Y-Y-V&=\Ad(n_1^{-1})(\Ad(n^{-1})Y-Y)+\Ad(n^{-1})Y-Y-V\\
&=\Ad(n^{-1})Y-Y+\Ad(n^{-1})Y-Y-V\quad\mod(\ngo^2)\\
&=0\quad\mod(\ngo^2).
\end{align*}

We obtain the lemma by induction.
\end{preuve}

\begin{lemma}\label{thm:18.7.3}
Let $X\in\ggo(\oc)$ be a generically semi-simple regular element. Then there exist integers $r$ and $n_0$ such that for every $n\geq n_0$ and every $Y\in\ggo(\oc)$ such that
$$
Y=X\quad\mod\eps^n
$$
there exists $g\in G(\oc)$ such that
\begin{enumerate}
\item $g=1\mod\eps^{n-r}$; \item $[\Ad(g)X,Y]=0$.
\end{enumerate}
If, moreover, $X$ and $Y$ are fixed by a power $\tau$, then we may require $g$ to be fixed by the same power of $\tau$.
\end{lemma}

\begin{preuve}
Let $Y\in\ggo(\oc)$. The functor which to every $\oc$-algebra $A$ associates the set
$$
\{g\in G(A)\mid [\Ad(g)X,Y]\}
$$
is representable by a scheme over $\oc$. Let us assume $Y=X\mod\eps^n$ for $n$ large enough. The generic fiber of this $\oc$-scheme is smooth, in particular if $Y=X$. This is not in general the case for its special fiber. The unit section provides a point in $\oc/\eps^n\oc$. By a variant of Lemma 1 of~\cite{Elkik}, using successive approximations, one can lift this section to an $\oc$-point which satisfies the required conditions.
\end{preuve}
\end{paragr}
\section{End of the proof of the weighted fundamental lemma}\label{sec:19}

\begin{paragr}\label{S:19.1}
Let $k$ be an algebraic closure of a finite field $\Fq$ with $q$ elements and $\tau$ the Frobenius of $k/\Fq$. Let $K_1=\Fq((\eps))$ and $K=k((\eps))$. Let $\oc_1=\Fq[[\eps]]$. Let $(G,T)$ be a pair consisting of a connected reductive group defined over $\Fq$ and a maximal sub-torus $T$ split over $\Fq$. Let $M\in\lc^G(T)$. Let $s\in\Tc$ and $\Mc'=\Mc_s$. We assume that we have $\Mc\neq\Mc'$ and $Z^0_{\Mc}=Z^0_{\Mc'}$. Let $M'$ be the group split over $\Fq$ dual to $\Mc'$. Then $M'$ is an ``unramified'' elliptic endoscopic group of $M$. For all $Y\in\mgo'(K_1)$ and $X\in\mgo(K_1)$ semi-simple and $G$-regular, Langlands and Shelstad (cf.\ \cite{Langlands-Shelstad}) have defined a transfer factor $\Delta_{M',M}(Y,X)$; more precisely, what is used here is a variant of their definition adapted to Lie algebras and ``deprived of the factor $\Delta_{\mathrm{IV}}$'', which Waldspurger introduced in \cite{Waldspurger-transfert}. This factor is normalized by the choice of a pinning: here we take a pinning of $M$ which is stable under the Frobenius $\tau$. Then the transfer factor does not depend on the choice of such a pinning (cf.\ \cite{Hales} II.7).

For every $Y\in\mgo'(K_1)$ semi-simple and $G$-regular, let
$$
\begin{aligned}
J^G_{M',M}(Y)=\sum_X &\Delta_{M',M}(Y,X)|D^G(X)|^{1/2}\\
&\times\int_{J_X(K_1)\back G(K_1)}\mathbf{1}_{\ggo(\oc_1)}(\Ad(g^{-1})X)v^G_M(g)dg,
\end{aligned}
$$
where the sum runs over the set of semi-simple $G$-regular elements of $\mgo(K_1)$ taken modulo $M(K_1)$-conjugacy. The weight $v^G_M$ is the Arthur weight considered previously. The conventions on Haar measures are also those considered previously. Let us note that $\Delta_{M'}(Y,X)$ is zero except on a finite subset of the summation set.

\begin{theorem}[Weighted fundamental lemma for Lie algebras over local fields of equal characteristic]\label{thm:19.1.1}
We keep the notation and the hypotheses above. Let $Y\in\mgo'(K_1)$ be semi-simple and $G$-regular. The map $J:\ec^G_{M,\el}(s)\cup\Gc\to\CC$ defined by
$$
J(\Gc)=J^G_{M',M}(Y)
$$
and for all $\Hc\in\ec^G_{M,\el}(s)$ by
$$
J(\Hc)=J^H_{M',M'}(Y)
$$
is stabilizing.
\end{theorem}

\begin{preuve}
It is entirely analogous to Ngô's proof of the ordinary fundamental lemma (cf.\ \cite[\S8.6]{Ngo-lemme}). Let $Y\in\mgo'(K_1)$ be semi-simple and $G$-regular. We assume that $\chi_{M'}(Y)\in\car_{M'}(\oc_1)$, since otherwise the theorem is trivially true ($J$ is then identically zero). The values of $J$ depend only on $b=\chi_{M'}(Y)\in\car_{M'}(\oc_1)$. Let $b_1\in\car_M(\oc_1)$ be deduced from $b$ by the canonical morphism $\car_{M'}\to\car_M$. For every choice of a pinning of $M$ stable under the Frobenius $\tau$, one has a Kostant section $\eps_M:\car_M\to\mgo$. Let $Y_1=\eps_M(b_1)$. By Kottwitz (cf.\ \cite{Kottwitz-transfer}), one has $\Delta_{M',M}(Y,Y_1)=1$, so that one checks that one has
\begin{equation}\label{eq:19.1.1}
J^G_{M',M}(Y)=J^{G,s,\tau}_M(Y_1),
\end{equation}
where the right-hand side is the one defined in~(\ref{eq:18.4.4}) relative to the Frobenius $\tau$.

Let $n\geq 1$ be an integer, $K_n=\FF_{q^n}((\eps))$ and $\oc_n=\FF_{q^n}[[\eps]]$. For every integer $n\geq 1$ and every $b\in\car_{M'}(\oc_n)$, we define $J_{n,b}$ as the map $J$ defined as above but associated with the field $K_n$ and with an arbitrary lift $Y_n\in\mgo'(K_n)$ of $b$ (that is, $\chi_{M'}(Y_n)=b$). We always normalize the transfer factors by a pinning stable under the Frobenius $\tau$. As before, one has $J_{n,b}=J^{G,s,\tau^n}_M(Y_1)$ where $Y_1$ is deduced from $b$ by the choice of a Kostant section attached to a pinning stable under the Frobenius $\tau$ and where the right-hand side is the one defined in~(\ref{eq:18.4.4}) but relative to the Frobenius $\tau^n$.

Let $b\in\car_{M'}(\oc_1)$. It follows from Theorem~\ref{thm:18.5.1} that there exist almost-zero families of complex numbers $(c_i)_{i\in\NN}$ and $(\la_i)_{i\in\NN}$ which depend on $b$ and $\Gc$ such that for every integer $n\in\NN$ one has
$$
J_{n,b}(\Gc)=\sum_{i\in\NN}c_i\la_i^n.
$$
The same holds for every $\Hc\in\ec^G_{M,\el}(s)$. Consequently, one deduces from a standard argument (the use of a Van der Monde determinant) that the map $J_{1,b}$ is stabilizing (which is our goal) as soon as there exists an integer $m$ such that for every integer $n\geq m$, the map $J_{n,b}$ is stabilizing.

Let $N$ be an integer such that for every $n\geq 1$ and every $b'\in\car_{M'}(\oc_n)$ which satisfies
\begin{equation}\label{eq:19.1.2}
b'\equiv b \mod \eps^N,
\end{equation}
the affine Springer fibers $\Xgo^{G,\xi}_M(b)$ and $\Xgo^{G,\xi}_M(b')$ on the one hand and $\Xgo^{H,\xi}_{M'}(b)$ and $\Xgo^{H,\xi}_{M'}(b')$ for every $\Hc\in\ec^G_{M,\el}(s)$ on the other hand satisfy the conclusions of Proposition~\ref{thm:18.7.1}. We also assume that, under the condition~(\ref{eq:19.1.2}), one has
\begin{equation}\label{eq:19.1.3}
D^G(b')\in D^G(b)\oc_n^\times.
\end{equation}
By Proposition~\ref{thm:18.7.1}, such an $N$ exists. For every $b'$ which satisfies~(\ref{eq:19.1.2}), Proposition~\ref{thm:18.7.1} and Theorem~\ref{thm:18.5.1} imply that for every $n\geq 1$, one has
\begin{equation}\label{eq:19.1.4}
J_{n,b'}=J_{n,b}.
\end{equation}

We fix a geometrically connected projective curve $C_0$ over $\Fq$ and two distinct points $v$ and $\infty$ in $C_0(\Fq)$. Let $g$ be its genus. We fix an isomorphism between $K_1$ and the completion $F_{0,v}$ of the function field of $C_0$ at the point $v$. Thus $b$ is identified with an element of $\car_{M'}(\oc_{0,v})$ where $\oc_{0,v}$ is the ring of integers of $F_{0,v}$. Let $D$ be a divisor of $C_0$ whose support contains neither $\infty$ nor $v$ and whose degree satisfies the following inequalities:
\begin{equation}\label{eq:19.1.5}
\val_v(D^G(b))\leq 2|W^G|^{-1}(\deg(D)-2g+2)
\end{equation}
and
\begin{equation}\label{eq:19.1.6}
N\rang(G)+2g<\deg(D).
\end{equation}
This last inequality implies that the restriction map
\begin{equation}\label{eq:19.1.7}
H^0(C_0,\car_{M',D})\to\car_{M'}(\oc_{0,v}/\varpi_v^N)
\end{equation}
is surjective. The closed subscheme $\Ac_{M'}(b,N)$ of $\Ac_{M'}$ formed of the $(h_a,t)$ such that $h_a$ has the same image as $b$ under the map~(\ref{eq:19.1.7}) is of codimension $N\rang(G)$ in $\Ac_{M'}$.

Let $\Ac_{M'}(b,N)^0$ be the open subset of $\Ac_{M'}(b,N)$ formed of the $a=(h_a,t)$ such that, for every point $c\neq v$ of $C_0$, if $h_a(c)$ belongs to $\mathfrak{R}^G_{M'}\cup\mathfrak{D}^{M'}$, this point is smooth and the intersection of $h_a(C)$ with $\mathfrak{R}^G_{M'}\cup\mathfrak{D}^{M'}$ is transverse there. By a variant of a theorem of Bertini (cf.\ \cite[\S8.6.6]{Ngo-lemme}), $\Ac_{M'}(b,N)^0$ is non-empty. Since the codimension of $\Ac_{M'}-\Ac^{\el}_{M'}$ in $\Ac_{M'}$ is $\geq\deg(D)$ and since the codimension of $\Ac_{M'}(b,N)^0$ in $\Ac_{M'}$ is $N\rang(G)<\deg(D)$ (cf.~(\ref{eq:19.1.6})), one has $\Ac_{M'}(b,N)^0\cap\Ac^{\el}_{M'}\neq\emptyset$. This intersection is even of dimension $>0$. By Deligne's theorem, there exists an integer $m$ such that for every $n\geq m$, there exists $a_M\in\Ac_{M'}(k)^{\tau^n}$ which belongs to this intersection. Now such a point satisfies the hypotheses of Theorem~\ref{thm:9.1.3}: the image of $a_M$ in $\Ac_G$ therefore belongs to the open subset $\Ac^{\mathrm{bon}}_G$. Theorem~\ref{thm:17.3.1} then implies that the map $J_{n,b'}$, where $b'$ is the image of $a_M$ in $\car_{M'}(\oc_n)$, is stabilizing. Now $J_{n,b'}=J_{n,b}$ by the equality~(\ref{eq:19.1.4}). The map $J_{n,b}$ is therefore stabilizing for every $n\geq m$. As we explained at the beginning of the proof, this entails that $J_{b,1}$ is stabilizing, which is what we had to prove.
\end{preuve}
\end{paragr}
\begin{thebibliography}{99}

\bibitem{Arthur-characters} J.~Arthur, The characters of discrete series as orbital integrals, \emph{Invent. Math.} \textbf{32} (1976), 205--261. \href{https://mathscinet.ams.org/mathscinet-getitem?mr=0412348}{MR 0412348}. \href{https://zbmath.org/?q=an:0359.22008}{Zbl 0359.22008}. \href{https://doi.org/10.1007/BF01425569}{doi:10.1007/BF01425569}.

\bibitem{Arthur-invariant-I} J.~Arthur, The invariant trace formula. I. Local theory, \emph{J. Amer. Math. Soc.} \textbf{1} (1988), 323--383. \href{https://mathscinet.ams.org/mathscinet-getitem?mr=0928262}{MR 0928262}. \href{https://zbmath.org/?q=an:0682.10021}{Zbl 0682.10021}. \href{https://doi.org/10.2307/1990920}{doi:10.2307/1990920}.

\bibitem{Arthur-transfer} J.~Arthur, On the transfer of distributions: weighted orbital integrals, \emph{Duke Math. J.} \textbf{99} (1999), 209--283. \href{https://mathscinet.ams.org/mathscinet-getitem?mr=1708030}{MR 1708030}. \href{https://zbmath.org/?q=an:0938.22019}{Zbl 0938.22019}. \href{https://doi.org/10.1215/S0012-7094-99-09909-X}{doi:10.1215/S0012-7094-99-09909-X}.

\bibitem{Arthur-stable-I} J.~Arthur, A stable trace formula. I. General expansions, \emph{J. Inst. Math. Jussieu} \textbf{1} (2002), 175--277. \href{https://mathscinet.ams.org/mathscinet-getitem?mr=1954821}{MR 1954821}. \href{https://zbmath.org/?q=an:1040.11038}{Zbl 1040.11038}. \href{https://doi.org/10.1017/S1474748002000051}{doi:10.1017/S1474748002000051}.

\bibitem{Beauville-Laszlo} A.~Beauville and Y.~Laszlo, Un lemme de descente, \emph{C. R. Acad. Sci. Paris S\'er. I Math.} \textbf{320} (1995), 335--340. \href{https://mathscinet.ams.org/mathscinet-getitem?mr=1320381}{MR 1320381}. \href{https://zbmath.org/?q=an:0852.13005}{Zbl 0852.13005}.

\bibitem{BBD} A.~A. Be\u{\i}linson, J.~Bernstein, and P.~Deligne, Faisceaux pervers, in \emph{Analysis and Topology on Singular Spaces, I} (Luminy, 1981), Ast\'erisque \textbf{100}, Soc. Math. France, Paris, 1982, pp.~5--171. \href{https://mathscinet.ams.org/mathscinet-getitem?mr=0751966}{MR 0751966}. \href{https://zbmath.org/?q=an:0536.14011}{Zbl 0536.14011}.

\bibitem{Bezrukavnikov} R.~Bezrukavnikov, The dimension of the fixed point set on affine flag manifolds, \emph{Math. Res. Lett.} \textbf{3} (1996), 185--189. \href{https://mathscinet.ams.org/mathscinet-getitem?mr=1386839}{MR 1386839}. \href{https://zbmath.org/?q=an:0874.20033}{Zbl 0874.20033}.

\bibitem{Biswas-Ramanan} I.~Biswas and S.~Ramanan, An infinitesimal study of the moduli of Hitchin pairs, \emph{J. London Math. Soc.} \textbf{49} (1994), 219--231. \href{https://mathscinet.ams.org/mathscinet-getitem?mr=1260109}{MR 1260109}. \href{https://zbmath.org/?q=an:0819.58007}{Zbl 0819.58007}. \href{https://doi.org/10.1112/jlms/49.2.219}{doi:10.1112/jlms/49.2.219}.

\bibitem{Bourbaki} N.~Bourbaki, \emph{\'El\'ements de Math\'ematique. Fasc. XXXIV. Groupes et Alg\`ebres de Lie. Chapitre IV: Groupes de Coxeter et Syst\`emes de Tits. Chapitre V: Groupes Engendr\'es par des R\'eflexions. Chapitre VI: Syst\`emes de Racines}, Actualit\'es Scientifiques et Industrielles \textbf{1337}, Hermann, Paris, 1968. \href{https://mathscinet.ams.org/mathscinet-getitem?mr=0240238}{MR 0240238}. \href{https://zbmath.org/?q=an:0186.33001}{Zbl 0186.33001}.

\bibitem{CL-homologie} P.-H. Chaudouard and G.~Laumon, Sur l'homologie des fibres de Springer affines tronqu\'ees, \emph{Duke Math. J.} \textbf{145} (2008), 443--535. \href{https://mathscinet.ams.org/mathscinet-getitem?mr=2462112}{MR 2462112}. \href{https://zbmath.org/?q=an:1206.11065}{Zbl 1206.11065}. \href{https://doi.org/10.1215/00127094-2008-057}{doi:10.1215/00127094-2008-057}.

\bibitem{LFPI} P.-H. Chaudouard and G.~Laumon, Le lemme fondamental pond\'er\'e. I. Constructions g\'eom\'etriques, \emph{Compos. Math.} \textbf{146} (2010), 1416--1506. \href{https://mathscinet.ams.org/mathscinet-getitem?mr=2735371}{MR 2735371}. \href{https://zbmath.org/?q=an:1206.14026}{Zbl 1206.14026}. \href{https://doi.org/10.1112/S0010437X10004756}{doi:10.1112/S0010437X10004756}.

\bibitem{Donagi-Gaitsgory} R.~Y. Donagi and D.~Gaitsgory, The gerbe of Higgs bundles, \emph{Transform. Groups} \textbf{7} (2002), 109--153. \href{https://mathscinet.ams.org/mathscinet-getitem?mr=1903115}{MR 1903115}. \href{https://zbmath.org/?q=an:1083.14519}{Zbl 1083.14519}. \href{https://doi.org/10.1007/s00031-002-0008-z}{doi:10.1007/s00031-002-0008-z}.

\bibitem{Elkik} R.~Elkik, Solutions d'\'equations \`a coefficients dans un anneau hens\'elien, \emph{Ann. Sci. \'Ecole Norm. Sup.} \textbf{6} (1973), 553--603 (1974). \href{https://mathscinet.ams.org/mathscinet-getitem?mr=0345966}{MR 0345966}. \href{https://zbmath.org/?q=an:0327.14001}{Zbl 0327.14001}.

\bibitem{GKM} M.~Goresky, R.~Kottwitz, and R.~Macpherson, Homology of affine Springer fibers in the unramified case, \emph{Duke Math. J.} \textbf{121} (2004), 509--561. \href{https://mathscinet.ams.org/mathscinet-getitem?mr=2040285}{MR 2040285}. \href{https://zbmath.org/?q=an:1162.14311}{Zbl 1162.14311}. \href{https://doi.org/10.1215/S0012-7094-04-12135-9}{doi:10.1215/S0012-7094-04-12135-9}.

\bibitem{EGAIII} A.~Grothendieck, \'El\'ements de g\'eom\'etrie alg\'ebrique. III. \'Etude cohomologique des faisceaux coh\'erents. I, \emph{Inst. Hautes \'Etudes Sci. Publ. Math.} \textbf{11} (1961), 5--167. \href{https://mathscinet.ams.org/mathscinet-getitem?mr=0217085}{MR 0217085}. \href{https://zbmath.org/?q=an:0118.36206}{Zbl 0118.36206}. \href{https://doi.org/10.1007/BF02684274}{doi:10.1007/BF02684274}.

\bibitem{EGAIV3} A.~Grothendieck, \'El\'ements de g\'eom\'etrie alg\'ebrique. IV. \'Etude locale des sch\'emas et des morphismes de sch\'emas. III, \emph{Inst. Hautes \'Etudes Sci. Publ. Math.} \textbf{28} (1966), 5--248. \href{https://mathscinet.ams.org/mathscinet-getitem?mr=0217086}{MR 0217086}. \href{https://zbmath.org/?q=an:0144.19904}{Zbl 0144.19904}. \href{https://doi.org/10.1007/BF02684343}{doi:10.1007/BF02684343}.

\bibitem{EGAIV4} A.~Grothendieck, \'El\'ements de g\'eom\'etrie alg\'ebrique. IV. \'Etude locale des sch\'emas et des morphismes de sch\'emas. IV, \emph{Inst. Hautes \'Etudes Sci. Publ. Math.} \textbf{32} (1967), 5--361. \href{https://mathscinet.ams.org/mathscinet-getitem?mr=0238860}{MR 0238860}. \href{https://zbmath.org/?q=an:0153.22301}{Zbl 0153.22301}. Available at \href{http://www.numdam.org/item?id=PMIHES_1967__32__5_0}{numdam.org}.

\bibitem{Hales} T.~C. Hales, A simple definition of transfer factors for unramified groups, in \emph{Representation Theory of Groups and Algebras}, Contemp. Math. \textbf{145}, Amer. Math. Soc., Providence, RI, 1993, pp.~109--134. \href{https://mathscinet.ams.org/mathscinet-getitem?mr=1216184}{MR 1216184}. \href{https://zbmath.org/?q=an:0828.22015}{Zbl 0828.22015}. \href{https://doi.org/10.1090/conm/145/1216184}{doi:10.1090/conm/145/1216184}.

\bibitem{Heinloth} J.~Heinloth, Uniformization of $\mathcal{G}$-bundles, \emph{Math. Ann.} \textbf{347} (2010), 499--528. \href{https://mathscinet.ams.org/mathscinet-getitem?mr=2640041}{MR 2640041}. \href{https://zbmath.org/?q=an:1193.14014}{Zbl 1193.14014}. \href{https://doi.org/10.1007/s00208-009-0443-4}{doi:10.1007/s00208-009-0443-4}.

\bibitem{Illusie} L.~Illusie, \emph{Complexe Cotangent et D\'eformations. I}, Lecture Notes in Math. \textbf{239}, Springer-Verlag, New York, 1971. \href{https://mathscinet.ams.org/mathscinet-getitem?mr=0491680}{MR 0491680}. \href{https://zbmath.org/?q=an:0224.13014}{Zbl 0224.13014}. \href{https://doi.org/10.1007/BFb0059052}{doi:10.1007/BFb0059052}.

\bibitem{Kazhdan-Lusztig} D.~Kazhdan and G.~Lusztig, Fixed point varieties on affine flag manifolds, \emph{Israel J. Math.} \textbf{62} (1988), 129--168. \href{https://mathscinet.ams.org/mathscinet-getitem?mr=0947819}{MR 0947819}. \href{https://zbmath.org/?q=an:0658.22005}{Zbl 0658.22005}. \href{https://doi.org/10.1007/BF02787119}{doi:10.1007/BF02787119}.

\bibitem{Kottwitz-isocrystals} R.~E. Kottwitz, Isocrystals with additional structure, \emph{Compositio Math.} \textbf{56} (1985), 201--220. \href{https://mathscinet.ams.org/mathscinet-getitem?mr=0809866}{MR 0809866}. \href{https://zbmath.org/?q=an:0597.20038}{Zbl 0597.20038}. Available at \href{http://www.numdam.org/item?id=CM_1985__56_2_201_0}{numdam.org}.

\bibitem{Kottwitz-basechange} R.~E. Kottwitz, Base change for unit elements of Hecke algebras, \emph{Compositio Math.} \textbf{60} (1986), 237--250. \href{https://mathscinet.ams.org/mathscinet-getitem?mr=0868140}{MR 0868140}. Available at \href{http://www.numdam.org/item?id=CM_1986__60_2_237_0}{numdam.org}.

\bibitem{Kottwitz-isocrystals-II} R.~E. Kottwitz, Isocrystals with additional structure. II, \emph{Compositio Math.} \textbf{109} (1997), 255--339. \href{https://mathscinet.ams.org/mathscinet-getitem?mr=1485921}{MR 1485921}. \href{https://zbmath.org/?q=an:0966.20022}{Zbl 0966.20022}. \href{https://doi.org/10.1023/A:1000102604688}{doi:10.1023/A:1000102604688}.

\bibitem{Kottwitz-transfer} R.~E. Kottwitz, Transfer factors for Lie algebras, \emph{Represent. Theory} \textbf{3} (1999), 127--138. \href{https://mathscinet.ams.org/mathscinet-getitem?mr=1703328}{MR 1703328}. \href{https://zbmath.org/?q=an:1044.22011}{Zbl 1044.22011}. \href{https://doi.org/10.1090/S1088-4165-99-00068-0}{doi:10.1090/S1088-4165-99-00068-0}.

\bibitem{Langlands-Shelstad} R.~P. Langlands and D.~Shelstad, On the definition of transfer factors, \emph{Math. Ann.} \textbf{278} (1987), 219--271. \href{https://mathscinet.ams.org/mathscinet-getitem?mr=0909227}{MR 0909227}. \href{https://zbmath.org/?q=an:0644.22005}{Zbl 0644.22005}. \href{https://doi.org/10.1007/BF01458070}{doi:10.1007/BF01458070}.

\bibitem{Laumon-Ngo} G.~Laumon and B.~C. Ng\^o, Le lemme fondamental pour les groupes unitaires, \emph{Ann. of Math.} \textbf{168} (2008), 477--573. \href{https://mathscinet.ams.org/mathscinet-getitem?mr=2434884}{MR 2434884}. \href{https://zbmath.org/?q=an:1179.22019}{Zbl 1179.22019}. \href{https://doi.org/10.4007/annals.2008.168.477}{doi:10.4007/annals.2008.168.477}.

\bibitem{Ngo-endoscopie} B.~C. Ng\^o, Fibration de Hitchin et endoscopie, \emph{Invent. Math.} \textbf{164} (2006), 399--453. \href{https://mathscinet.ams.org/mathscinet-getitem?mr=2218781}{MR 2218781}. \href{https://zbmath.org/?q=an:1098.14023}{Zbl 1098.14023}. \href{https://doi.org/10.1007/s00222-005-0483-7}{doi:10.1007/s00222-005-0483-7}.

\bibitem{Ngo-lemme} B.~C. Ng\^o, Le lemme fondamental pour les alg\`ebres de Lie, \emph{Publ. Math. Inst. Hautes \'Etudes Sci.} \textbf{111} (2010), 1--169. \href{https://mathscinet.ams.org/mathscinet-getitem?mr=2653248}{MR 2653248}. \href{https://zbmath.org/?q=an:1200.22011}{Zbl 1200.22011}. \href{https://doi.org/10.1007/s10240-010-0026-7}{doi:10.1007/s10240-010-0026-7}.

\bibitem{Waldspurger-transfert} J.-L. Waldspurger, Le lemme fondamental implique le transfert, \emph{Compositio Math.} \textbf{105} (1997), 153--236. \href{https://mathscinet.ams.org/mathscinet-getitem?mr=1440722}{MR 1440722}. \href{https://zbmath.org/?q=an:0871.22005}{Zbl 0871.22005}. \href{https://doi.org/10.1023/A:1000103112268}{doi:10.1023/A:1000103112268}.

\bibitem{Waldspurger-tordu} J.-L. Waldspurger, \`A propos du lemme fondamental pond\'er\'e tordu, \emph{Math. Ann.} \textbf{343} (2009), 103--174. \href{https://mathscinet.ams.org/mathscinet-getitem?mr=2448443}{MR 2448443}. \href{https://zbmath.org/?q=an:1158.22011}{Zbl 1158.22011}. \href{https://doi.org/10.1007/s00208-008-0267-7}{doi:10.1007/s00208-008-0267-7}.

\bibitem{Waldspurger-endoscopie} J.-L. Waldspurger, Endoscopie et changement de caract\'eristique: int\'egrales orbitales pond\'er\'ees, \emph{Ann. Inst. Fourier} (Grenoble) \textbf{59} (2009), 1753--1818. \href{https://mathscinet.ams.org/mathscinet-getitem?mr=2573190}{MR 2573190}. \href{https://zbmath.org/?q=an:1184.22003}{Zbl 1184.22003}. \href{https://doi.org/10.5802/aif.2476}{doi:10.5802/aif.2476}.

\bibitem{Whitehouse} D.~Whitehouse, The twisted weighted fundamental lemma for the transfer of automorphic forms from $\mathrm{GSp}(4)$, in \emph{Formes Automorphes. II. Le cas du Groupe} $\mathrm{GSp}(4)$, Ast\'erisque \textbf{302}, Soc. Math. France, Paris, 2005, pp.~291--436. \href{https://mathscinet.ams.org/mathscinet-getitem?mr=2234863}{MR 2234863}. \href{https://zbmath.org/?q=an:1097.11025}{Zbl 1097.11025}.

\end{thebibliography}

\bigskip
\begin{flushleft}\small
(Received: May 2, 2010)\\[1ex]
Institut de Math\'ematiques de Jussieu et\\
Universit\'e Paris Diderot (Paris 7), Paris, France\\
E-mail: \href{mailto:Chaudouard@math.jussieu.fr}{Chaudouard@math.jussieu.fr}\\[1ex]
CNRS et Universit\'e Paris-Sud, Orsay, France\\
E-mail: \href{mailto:Gerard.Laumon@math.u-psud.fr}{Gerard.Laumon@math.u-psud.fr}
\end{flushleft}

\end{document}
